\documentclass[11pt,a4paper,twoside]{book}

% ===================================================================
%  1. ΒΑΣΙΚΑ ΠΑΚΕΤΑ & ΓΛΩΣΣΑ
% ===================================================================
\usepackage[english,greek]{babel}
% ΣΗΜΕΙΩΣΗ: Χρήση του parskip για αυτόματο κενό μεταξύ παραγράφων,
% αντί για χειροκίνητο μηδενισμό της εσοχής (\parindent).
\usepackage{parskip}

% ===================================================================
%  2. ΓΡΑΜΜΑΤΟΣΕΙΡΕΣ (Απαιτείται XeLaTeX ή LuaLaTeX)
% ===================================================================
% ΣΗΜΕΙΩΣΗ: Τα inputenc και fontenc είναι περιττά με το XeLaTeX.
% Το xunicode είναι παρωχημένο και η λειτουργικότητά του
% είναι ενσωματωμένη στο fontspec.
\usepackage[no-math]{fontspec}
\usepackage{xgreek}
\let\hbar\relax % Για συμβατότητα με το πακέτο mtpro2
\defaultfontfeatures{Mapping=tex-text,Scale=MatchLowercase}

% --- Κύριες γραμματοσειρές ---
\setmainfont[Mapping=tex-text,Numbers=Lining,Scale=1.0,BoldFont={Minion Pro Bold}]{Minion Pro}
\newfontfamily\titlefont{Century Gothic Bold}
\newfontfamily\scfont{Nimbus Roman}
\newfontfamily\notoserif{Noto Serif}
\newfontfamily\urwgothic{URW Gothic}

% --- Προσαρμοσμένες εντολές για γραμματοσειρές ---
% ΣΗΜΕΙΩΣΗ: Αντικατάσταση της παλιάς εντολής \font με σύγχρονες
% εντολές του LaTeX για καλύτερη διαχείριση μεγέθους.
\newcommand{\kefalaio}[1]{{\notoserif\bfseries\fontsize{36}{42}\selectfont #1}}
\newcommand{\ArKef}[1]{{\notoserif\bfseries\itshape\fontsize{72}{80}\selectfont #1}}
\newcommand{\OnKef}[1]{{\notoserif\fontsize{40}{46}\selectfont #1}}
\newcommand{\OnPar}[1]{{\notoserif\bfseries\fontsize{18}{22}\selectfont #1}}
\newcommand{\onoma}[1]{{\urwgothic\bfseries\normalsize #1}}
\newcommand{\Askhseis}[1]{{\urwgothic\bfseries\fontsize{16}{19}\selectfont #1}}

% ===================================================================
%  3. ΓΕΩΜΕΤΡΙΑ ΣΕΛΙΔΑΣ
% ===================================================================
\usepackage[
    a4paper,
    twoside,
    inner=1.5cm,
    top=3cm,
    bottom=2cm,
    textwidth=18cm,
%    marginparsep=5mm,
%    marginparwidth=5.5cm,
    ]{geometry}

% ===================================================================
%  4. ΧΡΩΜΑΤΑ
% ===================================================================
\usepackage[dvipsnames,table,x11names]{xcolor}
\definecolor{maincolor}{RGB}{35,150,210}
\definecolor{secondarycolor}{RGB}{210,20,210}
\definecolor{thirdcolor}{RGB}{110,195,7}

% ===================================================================
%  5. ΜΑΘΗΜΑΤΙΚΑ
% ===================================================================
\usepackage{mathtools} % Το mathtools φορτώνει αυτόματα το amsmath
\let\myBbbk\Bbbk
\let\Bbbk\relax
\usepackage[amsbb,subscriptcorrection,zswash,mtpcal,mtphrb,mtpfrak]{mtpro2}
\usepackage{mathimatika,venndiagram,diffcoeff}
\allowdisplaybreaks

\providecommand\PARENS[1]{\left(#1\right)}

\makeatletter
\newsavebox{\mtp@matrixbox}
\renewenvironment{pmatrix}{%
  \matrix@check\pmatrix
  \setbox\mtp@matrixbox=\hbox\bgroup$\begin{matrix}%
}{%
  \end{matrix}$\egroup
  \PARENS{\copy\mtp@matrixbox}%
}
\renewenvironment{pmatrix*}[1][c]{%
  \matrix@check\pmatrix
  \setbox\mtp@matrixbox=\hbox\bgroup$\begin{matrix*}[#1]%
}{%
  \end{matrix*}$\egroup
  \PARENS{\copy\mtp@matrixbox}%
}
\makeatother

\DeclareMathSizes{10.95}{10.95}{7}{5}
\DeclareMathSizes{6}{6}{3.8}{2.7}
\DeclareMathSizes{8}{8}{5.1}{3.6}
\DeclareMathSizes{9}{9}{5.8}{4.1}
\DeclareMathSizes{10}{10}{6.4}{4.5}
\DeclareMathSizes{12}{12}{7.7}{5.5}
\DeclareMathSizes{14.4}{14.4}{9.2}{6.5}
\DeclareMathSizes{17.28}{17.28}{11}{7.9}
\DeclareMathSizes{20.74}{20.74}{13.3}{9.4}
\DeclareMathSizes{24.88}{24.88}{16}{11.3}

\makeatletter
\AtBeginDocument{
\check@mathfonts
\fontdimen16\textfont2=2.5pt
\fontdimen17\textfont2=2.5pt
\fontdimen14\textfont2=4.5pt
\fontdimen13\textfont2=4.5pt}
\makeatother
\DeclareMathOperator{\tr}{tr}
\DeclareMathOperator{\adj}{adj}
\DeclareMathOperator{\diag}{diag}

% ===================================================================
%  6. ΓΡΑΦΙΚΑ & ΣΧΗΜΑΤΑ
% ===================================================================
\usepackage{graphicx,tikz,pgfplots,tkz-euclide,tkz-tab}
\usetikzlibrary{shadows,calc,fadings,patterns}
\pgfplotsset{compat=newest}
\tkzSetUpPoint[size=2.8,fill=white]
\makeatletter
\xpatchcmd{\tkzTabLine}
{\node at (Z\thetkz@cnt@impair\thetkz@cnt@lg){$0$};} % search
{\node[fill=white,inner sep=.5mm] at (Z\thetkz@cnt@impair\thetkz@cnt@lg){$0$};} % replace
{}{}
\makeatother
% ΣΗΜΕΙΩΣΗ: Χρήση της σύγχρονης εντολής \tikzset αντί για \tikzstyle.
\tikzset{
    >=latex,
    l3/.style={line width=0.3mm},
    l4/.style={line width=0.4mm},
    labelbox/.style args={#1}{align=center,draw=#1,fill=#1!30!white,rectangle,rounded corners=1,font=\tiny\linespread{0.8}\selectfont}
}
\usetikzlibrary{positioning, shapes.geometric, arrows.meta}

\usepackage{tikz-3dplot}
\newcommand{\DrawSphereInCylinder}[3]{\begin{tikzpicture}[declare function={R=#2;
        r=#3;
        theta=#1;
    }]
    \pgfmathsetmacro{\myalpha}{asin(r/R)}
    \pgfmathsetmacro{\myh}{R*cos(\myalpha)}
    \pgfmathtruncatemacro{\iflag}{ifthenelse(theta>90-\myalpha,1,0)}
    \tdplotsetmaincoords{theta}{0}
    \ifnum\iflag=0\relax
        \path[ball color=gray] (0,0) circle[radius=R];          
    \else
        \pgfmathsetmacro{\myhprime}{\myh*sin(theta)+r*cos(theta)}
        \pgfmathsetmacro{\mytcrit}{-asin((r*cos(theta)*\myhprime)/(r*r-\myhprime*\myhprime))}
        \pgfmathsetmacro{\mybeta}{90-atan2(\myhprime+r*cos(theta)*sin(\mytcrit),r*cos(\mytcrit))}
        \pgfmathsetmacro{\mytcrit}{asin((cos(theta)*\myhprime)/(r*sin(theta)*sin(theta)))}
        \path[ball color=orange]  (90-\mybeta:R) arc[start angle=90-\mybeta,end angle=-90+\mybeta,radius=R]
        [tdplotmaincoords] 
        -- ({r*cos(-\mytcrit)},{-r*sin(\mytcrit)},-\myh)
        arc[start angle=-\mytcrit,end angle=-180+\mytcrit,radius=r]
        [tdplotscreencoords]
        -- (-90-\mybeta:R) arc[start angle=-90-\mybeta,end angle=-270+\mybeta,radius=R];
    \fi
    \begin{scope}[tdplotmaincoords]
        \ifnum\iflag=0\relax
            \draw[dashed] (0,0,-\myh) circle[radius=r] (-r,0,\myh) edge (-r,0,-\myh) (r,0,\myh) edge (r,0,-\myh);
        \else 
            \draw[dashed] ({r*cos(-\mytcrit)},{-r*sin(\mytcrit)},-\myh)
            arc[start angle=-\mytcrit,end angle=180+\mytcrit,radius=r]
            (-r,0,\myh) edge (-r,0,-\myh) (r,0,\myh) edge (r,0,-\myh);
        \fi    
        \clip (0,0,\myh) circle[radius=r];
        \path[left color=orange!50!black,right color=orange,middle color=orange!20] (r/2,0,\myh) circle[radius=1.5*r];
        \fill[white] (0,0,-\myh) circle[radius=r];  
    \end{scope}
    \path (current bounding box.north) node[above]{};\end{tikzpicture}}
% ===================================================================
%  7. ΔΙΑΦΟΡΑ ΠΑΚΕΤΑ
% ===================================================================
\usepackage{gensymb,fontawesome5,eurosym,titletoc,sidenotes,multicol,adjmulticol,siunitx,fullwidth,svg,wrap-rl,wrapfig}
\usepackage[framemethod=tikz]{mdframed}
\usepackage[explicit]{titlesec}

% ===================================================================
%  8. ΚΕΦΑΛΙΔΑ & ΥΠΟΣΕΛΙΔΟ (fancyhdr)
% ===================================================================
\usepackage{fancyhdr}
\pagestyle{fancy}
%\fancyheadoffset[LE]{\marginparwidth+\marginparsep}
%\fancyheadoffset[RO]{\marginparwidth+\marginparsep}
\renewcommand{\headrulewidth}{\iftopfloat{0pt}{.5pt}}
\renewcommand{\chaptermark}[1]{\markboth{#1}{}}
\renewcommand{\sectionmark}[1]{\markright{\it\thesection\ #1}}
\fancyhf{}
\fancyhead[LE]{\thepage\ $\cdot$\ \nouppercase{\leftmark}}
\fancyhead[RO]{\nouppercase{\rightmark} $\cdot$\ \thepage}
\fancypagestyle{plain}{%
\fancyhead{}
\renewcommand{\headrulewidth}{0pt}}

% ===================================================================
%  9. ΟΡΙΣΜΟΙ - ΘΕΩΡΗΜΑΤΑ - ΔΙΑΦΟΡΑ ΠΛΑΙΣΙΑ (tcolorbox)
% ===================================================================
% ΣΗΜΕΙΩΣΗ: Ολόκληρη η ενότητα έχει αναδομηθεί με τη βιβλιοθήκη 'theorems'
% του tcolorbox. Αυτό κάνει τον κώδικα δραματικά πιο καθαρό,
% ευέλικτο και εύκολο στη συντήρηση.
\usepackage[most]{tcolorbox}
\tcbuselibrary{skins,theorems,breakable}

% ΣΤΥΛ ΛΥΣΗΣ - ΑΠΟΔΕΙΞΗΣ
\newcounter{lysh}[chapter]
\renewcommand{\thelysh}{
	\bf{\titlefont{\arabic{lysh}}}
}
\newcommand{\lysh}[2]{\refstepcounter{lysh}{\vspace{4mm}\titlefont{\faCheck\ {\large Λύση άσκησης 1}.\!\thelysh\ - \textcolor{maincolor!80!black}{\titlefont{Θέμα #2ο (#1)}}}}}
\newcommand{\apodeiksi}{\textcolor{secondarycolor!80!black}{\titlefont{\faPen\ \textbf{ΑΠΟΔΕΙΞΗ}\\}}}

% Κάρτα πληροφοριών με φωτογραφία
\newcommand{\Info}[5]{
\begin{tcolorbox}[
    enhanced,
    width=5.5cm,
    arc=2mm,
    colback=white,
    colframe=black,
    boxrule=0pt,
    bottom=0pt, top=0pt, left=0pt, right=0pt,
	enhanced standard,
	fuzzy shadow={2mm}{-2mm}{0mm}{0.4mm}{black!50!white},
	boxsep=0mm
]

% --- Εικόνα χωρίς καθόλου περιθώρια ---
\begin{tikzpicture}\node[
    inner sep=0pt,
    outer sep=0pt,
    rounded corners=2mm,
    clip
] {\includegraphics[width=\linewidth,height=4cm]{#5}};\end{tikzpicture}

% --- Κείμενο ---
\begin{tcolorbox}[
    colback=white,
    boxrule=0pt,
    top=2mm,bottom=0mm,left=3mm,right=3mm,
    sharp corners=all
]
    \begin{center}
    {\bfseries\LARGE \titlefont{#1}} \\[4pt]
	{\color{secondarycolor!80!black} \titlefont{#2}}
	\end{center}
    {\small #3}
\end{tcolorbox}

% --- Kάτω μπάρα ---
\begin{tcolorbox}[
    colback=maincolor!90!black,
    colupper=white,
    boxrule=0pt,
    bottom=3pt,top=3pt,left=2mm,right=2mm,
    rounded corners=south,
	sharp corners=north,
    arc=2mm
]
\titlefont{#4}
\end{tcolorbox}
\end{tcolorbox}}

% ===================================================================
%  10. ΛIΣTEΣ (enumitem)
% ===================================================================
\usepackage{enumitem,moreenum}
\newlist{rlist}{enumerate}{3} \setlist[rlist]{itemsep=0mm,label=\roman*.}
\newlist{alist}{enumerate}{3} \setlist[alist]{itemsep=0mm,label=\alph*.}
\newlist{balist}{enumerate}{3} \setlist[balist]{itemsep=0mm,label=\bf\alph*.}
\newlist{Alist}{enumerate}{3} \setlist[Alist]{itemsep=0mm,label=\Alph*.}
\newlist{bAlist}{enumerate}{3} \setlist[bAlist]{itemsep=0mm,label=\bf\Alph*.}
\newlist{askhseis}{enumerate}{3} \setlist[askhseis]{label={\Large\thesection}.\arabic*.}
\renewcommand{\textstigma}{\textsigma\texttau}
%\renewcommand{\textdexiakeraia}{}
\AtBeginDocument{\renewcommand{\textdexiakeraia}{}}
\setlist[itemize]{itemsep=0mm}

% --- Στυλ άσκησης ---
\newcounter{askhsh}[chapter]
\renewcommand{\theaskhsh}{
	\bf{\titlefont{\Large{\thechapter}}.\arabic{askhsh}}
}
\newcommand{\Askhsh}[2]{\refstepcounter{askhsh}{\vspace{4mm}\titlefont{{\large Άσκηση}\theaskhsh} - \textcolor{maincolor}{ΘEMA #2}} ({\small{\textbf{Κωδικός:} #1}})}{}

% --- Περιβάλλον Ύλης ---
\newenvironment{ylh}[1]{\raggedright
\begin{tcolorbox}[
breakable,
halign=flush left,
enhanced standard,
titlerule=0pt,
toptitle=1mm,
toprule=0pt, 
rightrule=0pt, 
bottomrule=0pt,
title={\titlefont{Ύλη άσκησης}},
colback=white,
opacityback=0,
left=1mm,
top=0mm,
bottom=2mm,
boxrule=0pt,
colframe=white,
%borderline west={0mm}{0pt}{maincolor},
borderline north={0.2mm}{0pt}{maincolor},
borderline south={0.2mm}{0pt}{maincolor},
leftrule=2mm,
sharp corners,
coltitle=black
]
\begin{multicols}{#1}
\begin{itemize}[leftmargin=5mm, itemsep=0mm, label=\textcolor{maincolor}{{\small\faBook}}]
}{
\end{itemize}
\end{multicols}
\end{tcolorbox}
}


% ===================================================================
%  11. BΙBΛΙOΓΡAΦΙA (biblatex)
% ===================================================================
\usepackage{csquotes} % Για σωστή διαχείριση εισαγωγικών, συνιστάται από το biblatex
\usepackage[backend=biber,style=alphabetic,sorting=ynt]{biblatex}

% ===================================================================
%  12. ΠΙNAKEΣ (tabularray, longtable, caption)
% ===================================================================
\usepackage{tabularx,longtable,tabularray,capt-of,caption,multirow}
\DeclareTblrTemplate{caption}{nocaptemplate}{}
\DeclareTblrTemplate{capcont}{nocaptemplate}{}
\DeclareTblrTemplate{contfoot}{nocaptemplate}{}
\NewTblrTheme{mytabletheme}{
\SetTblrTemplate{caption}{nocaptemplate}{}
\SetTblrTemplate{capcont}{nocaptemplate}{}
\SetTblrTemplate{contfoot}{nocaptemplate}{}}

\NewTblrEnviron{myhtblr}
\SetTblrOuter[myhtblr]{theme=mytabletheme}
\SetTblrInner[myhtblr]{
    hlines = {1pt,white},
	vlines = {0.5pt,white},
	column{odd} = {maincolor!40,c},
	column{even} = {maincolor!10,c},
	column{1} = {maincolor,c,fg={white}}
}

\NewTblrEnviron{mytblr}
\SetTblrStyle{firsthead}{font=\bfseries}
\SetTblrStyle{firstfoot}{fg=red2}
\SetTblrOuter[mytblr]{theme=mytabletheme}
\SetTblrInner[mytblr]{
    rowspec={t{7mm}},
    columns = {c},
    width = 0.85\linewidth,
    row{odd} = {bg=maincolor!20!white,fg=black,ht=8mm},
    row{even} = {bg=gray!10!white,fg=black,ht=8mm},
    hlines={white},
    vlines={white},
    row{1} = {bg=maincolor, fg=white, font=\large\bfseries\titlefont},
    rowhead = 1,
    hline{2} = {.7mm},
}
\captionsetup[figure]{format=hang,labelsep=period,name={\titlefont{\textbf{Σχήμα}}}}
\renewcommand\thefigure{{\bf\titlefont{\thechapter.\arabic{figure}}}}

% ===================================================================
%  13. ΠEΡΙEΧOMENA (etoc & titletoc)
% ===================================================================
% 2. ΝΕΟ: Μορφοποίηση για άτιτλα κεφάλαια (Περιεχόμενα, Βιβλιογραφία)
\titleformat{name=\chapter,numberless}[block]
  {\normalfont\sffamily}
  {}
  {0pt}
  {%
    \begin{tikzpicture}[remember picture, overlay]
      \fill[maincolor] (current page.north west) rectangle ($(current page.north east) + (0,-4.5cm)$);%
      \fill[maincolor!50!black] (current page.north west) rectangle ($(current page.north west) + (2cm,-4.5cm)$);%
      \fill[maincolor!50] ($(current page.north west) + (0,-4.5cm)$) rectangle ($(current page.north east) + (0,-4.65cm)$);%
      \node[anchor=north west, text=white, font=\titlefont\Huge, text width=\paperwidth-5cm] 
        at ($(current page.north west) + (2.4cm,-2.5cm)$) {#1};\end{tikzpicture}%
  }

\titlespacing*{\chapter}{0pt}{-1cm}{3.5cm}
\titlespacing*{name=\chapter,numberless}{0pt}{-1cm}{3.5cm}

\titleformat{\section}
  {\normalfont\titlefont\Large\color{maincolor}}{\thesection}{1em}{}
\titlecontents{chapter}
  [0em] % Αριστερό περιθώριο
  {\addvspace{1.5ex}\titlefont\large\bfseries\color{maincolor}} % Κενό και γραμματοσειρά
  {\color{maincolor!50}\thecontentslabel\hspace*{1em}\color{maincolor}} % Αριθμός κεφαλαίου σε cyan
  {} % Για τα άτιτλα (δεν χρειάζεται κάτι)
  {\hfill\color{maincolor!50}\bfseries\contentspage} % Αριθμός σελίδας δεξιά

% Εμφάνιση των ενοτήτων στον πίνακα περιεχομένων
\titlecontents{section}
  [2.5em] % Εσοχή
  {\titlefont\color{maincolor!50!black}}
  {\thecontentslabel\hspace*{1em}}
  {}
  {\titlerule*[0.5pc]{\color{gray}.}\color{maincolor!50!black}\contentspage} % Τελίτσες μέχρι τον αριθμό σελίδας

\usepackage{tikzpagenodes,eso-pic}
\newcommand{\chaptertoc}{%
\AddToShipoutPictureFG*{%
\begin{tikzpicture}[overlay]
  \coordinate (NE) at ($(current page text area.north east)+(-59mm,-14.5cm)$);
  \coordinate (SE) at ($(current page text area.south east)+(-59mm,0)$);
  \draw[magenta!80!black] ([xshift=0.5\marginparsep]NE) -- ([xshift=0.5\marginparsep]SE);\end{tikzpicture}%
}%
\marginpar{%
\vspace*{8mm}
\begin{minipage}{\dimexpr\marginparwidth+4mm}
\begin{flushleft}
\localtableofcontents
\end{flushleft}
\end{minipage}%
}%
}

% ===================================================================
%  14. ΣTYΛ KEΦAΛAΙOY & ENOTHTΩN
% ===================================================================

\newcommand*\chapterlabel{}
\newcommand{\fancychapter}{%
\titleformat{\chapter}
{\normalfont\huge }
{\gdef\chapterlabel{\thechapter\ }}{-2cm}
{    \begin{tikzpicture}[remember picture, overlay]
      \fill[maincolor] (current page.north west) rectangle ($(current page.north east) + (0,-4.5cm)$);%
      \fill[maincolor!50!black] (current page.north west) rectangle ($(current page.north west) + (2cm,-4.5cm)$);%
      \fill[maincolor!50] ($(current page.north west) + (0,-4.5cm)$) rectangle ($(current page.north east) + (0,-4.65cm)$);%
      \node[anchor=south west, text=maincolor!50, font=\titlefont\Large] 
        at ($(current page.north west) + (2.5cm,-2.4cm)$) {ΚΕΦΑΛΑΙΟ \thechapter};%
      \node[anchor=north west, text=white, font=\titlefont\Huge, text width=\paperwidth-5cm] 
        at ($(current page.north west) + (2.4cm,-2.5cm)$) {##1};\end{tikzpicture}%
}
}

\titleformat{\section}[hang]
  {\normalfont\titlefont\Large\color{maincolor}}
  {\textcolor{maincolor!50}{\rule[-0.15em]{3pt}{1.1em}}\hspace{0.5em}\textcolor{maincolor!50!black}{\thesection}}
  {0.5em}
  {#1}

% --- 4. Υπο-ενότητες (Subsections) ---
% Πιο διακριτικές, με τον αριθμό σε cyan και τον τίτλο σε σκούρο μπλε
\titleformat{\subsection}[hang]
  {\normalfont\titlefont\large\color{maincolor!50!black}}
  {\textcolor{maincolor!50}{\thesubsection}}
  {0.5em}
  {#1}

\apptocmd{\mainmatter}{\fancychapter}{}{}
\titlespacing{\chapter}{0pt}{0cm}{1cm}

\titleformat{\paragraph}{\large}{}{0em}{\textcolor{maincolor}{\faSquare\ \ \titlefont{\bmath{#1}}}}
\titlespacing{\paragraph}{0mm}{2mm}{2mm}

% ===================================================================
%  15. BOHΘHTΙKEΣ ENTOΛEΣ
% ===================================================================
\newcommand{\eng}[1]{\selectlanguage{english}#1\selectlanguage{greek}}
\newcommand{\tss}[1]{\textsuperscript{#1}}
\newcommand{\tssL}[1]{\MakeLowercase{\textsuperscript{#1}}}
\newcommand{\full}[1]{\begin{mdframed}[outermargin=\dimexpr-\marginparwidth-\marginparsep\relax,innerleftmargin=0mm,innerrightmargin=0mm,hidealllines=true] #1 \end{mdframed}}
\newcommand{\fulltwoc}[1]{\begin{adjmulticols}{2}{5mm}{\dimexpr-\marginparwidth-\marginparsep+5mm\relax} #1 \end{adjmulticols}}

\newcommand{\monades}[1]{
\hspace*{\fill}
\textbf{\textit{\textcolor{maincolor}{Μονάδες #1}}}}


% ===================================================================
%  TEΛOΣ PREAMBLE
% ===================================================================



\begin{document}
\pagestyle{empty}
\frontmatter
\pagenumbering{gobble}
\begin{tikzpicture}[remember picture, overlay]
    % --- 1. TOP SECTION: MATH ART BACKGROUND (Placeholder for Image) ---
    \begin{scope}
        \clip (current page.north west) rectangle ([yshift=-15cm]current page.east);
        \node[anchor=north west, inner sep=0pt] at (current page.north west) {
            \includegraphics[width=\paperwidth, height=17cm, keepaspectratio=false]{/home/spyros/Εικόνες/wallpapers/sebastian-boring-8zD7rs8UpxU-unsplash.jpg}};
    \end{scope}
        % --- 2. MIDDLE GEOMETRY: ORANGE ACCENTS ---
        % Orange Main Polygon
    \path[right color=maincolor!55!black, left color=maincolor] 
        ([xshift=10cm, yshift=-14cm]current page.north west) -- 
        ([xshift=22cm, yshift=-14cm]current page.north west) -- 
        ([xshift=22cm, yshift=-17cm]current page.north west) -- 
        ([xshift=13.7cm, yshift=-17cm]current page.north west) -- cycle;
    \node[anchor=south, inner sep=0pt] at (current page.south){
    \includesvg[width=\paperwidth, height=\paperheight]{cover.svg}
    };
    % --- 5. LOGO & YEAR ---
    % Fake Logo
    \node[anchor=north west] at ([xshift=1cm, yshift=-14cm]current page.north west) {
	\begin{tikzpicture}[scale=0.7]
        \fill[maincolor,yshift=-2mm](0,0) rectangle (2cm,2.4cm);
		\node (fig1) at ($(0,0)+(0.21cm,1.9cm)$)
		{
\includegraphics[width=0.049\linewidth]{/home/spyros/Μαθηματικά/Φροντιστήριο ΦΙΛΟΜΑΘΕΙΑ/Βαθμοί - Φυλλάδιο ύλης - Διάφορα/Λογότυπα/logo.png}};
		\node[white] at ($(0,0)+(2.4cm,2.3cm)$) {\titlefont{\fontsize{10}{12}\selectfont ΦΡΟΝΤΙΣΤΗΡΙΟ ΜΕΣΗΣ ΕΚΠΑΙΔΕΥΣΗΣ}};
		\node[white] at ($(0,0)+(2.3cm,1.7cm)$) {\titlefont{ {\fontsize{26}{33.6}\selectfont \textbf{ΦΙΛΟΜΑΘΕΙΑ}}}};
		\draw[white] (2.5,0.45)--($(2.4,0.45)+(7.7cm,0cm)$);
		\node[white] at ($(0,0)+(2.5cm,0.5cm)$) {\textcolor{white}{\fontsize{10}{12}\selectfont\faMapMarker*\ \ Ιακώβου Πολυλά 24, Πεζόδρομος}};\end{tikzpicture}
    };

    % Year
    \node[anchor=north west, text=white, align=left] at ([xshift=1cm, yshift=-17cm]current page.north west) {
        \fontsize{28}{40}\selectfont \titlefont{\textbf{Τράπεζα Θεμάτων}}\\[1mm]  \fontsize{32}{40}\selectfont \titlefont{\textcolor{maincolor}{Α' Λυκείου}}
    };

    % --- 6. MAIN TITLE ---
    \node[anchor=north west, text=white, align=left] at ([xshift=1cm, yshift=-7cm]current page.west) {
        \fontsize{14}{16}\selectfont \titlefont{ΘΕΜΑΤΑ -  \textcolor{secondarycolor}{ΛΥΣΕΙΣ ΘΕΜΑΤΩΝ}}\\[0.8em]
        \fontsize{38}{40}\selectfont \textbf{\textcolor{maincolor}{ΑΛΓΕΒΡΑ}}\\[0.4em]
        \fontsize{14}{16}\selectfont {\titlefont{\textcolor{white}{\textbf{ΣΧΟΛΙΚΟ ΕΤΟΣ 2025 - 2026}}}}
    };

    % --- 8. FOOTER / CONTACT INFO ---
    % Vertical Yellow Bar
    \path[bottom color=maincolor!80!black, top color=secondarycolor!80] ([xshift=15cm, yshift=0cm]current page.south west) rectangle ([xshift=15.5cm, yshift=3.7cm]current page.south west);
    
    % Contact Text
    \node[anchor=south west, text=white, align=left] at ([xshift=14.75cm, yshift=.5cm]current page.south west) {
        \fontsize{9}{12}\selectfont
		\begin{tabular}{cl}
		 & \textbf{ΕΠΙΚΟΙΝΩΝΙΑ} \\[1mm]
		\faPhone* & 26610 20144\\[1mm]
		\faEnvelope & {frontistirio.filomatheia@gmail.com} \\[1mm]
		\faMobile* & 693 232 7283\ \ \ \faTelegram\ \faViber\\[1mm]
		\faFacebook & Φροντιστήριο Φιλομάθεια\\[1mm]
		\faInstagram & {filomatheia\_kerkyra}
		\end{tabular}        
    };
\end{tikzpicture}
\newpage
\newgeometry{inner=2.00cm, top=3.00cm, bottom=2.00cm,outer=1.50cm,headheight=14pt}
\tableofcontents
\restoregeometry % restores the geometry
\mainmatter
\pagestyle{fancy}
%\setchapterimage{<image path>}
\chapter{Tράπεζα Θεμάτων}




\restoregeometry
% ===================================================================
% ΑΣΚΗΣΕΙΣ ΤΡΑΠΕΖΑΣ ΘΕΜΑΤΩΝ – Αυτόματη εισαγωγή (275 ασκήσεις)
% ===================================================================

% --- Άσκηση 14586 (14586.tex) ---
\Askhsh{0000}{2}
Δίνονται τα σημεία $A(1,2)$, $B(3,4)$ και $\Gamma(5,-2)$.

\begin{alist}    \item Να βρείτε τις συντεταγμένες των διανυσμάτων $\vec{AB}$, $\vec{A\Gamma}$ και να αποδείξετε ότι η γωνία $\widehat{A}$ είναι ορθή. \monades{9}
    \item Αν $M$ είναι το μέσο του $B\Gamma$, να βρείτε τα μέτρα των $\vec{AM}$ και $\vec{B\Gamma}$. \monades{8}
    \item Να γραφεί το $\vec{B\Gamma}$ ως γραμμικός συνδυασμός των $\vec{A\Gamma}$ και $\vec{AM}$. \monades{8}\end{alist}

\vspace{4mm}

% --- Άσκηση 14666 (14666.tex) ---
\Askhsh{}{2}
Θεωρούμε τα διανύσματα $\vec{a} = (1,-3)$ και $\vec{b} = (-2,-1)$ και ορίζουμε τα διανύσματα
\[ \vec{u} = 3\vec{a} - 5\vec{b} \quad \text{και} \quad \vec{v} = 5\vec{a} - 9\vec{b}. \]

\begin{alist}    \item Να βρείτε τις συντεταγμένες των διανυσμάτων $\vec{u}$ και $\vec{v}$. \monades{9}
    \item Αν $\vec{w} = 2\vec{u} - \vec{v}$, να γράψετε το $\vec{w}$ ως γραμμικό συνδυασμό των $\vec{a}, \vec{b}$. \monades{9}
    \item Αν τα $\vec{b}, \vec{v}, \vec{u}$ είναι τα διανύσματα θέσης των σημείων $K, \Lambda, \text{και } M$ αντίστοιχα, να αποδείξετε ότι τα σημεία αυτά είναι συνευθειακά. \monades{7}\end{alist}

\vspace{4mm}

% --- Άσκηση 14953 (14953.tex) ---
\Askhsh{}{2}

Θεωρούμε τρίγωνο $AB\Gamma$ με $A(-2,5)$, $B(7,8)$, $\Gamma(1,-4)$.

\begin{alist}    \item Να βρείτε τις συντεταγμένες των διανυσμάτων $\vec{AB}$ και $\vec{A\Gamma}$. \monades{10}
    \item Να υπολογίσετε το εσωτερικό γινόμενο $\vec{AB} \cdot \vec{A\Gamma}$. \monades{10}
    \item Να βρείτε, σε μοίρες, τη γωνία $\widehat{BA\Gamma}$. \monades{5}\end{alist}

\vspace{4mm}

% --- Άσκηση 14954 (14954.tex) ---
\Askhsh{M1234}{4}
Θεωρούμε τις εξισώσεις \((\varepsilon1)\): \(\mu x - y - \mu = 0\) και \((\varepsilon2)\): \((\mu + 1)x + (\mu - 1) y - \mu + 1 = 0\), με \(\mu \in \mathbb{R}\).

\begin{alist}    \item Να αποδείξετε ότι οι \((\varepsilon1)\) και \((\varepsilon2)\) παριστάνουν εξισώσεις ευθειών για κάθε τιμή της παραμέτρου \(\mu\).
    \item Να αποδείξτε ότι η οξεία γωνία των ευθειών \((\varepsilon1)\) και \((\varepsilon2)\) είναι \(45^\circ\) για κάθε τιμή της παραμέτρου \(\mu\).  \monades{6}
    \item Να αποδείξετε ότι τα σημεία τομής των ευθειών \((\varepsilon1)\) και \((\varepsilon2)\) ανήκουν στον κύκλο με κέντρο την αρχή των αξόνων και ακτίνα 1.  \monades{10}
 \monades{9}\end{alist}

\vspace{4mm}

% --- Άσκηση 14970 (14970.tex) ---
\Askhsh{}{4}
Δίνεται η ευθεία $y = \lambda(x-2) + \lambda - 2$, $\lambda \in \mathbb{R}$ \textnormal{(1)}.

\begin{alist}    \item Να βρείτε τις ευθείες που προκύπτουν όταν $\lambda=1$ και όταν $\lambda=2$. Κατόπιν να βρείτε το κοινό σημείο $M$ των δυο ευθειών.
    \item Να αποδείξετε ότι όλες οι ευθείες που προκύπτουν από την \textnormal{(1)} για τις διάφορες τιμές του $\lambda$, διέρχονται από το $M$.  \monades{7} Έστω $M(1, -2)$.
    \item Να βρείτε:  \monades{5}
    \begin{rlist}
        \item τα σημεία τομής $A$, $B$ της ευθείας \textnormal{(1)} με τους άξονες $x'x$ και $y'y$.

        \item για ποιες τιμές του $\lambda$ το εμβαδόν του τριγώνου $OAB$ είναι ίσο με $\frac{1}{2}$.  \monades{6}
 \monades{7}
    \end{rlist}\end{alist}

\vspace{4mm}

% --- Άσκηση 14978 (14978.tex) ---
\Askhsh{}{4}
Δίνονται τα σημεία $A(1,1)$, $B(3,3)$.

\begin{alist}    \item Αν $M(x,y)$ σημείο του επιπέδου, να βρείτε τις αποστάσεις $d_1, d_2$ του $M$ από τα $A$ και $B$ αντίστοιχα.
    \item Να γράψετε τη σχέση που πρέπει να πληρούν οι $d_1, d_2$, ώστε το σημείο $M$ να ανήκει στη μεσοκάθετο του $AB$.  \monades{6}
    \item Να βρείτε την εξίσωση της μεσοκαθέτου του $AB$.  \monades{4}
    \item Να βρείτε σημείο $\Sigma$ τέτοιο ώστε το τρίγωνο $\Sigma AB$ να είναι ισόπλευρο.  \monades{8}
 \monades{7}\end{alist}

\vspace{4mm}

% --- Άσκηση 14984 (14984.tex) ---
\Askhsh{XXXX}{4}
Θεωρούμε τα σημεία $A(-2, -3)$ και $B(7, 9)$. Έστω $S$ το σύνολο των σημείων $M$ που είναι κορυφές των τριγώνων $AMB$ ώστε $(AMB) = 12$ τ.μ.

\begin{alist}    \item Να αποδείξτε ότι το $S$ αποτελείται από τα σημεία των παραλλήλων ευθειών
    \item Να αποδείξετε ότι η ευθεία $AB$ είναι η μεσοπαράλληλη των $(\varepsilon1)$ και $(\varepsilon2)$.  \[ (\varepsilon1): 4x - 3y - 9 = 0 \quad \text{και} \quad (\varepsilon2): 4x - 3y + 7 = 0. \] \monades{9}
    \item Θεωρούμε ένα σημείο $M_1$ στην $(\varepsilon1)$ και ένα σημείο $M_2$ στην $(\varepsilon2)$ ώστε να σχηματίζεται το τετράπλευρο $AM_1 BM_2$. Πόσο είναι το εμβαδόν του; Πόσα τετράπλευρα $AXBY$ υπάρχουν, αν το $X$ πρέπει να είναι σημείο της $(\varepsilon1)$ και το $Y$ σημείο της $(\varepsilon2)$, που έχουν το ίδιο εμβαδόν με το $AM_1 BM_2$; Εξηγήστε.  \monades{9}
 \monades{7}\end{alist}

\vspace{4mm}

% --- Άσκηση 15002 (15002.tex) ---
\OnPar{ΘΕΜΑ 2}
Δίνονται τα σημεία $A(0,5)$ και $\Delta(4,5)$ και τα διανύσματα $\vec{AB} = (3, -3)$ και $\vec{A\Gamma} = (3,1)$.

\begin{alist}    \item Να αποδείξετε ότι το σημείο $\Gamma$ έχει συντεταγμένες $\Gamma(3,6)$.
    \item  \monades{11}
    \begin{rlist}
        \item Να βρείτε τις συντεταγμένες του διανύσματος $\vec{\Gamma\Delta}$.

        \item Να αποδείξετε ότι $\vec{AB}//\vec{\Gamma\Delta}$.  \monades{6}
 \monades{8}
    \end{rlist}\end{alist}

\vspace{4mm}

% --- Άσκηση 15004 (15004.tex) ---
\textbf{ΘΕΜΑ 4}

\begin{alist}    \item Να βρείτε την εξίσωση της ευθείας $\varepsilon1$ που διέρχεται από τα σημεία $A(4,2)$ και $B(8,5)$. \monades{5}
    \item Αν $\varepsilon1: 3x - 4y - 4 = 0$, να δείξετε ότι σχηματίζει με την ευθεία $\varepsilon2: 7x - y - 1 = 0$ γωνία $\phi = 45^\circ$. \monades{8}
    \item Να βρείτε το σημείο τομής των $\varepsilon1$ και $\varepsilon2$. \monades{4}
    \item Να βρείτε την εξίσωση ευθείας $\varepsilon3$ τέτοιας ώστε η $\varepsilon2$ να διχοτομεί τη γωνία που σχηματίζουν οι ευθείες $\varepsilon1$ και $\varepsilon3$. \monades{8}
\end{alist}

\vspace{4mm}

% --- Άσκηση 15010 (15010.tex) ---
\Askhsh{}{2}
Δίνονται τα μη συνευθειακά σημεία του επιπέδου $A,B,\Gamma$ και τα διανύσματα $\vec{B\Delta}$ και $\vec{\Gamma E}$ τέτοια ώστε $\vec{B\Delta} = \vec{BA}+\vec{B\Gamma}$ και $\vec{\Gamma E} = \vec{\Gamma A}+\vec{\Gamma B}$.

\begin{alist}    \item
    \begin{rlist}
        \item Να δείξετε ότι $\vec{A\Delta}= \vec{B\Gamma}$ και $\vec{AE} = \vec{\Gamma B}$. \monades{8}
        \item Να δείξετε ότι τα διανύσματα $\vec{A\Delta}$ και $\vec{AE}$ είναι αντίθετα. \monades{8}
    \end{rlist}
    \item Να δικαιολογήσετε γιατί τα σημεία $A, \Delta$ και $E$ είναι συνευθειακά. \monades{9}\end{alist}

\vspace{4mm}

% --- Άσκηση 15027 (15027.tex) ---
\Askhsh{M21}{2}
Δίνονται τα σημεία $A(1,-1)$ και $B(3,5)$ όπως φαίνεται στο παρακάτω σχήμα.

\begin{center}
\begin{tikzpicture}[scale=1]
    % Ορισμός των σημείων σχηματικά, όπως στην εκφώνηση
    \tkzDefPoint(-2,0){Adiag}
    \tkzDefPoint(2,0){Bdiag}
    \tkzDefMidPoint(Adiag,Bdiag) \tkzGetPoint{Mdiag}
    % Σχεδίαση του ευθύγραμμου τμήματος Adiag Bdiag
    \tkzDrawSegment[l3](Adiag,Bdiag)
    % Σχεδίαση της κάθετης ευθείας που διέρχεται από το Mdiag
    \coordinate (Pup) at ($(Mdiag) + (0,2.5)$); % Σημείο πάνω από το Mdiag
    \coordinate (Pdown) at ($(Mdiag) + (0,-2.5)$); % Σημείο κάτω από το Mdiag
    \tkzDrawSegment[l3](Pup,Pdown)
    % Σήμανση της ορθής γωνίας στο Mdiag
    \tkzMarkRightAngle[size=0.3,fill=thirdcolor!20,draw=thirdcolor](Pup,Mdiag,Bdiag)
    % Σχεδίαση και επισήμανση των σημείων
    \tkzDrawPoints[color=maincolor](Adiag,Bdiag) % Τα σημεία A και B είναι μπλε
    \tkzDrawPoint(Mdiag) % Το σημείο M είναι μαύρο
    % Ετικέτες των σημείων
    \tkzLabelPoint[below left](Adiag){A}
    \tkzLabelPoint[below right](Bdiag){B}
    \tkzLabelPoint[below](Mdiag){M}\end{tikzpicture}
\end{center}

\begin{alist}    \item Να βρείτε το συντελεστή διεύθυνσης της ευθείας $AB$. \monades{7}
    \item Να βρείτε τις συντεταγμένες του μέσου $M$ του τμήματος $AB$. \monades{8}
    \item Να βρείτε την εξίσωση της μεσοκαθέτου του τμήματος $AB$. \monades{10}\end{alist}

\vspace{4mm}

% --- Άσκηση 15028 (15028.tex) ---
\Askhsh{}{2}
Έστω κύκλος $C$ με κέντρο $K(1,2)$ και ακτίνα $\rho = 2$ και ευθεία $(\varepsilon)$ με εξίσωση $3x+4y-1=0$.

\begin{alist}    \item Να γράψετε την εξίσωση του κύκλου $C$.
    \item Να δείξετε ότι η απόσταση του κέντρου $K(1,2)$ από την ευθεία $(\varepsilon)$ είναι ίση με $2$.  \monades{8}
    \item Να δείξετε ότι η ευθεία $(\varepsilon)$ εφάπτεται στον κύκλο $C$.  \monades{9}
 \monades{8}\end{alist}

\vspace{4mm}

% --- Άσκηση 15029 (15029.tex) ---
\Askhsh{}{4}
Στο παρακάτω σχήμα δίνονται τα σημεία $O(0,0)$, $A(1,\sqrt{3})$, $B(\sqrt{3}+1,\sqrt{3}-1)$.

\begin{center}
\begin{tikzpicture}[scale=1.5]
    % Axes
    \draw[->] (-1,0) -- (4,0) node[right] {$x'$x};
    \draw[->] (0,-0.5) -- (0,2.5) node[above] {$y'$y};
    % Define points using exact coordinates for later calculations/consistency
    % Use numerical approximations for drawing, but define points for tkz-euclide
    \tkzDefPoint(0,0){Otkz}
    \tkzDefPoint(1,{sqrt(3)}){Atkz}
    \tkzDefPoint({1+sqrt(3)},{sqrt(3)-1}){Btkz}
    % Draw lines OA and AB
    \tkzDrawLine[add=0 and 0.5, l3, maincolor](Otkz,Atkz) % Line OA
    \tkzDrawLine[add=0 and 0.5, l3, secondarycolor](Atkz,Btkz) % Line AB
    % Draw points
    \tkzDrawPoints(Otkz,Atkz,Btkz)
    % Labels for points
    \tkzLabelPoint[below left](Otkz){O}
    \tkzLabelPoint[above left](Atkz){A}
    \tkzLabelPoint[above right](Btkz){B}
\end{tikzpicture}
\end{center}

\begin{alist}    \item Να βρείτε την εξίσωση της ευθείας ΟΑ καθώς και τη γωνία $\omega$ που σχηματίζει με τον άξονα $x'x$.
    \item Να βρείτε την εξίσωση της ευθείας ΑΒ καθώς και τη γωνία $\phi$ που σχηματίζει με τον άξονα $x'x$.  \monades{6}
    \item Να δείξετε ότι το τρίγωνο ΟΑΒ είναι ορθογώνιο και ισοσκελές με $\hat{A} = 90^\circ$.  \monades{6}
    \item Να δείξετε ότι $\text{εφ}15^\circ = \frac{\sqrt{3}-1}{\sqrt{3}+1}$.  \monades{7}
 \monades{6}\end{alist}

\vspace{4mm}

% --- Άσκηση 15030 (15030.tex) ---
\section*{ΘΕΜΑ 4}

Δίνεται ο κύκλος $C: (x-2)^2+(y+3)^2 = 5$ και η ευθεία $\varepsilon: 2x+y+5=0$.

\begin{alist}    \item Να βρείτε το κέντρο και την ακτίνα του κύκλου $C$.
    \item Να δείξετε ότι ο κύκλος $C$ και η ευθεία $(\varepsilon)$ δεν έχουν κοινά σημεία.  \monades{6}
    \item Να δείξετε ότι υπάρχουν δύο ευθείες $(\eta_1),(\eta_2)$ που είναι παράλληλες στην ευθεία $(\varepsilon)$ και εφάπτονται του κύκλου $C$ και να βρείτε τις εξισώσεις τους.  \monades{6}
    \item Να βρείτε τη μεσοπαράλληλη των ευθειών $(\eta_1),(\eta_2)$.  \monades{7}
 \monades{6}\end{alist}

\vspace{4mm}

% --- Άσκηση 15038 (15038.tex) ---
\Askhseis{ΘΕΜΑ 2}

Θεωρούμε διανύσματα $\vec{a}$, $\vec{b}$ τέτοια ώστε $|\vec{a}| = 3$, $|\vec{b}| = 4$ και $\left(\widehat{\vec{a}, \vec{b}}\right) = \frac{\pi}{3}$.

\begin{alist}    \item Να βρείτε το εσωτερικό γινόμενο των διανυσμάτων $\vec{a}$, $\vec{b}$. \par\monades{9}
    \item Να βρείτε τα $\vec{a}^2$ και $\vec{b}^2$. \par\monades{6}
    \item Να αποδείξετε ότι $(3\vec{a} - \vec{b}) \cdot (\vec{a} - 3\vec{b}) = 15$. \par\monades{10}\end{alist}

\vspace{4mm}

% --- Άσκηση 15042 (15042.tex) ---
\noindent\textbf{ΘΕΜΑ 4}

Δίνεται τρίγωνο ΑΒΓ και σημείο του επιπέδου Μ, τέτοιο ώστε:
\[ \vec{AB} - 2\vec{AM} + \vec{A\Gamma} = \vec{0} \]

\begin{alist}    \item Να αποδείξετε ότι τα σημεία Β, Γ, Μ είναι συνευθειακά. \monades{8}
    \item Να αποδείξετε ότι το Μ είναι το μέσο του ΒΓ. \monades{2}
    \item Έστω πραγματικοί αριθμοί $\kappa, \lambda$ τέτοιοι ώστε $\vec{AB} \cdot \vec{A\Gamma} = \kappa$ και $\vec{AM} \cdot \vec{B\Gamma} = \lambda$.
 Αν επιπλέον είναι γνωστό ότι για τα μη παράλληλα διανύσματα $\vec{A\Gamma}$, $\vec{AB}$ ισχύει ότι $\kappa \vec{A\Gamma} = \lambda \vec{AB}$, τότε:
    \begin{rlist}
        \item Να αποδείξετε ότι $\kappa = \lambda = 0$. \monades{7}
        \item Να αποδείξετε ότι το τρίγωνο ΑΒΓ είναι ορθογώνιο και ισοσκελές. Να προσδιορίσετε την ορθή γωνία και τις πλευρές που είναι ίσες. \monades{8}
    \end{rlist}\end{alist}

\vspace{4mm}

% --- Άσκηση 15043 (15043.tex) ---
\OnPar{ΘΕΜΑ 2}

Ένα γραφείο μελετών έχει αναλάβει την αναμόρφωση μιας οικιστικής περιοχής, η οποία
αποτυπώνεται σε τοπογραφικό σχέδιο με ένα καρτεσιανό σύστημα συντεταγμένων. Τα
σημεία $A(0,5)$, $B(4,1)$ και $\Gamma(6, -1)$ παριστάνουν τη θέση τριών οικισμών στο χάρτη.

\begin{center}
\begin{tikzpicture}\begin{axis}[
    axis lines = middle,
    xlabel = {$x \text{ (km)}$},
    ylabel = {$y \text{ (km)}$},
    xmin = -2, xmax = 9,
    ymin = -2, ymax = 8,
    xtick = {-2,-1,0,1,2,3,4,5,6,7,8,9},
    ytick = {-2,-1,0,1,2,3,4,5,6,7,8},
    grid = both,
    grid style = {gray!30},
    width = 10cm, height = 7cm,
    tick label style = {font=\scriptsize},
    every axis x label/.style = {at={(current axis.right of origin)}, anchor=north west},
    every axis y label/.style = {at={(current axis.above origin)}, anchor=south},
    clip=false 
]
    % Plot points A, B, Gamma
    \addplot[only marks, mark=*, mark size=2pt, maincolor] coordinates {(0,5) (4,1) (6,-1)};
    % Label points
    \node[maincolor, above right] at (axis cs:0,5) {$A$};
    \node[maincolor, below right] at (axis cs:4,1) {$B$};
    \node[maincolor, below right] at (axis cs:6,-1) {$\Gamma$};
\end{axis}\end{tikzpicture}
\end{center}

\begin{alist}    \item % Part α)
    \begin{rlist}
        \item Να βρείτε τις συντεταγμένες των διανυσμάτων $\vec{AB}$ και $\vec{B\Gamma}$. \monades{06}
        \item Να αποδείξετε ότι τα σημεία $A$, $B$, $\Gamma$ είναι συνευθειακά και ως εκ τούτου υπάρχει η δυνατότητα να σχεδιασθεί ένας ευθύγραμμος δρόμος που να συνδέει τους τρεις οικισμούς. \monades{07}
    \end{rlist}
    \item % Part β)
 Να αποδείξετε ότι η απόσταση του οικισμού $B$ από τον οικισμό $A$ είναι διπλάσια από την απόσταση του οικισμού $B$ από τον οικισμό $\Gamma$. \monades{12}\end{alist}

\vspace{4mm}

% --- Άσκηση 15044 (15044.tex) ---
\Askhsh{}{\textbf{2}}
Δίνονται τα σημεία $A(0,5)$ και $B(6, -1)$.

\begin{alist}   
\item 
\begin{rlist}
        \item Να βρείτε τον συντελεστή διεύθυνσης της ευθείας που διέρχεται από τα σημεία $A$ και $B$.
        \item Να αποδείξετε ότι το μέσο του ευθύγραμμου τμήματος $AB$, είναι το σημείο $M(3,2)$.  \monades{5}
 \monades{5}
    \end{rlist}
    \item Να βρείτε την εξίσωση της μεσοκάθετης ευθείας $(\varepsilon)$ του ευθύγραμμου τμήματος $AB$.
 \monades{15}\end{alist}

\vspace{4mm}

% --- Άσκηση 15073 (15073.tex) ---
\Askhsh{}{2}
Δίνονται τα διανύσματα $\vec{a} = (1,2)$ και $\vec{b} = (2,3)$.

\begin{alist}    \item Να βρείτε τις συντεταγμένες του διανύσματος $\vec{c} = 2\vec{a} + \vec{b}$. \monades{8}
    \item Να βρείτε το μέτρο του διανύσματος $\vec{c}$. \monades{8}
    \item Να βρείτε το εσωτερικό γινόμενο $\vec{a} \cdot \vec{c}$. \monades{9}\end{alist}

\vspace{4mm}

% --- Άσκηση 15080 (15080.tex) ---
\Askhsh{}{4}
Δίνονται οι εξισώσεις
\[
C_1: x^2 + y^2 - 2x - 8 = 0 \quad (1)
\]
και
\[
C_2: x^2 + y^2 - 6x + 8 = 0 \quad (2).
\]

\begin{alist}    \item Να δείξετε ότι οι (1) και (2) είναι εξισώσεις κύκλων, με κέντρα $K(1,0)$, $\Lambda(3,0)$ και ακτίνες $\rho_1 =3$, $\rho_2 =1$ αντίστοιχα. \monades{6}
    \item
    \begin{rlist}
        \item Να βρείτε το μήκος της διακέντρου $(K\Lambda)$. \monades{5}
        \item Να δείξετε ότι ο κύκλος $C_2$ εφάπτεται εσωτερικά του κύκλου $C_1$. \monades{5}
    \end{rlist}
    \item Να βρείτε τις εξισώσεις των ακτίνων του κύκλου $C_1$ που εφάπτονται στον κύκλο $C_2$. \monades{9}\end{alist}

\vspace{4mm}

% --- Άσκηση 15081 (15081.tex) ---
\Askhsh{A4MATHGEO01}{4}

Δίνονται οι κύκλοι
\[
C_1 \colon x^2 + y^2 + 2\sqrt{2}x + 1 = 0
\]

και
\[
C_2 \colon x^2 + y^2 - 6\sqrt{2}x + 9 = 0.
\]

\begin{alist}    \item Να δείξετε ότι οι κύκλοι $C_1$ και $C_2$ έχουν κέντρα $K(-\sqrt{2},0)$, $\Lambda(3\sqrt{2},0)$ και ακτίνες $\rho_1 = 1$, $\rho_2 = 3$ αντίστοιχα.
    \item  \monades{8}
    \begin{rlist}
        \item Να δείξετε ότι από την αρχή των αξόνων διέρχονται δύο κοινές εφαπτόμενες των κύκλων $C_1$ και $C_2$.

        \item Να σχεδιάσετε ένα πρόχειρο σχήμα όπου να φαίνονται οι κύκλοι και οι δύο αυτές εφαπτόμενες.  \monades{10}
        \begin{center}
        \begin{tikzpicture}            
\begin{axis}[
 axis lines=middle, xlabel=$x$, ylabel=$y$, xmin=-3.5, xmax=7.5, ymin=-4, ymax=4, xtick=\empty, ytick=\empty, grid=none, width=10cm, height=7cm, axis equal, every axis x label/.style={at={(current axis.right of origin)}, anchor=north west}, every axis y label/.style={at={(current axis.above origin)}, anchor=south}, declare function={ sqrt2 = sqrt(2); Kx = -sqrt2; Lx = 3*sqrt2; rho1 = 1; rho2 = 3; } ]
                % Draw circle C1
 \draw[thick, maincolor] (axis cs:Kx, 0) circle (rho1); \node[maincolor, anchor=south west, inner sep=1pt] at (axis cs:Kx+rho1*0.7, rho1*0.7) {$C1$};
                % Draw circle C2
 \draw[thick, secondarycolor] (axis cs:Lx, 0) circle (rho2); \node[secondarycolor, anchor=south east, inner sep=1pt] at (axis cs:Lx-rho2*0.7, rho2*0.7) {$C2$};
                % Draw centers K and Lambda
 \fill[black] (axis cs:Kx, 0) circle (1.5pt); % Use fill for points \fill[black] (axis cs:Lx, 0) circle (1.5pt); \node[below, inner sep=1pt] at (axis cs:Kx, 0) {$K(-\sqrt{2},0)$}; \node[below, inner sep=1pt] at (axis cs:Lx, 0) {$\Lambda(3\sqrt{2},0)$};
                % Draw Origin
 \fill[black] (axis cs:0, 0) circle (1.5pt); \node[below left, inner sep=1pt] at (axis cs:0, 0) {$O$};
                % Draw tangents y=x and y=-x
 \addplot[very thick, thirdcolor, domain=-2.5:6.5] {x}; \addplot[very thick, thirdcolor, domain=-2.5:6.5] {-x};
                % Position labels for tangents
 \node[thirdcolor, below right, inner sep=1pt] at (axis cs: 6, 6) {$y=x$}; \node[thirdcolor, above right, inner sep=1pt] at (axis cs: 6, -6) {$y=-x$};
            \end{axis}
\end{tikzpicture}
        \end{center}
 \monades{7}
    \end{rlist}\end{alist}

\vspace{4mm}

% --- Άσκηση 15082 (15082.tex) ---
\OnPar{ΘΕΜΑ 4}
Δίνονται δύο κύκλοι με εξισώσεις:
\[
C_1 :(x-2)^2 + (y-3)^2 = 8 \quad \text{και} \quad C_2:(x-7)^2 + (y+2)^2 = 18.
\]
\begin{alist}    \item Να υπολογίσετε το μήκος της διακέντρου $(K\Lambda)$, όπου $K, \Lambda$ τα κέντρα των κύκλων $C_1, C_2$ αντίστοιχα. Ακολούθως να δείξετε ότι οι δύο κύκλοι εφάπτονται εξωτερικά. \monades{5}
    \item
    \begin{rlist}
        \item Να βρείτε την εξίσωση της ευθείας $K\Lambda$. \monades{5}
        \item Να βρείτε τα σημεία τομής της ευθείας $K\Lambda$ με τον κύκλο $C_1$ και το σημείο επαφής των δύο κύκλων. \monades{7}
    \end{rlist}
    \item Να βρείτε την εξίσωση της κοινής εσωτερικής εφαπτομένης των κύκλων. \monades{8}\end{alist}

\vspace{4mm}

% --- Άσκηση 15152 (15152.tex) ---
\Askhsh{}{3}
Δίνονται τα σημεία $A(1,3)$, $B(-2,2)$ και η ευθεία $\varepsilon: 3x + y + \alpha = 0$ με $\alpha \in \mathbb{R}$.
\begin{alist}    \item Να βρεθεί η απόσταση του σημείου $A$ από το σημείο $B$. \monades{5}
    \item Για ποιες τιμές του $\alpha$, η απόσταση $AB$ είναι ίση με την απόσταση του σημείου $A$ από την ευθεία $\varepsilon$. \monades{8}
    \item Για $\alpha = 4$ να βρεθεί το εμβαδόν του τριγώνου $AB\Gamma$, όπου $\Gamma$ το σημείο τομής της ευθείας $\varepsilon$ με τον άξονα $y'y$. \monades{12}\end{alist}

\vspace{4mm}

% --- Άσκηση 15177 (15177.tex) ---
\section*{\centering ΘΕΜΑ 4}
Δίνονται τα σημεία $A(1,0)$ και $B(0, -1)$ και ο κύκλος $C_1$ με εξίσωση
\[ C_1: \left(x-\frac{1}{2}\right)^2 + \left(y+\frac{1}{2}\right)^2 = 2. \]

\begin{alist}    \item Να αποδείξετε ότι το σύνολο των σημείων $N(x, y)$ του επιπέδου, τα οποία ικανοποιούν
    \item Να αποδείξετε ότι το σύνολο των σημείων $P$ του επιπέδου, τα οποία ικανοποιούν την  τη σχέση $NA^2 - NB^2 = 4$, ανήκουν στην ευθεία $(\varepsilon)$ με εξίσωση $y = -x - 2$. \monades{07}
    \item  εξίσωση $2x^2 + 2y^2 + 10x + 14y + 21 = 0$, ανήκουν σε κύκλο $C_2$ κέντρου $\Lambda\left(-\frac{5}{2}, -\frac{7}{2}\right)$ και ακτίνας $R = 2\sqrt{2}$. \monades{06}
    \begin{rlist}
        \item Να αποδείξετε ότι οι δύο κύκλοι $C_1$ και $C_2$ εφάπτονται εξωτερικά και στη συνέχεια να

        \item Να αποδείξετε ότι η ευθεία $(\varepsilon)$ είναι κοινή εφαπτομένη των κύκλων $C_1$ και $C_2$.  βρείτε την ελάχιστη και τη μέγιστη απόσταση των σημείων τους. \monades{06}
 \monades{06}
    \end{rlist}\end{alist}

\vspace{4mm}

% --- Άσκηση 15178 (15178.tex) ---
\Askhsh{}{3}
Δίνεται η εξίσωση $(\mu + 1)x + (\mu + 2)y = 0$, $\mu \in \mathbb{R}$ (1).

\begin{alist}    \item Να αποδείξετε ότι η εξίσωση (1) παριστάνει ευθεία για κάθε $\mu \in \mathbb{R}$. \monades{08}
    \item Να αποδείξετε ότι όλες οι ευθείες που ορίζονται από την εξίσωση (1) διέρχονται από την αρχή των αξόνων. \monades{02}\end{alist}

\begin{alist}    \item Πότε η εξίσωση (1) παριστάνει ευθεία που έχει συντελεστή διεύθυνσης 0; Ποια είναι η εξίσωσή της; \monades{03}
    \item Πότε η εξίσωση (1) παριστάνει ευθεία για την οποία δεν ορίζεται συντελεστής διεύθυνσης; Ποια είναι η εξίσωσή της; \monades{03}\end{alist}

\begin{alist}    \item Να βρείτε για ποια τιμή του πραγματικού αριθμού $\mu$, προκύπτει ευθεία η οποία σχηματίζει γωνία $45^\circ$ με τον άξονα $\chi'\chi$. Ποια είναι η εξίσωσή της; \monades{09}\end{alist}

\vspace{4mm}

% --- Άσκηση 15186 (15186.tex) ---
\Askhsh{}{2} % Assuming an empty code and Theme 2 based on the PDF content
Δίνονται τα σημεία $A(2,1)$, $B(6,3)$, $\Delta(1, -2)$ και $\Gamma'(9,2)$.
Να αποδείξετε ότι:
\begin{alist}    \item Το μέσο $M$ του τμήματος $AB$ έχει συντεταγμένες $(4,2)$ και το μέσο $N$ του τμήματος $\Gamma'\Delta$ έχει συντεταγμένες $(5,0)$.
    \item $\vec{MN} = (1, -2)$ και $\vec{\Delta\Gamma'} = (8,4)$.  \monades{8}
    \item $\vec{MN} \perp \vec{\Delta\Gamma'}$.  \monades{8}
 \monades{9}\end{alist}

\vspace{4mm}

% --- Άσκηση 15189 (15189.tex) ---
\Askhsh{GEO4}{4}
Δίνονται τα σημεία $A(-2,0)$ και $B(2, -2)$.
\begin{alist}    \item Να βρείτε τις συντεταγμένες του μέσου $K$ και το μήκος του ευθυγράμμου τμήματος $AB$.
    \item Να δείξετε ότι ο κύκλος $C$ με διάμετρο $AB$ έχει εξίσωση $C: x^2 + (y + 1)^2 = 5$.  \monades{6}
    \item Να δείξετε ότι τα σημεία $M(x, y)$ του επιπέδου για τα οποία $(AMB) = 5$ ανήκουν στις ευθείες $\varepsilon1: x + 2y - 3 = 0$ και $\varepsilon2: x + 2y + 7 = 0$.  \monades{6}
    \item Να δείξετε ότι οι ευθείες $\varepsilon1$ και $\varepsilon2$ εφάπτονται του κύκλου $C$.  \monades{7}
 \monades{6}\end{alist}

\vspace{4mm}

% --- Άσκηση 15194 (15194.tex) ---
ΘΕΜΑ 4
Δίνονται τα σημεία του επιπέδου $A(1,1)$, $B(4,4)$ και $\Gamma(3,1)$.
\begin{alist}    \item Να δείξετε ότι τα σημεία αυτά σχηματίζουν τρίγωνο.
    \item Να δείξετε ότι η μεσοκάθετος του τμήματος $B\Gamma$ είναι η ευθεία $(\varepsilon)$: $y = -\frac{1}{3}x + \frac{11}{3}$.  \monades{7}
    \item Να βρείτε σημείο $K$ της ευθείας $(\varepsilon)$ του β) ερωτήματος τέτοιο ώστε $(KA) = (KB)$. Τι ιδιότητα έχει το σημείο $K$;  \monades{9}
 \monades{9}\end{alist}

\vspace{4mm}

% --- Άσκηση 15252 (15252.tex) ---
\textbf{ΘΕΜΑ 2}

Δίνονται τα διανύσματα $\vec{AB} = (2,1)$ και $\vec{A\Gamma} = (3,-1)$.

\begin{alist}    \item Να δείξετε ότι $\vec{B\Gamma} = (1,-2)$. \monades{8}
    \item Να δείξετε ότι το τρίγωνο ΑΒΓ είναι ορθογώνιο και ισοσκελές με υποτείνουσα την ΑΓ. \monades{9}
    \item Να βρείτε το εμβαδόν του τριγώνου ΑΒΓ. \monades{8}\end{alist}

\vspace{4mm}

% --- Άσκηση 15253 (15253.tex) ---
\Askhsh{}{4}
Δίνεται η εξίσωση
\[ (\mu^2 - 1)x + (3\mu^2 - 2\mu - 1)y - 5\mu^2 + 4\mu + 1 = 0 \quad (1) \]
όπου $\mu \in \mathbb{R}$.

\begin{alist}    \item Να βρείτε για ποιες τιμές του $\mu$ η (1) παριστάνει ευθεία $\varepsilon$. \monades{5}
    \item Να βρείτε για ποιες τιμές του $\mu$ οι ευθείες $\varepsilon$:
    \begin{rlist}
        \item είναι παράλληλες στον $xx'$. \monades{4}
        \item είναι παράλληλες στον $yy'$. \monades{4}
        \item διέρχονται από το $(0,0)$. \monades{4}
    \end{rlist}
    \item Να δείξετε ότι όλες οι ευθείες $\varepsilon$ που προκύπτουν από την (1) διέρχονται από σταθερό σημείο. \monades{8}\end{alist}

\vspace{4mm}

% --- Άσκηση 15271 (15271.tex) ---
\OnPar{ΘΕΜΑ 2}

Δίνονται τα σημεία $A(-3, 2)$, $B(1, 6)$ και $\Gamma(-13, -7)$.

\begin{alist}    \item Να βρείτε το συντελεστή διεύθυνσης της ευθείας που διέρχεται από τα $A$, $B$. \monades{8}
    \item Να αποδείξετε ότι η ευθεία που διέρχεται από τα $A$, $B$ έχει εξίσωση $y=x+5$. \monades{7}
    \item Να αιτιολογήσετε γιατί το σημείο $\Gamma$ δεν είναι πάνω στην $AB$. \monades{10}\end{alist}

\vspace{4mm}

% --- Άσκηση 15272 (15272.tex) ---
\Askhseis{ΘΕΜΑ 4}

Δίνεται η εξίσωση $x^2 + y^2 - 2x + 4y = -1$.

\begin{alist}    \item Να αποδείξετε ότι παριστάνει κύκλο του οποίου να βρείτε το κέντρο και την ακτίνα. \monades{6}
    \item Να αποδείξετε ότι το σημείο $M(3, 2)$ βρίσκεται έξω από τον κύκλο. \monades{7}
    \item Να βρείτε τις εφαπτόμενες του κύκλου που διέρχονται από το $M$. \monades{12}\end{alist}

\vspace{4mm}

% --- Άσκηση 15273 (15273.tex) ---
\OnPar{ΘΕΜΑ 4}

Θεωρούμε τα σταθερά σημεία $A(3,4)$, $B(2,5)$ και $\Gamma(-2,2)$ και το μεταβλητό σημείο $M(4\alpha - 1, 3\alpha + 1)$, $\alpha \in \mathbb{R}$.

\begin{alist}    \item Να αποδείξετε ότι τα $A$, $B$, $\Gamma$ σχηματίζουν τρίγωνο.
    \item Να βρείτε την εξίσωση της ευθείας $B\Gamma$.  \monades{5}
    \item Να αποδείξετε ότι τα σημεία $M$ κινούνται στην ευθεία που διέρχεται από το $A$ και είναι παράλληλη στην $B\Gamma$.  \monades{5}
    \item Να αποδείξετε ότι για οποιαδήποτε θέση του σημείου $M$ ισχύει $(MB\Gamma) = (AB\Gamma)$. Πως αιτιολογείται αυτό γεωμετρικά;  \monades{7}
 \monades{8}\end{alist}

\vspace{4mm}

% --- Άσκηση 15275 (15275.tex) ---
\Askhsh{}{4}

Σε ορθοκανονικό σύστημα συντεταγμένων θεωρούμε το σημείο $M(2, 1)$.

\begin{alist}    \item Μια ευθεία $(\varepsilon)$ με συντελεστή διεύθυνσης $\lambda$ διέρχεται από το $M$. Να βρείτε:
    \begin{rlist}
        \item Την εξίσωση της. \monades{2}
        \item Για ποιες τιμές του $\lambda$ η ευθεία σχηματίζει τρίγωνο με τους άξονες. \monades{5}
    \end{rlist}
    \item Έστω ότι η ευθεία $(\varepsilon)$ τέμνει τους άξονες $x'x$ και $y'y$ στα σημεία $A$, $B$ αντίστοιχα.
    \begin{rlist}
        \item Να βρείτε, με τη βοήθεια του $\lambda$, τα μήκη των τμημάτων $OA$, $OB$. \monades{6}
        \item Να βρείτε για ποιες τιμές του $\lambda$ η ευθεία σχηματίζει με τους άξονες ισοσκελές τρίγωνο. \monades{6}
        \item Να υπολογίσετε, σε κάθε περίπτωση, το εμβαδόν του ισοσκελούς τριγώνου που σχηματίζεται. \monades{6}
    \end{rlist}\end{alist}

\vspace{4mm}

% --- Άσκηση 15317 (15317.tex) ---
\Askhsh{ALGB201}{2}

Δίνονται τα διανύσματα $\vec{v} = (3,0)$ και $\vec{w} = (-3,4)$.

\begin{alist}    \item Να δείξετε ότι τα διανύσματα δεν είναι παράλληλα. \monades{12}
    \item
    \begin{rlist}
        \item Στο παρακάτω σύστημα συντεταγμένων να σχεδιάσετε τα διανύσματα $\vec{v}$ και $\vec{w}$.
        \begin{center}
        \begin{tikzpicture}        \begin{axis}[
 axis lines = middle, xmin = -3, xmax = 5, ymin = -5, ymax = 5, xtick = {-2,-1,0,1,2,3,4}, ytick = {-4,-3,-2,-1,0,1,2,3,4}, grid = both, grid style = {gray!30}, width = 10cm, height = 10cm, tick label style = {font=\scriptsize},
 ]
            % Draw vector v
 \draw[->, very thick, maincolor] (0,0) -- (3,0) node[above right, maincolor] {$\vec{v}$};
            % Draw vector w
 \draw[->, very thick, secondarycolor] (0,0) -- (-3,4) node[above left, secondarycolor] {$\vec{w}$};
        \end{axis}\end{tikzpicture}
        \end{center}

        \item Να προσδιορίσετε το είδος της γωνίας $\theta$ που σχηματίζουν τα διανύσματα. \monades{3}  \monades{10}
    \end{rlist}\end{alist}

\vspace{4mm}

% --- Άσκηση 15320 (15320.tex) ---
\section*{ΘΕΜΑ 4}

Δίνεται παραλληλόγραμμο ΟΑΓΒ με $\vec{OA}=\vec{a}$ και $\vec{OB}=\vec{\beta}$, όπου $\vec{a}$ και $\vec{\beta}$ είναι μη μηδενικά διανύσματα.

\begin{alist}    \item Να δείξετε ότι:
    \begin{rlist}
        \item $\left|\vec{O\Gamma}\right|^2 = \left|\vec{a}\right|^2 + 2\vec{a}\cdot\vec{\beta} + \left|\vec{\beta}\right|^2$. \monades{9}
        \item $\left|\vec{AB}\right|^2 = \left|\vec{a}\right|^2 - 2\vec{a}\cdot\vec{\beta} + \left|\vec{\beta}\right|^2$. \monades{9}
    \end{rlist}
    \item Αν $\left|\vec{O\Gamma}\right| = \left|\vec{AB}\right|$, να δείξετε ότι το ΟΑΓΒ είναι ορθογώνιο. \monades{7}\end{alist}

\vspace{4mm}

% --- Άσκηση 15379 (15379.tex) ---
\Askhsh{}{2}
Δίνονται τα διανύσματα $\vec{a} = (1,3)$, $\vec{b} = (3,-1)$.
Να υπολογίσετε:
\begin{alist}    \item το εσωτερικό γινόμενο $\vec{a} \cdot \vec{b}$ και την γωνία μεταξύ των δύο διανυσμάτων $\vec{a}$, $\vec{b}$. \monades{13}
    \item το διάνυσμα $\vec{\gamma} = 2\vec{a} - \vec{b}$. \monades{12}\end{alist}

\vspace{4mm}

% --- Άσκηση 15380 (15380.tex) ---
\Askhsh{4}{4}
Δίνονται τα σημεία $A(1,3)$, $B(-2,2)$ και η ευθεία $\epsilon: 3x + y + \alpha = 0$, με $\alpha \in \mathbb{R}$.

\begin{alist}    \item Να βρείτε για ποια τιμή του $\alpha$, η απόσταση του σημείου $A$ από το σημείο $B$ είναι ίση με την απόσταση του σημείου $A$ από την ευθεία $\epsilon$.
    \item Για $\alpha = 4$  \monades{8}
    \begin{rlist}
        \item Να βρείτε το εμβαδόν του τριγώνου $AB\Gamma$, όπου $\Gamma$ το σημείο τομής της ευθείας $\epsilon$ με τον άξονα $y'y$.
        \item Να βρείτε το σημείο της ευθείας $\epsilon$ που απέχει την μικρότερη απόσταση από την αρχή των αξόνων.  \monades{8}
 \monades{9}
    \end{rlist}\end{alist}

\vspace{4mm}

% --- Άσκηση 15394 (15394.tex) ---
\paragraph{\textbf{ΘΕΜΑ 4}}

Στο παρακάτω σχήμα δίνεται η παραβολή $C: y^2 = 12x$ με εστία $E$ και η εφαπτομένη ευθεία $(\varepsilon)$ της $(C)$ στο σημείο της $M(1,2\sqrt{3})$, η οποία τέμνει τον άξονα $x'x$ στο σημείο $B$. Από το σημείο $M$ φέρνουμε ευθεία $t't$ παράλληλη προς τον άξονα $x'x$, η οποία τέμνει την διευθετούσα $(\delta)$ στο σημείο $H$.

\begin{alist}    \item Να αποδείξετε ότι η $(\varepsilon)$ έχει εξίσωση $y = \sqrt{3}x + \sqrt{3}$. \monades{6}
    \item Να βρείτε τις συντεταγμένες των σημείων $B$, $H$, $E$. \monades{6}
    \item Να αποδείξετε ότι το τετράπλευρο $MEBH$ είναι ρόμβος. \monades{7}
    \item Να βρείτε την εξίσωση της ευθείας $(\zeta)$ η οποία διχοτομεί την γωνία $E\hat{M}t$. \monades{6}\end{alist}

\begin{center}
\begin{tikzpicture}[scale=0.8]
    % Define the parabola parameters
    \pgfmathsetmacro{\p}{6} % 2p = 12 => p = 6
    \pgfmathsetmacro{\focx}{\p/2} % Focus E: (p/2, 0) = (3,0)
    \pgfmathsetmacro{\dirx}{-\p/2} % Directrix (delta): x = -p/2 = -3
    % Define points
    \tkzDefPoint(0,0){O}
    \tkzDefPoint(\focx,0){E} % E(3,0)
    \tkzDefPoint(1,{2*sqrt(3)}){M} % M(1, 2√3)
    % Point B: intersection of tangent y=sqrt(3)x+sqrt(3) with y=0 => x=-1
    \tkzDefPoint(-1,0){B}
    % Point H: intersection of line y=2*sqrt(3) with directrix x=-3
    \tkzDefPoint(\dirx,{2*sqrt(3)}){H}
    % Reference points for labels and lines
    \tkzDefPoint(\dirx,4){deltalabelpoint}
    \tkzDefPoint(3,{sqrt(3)*3 + sqrt(3)}){epsilonlabelpoint} % A point on tangent for label
    \tkzDefPoint(5,{2*sqrt(3)}){tlabelpoint} % A point on line t't for label t
    \tkzDefPoint(-5,{2*sqrt(3)}){tprimelabelpoint} % A point on line t't for label t'
    \tkzDefPoint(2,{2*sqrt(3)}){Tangleref} % A point to define the direction Mt for angle
    % Draw axes x'x and y'y
    \draw[l3,->] (-4.5,0) node[below left] {$x'$} -- (5.5,0) node[below right] {$x$};
    \draw[l3,->] (0,-3.5) node[below left] {$y'$} -- (0,4.5) node[above left] {$y$};
    \tkzLabelPoint[below left](O){O}
    % Draw Parabola C: y^2 = 12x
    \draw[maincolor, thick, samples=100, domain=0:5.5] plot (\x, {sqrt(12*\x)});
    \draw[maincolor, thick, samples=100, domain=0:5.5] plot (\x, {-sqrt(12*\x)});
    \node[maincolor, above right, xshift=1mm, yshift=1mm] at (2.5, {sqrt(12*2.5)}) {$(C): y^2 = 12x$};
    % Draw Directrix (delta)
    \draw[red, l3] (\dirx,-4) -- (\dirx,4);
    \node[red, above left, yshift=2mm] at (deltalabelpoint) {$(\delta)$};
    % Draw tangent (epsilon)
    \draw[secondarycolor, l3] ($(B)!-0.8!(M)$) -- ($(M)!1.5!(B)$); % Draw extending beyond M and B
    \node[secondarycolor, above right] at (epsilonlabelpoint) {$(\varepsilon)$};
    % Draw line t't
    \draw[thirdcolor, l3] ($(H)!-0.5!(M)$) -- ($(M)!1.5!(H)$); % Draw extending beyond H and M
    \node[thirdcolor, below left] at (tprimelabelpoint) {$t'$};
    \node[thirdcolor, below right] at (tlabelpoint) {$t$};
    % Draw points
    \tkzDrawPoints(O,E,M,B,H)
    % Label points
    \tkzLabelPoint[below](E){E}
    \tkzLabelPoint[left](B){B}
    \tkzLabelPoint[above left](H){H}
    \tkzLabelPoint[above right](M){M} 
    % Draw segments for the rhombus MEBH
    \draw[l3] (M) -- (E);
    \draw[l3] (E) -- (B);
    \draw[l3] (B) -- (H);
    % Angle mark for EMt
    \tkzMarkAngle[fill=secondarycolor!10, size=0.6cm, opacity=0.8, color=secondarycolor](E,M,Tangleref)\end{tikzpicture}
\end{center}

\vspace{4mm}

% --- Άσκηση 15432 (15432.tex) ---
\Askhsh{}{}
\OnPar{ΘΕΜΑ 4}
Δίνεται η εξίσωση $x^2 + y^2 - 4\kappa x - 2\kappa y + 4 = 0$ (1) με $\kappa \in \mathbb{R}$.
\begin{alist}    \item Να βρείτε τις τιμές του $\kappa \in \mathbb{R}$ ώστε η εξίσωση (1) να παριστάνει κύκλο. \monades{7}
    \item Να βρείτε τις συντεταγμένες του κέντρου και την ακτίνα του κάθε κύκλου. \monades{3}
    \item Να βρείτε την ευθεία στην οποία ανήκουν τα κέντρα των παραπάνω κύκλων. \monades{7}
    \item Για $\kappa = 1$ να βρείτε την εξίσωση εφαπτομένης του αντίστοιχου κύκλου της εξίσωσης (1) στο σημείο $\Gamma(2,2)$. \monades{8}\end{alist}

\vspace{4mm}

% --- Άσκηση 15433 (15433.tex) ---
\Askhsh{}{4}
Δύο οικισμοί Α και Β βρίσκονται στις θέσεις που ορίζουν τα σημεία $A(-1, -2)$ και $B(3,1)$. Εξωτερικά των οικισμών υπάρχει ευθύγραμμος δρόμος με εξίσωση $\delta: x + y - 1 = 0$.

\begin{alist}    \item Να βρείτε σε ποια θέση του δρόμου $\delta$:
    \begin{rlist}
        \item Ο οικισμός Α έχει τη μικρότερη απόσταση από τον δρόμο. \monades{8}
        \item Υπάρχει το Κέντρο Υγείας της περιοχής, αν είναι γνωστό ότι ισαπέχει από τους δύο οικισμούς. \monades{7}
    \end{rlist}
    \item Να βρείτε τη θέση $\Gamma$ ενός αυτοκινήτου πάνω στο δρόμο, αν είναι γνωστό, ότι το εμβαδόν του τριγώνου που σχηματίζουν τα τρία σημεία Α, Β και $\Gamma$ είναι ίσο με 8. \monades{10}\end{alist}

\vspace{4mm}

% --- Άσκηση 15439 (15439.tex) ---
\Askhsh{A2000004}{4}
Μία φωτεινή ακτίνα διερχόμενη από το σημείο $A(2,3)$ και προσπίπτουσα στην ευθεία $(\varepsilon)$
με εξίσωση $x + y + 1 = 0$, μετά την ανάκλασή της διέρχεται από το σημείο $M(1,1)$.

\begin{center}
\begin{tikzpicture}[scale=0.8]
    % Define points
    \tkzDefPoint(0,0){O}
    \tkzDefPoint(2,3){A}
    \tkzDefPoint(1,1){M}
    \tkzDefPoint(-1,0){P} % Point Pi
    \tkzDefPoint(-0.5,-0.5){S} % Point Sigma (visually estimated from PDF)
    \tkzDefPoint(-4,-3){Ap} % Point A' from question a.ii.
    % Draw axes
    \draw[l3,->] (-5,0) -- (4,0) node[below right] {$x$};
    \draw[l3,->] (0,-4) -- (0,4) node[above left] {$y$};
    \tkzLabelPoint[below left](O){O}
    % Draw line (epsilon)
    % Find another point on x+y+1=0 to draw the line. E.g., (0,-1)
    \tkzDefPoint(0,-1){Py}
    \tkzDrawLine[add=1 and 1.5,l3](P,Py) % Extend line in both directions
    \tkzLabelLine[below right,pos=1.1](P,Py){$(\varepsilon): x+y+1=0$}
    % Draw incident ray (epsilon1)
    \tkzDrawSegment[l3,->,color=black](A,S)
    \tkzLabelLine[above left,pos=0.2](A,S){$(\varepsilon1)$}
    % Draw reflected ray (epsilon2)
    \tkzDrawSegment[l3,->,color=black](S,M)
    \tkzLabelLine[above right,pos=0.4](S,M){$(\varepsilon2)$}
    % Draw dashed line from A' to S
    \tkzDrawSegment[dashed,l3](Ap,S)
    % Draw points
    \tkzDrawPoints(A,M,P,S,Ap)
    % Label points
    \tkzLabelPoint[above left](A){$A(2,3)$}
    \tkzLabelPoint[below right](M){$M(1,1)$}
    \tkzLabelPoint[below left](P){$\Pi(-1,0)$}
    \tkzLabelPoint[above right](S){$\Sigma$}
    \tkzLabelPoint[below left](Ap){$A'$}\end{tikzpicture}
\end{center}

\begin{alist}
\item[α)]
    \begin{rlist}        \item Να αποδείξετε ότι η προβολή του σημείου $A$ πάνω στην ευθεία $(\varepsilon)$ είναι το σημείο $\Pi(-1,0)$. \monades{7}
        \item Να αποδείξετε ότι το συμμετρικό του σημείου $A$ ως προς την ευθεία $(\varepsilon)$, είναι το σημείο $A'(-4, -3)$. \monades{5}\end{rlist}
\item[β)]
    \begin{rlist}        \item Αν γνωρίζετε ότι η ανακλώμενη ακτίνα είναι η ευθεία $(\varepsilon2)$, η οποία διέρχεται από τα σημεία $A', \Sigma, M$, τότε να βρείτε την εξίσωσή της. \monades{4}
        \item Να βρείτε τις συντεταγμένες του σημείου πρόσπτωσης $\Sigma$ της φωτεινής ακτίνας $(\varepsilon1)$ πάνω στην ευθεία $(\varepsilon)$. \monades{5}\end{rlist}
\item[γ)] Αν $\Sigma\left(-\frac{2}{3}, -\frac{1}{3}\right)$, τότε να βρείτε την εξίσωση της προσπίπτουσας ακτίνας $(\varepsilon1)$. \monades{4}
\end{alist}

\vspace{4mm}

% --- Άσκηση 15440 (15440.tex) ---
\Askhsh{}{2}
Δίνονται τα σημεία $A(0,2)$, $B(3,0)$ και $\Gamma(1,1)$.
\begin{alist}    \item Να βρείτε τις συντεταγμένες των διανυσμάτων $\overrightarrow{AB}$, $\overrightarrow{A\Gamma}$. \monades{9}
    \item
    \begin{rlist}
        \item Να εξετάσετε αν τα σημεία $A$, $B$ και $\Gamma$ ορίζουν τρίγωνο. \monades{8}
        \item Να υπολογίσετε το εμβαδόν του τριγώνου $AB\Gamma$. \monades{8}
    \end{rlist}\end{alist}

\vspace{4mm}

% --- Άσκηση 15463 (15463.tex) ---
\textbf{ΘΕΜΑ 2}

Δίνονται τα διανύσματα $\overrightarrow{\text{AB}} = (2,1)$ και $\overrightarrow{A\Gamma} = (3,-1)$.

\begin{alist}    \item Να αποδείξετε ότι $\overrightarrow{B\Gamma} = (1,-2)$. \monades{8}
    \item Να αποδείξετε ότι $\overrightarrow{\text{AB}} \perp \overrightarrow{B\Gamma}$. \monades{9}
    \item Να αποδείξετε ότι $|\overrightarrow{\text{AB}}| = |\overrightarrow{B\Gamma}|$. \monades{8}\end{alist}

\vspace{4mm}

% --- Άσκηση 15475 (15475.tex) ---
\textbf{ΘΕΜΑ 4}

Δύο εργοστάσια $A$ και $B$ τα οποία σε ένα σύστημα συντεταγμένων έχουν συντεταγμένες $A(2,1)$ και $B(4,3)$, βρίσκονται κοντά σε μια ακτή που πρόκειται να κατασκευαστεί μια αποβάθρα και θα εξυπηρετεί τα δύο εργοστάσια.

\begin{center}
\begin{tikzpicture}[scale=0.8] % Προσαρμόστε την κλίμακα όπως χρειάζεται
    % Φόντο (λόφοι) - απλοποιημένη αναπαράσταση
    \fill[thirdcolor!30] (-2,0.5) .. controls (-1,1.5) and (0,1.0) .. (1,2.0) .. controls (2,2.5) and (3,1.8) .. (4,2.2) .. controls (5,2.0) and (6,1.0) .. (7,0.5) -- (7,0) -- (-2,0) -- cycle;
    \fill[thirdcolor!20] (-2,0) rectangle (7,0.5); % Επίπεδο έδαφος
    % Περιοχή νερού
    \fill[maincolor!10] (-2,-2) rectangle (7,0);
    % Γραμμή ακτής (ΑΚΤΗ)
    \draw[line width=1.5pt, draw=maincolor!80!black] (-2,0) .. controls (-1,0.3) and (0,-0.1) .. (1,0.2) .. controls (2,0.4) and (3,-0.2) .. (4,0.1) .. controls (5,0.3) and (6,-0.1) .. (7,0);
    \node[below, maincolor!80!black, font=\bfseries] at (2.5,-0.7) {ΑΚΤΗ};
    % Εργοστάσιο A
    \colorlet{factoryAcolor}{maincolor!60!black}
    % Κεντρικό κτίριο (απλό 3D εφέ)
    \draw[fill=factoryAcolor, draw=factoryAcolor!80!black, line width=0.5pt] (0.5,0.5) -- (1.5,0.5) -- (1.5,1.5) -- (0.5,1.5) -- cycle; % Μπροστινή όψη
    \draw[fill=factoryAcolor!40!black, draw=factoryAcolor!80!black, line width=0.5pt] (0.5,0.5) -- (0.2,0.7) -- (0.2,1.7) -- (0.5,1.5) -- cycle; % Πλαϊνή όψη (αριστερά)
    \draw[fill=factoryAcolor!80!black, draw=factoryAcolor!80!black, line width=0.5pt] (0.5,1.5) -- (0.2,1.7) -- (1.2,1.9) -- (1.5,1.5) -- cycle; % Πάνω όψη
    % Καμινάδα για το A
    \draw[fill=factoryAcolor!70!black, draw=factoryAcolor!90!black, line width=0.5pt] (1.2,1.5) -- (1.3,1.5) -- (1.3,2.0) -- (1.2,2.0) -- cycle;
    \draw[fill=factoryAcolor!50!black, draw=factoryAcolor!90!black, line width=0.5pt] (1.2,1.5) -- (1.1,1.6) -- (1.1,2.1) -- (1.2,2.0) -- cycle; % Πλαϊνή όψη καμινάδας
    \draw[fill=factoryAcolor!90!black, draw=factoryAcolor!90!black, line width=0.5pt] (1.2,2.0) -- (1.1,2.1) -- (1.25,2.2) -- (1.35,2.1) -- cycle; % Πάνω όψη καμινάδας
    \node[above=0pt, font=\bfseries, color=maincolor!80!black] at (1.0,1.8) {A}; % Ετικέτα για το A
    % Εργοστάσιο B
    \colorlet{factoryBcolor}{secondarycolor!60!black}
    % Κεντρικό κτίριο (απλό 3D εφέ)
    \draw[fill=factoryBcolor, draw=factoryBcolor!80!black, line width=0.5pt] (4.5,0.5) -- (5.5,0.5) -- (5.5,1.5) -- (4.5,1.5) -- cycle; % Μπροστινή όψη
    \draw[fill=factoryBcolor!40!black, draw=factoryBcolor!80!black, line width=0.5pt] (4.5,0.5) -- (4.2,0.7) -- (4.2,1.7) -- (4.5,1.5) -- cycle; % Πλαϊνή όψη (αριστερά)
    \draw[fill=factoryBcolor!80!black, draw=factoryBcolor!80!black, line width=0.5pt] (4.5,1.5) -- (4.2,1.7) -- (5.2,1.9) -- (5.5,1.5) -- cycle; % Πάνω όψη
    % Καμινάδα για το B
    \draw[fill=factoryBcolor!70!black, draw=factoryBcolor!90!black, line width=0.5pt] (5.2,1.5) -- (5.3,1.5) -- (5.3,2.0) -- (5.2,2.0) -- cycle;
    \draw[fill=factoryBcolor!50!black, draw=factoryBcolor!90!black, line width=0.5pt] (5.2,1.5) -- (5.1,1.6) -- (5.1,2.1) -- (5.2,2.0) -- cycle; % Πλαϊνή όψη καμινάδας
    \draw[fill=factoryBcolor!90!black, draw=factoryBcolor!90!black, line width=0.5pt] (5.2,2.0) -- (5.1,2.1) -- (5.25,2.2) -- (5.35,2.1) -- cycle; % Πάνω όψη καμινάδας
    \node[above=0pt, font=\bfseries, color=secondarycolor!80!black] at (5.0,1.8) {B}; % Ετικέτα για το B
\end{tikzpicture}
\end{center}

\begin{alist}    
\item Να βρείτε την εξίσωση της ευθείας που συνδέει τα δύο εργοστάσια. \monades{8}
    \item Αν η ακτή είναι ευθύγραμμη με εξίσωση $\varepsilon:y = 2x - 7$, να βρείτε τις συντεταγμένες του σημείου της ακτής στο οποίο πρέπει να τοποθετηθεί η αποβάθρα ώστε να απέχει εξ ίσου από τα δύο εργοστάσια. \monades{10}
    \item Αν το ζητούμενο σημείο του ερωτήματος β) είναι $N(4,1)$, να βρείτε πόσο απέχει το κάθε εργοστάσιο από το σημείο αυτό. \monades{7}
\end{alist}

\vspace{4mm}

% --- Άσκηση 15628 (15628.tex) ---
\Askhsh{}{4}
Δίνεται η εξίσωση $x^2 + y^2 + (4 - 2k)x - 2(1 + k)y + 5 - 2k = 0$ (I), όπου $k \in (0, +\infty)$.
\begin{alist}    \item Να αποδείξετε ότι η (I) παριστάνει κύκλο με κέντρο $M(k - 2, k + 1)$ και ακτίνα $R = k\sqrt{2}$ για κάθε $k > 0$. \monades{10}
    \item Να αποδείξετε ότι το σημείο $M$ ανήκει σε μια σταθερή ευθεία για κάθε $k > 0$. \monades{7}
    \item Να αποδείξετε ότι η ευθεία $(\varepsilon): y = -x - 1$ είναι εφαπτομένη του παραπάνω κύκλου για κάθε $k > 0$. \monades{8}\end{alist}

\vspace{4mm}

% --- Άσκηση 15646 (15646.tex) ---
\Askhsh{}{4}
Δίνονται οι κύκλοι $C_1: (x-1)^2 + (y-1)^2 = 9$ και $C_2: (x-4)^2 + (y-4)^2 = 9$.

\begin{alist}    \item Να δείξετε ότι τα κέντρα $K,\Lambda$ των κύκλων $C_1$ και $C_2$ αντίστοιχα βρίσκονται στην διχοτόμο της γωνίας $xOy$ του συστήματος συντεταγμένων. \monades{8}
    \item Να βρείτε τα σημεία τομής $B,\Gamma$ των κύκλων $C_1$ και $C_2$. \monades{7}
    \item Να βρείτε τα σημεία της ευθείας $y = x$ ώστε το τρίγωνο που σχηματίζεται με τα $B,\Gamma$ να έχει εμβαδόν $\frac{21}{2}$ τ.μ. \monades{10}\end{alist}

\vspace{4mm}

% --- Άσκηση 15657 (15657.tex) ---
\textbf{ΘΕΜΑ 2}

Δίνονται οι ευθείες: $\varepsilon1 : 2x + y = 6$ και $\varepsilon2 : x-2y = -2$.

\begin{alist}    \item Να βρείτε το κοινό τους σημείο M. \monades{12}
    \item Να δείξετε ότι οι ευθείες $\varepsilon1$, $\varepsilon2$ και $\varepsilon3 : 3x - y = 4$ διέρχονται από το ίδιο σημείο. \monades{13}\end{alist}

\vspace{4mm}

% --- Άσκηση 15658 (15658.tex) ---
\Askhsh{}{4}
Δίνονται τα διανύσματα $\vec{\alpha} = (2,-2)$ και $\vec{\beta} = (1,1)$ τα οποία έχουν κοινή αρχή το σημείο $K(2,1)$.

\begin{alist}    \item Να αποδείξετε ότι τα διανύσματα $\vec{\alpha}$ και $\vec{\beta}$ είναι κάθετα. \monades{4}
    \item Αν το σημείο Α είναι το πέρας του διανύσματος $\vec{\alpha}$, Β είναι το πέρας του διανύσματος $\vec{\beta}$ και $\Gamma(x_\Gamma, y_\Gamma)$ ένα τυχαίο σημείο της ευθείας ΑΒ,
    \begin{rlist}
        \item να δείξετε ότι οι συντεταγμένες των σημείων Α και Β είναι $Α(4,-1)$ και $B(3,2)$. \monades{5}
        \item να δείξετε ότι $3x_\Gamma + y_\Gamma = 11$. \monades{6}
        \item να βρείτε τις συντεταγμένες του σημείου $\Gamma(x_\Gamma, y_\Gamma)$, αν ισχύει ότι το $\Gamma$ είναι εσωτερικό σημείο του ευθύγραμμου τμήματος ΑΒ και $\left|\vec{K\Gamma}\right| = \frac{1}{2}\left|\vec{AB}\right|$. \monades{10}
    \end{rlist}\end{alist}

\vspace{4mm}

% --- Άσκηση 15680 (15680.tex) ---
\section*{ΘΕΜΑ 2}
Δίνεται ο κύκλος $C: x^2 + y^2 -2x-4y+1=0$ με κέντρο $K(1,2)$ και η ευθεία $\varepsilon: 3x+4y+1=0$.

\begin{alist}    \item Να αποδείξετε ότι η ακτίνα του κύκλου $C$ είναι $\rho = 2$. \monades{10}
    \item Να αποδείξετε ότι η απόσταση του κέντρου $K$ από την ευθεία $\varepsilon$ είναι $\frac{12}{5}$. \monades{10}
    \item Να αιτιολογήσετε γιατί η ευθεία $\varepsilon$ και ο κύκλος $C$ δεν έχουν κοινά σημεία. \monades{5}\end{alist}

\vspace{4mm}

% --- Άσκηση 15681 (15681.tex) ---
\Askhsh{}{4}

Δίνονται τα σημεία $O(0,0)$, $A(\alpha,0)$, $B\left(\frac{\alpha}{2},\beta\right)$ και $M\left(\frac{\alpha}{2},0\right)$, όπου $\alpha,\beta$ σταθεροί θετικοί πραγματικοί αριθμοί.

\begin{alist}    \item Να μεταφέρετε τα παραπάνω σημεία σε ορθοκανονικό σύστημα συντεταγμένων. Κατόπιν, να αποδείξετε ότι το τρίγωνο $OAB$ είναι ισοσκελές και το σημείο $M$ είναι το μέσο της βάσης του $OA$.
    \item Να αποδείξετε ότι οι εξισώσεις των ευθειών $OB$ και $AB$ είναι $OB: 2\beta x - \alpha y = 0$ και $AB: 2\beta x + \alpha y - 2\alpha\beta = 0$ αντίστοιχα.  \monades{6}
    \item Αν $d_1$ είναι η απόσταση του σημείου $M$ από την ευθεία $OB$ και $d_2$ η απόσταση του σημείου $M$ από την ευθεία $AB$, να αποδείξετε ότι $d_1 = d_2$.  \monades{8}
    \item Ποια πρόταση της Ευκλείδειας Γεωμετρίας έχει αποδειχθεί;  \monades{8}
 \monades{3}\end{alist}

\vspace{4mm}

% --- Άσκηση 15692 (15692.tex) ---
\textbf{ΘΕΜΑ 4}

Δίνεται η εξίσωση
\[ x^2 + y^2 + y = x + 2xy + 6 \]
\begin{alist}    \item Να αποδείξετε ότι:
    \begin{rlist}
        \item
        \item Η εξίσωση παριστάνει ένα ζεύγος παράλληλων ευθειών, τις οποίες να βρείτε. \monades{4}  \[ (x-y)^2 - (x-y) - 6 = 0 \] \monades{4}
    \end{rlist}
    \item Να αποδείξετε ότι όλα τα σημεία $M\left(\alpha, \alpha - \frac{1}{2}\right)$, $\alpha \in \mathbb{R}$ ισαπέχουν από τις δυο ευθείες. \monades{10}  Έστω $\varepsilon1: x-y-3=0$ και $\varepsilon2: x-y+2=0$ οι δυο παράλληλες ευθείες.
    \item Να βρείτε την μεσοπαράλληλη των δυο ευθειών. \monades{7}\end{alist}

\vspace{4mm}

% --- Άσκηση 15791 (15791.tex) ---
\Askhsh{}{4}
Στο παρακάτω σχήμα έχουμε σχεδιάσει κύκλο $C_1$ κέντρου $A$ και την ευθεία $(\varepsilon): x=5$.

\begin{tikzpicture}\begin{axis}[
    axis lines = middle,
    xmin = -4.5, xmax = 8.5,
    ymin = -1.5, ymax = 10.5,
    xtick = {-4,-3,-2,-1,0,1,2,3,4,5,6,7,8},
    ytick = {-1,0,1,2,3,4,5,6,7,8,9,10},
    grid = both,
    grid style = {gray!30},
    width = 10cm, height = 7cm,
    tick label style = {font=\scriptsize},
    axis equal image, 
]
    % Σχεδίαση του κύκλου C1 με κέντρο A(0,7) και ακτίνα 2
    \addplot[thick, black, domain=0:360, samples=100] ({0 + 2*cos(x)}, {7 + 2*sin(x)});
    % Σημείο A
    \node[black] at (axis cs:0.1, 6.7) {A};
    % Σχεδίαση της ευθείας (ε): x=5
    \draw[thick, black] (axis cs:5, \pgfkeysvalueof{/pgfplots/ymin}) -- (axis cs:5, \pgfkeysvalueof{/pgfplots/ymax});
\end{axis}\end{tikzpicture}

\begin{alist}    \item Να βρείτε την εξίσωση του κύκλου $C_1$. \monades{3}
    \item Έστω ένα σημείο του επιπέδου $B(x_1,y_1)$.
    \begin{rlist}
        \item Να βρείτε την εξίσωση του κύκλου με κέντρο $B(x_1, y_1)$ και ακτίνα 2. \monades{6}
        \item Να βρείτε το μήκος της διακέντρου $AB$ σε συνάρτηση με τις συντεταγμένες του σημείου $B$. \monades{6}
    \end{rlist}
    \item Να βρείτε όλους τους κύκλους του ερωτήματος β)ί. με ακτίνα 2, που εφάπτονται εξωτερικά στον $C_1$ και στην ευθεία $(\varepsilon)$. \monades{10}\end{alist}

\vspace{4mm}

% --- Άσκηση 15825 (15825.tex) ---
\noindent \textbf{ΘΕΜΑ 2}

Θεωρούμε τα διανύσματα $\vec{\alpha}, \vec{\beta}$ με $|\vec{\alpha}|=2$, $|\vec{\beta}|=4$, $\angle(\vec{\alpha}, \vec{\beta})=\frac{\pi}{3}$ και το $\vec{\gamma}=\vec{\alpha}-\vec{\beta}$.

\begin{alist}    \item Να αποδείξετε ότι $\vec{\alpha}\cdot \vec{\beta}=4$. \monades{8}
    \item Να αποδείξετε ότι $\vec{\alpha}\cdot \vec{\gamma}=0$. \monades{10}
    \item Να βρείτε τη $\angle(\vec{\alpha}, \vec{\gamma})$. \monades{7}\end{alist}

\vspace{4mm}

% --- Άσκηση 15826 (15826.tex) ---
\Askhsh{}{4}
Δίνεται η εξίσωση
\[ x^2 + y^2 -2(\lambda+1)x-2\lambda y+2\lambda+1=0 \quad (1) \]
όπου $\lambda \in \mathbb{R}$.

\begin{alist}    \item Να βρείτε για ποιες τιμές του $\lambda$ η εξίσωση (1) παριστάνει κύκλο και να γράψετε ως συνάρτηση του $\lambda$ τις συντεταγμένες του κέντρου $K$ και την ακτίνα $\rho$. \monades{7}
    \item Τι παριστάνει η εξίσωση (1) για $\lambda=0$; \monades{3}
    \item Στο παρακάτω σχήμα φαίνονται 4 κύκλοι με τα αντίστοιχα κέντρα τους $K_1, K_2, K_3, K_4$ που προκύπτουν από την (1) για 4 αντίστοιχες τιμές του $\lambda$. Αξιοποιώντας το σχήμα,
    \begin{rlist}
        \item να αποδείξετε ότι τα κέντρα όλων των κύκλων που προκύπτουν από την (1) βρίσκονται πάνω σε μια ευθεία της οποίας να βρείτε την εξίσωση. \monades{5}

        \item να αποδείξετε ότι όλοι οι κύκλοι που προκύπτουν από την (1) διέρχονται από σταθερό σημείο του οποίου να βρείτε τις συντεταγμένες. \monades{5}

        \item να αποδείξετε ότι η ευθεία $\varepsilon: x+y-1=0$ είναι κοινή εφαπτομένη όλων των κύκλων που προκύπτουν από την (1). \monades{5}
    \end{rlist}\end{alist}

\begin{center}
\begin{tikzpicture}\begin{axis}[
    axis lines = middle,
    xmin = -6.5, xmax = 8.5,
    ymin = -4.5, ymax = 5.5,
    xtick = {-6,-5,-4,-3,-2,-1,0,1,2,3,4,5,6,7,8},
    ytick = {-4,-3,-2,-1,0,1,2,3,4,5},
    grid = both,
    grid style = {gray!30},
    width = 0.95\textwidth,
    height = 9cm,
    tick label style = {font=\small},
    xlabel={}, ylabel={}, 
]
    % Circle K1 at (2,3) with radius 1
    \addplot[thick, maincolor, samples=100, domain=0:360] ({2+cos(x)}, {3+sin(x)});
    \node[above right, font=\small] at (axis cs:2,3) {$K1$};
    % Circle K2 at (3,2) with radius 1
    \addplot[thick, maincolor, samples=100, domain=0:360] ({3+cos(x)}, {2+sin(x)});
    \node[above right, font=\small] at (axis cs:3,2) {$K2$};
    % Circle K3 at (-1,-2) with radius 1
    \addplot[thick, maincolor, samples=100, domain=0:360] ({-1+cos(x)}, {-2+sin(x)});
    \node[below left, font=\small] at (axis cs:-1,-2) {$K3$};
    % Circle K4 at (-2,-1) with radius 1
    \addplot[thick, maincolor, samples=100, domain=0:360] ({-2+cos(x)}, {-1+sin(x)});
    \node[below left, font=\small] at (axis cs:-2,-1) {$K4$};
\end{axis}\end{tikzpicture}
\end{center}

\vspace{4mm}

% --- Άσκηση 15852 (15852.tex) ---
\Askhsh{}{2}

Δίνονται τα διανύσματα $\vec{a} = (3,2)$, $\vec{b} = (-2,1)$.
Να υπολογίσετε:
\begin{alist}    \item το διάνυσμα $\vec{v} = 2\vec{a} + 3\vec{b}$. \monades{7}
    \item το εσωτερικό γινόμενο $\vec{a} \cdot \vec{b}$ και το μέτρο του διανύσματος $\vec{a}$. \monades{6}
    \item το εσωτερικό γινόμενο $\vec{a} \cdot \vec{v}$. \monades{12}\end{alist}

\vspace{4mm}

% --- Άσκηση 15854 (15854.tex) ---
\OnPar{ΘΕΜΑ 2}

Δίνονται τα διανύσματα $\vec{\alpha} = (2, 1)$ και $\vec{\beta} = (-8, -4)$.

\begin{alist}    \item Να δείξετε ότι $\vec{\alpha}//\vec{\beta}$.
    \item Να δείξετε ότι για τα διανύσματα $\vec{\alpha}$, $\vec{\beta}$ ισχύει $\vec{\beta} = -4\vec{\alpha}$.  \monades{10}
    \item Να αποδείξετε ότι το μέτρο του διανύσματος $\vec{\beta}$ είναι τετραπλάσιο του μέτρου του διανύσματος $\vec{\alpha}$.  \monades{7}
 \monades{8}\end{alist}

\vspace{4mm}

% --- Άσκηση 15986 (15986.tex) ---
\Askhsh{B21}{2} % Χρησιμοποιούμε μια αναγνωριστική τιμή (π.χ. Β2_1) και το "2" για το ΘΕΜΑ 2
Δίνονται τα σημεία $A(1,1)$ και $B(2,3)$.

\begin{alist}    \item
    \begin{rlist}
        \item Να βρείτε το συντελεστή διεύθυνσης της ευθείας που διέρχεται από τα $A$, $B$.
        \item Να αποδείξετε ότι η εξίσωση της ευθείας $AB$ είναι η $(\varepsilon): y=2x-1$.
    \end{rlist}
    \item Να εξετάσετε αν το σημείο $\Gamma(2^{100}, 5)$ ανήκει στην ευθεία $(\varepsilon)$.  \monades{12}
 \monades{13}\end{alist}

\vspace{4mm}

% --- Άσκηση 15987 (15987.tex) ---
\Askhsh{12345}{4}
Δίνονται τα σημεία $A(1,1)$ και $B(2,3)$.

\begin{alist}    \item Να αποδείξετε ότι η εξίσωση της ευθείας $AB$ είναι η $(\varepsilon): y=2x-1$.
    \item Να αιτιολογήσετε αν το σημείο $\Gamma' (2^{100},5)$ ανήκει ή όχι στο ημιεπίπεδο που ορίζεται από την ευθεία $(\varepsilon)$ και την αρχή των αξόνων $O(0,0)$.  \monades{8}
    \item Να αιτιολογήσετε αν το εμβαδόν του τριγώνου $AB\Gamma$ είναι μεγαλύτερο ή μικρότερο από το εμβαδόν του τριγώνου $AOB$.  \monades{8}
 \monades{9}\end{alist}

\vspace{4mm}

% --- Άσκηση 15993 (15993.tex) ---
\Askhsh{21160}{4}
Δίνεται η εξίσωση $(x-2)^2 + (y-\lambda)^2 = \lambda^2+1$ (1), όπου $\lambda \in \mathbb{R}$.
\begin{alist}    \item Να αποδείξετε ότι για κάθε τιμή του $\lambda$ η (1) παριστάνει κύκλο, του οποίου να βρείτε το κέντρο και την ακτίνα. \monades{03}
    \item Να αποδείξετε ότι όλοι οι κύκλοι που ορίζονται από την (1) για τις διάφορες τιμές του $\lambda$ διέρχονται από δύο σταθερά σημεία. \monades{10}
    \item Αν $A(1,0)$ και $B(3,0)$ είναι τα μοναδικά σημεία από τα οποία διέρχονται όλοι οι κύκλοι, τότε να βρείτε την εξίσωση της κοινής χορδής τους και να αποδείξετε ότι είναι κάθετη στην ευθεία που διέρχεται από τα κέντρα των κύκλων. \monades{07}
    \item Αν ένα σημείο $M(\alpha, \beta)$ επαληθεύει την (1) για κάθε $\lambda \in \mathbb{R}$, τότε να αποδείξετε ότι $\alpha \cdot \beta = 0$. \monades{05}\end{alist}

\vspace{4mm}

% --- Άσκηση 15994 (15994.tex) ---
\Askhsh{LK2BMATH001}{2}
Δίνεται η εξίσωση $x^2 + y^2 - 4x + 3 = 0$ (1).

\begin{alist}    \item Να αποδείξετε ότι η εξίσωση (1) παριστάνει κύκλο του οποίου να βρείτε το κέντρο και την ακτίνα. \monades{13}
    \item Να σχεδιάσετε τον κύκλο (c) και να βρείτε, χρησιμοποιώντας το σχήμα ή με οποιανδήποτε άλλον τρόπο, τα κοινά του σημεία με τους άξονες. \monades{12}
    \begin{center}
    \begin{tikzpicture}[scale=1]
        % Ορισμός του κέντρου και της ακτίνας
 \tkzDefPoint(2,0){K} % Κέντρο του κύκλου \def\radius{1}       % Ακτίνα του κύκλου
        % Σχεδίαση των αξόνων
 \draw[->, gray!50] (-0.5,0) -- (4.5,0) node[right] {$x$}; \draw[->, gray!50] (0,-1.5) -- (0,1.5) node[above] {$y$};
        % Σημείο αρχής των αξόνων
 \node[below left, gray!50] at (0,0) {O};
        % Σχεδίαση του κύκλου (c)
 \draw[maincolor, line width=0.8pt] (K) circle (\radius);
        % Σημείωση του κέντρου K
 \tkzDrawPoint(K) \tkzLabelPoint[below left](K){K$(2,0)$}
        % Ορισμός και σημείωση των σημείων τομής με τον άξονα x
 \tkzDefPoint(1,0){A} \tkzDefPoint(3,0){B} \tkzDrawPoints(A,B) \tkzLabelPoint[below](A){$(1,0)$} \tkzLabelPoint[below](B){$(3,0)$}
        % Προσθήκη σημάνσεων (ticks) στους άξονες
 \foreach \x in {1,2,3,4} \draw (\x, -0.05) -- (\x, 0.05) node[below] {\x}; \foreach \y in {-1,1} \draw (-0.05, \y) -- (0.05, \y) node[left] {\y};\end{tikzpicture}
    \end{center}\end{alist}

\vspace{4mm}

% --- Άσκηση 15996 (15996.tex) ---
\Askhsh{2}{2}{}

Δίνονται τα σημεία $A(-6, -1)$, $B(8,7)$, $\Gamma(-10,6)$, τα οποία ορίζουν τρίγωνο $AB\Gamma$.

\begin{center}
\begin{tikzpicture}\begin{axis}[
    axis lines=middle,
    xlabel={$x$},
    ylabel={$y$},
    xmin=-12.5, xmax=10.5,
    ymin=-2.5, ymax=10.5,
    xtick={-12,-10,...,10}, 
    ytick={-2,0,...,10},     
    major grid style={gray!30, line width=0.2pt}, 
    minor x tick num=1,      
    minor y tick num=1,      
    minor grid style={gray!15, dotted, line width=0.1pt}, 
    grid=both, 
    tick label style={font=\scriptsize},
    every axis x label/.style={at={(current axis.right of origin)}, anchor=north west},
    every axis y label/.style={at={(current axis.above origin)}, anchor=south east},
    width=10cm, height=8cm,
    axis on top=true, 
]
    % Define coordinates for the points
    \coordinate (A) at (-6,-1);
    \coordinate (B) at (8,7);
    \coordinate (Gamma) at (-10,6);
    % Fill the triangle region with a dotted pattern
    \addplot[fill=maincolor!10, pattern=dots, pattern color=gray!50] coordinates {(A) (B) (Gamma)} -- cycle;
    % Draw the triangle outline
    \draw[l4, maincolor] (A) -- (B) -- (Gamma) -- cycle;
    % Plot the points (as small dots, using thirdcolor)
    \addplot[only marks, mark=*, mark size=1.5pt, thirdcolor] coordinates {(A) (B) (Gamma)};
    % Labels for points
    \node[below] at (A) {A};
    \node[above right=2pt and 0pt] at (B) {B};
    \node[above left=2pt and 0pt] at (Gamma) {$\Gamma$};
\end{axis}\end{tikzpicture}
\end{center}

\begin{alist}    \item Να βρείτε τις συντεταγμένες των διανυσμάτων $\vec{AB}$, $\vec{B\Gamma}$ και του αθροίσματος τους $\vec{AB} + \vec{B\Gamma}$. \monades{10}
    \item Ένας μαθητής βλέποντας το τρίγωνο $AB\Gamma$ ισχυρίστηκε ότι είναι ορθογώνιο. Να ελέγξετε την αλήθεια του ισχυρισμού. \monades{15}\end{alist}

\vspace{4mm}

% --- Άσκηση 16002 (16002.tex) ---
\Askhsh{}{2}
Σε τρίγωνο ΑΒΓ είναι $A(3, -2)$ και $\Gamma(5, 2)$. Αν το σημείο $M \left(3, \frac{1}{2}\right)$ είναι το μέσο της $B\Gamma$, τότε:

\begin{alist}    \item Να αποδείξετε ότι $B(1, -1)$. \monades{9}
    \item Να βρείτε το μήκος της πλευράς $B\Gamma$. \monades{6}
    \item Να βρείτε την εξίσωση της ευθείας $A\Gamma$. \monades{10}\end{alist}

\vspace{4mm}

% --- Άσκηση 16003 (16003.tex) ---
\Askhsh{}{4}
Θεωρούμε την οικογένεια των ευθειών $\varepsilon_\alpha : (\alpha-4)x-2\alpha y+\alpha+4=0, \alpha \in \mathbb{R}$.

\begin{alist}    \item Να βρείτε τις ευθείες που προκύπτουν όταν $\alpha=0$ και όταν $\alpha=1$ και κατόπιν να προσδιορίσετε το κοινό τους σημείο Μ. \monades{8}
    \item Να αποδείξετε ότι όλες οι ευθείες της οικογένειας διέρχονται από το Μ. \monades{6}
    \item Έστω ότι μια ευθεία της παραπάνω οικογένειας τέμνει τους θετικούς ημιάξονες $Ox, O\gamma$ στα σημεία Α και Β αντίστοιχα.
    \begin{rlist}
        \item Να αποδείξετε ότι $0 < \alpha < 4$. \monades{6}
        \item Να βρείτε για ποια τιμή του $\alpha$ ισχύει $(OA) = 2(OB)$. \monades{5}
    \end{rlist}\end{alist}

\vspace{4mm}

% --- Άσκηση 16057 (16057.tex) ---
\Askhsh{BLYKALGT4001}{4}
Δίνονται τα σημεία $A(2,0)$, $B(3,4)$ και $\lambda \in \mathbb{R}$.
\begin{alist}    \item
    \begin{rlist}
        \item Να βρείτε την εξίσωση που περιγράφει όλες τις ευθείες που διέρχονται από το σημείο $A$ και έχουν κλίση $\lambda$.
        \item Να αποδείξετε οτι η ευθεία, η οποία διέρχεται από το σημείο $A$, έχει κλίση $\lambda$ και απέχει απόσταση ίση με $1$ από το σημείο $B$, έχει εξίσωση $(\varepsilon): 15x - 8y - 30 = 0$.  \monades{5}
 \monades{8}
    \end{rlist}
    \item Να αποδείξετε οτι υπάρχει και άλλη ευθεία $(\zeta)$, εκτός από την $(\varepsilon)$, η οποία διέρχεται από το σημείο $A$ και απέχει απόσταση ίση με $1$ από το σημείο $B$.
    \item Να βρείτε τις εξισώσεις των διχοτόμων των γωνιών που σχηματίζουν οι ευθείες $(\varepsilon)$ και $(\zeta)$.  \monades{5}
 \monades{7}\end{alist}

\vspace{4mm}

% --- Άσκηση 16128 (16128.tex) ---
\Askhsh{}{2}
Δίνεται η υπερβολή (C): \[ \frac{x^2}{16} - \frac{y^2}{9} = 1 \]
\begin{alist}    \item Να βρείτε τις συντεταγμένες των εστιών $E'$ και $E$. \monades{10}
    \item Αν το $N$ είναι τυχαίο σημείο της (C), να βρείτε την τιμή της διαφοράς $|(NE') - (NE)|$. \monades{5}
    \item Να σχεδιάσετε την υπερβολή (C). \monades{10}
    \begin{center}
    \begin{tikzpicture}    \begin{axis}[
 axis lines = middle, xlabel = {$x$}, ylabel = {$y$}, xmin = -8, xmax = 8, ymin = -6, ymax = 6, xtick = {-5,-4,-3,-2,-1,0,1,2,3,4,5}, ytick = {-5,-4,-3,-2,-1,0,1,2,3,4,5}, grid = both, grid style = {gray!30}, samples = 100, width = 10cm, height = 7cm, tick label style = {font=\scriptsize}, every axis x label/.style = {at={(current axis.right of origin)}, anchor=north west}, every axis y label/.style = {at={(current axis.above origin)}, anchor=south}, clip = false, ]
        % Asymptotes: y = +/- (b/a)x = +/- (3/4)x
 \addplot[dashed, thirdcolor, domain=-8:8] {3/4 * x}; \addplot[dashed, thirdcolor, domain=-8:8] {-3/4 * x};
        % Hyperbola branches: y = +/- (3/4) * sqrt(x^2 - 16)
        % Upper branch
 \addplot[very thick, maincolor, domain=4:8] {3/4 * sqrt(x^2 - 16)}; \addplot[very thick, maincolor, domain=-8:-4] {3/4 * sqrt(x^2 - 16)};
        % Lower branch
 \addplot[very thick, maincolor, domain=4:8] {-3/4 * sqrt(x^2 - 16)}; \addplot[very thick, maincolor, domain=-8:-4] {-3/4 * sqrt(x^2 - 16)};
        % Foci E'(-5,0) and E(5,0) (c^2 = a^2 + b^2 = 16 + 9 = 25 -> c=5)
 \addplot[only marks, mark=*, mark size=1.5pt, black] coordinates {(-5,0) (5,0)}; \node[below left] at (axis cs:-5,0) {$E'$}; \node[below right] at (axis cs:5,0) {$E$};
        % Vertices V1(-4,0) and V2(4,0) (a=4)
 \addplot[only marks, mark=*, mark size=1.5pt, black] coordinates {(-4,0) (4,0)}; \node[above] at (axis cs:-4,0) {$V1$}; \node[above] at (axis cs:4,0) {$V2$};
    \end{axis}\end{tikzpicture}
    \end{center}\end{alist}

\vspace{4mm}

% --- Άσκηση 16141 (16141.tex) ---
ΘΕΜΑ 2
Δίνεται ισόπλευρο τρίγωνο $AB\Gamma$ πλευράς $10$ και το μέσο $M$ της πλευράς $B\Gamma$.

\begin{alist}    \item Να βρεθούν τα μέτρα των γωνιών: \monades{10}
    \begin{rlist}
        \item $\left(\vec{AB}, \vec{A\Gamma}\right)$
        \item $\left(\vec{AM}, \vec{B\Gamma}\right)$
        \item $\left(\vec{AM}, \vec{\Gamma A}\right)$
        \item $\left(\vec{BM}, \vec{\Gamma M}\right)$
        \item $\left(\vec{\Gamma M}, \vec{\Gamma B}\right)$
    \end{rlist}
    \item Να υπολογιστούν τα εσωτερικά γινόμενα: \monades{15}
    \begin{rlist}
        \item $\vec{AM} \cdot \vec{B\Gamma}$
        \item $\vec{AM} \cdot \vec{\Gamma A}$
        \item $\vec{\Gamma M} \cdot \vec{\Gamma B}$
    \end{rlist}\end{alist}

\begin{center}
\begin{tikzpicture}    % Ορισμός κορυφών τριγώνου
    \tkzDefPoint(0,0){B}
    \tkzDefPoint(10,0){C}
    \tkzDefPointWith[euclidian](B,C){A}{sqrt(3)/2*10} % ίσοπλευρο τρίγωνο πλευράς 10
    \tkzGetPoint{A}
    % Ορισμός του μέσου Μ
    \tkzDefMidPoint(B,C) \tkzGetPoint{M}
    % Σχεδίαση του περιγράμματος
    \tkzDrawPolygon[l3](A,B,C)
    % Σχεδίαση του τμήματος AM
    \tkzDrawSegment[l3,dashed](A,M)
    % Σημεία
    \tkzDrawPoints(A,B,C,M)
    % Ετικέτες
    \tkzLabelPoint[above](A){A}
    \tkzLabelPoint[below left](B){B}
    \tkzLabelPoint[below](M){M}
    \tkzLabelPoint[below right](C){$\Gamma$}\end{tikzpicture}
\end{center}

\vspace{4mm}

% --- Άσκηση 16144 (16144.tex) ---
\OnPar{ΘΕΜΑ 2}
\vspace{4mm} % Add some vertical space to match the PDF layout

Δίνεται ρόμβος ΑΒΓΔ με κέντρο Ο, πλευρά $4$ και $\hat{A} = 60^\circ$. Να υπολογίσετε τα εσωτερικά γινόμενα : \monades{25}

\begin{alist}    \item $\vec{AB} \cdot \vec{A\Delta}$
    \item $\vec{A\Delta} \cdot \vec{B\Gamma}$
    \item $\vec{O\Delta} \cdot \vec{AO}$
    \item $\vec{O\Delta} \cdot \vec{OB}$
    \item $\vec{A\Delta} \cdot \vec{\Gamma\Delta}$\end{alist}

\begin{center}
\begin{tikzpicture}[scale=1.2]
    % Coordinates calculation based on problem text: Rhombus with side 4, angle A = 60 degrees.
    % If the full angle $\hat{A} = 60^\circ$, then the half angle $\angle OAD = 30^\circ$.
    % In the right-angled triangle AOD:
    % Hypotenuse AD = 4.
    % OD = AD \cdot \sin(\angle OAD) = 4 \cdot \sin(30^\circ) = 4 \cdot 0.5 = 2$.
    % AO = AD \cdot \cos(\angle OAD) = 4 \cdot \cos(30^\circ) = 4 \cdot \frac{\sqrt{3}}{2} = 2\sqrt{3}$.
    % Set O at the origin (0,0)
    \tkzDefPoint(0,0){O}
    % Define points D, B, A, C based on calculated lengths
    \tkzDefPoint(2,0){D} % D is at (OD, 0)
    \tkzDefPoint(-2,0){B} % B is at (-OD, 0)
    \tkzDefPoint(0, {2*sqrt(3)}){A} % A is at (0, AO) approx (0, 3.46)
    \tkzDefPoint(0, {-2*sqrt(3)}){C} % C is at (0, -AO) approx (0, -3.46)
    % Draw the rhombus ABCD
    \tkzDrawPolygon(A,B,C,D)
    % Draw diagonals AC and BD
    \tkzDrawSegment(A,C)
    \tkzDrawSegment(B,D)
    % Draw points A, B, C, D, O
    \tkzDrawPoints(A,B,C,D,O)
    % Labels for points
    \tkzLabelPoint[above](A){A}
    \tkzLabelPoint[left](B){B}
    \tkzLabelPoint[below](C){$\Gamma$}
    \tkzLabelPoint[above left, xshift=-2pt, yshift=2pt](O){O} % Adjust position of O label
    \tkzLabelPoint[right](D){$\Delta$}
    % Label side length '4' on segment AD, positioned as seen in the PDF.
    \node at ($(A)!0.6!(D)$) [right] {4}; % 60% of the way from A to D, text to the right
    % Highlight segment AO in 'maincolor' (green in the PDF image)
    \tkzDrawSegment[l3, maincolor](A,O) % l3 is 0.3mm linewidth from the template
    % Draw the angle arc for OAD and label it '60 degrees' in 'maincolor', as seen in the PDF image.
    % NOTE: The problem text states $\hat{A} = 60^\circ$, which implies $\angle OAD = 30^\circ$.
    % However, to faithfully reproduce the given PDF image, the label '60°' is used for $\angle OAD$.
    \tkzMarkAngle[fill=maincolor!10, size=0.5, arc=arc arrow](D,A,O)
    \tkzLabelAngle[pos=0.8, text=maincolor](D,A,O){$60^\circ$}\end{tikzpicture}
\end{center}

\vspace{4mm}

% --- Άσκηση 16147 (16147.tex) ---
\Askhsh{}{2}

Δίνονται τα διανύσματα:
\[
    \vec{a} = 3\vec{i} + 3\sqrt{3}\vec{j}, \quad \vec{b} = \sqrt{2}\vec{i}, \quad \vec{\gamma} = -3\vec{j} \quad \text{και} \quad \vec{\delta} = (-1, 1).
\]
\begin{alist}    \item Να βρείτε το συντελεστή διεύθυνσης καθενός από τα διανύσματα $\vec{a}$, $\vec{b}$ και $\vec{\delta}$. \monades{09}
    \item Να γράψετε τη γωνία που σχηματίζει καθένα από τα διανύσματα $\vec{a}$, $\vec{b}$, $\vec{\gamma}$ και $\vec{\delta}$ με το θετικό ημιάξονα $Ox$. \monades{10}
    \item Να βρείτε τα μέτρα των διανυσμάτων $\vec{a}$ και $\vec{\gamma}$. \monades{06}\end{alist}

\vspace{4mm}

% --- Άσκηση 16151 (16151.tex) ---
\Askhsh{}{2}
Δίνονται τα διανύσματα $\vec{\alpha} = (3,3)$ και $\vec{\beta} = (-\sqrt{3}, 1)$.
\begin{alist}    \item Να βρείτε το συντελεστή διεύθυνσής των διανυσμάτων $\vec{\alpha}$ και $\vec{\beta}$ καθώς και τη γωνία που σχηματίζει καθένα από αυτά με τον άξονα $\chi'\chi$. \monades{16}
    \item Να βρείτε τη γωνία $(\vec{\alpha},\vec{\beta})$. \monades{09}\end{alist}

\vspace{4mm}

% --- Άσκηση 16191 (16191.tex) ---
\Askhsh{}{4}
Δίνονται τα σημεία $A(1,1)$, $B(5,5)$.

\begin{alist}    \item Αν για το σημείο $M(x,\psi)$ ισχύει $\overrightarrow{AM}^2 + \overrightarrow{BM}^2 = 32$, να αποδείξετε ότι:
    \begin{rlist}
        \item Το σημείο $M$ βρίσκεται πάνω στην καμπύλη με εξίσωση $x^2 + \psi^2 - 6\psi - 6x + 10 = 0$ (1)
        \item Η εξίσωση (1) παριστάνει κύκλο.  \monades{8}
 \monades{3}
    \end{rlist}
    \item Αν το κέντρο του κύκλου είναι το $K(3,3)$ και η ακτίνα του $\rho = 2\sqrt{2}$:
    \begin{rlist}
        \item Να διερευνήσετε για ποιες τιμές του $\lambda$ η ευθεία $(\varepsilon)$: $\lambda x + \psi = 2$ εφάπτεται του κύκλου (1).
        \item Υπάρχει τιμή του $\lambda$ για την οποία η ευθεία $(\varepsilon)$ σχηματίζει με την $AB$ γωνία $45\textdegree$;  \monades{7}
 \monades{7}
    \end{rlist}\end{alist}

\vspace{4mm}

% --- Άσκηση 16194 (16194.tex) ---
\Askhseis{ΜΑΘΗΜΑΤΙΚΑ Β' ΛΥΚΕΙΟΥ}{2}

Δίνονται οι ευθείες $(\varepsilon1) : 8x + \psi - 28 = 0$, $(\varepsilon2) : x - \psi + 1 = 0$, $(\varepsilon3) : 3x + 4\psi + 5 = 0$.

\begin{alist}    \item Να βρείτε τις συντεταγμένες του σημείου τομής $M$ των $(\varepsilon1)$ και $(\varepsilon2)$.
    \item Αν το σημείο τομής είναι το $M(3,4)$ να υπολογίσετε:  \monades{09}
    \begin{rlist}
        \item Το μέτρο του διανύσματος $\overrightarrow{OM}$, όπου $O$ η αρχή των αξόνων.

        \item Την απόσταση του σημείου $M$ από την ευθεία $(\varepsilon3)$.  \monades{08}
 \monades{08}
    \end{rlist}\end{alist}

\vspace{4mm}

% --- Άσκηση 16425 (16425.tex) ---
\section*{ΘΕΜΑ 2}

Δίνονται οι ευθείες: $\varepsilon1: y = -\frac{2}{3}x + 1$ και $\varepsilon2: x = \frac{3}{2}y + 9$.

\begin{alist}    \item Να αποδείξετε ότι: $\varepsilon1 \parallel \varepsilon2$. \monades{12}
    \item Να υπολογίσετε την απόσταση των ευθειών $\varepsilon1$ και $\varepsilon2$. \monades{13}\end{alist}

\vspace{4mm}

% --- Άσκηση 16426 (16426.tex) ---
\OnPar{ΘΕΜΑ 2}

Δίνονται τα διανύσματα $\vec{\alpha} = (2, -1)$ και $\vec{\beta} = (-3, 2)$.
\begin{alist}    \item Να υπολογίσετε το γινόμενο $\vec{\alpha} \cdot (2\vec{\alpha} - \vec{\beta})$. \monades{10}
    \item Να βρείτε το διάνυσμα $\vec{\gamma} = (x, y)$ όταν $\vec{\gamma} \perp \vec{\alpha}$ και $|\vec{\gamma}| = \sqrt{5}$. \monades{15}\end{alist}

\vspace{4mm}

% --- Άσκηση 16427 (16427.tex) ---
\Askhsh{M2GEL24}{2}
Δίνονται τα σημεία $A(-2, 3)$, $B(0, 8)$, $\Gamma(5, 3)$ και $\Delta(10, 5)$. Να υπολογίσετε:
\begin{alist}    \item το εσωτερικό γινόμενο $\overrightarrow{AB} \cdot \overrightarrow{\Gamma\Delta}$. \monades{12}
    \item τη γωνία που σχηματίζει το διάνυσμα $\vec{u} = \overrightarrow{AB} + \overrightarrow{\Gamma\Delta}$ με τον άξονα $x'x$. \monades{13}\end{alist}

\vspace{4mm}

% --- Άσκηση 16428 (16428.tex) ---
\OnPar{ΘΕΜΑ 2}

Δίνονται τα διανύσματα $\vec{\alpha}$ και $\vec{\beta}$ με: $|\vec{\alpha}| = \frac{\sqrt{3}}{2}$, $|\vec{\beta}| = \frac{1}{2}$ και $|3\vec{\alpha} + 2\vec{\beta}| = |\vec{\alpha} - 2\vec{\beta}|$.

\begin{alist}    \item Να αποδείξετε ότι: $\vec{\alpha} \cdot \vec{\beta} = -\frac{3}{8}$. \monades{15}
    \item Να υπολογίσετε τη γωνία των διανυσμάτων $\vec{\alpha}$ και $\vec{\beta}$. \monades{10}\end{alist}

\vspace{4mm}

% --- Άσκηση 16477 (16477.tex) ---
\Askhsh{}{4}
Σε καρτεσιανό σύστημα συντεταγμένων $Oxy$, η εξίσωση ευθείας $l_\lambda: \lambda x+(1-\lambda)y+2=0$, όπου $\lambda$ αριθμός που μεταβάλλεται στο $\mathbb{R}$, παριστάνει τη φωτεινή ακτίνα που εκπέμπει ένας περιστρεφόμενος φάρος $\Phi$. Ακόμη δίνεται ότι ένα φορτηγό πλοίο είναι αγκυροβολημένο στο σημείο $O(0,0)$.

\begin{alist}    \item
    \begin{rlist}
        \item Να βρείτε τις συντεταγμένες του φάρου $\Phi$. \monades{10}
        \item Να εξετάσετε αν υπάρχει φωτεινή ακτίνα που εκπέμπεται από το φάρο προς το αγκυροβολημένο πλοίο. \monades{5}
    \end{rlist}
    \item Ένα ρυμουλκό πλοίο $P$ βρίσκεται βόρεια του φάρου $\Phi$. Η φωτεινή ακτίνα που φωτίζει το $P$ έχει εξίσωση $x+y+4=0$. Να βρείτε τις συντεταγμένες του σημείου $P$ όταν είναι γνωστό ότι η συντομότερη διαδρομή που πρέπει να διανύσει το ρυμουλκό πλοίο για να πάει προς το αγκυροβολημένο φορτηγό πλοίο είναι ίση με 4 μονάδες μήκους. \monades{10}\end{alist}

\begin{center}
\begin{tikzpicture}\begin{axis}[
    axis lines = middle,
    xmin = -7, xmax = 7,
    ymin = -5, ymax = 5,
    xtick = {-6,-4,-2,0,2,4,6},
    ytick = {-4,-2,0,2,4},
    major grid style = {draw=gray!30},
    grid = both,
    minor tick num = 1, 
    tick label style = {font=\scriptsize},
    width = 10cm, height = 10cm,
    axis equal, 
    x axis line style={maincolor, very thick},
    y axis line style={maincolor, very thick},
    xlabel={}, ylabel={}, 
]
    % Plot the line x+y+4=0, which is y = -x - 4
    % The line segment in the PDF image appears to go from x=-6 to x=0.
    \addplot[thick, secondarycolor, domain=-6:0] {-x - 4};
    % Compass-like element. Placed using axis cs to position it relative to the plot area.
    % Visual inspection suggests it's in the bottom right corner, slightly inside the bounds.
    \node[anchor=south east, xshift=-0.5cm, yshift=0.5cm] at (rel axis cs:1,0) { % rel axis cs for bottom right of plot area
        \begin{tikzpicture}[scale=0.35, thick, thirdcolor, >={Stealth[length=2mm, width=2mm]}]
            % Central point for arrows
            \coordinate (C) at (0,0);
            % North (B) - longer arrow, upright text
            \draw[->] (C) -- ++(0,1.8);
            \node[above=0.1em of C] {B};
            % East (A) - shorter arrow, upright text
            \draw[->] (C) -- ++(1.2,0);
            \node[right=0.1em of C] {A};
            % West (Z) - shorter arrow, upright text
            \draw[->] (C) -- ++(-1.2,0);
            \node[left=0.1em of C] {Z};
            % South - shorter arrow, no label
            \draw[->] (C) -- ++(0,-1.2);
            % Label 1 - positioned somewhere around south-west
            \node[font=\tiny] at (-0.5,-0.5) {1};\end{tikzpicture}
    };
\end{axis}
\end{tikzpicture}
\end{center}

\vspace{4mm}

% --- Άσκηση 16579 (16579.tex) ---
\Askhsh{}{2}
Δίνονται τα σημεία $A(2,1)$ και $B(6,7)$ του καρτεσιανού επιπέδου $Oxy$.

\begin{alist}    \item Να σχεδιάσετε το διάνυσμα $\vec{AB}$. \monades{07}
    \begin{center}
        \begin{tikzpicture} \tkzInit[xmin=-0.5,xmax=7.5,ymin=-0.5,ymax=8.5,xstep=1,ystep=1] \tkzGrid[color=gray!30] \tkzDrawXY[label style={font=\normalsize}] % Draw axes with labels x and y \tkzDefPoint(0,0){O} % Origin \tkzDefPoint(2,1){A} \tkzDefPoint(6,7){B} \tkzDrawPoints(O,A,B) % Draw the points \tkzLabelPoint[below left](O){$O$} \tkzLabelPoint[below left](A){$A(2,1)$} \tkzLabelPoint[above right](B){$B(6,7)$} \tkzDrawSegment[vector,maincolor,l4](A,B) % Draw the vector
\end{tikzpicture}
    \end{center}
    \item Αν $\vec{v} = \vec{AB}$ να βρείτε τις συντεταγμένες του διανύσματος $\vec{v}$. \monades{08}
    \item Να αποδείξετε ότι τα διανύσματα $\vec{u} = (-8, -12)$ και $\vec{v}$ του β) ερωτήματος είναι αντίρροπα. \monades{10}\end{alist}

\vspace{4mm}

% --- Άσκηση 16580 (16580.tex) ---
ΘΕΜΑ 2
Σε καρτεσιανό επίπεδο $Oxy$ δίνονται τα σημεία $A(2,4)$, $B(11,5)$, $\Gamma(3, 7)$ και ένα σημείο $\Delta$ ώστε το $\vec{A\Delta}$ να είναι ίσο με το άθροισμα των $\vec{AB}$ και $\vec{A\Gamma}$.

Να υπολογίσετε τις συντεταγμένες:
\begin{alist}    \item των διανυσμάτων $\vec{AB}$ και $\vec{A\Gamma}$. \monades{12}
    \item του διανύσματος $\vec{A\Delta}$. \monades{08}
    \item του σημείου $\Delta$. \monades{05}\end{alist}

\begin{center}
\begin{tikzpicture}[scale=1]
    % Ορισμός σημείων A, B, Γ, Δ (Το Δ ορίζεται από το άθροισμα διανυσμάτων)
    % Χρησιμοποιούμε γενικές συντεταγμένες για την αναπαράσταση του σχήματος.
    \tkzDefPoint(0,0){A}
    \tkzDefPoint(4,0.5){B} % Διάνυσμα AB
    \tkzDefPoint(1,3){C}  % Διάνυσμα AΓ (χρησιμοποιούμε C για το Γ)
    % Ορισμός του Δ ώστε να ισχύει ΑΔ = ΑΒ + ΑΓ
    % Αυτό σημαίνει ότι A, B, Δ, Γ αποτελούν παραλληλόγραμμο με διαγώνιο ΑΔ.
    \tkzDefPointWith[vector from=A to B, vector from=A to C]{D}
    \tkzGetPoint{D}
    % Σχεδίαση των τμημάτων
    \tkzDrawSegment[l3](A,B) % Τμήμα AB
    \tkzDrawSegment[l3](A,C) % Τμήμα AΓ
    \tkzDrawSegment[l3,maincolor](A,D) % Τμήμα AΔ (η διαγώνιος, με χρώμα)
    \tkzDrawSegment[dashed,l3](B,D) % Τμήμα BΔ (διακεκομμένη γραμμή)
    \tkzDrawSegment[dashed,l3](C,D) % Τμήμα ΓΔ (διακεκομμένη γραμμή)
    % Σχεδίαση των σημείων
    \tkzDrawPoints(A,B,C,D)
    % Ετικέτες των σημείων
    \tkzLabelPoint[below left](A){A}
    \tkzLabelPoint[below right](B){B}
    \tkzLabelPoint[above left](C){$\Gamma$}
    \tkzLabelPoint[above right](D){$\Delta$}\end{tikzpicture}
\end{center}

\vspace{4mm}

% --- Άσκηση 16581 (16581.tex) ---
\Askhsh{}{2}
Σε καρτεσιανό επίπεδο $Oxy$ δίνονται τα σημεία $A(-1, 6)$, $B(1,2)$ και $\Gamma(3, -2)$.
\begin{alist}    \item Να υπολογίσετε τις συντεταγμένες των διανυσμάτων $\vec{AB}$ και $\vec{B\Gamma}$. \monades{12}
    \item Να αποδείξετε ότι τα σημεία $A$, $B$ και $\Gamma$ είναι συνευθειακά. \monades{06}
    \item Να αποδείξετε ότι το $B$ είναι μέσο του ευθύγραμμου τμήματος $A\Gamma$. \monades{07}\end{alist}

\vspace{4mm}

% --- Άσκηση 16759 (16759.tex) ---
\Askhsh{TH2}{2}
Δίνονται οι ευθείες $(\varepsilon1)$, $(\varepsilon2)$ και $(\varepsilon3)$ με εξισώσεις $x - 2y = -1$, $2x + y = 4$ και $y = -1$ αντίστοιχα.

\begin{alist}    \item Να αποδείξετε ότι οι ευθείες $(\varepsilon1)$ και $(\varepsilon2)$ είναι κάθετες.
    \item Να αποδείξετε ότι οι ευθείες $(\varepsilon1)$ και $(\varepsilon2)$ τέμνονται στο σημείο $A\left(\frac{7}{5}, \frac{6}{5}\right)$.  \monades{8}
    \item Να υπολογίσετε την απόσταση του σημείου $A$ από την ευθεία $(\varepsilon3)$.  \monades{9}
 \monades{8}\end{alist}

\vspace{4mm}

% --- Άσκηση 16766 (16766.tex) ---
\Askhsh{}{2}
Δίνονται οι ευθείες $(\varepsilon1)$ και $(\varepsilon2)$ με εξισώσεις $x - 3y = 4$ και $9x + 3y = 6$ αντίστοιχα.

\begin{alist}    \item Να αποδείξετε ότι οι ευθείες $(\varepsilon1)$ και $(\varepsilon2)$ είναι κάθετες.
    \item Να αποδείξετε ότι οι ευθείες $(\varepsilon1)$ και $(\varepsilon2)$ τέμνονται στο σημείο A$(1, -1)$.  \monades{8}
    \item Να βρεθεί η εξίσωση της ευθείας η οποία διέρχεται από το σημείο A και είναι κάθετη στον άξονα $x'x$.  \monades{8}
 \monades{9}\end{alist}

\vspace{4mm}

% --- Άσκηση 16769 (16769.tex) ---
\Askhsh{}{2}
Δίνεται τρίγωνο ΑΒΓ με κορυφές $A(1,7)$, $B(-1,5)$ και $\Gamma(3,3)$.
\begin{alist}    \item Να υπολογίσετε το εμβαδόν του τριγώνου ΑΒΓ. \monades{09}
    \item Αν Μ είναι το μέσο της πλευράς ΒΓ, τότε να υπολογίσετε:
    \begin{rlist}
        \item Τις συντεταγμένες του Μ.
        \item Την εξίσωση της διαμέσου ΑΜ.
    \end{rlist}
 \monades{16}\end{alist}

\vspace{4mm}

% --- Άσκηση 16771 (16771.tex) ---
\Askhsh{}{2}
Δίνονται τα σημεία $A(2,1)$, $\Gamma(4,-1)$ και το διάνυσμα $\vec{AB} = (3,-1)$.
\begin{alist}    \item Να βρεθεί το σημείο Β. \monades{09}
    \item Αν $B(5,0)$:
    \begin{rlist}
        \item Να δείξετε ότι τα σημεία Α, Β και Γ σχηματίζουν τρίγωνο. \monades{08}
        \item Να υπολογίσετε το εμβαδόν του τριγώνου ΑΒΓ. \monades{08}
    \end{rlist}\end{alist}

\vspace{4mm}

% --- Άσκηση 16773 (16773.tex) ---
\Askhsh{}{2} % Askhsh{Κωδικός}{ΘΕΜΑ}
\begin{alist}    \item Να βρεθεί η εξίσωση του κύκλου που έχει κέντρο το $O(0,0)$ και διέρχεται από το σημείο $A(1,2)$. \monades{08}
    \item Δίνεται ο κύκλος $x^2 + y^2 = 5$.
    \begin{rlist}
        \item Να βρεθεί η εξίσωση της εφαπτομένης του στο σημείο $A$. \monades{09}
        \item Να βρεθεί το σημείο $B$, το οποίο είναι αντιδιαμετρικό του $A$ σε αυτόν τον κύκλο. \monades{08}
    \end{rlist}\end{alist}

\vspace{4mm}

% --- Άσκηση 16774 (16774.tex) ---
\Askhsh{}{2}
Δίνεται τρίγωνο ΑΒΓ με κορυφές τα σημεία $A(2,5)$, $B(3,6)$ και $\Gamma(-1,-2)$.

\begin{alist}    \item Να βρείτε τον συντελεστή διεύθυνσης της ευθείας ΒΓ. \monades{07}
    \item Να βρείτε την εξίσωση του ύψους που άγεται από το Α. \monades{09}
    \item Να βρείτε την γωνία που σχηματίζει η ευθεία ΑΒ με τον άξονα $x'x$. \monades{09}\end{alist}

\vspace{4mm}

% --- Άσκηση 16808 (16808.tex) ---
\Askhsh{}{2} % No code provided, Theme 2
Δίνονται τα σημεία του επιπέδου $A(-8, 1)$, $B(4, 5)$ και $\Gamma(-4, 9)$.
\begin{alist}    \item Να βρείτε τις συντεταγμένες του μέσου $K$ του ευθυγράμμου τμήματος $AB$.
    \item Να δείξετε ότι ο κύκλος $(C)$ που έχει κέντρο το σημείο $K$ και διάμετρο το τμήμα $AB$ διέρχεται από το σημείο $\Gamma$.  \monades{8}
    \item Να βρείτε την εξίσωση του κύκλου $(C)$.  \monades{9}
 \monades{8}\end{alist}

\vspace{4mm}

% --- Άσκηση 16810 (16810.tex) ---
\Askhsh{K0123}{2}
Στο ορθοκανονικό σύστημα συντεταγμένων θεωρούμε τα σημεία $A(1,1)$, $B(5,2)$, $\Gamma(0,-2)$ και $\Delta(8,0)$.

\begin{alist}    \item Να τοποθετήσετε τα παραπάνω σημεία του επιπέδου σε ένα πρόχειρο σχήμα και να αποδείξετε ότι το τετράπλευρο με κορυφές τα σημεία αυτά είναι τραπέζιο. \monades{10}
    \begin{center}
        \begin{tikzpicture}[scale=0.8]
            % Initialize coordinate system
 \tkzInit[xmin=-1,xmax=9,ymin=-3,ymax=3,xstep=1,ystep=1] \tkzGrid[sub,color=gray!30] \draw[l3,->] (-5,0) -- (5,0) node[below right] {$x$}; \draw[l3,->] (0,-5) -- (0,5) node[above left] {$y$}; \tkzLabelX[font=\scriptsize] \tkzLabelY[font=\scriptsize]
            % Define points
 \tkzDefPoint(1,1){A} \tkzDefPoint(5,2){B} \tkzDefPoint(0,-2){C} \tkzDefPoint(8,0){D}
            % Draw polygon (trapezoid) A-B-D-C, as AB and CD are parallel
 \tkzDrawPolygon[l3,maincolor](A,B,D,C)
            % Draw points
 \tkzDrawPoints[color=secondarycolor](A,B,C,D)
            % Label points
 \tkzLabelPoint[above left](A){$A$} \tkzLabelPoint[above right](B){$B$} \tkzLabelPoint[below left](C){$\Gamma$} \tkzLabelPoint[below right](D){$\Delta$}\end{tikzpicture}
    \end{center}
    \item Να υπολογίσετε το εμβαδόν του τραπεζίου του ερωτήματος α). \monades{15}\end{alist}

\vspace{4mm}

% --- Άσκηση 17070 (17070.tex) ---
\OnPar{ΘΕΜΑ 2}

Στο καρτεσιανό επίπεδο $Oxy$ δίνονται τα σημεία $A(3,4)$, $B(2,1)$, $\Gamma(3, -1)$ και $\Delta(4,2)$.

\begin{alist}    \item Να σχεδιάσετε τα παραπάνω σημεία $A$, $B$, $\Gamma$ και $\Delta$. \monades{8}
    \begin{center}
    \begin{tikzpicture}    \begin{axis}[
 axis lines=middle, xlabel={$x$}, ylabel={$y$}, xmin=-1, xmax=5, ymin=-2, ymax=5, xtick={-1,0,1,2,3,4,5}, ytick={-2,-1,0,1,2,3,4,5}, grid=both, grid style={gray!30}, width=8cm, height=8cm, tick label style={font=\scriptsize}, every axis x label/.style={at={(current axis.right of origin)}, anchor=north west}, every axis y label/.style={at={(current axis.above origin)}, anchor=south}, clip=false, ]
        % Plot the points
 \addplot[only marks, mark=*, mark size=2.5pt, black] coordinates { (3,4) % A (2,1) % B (3,-1) % Gamma (4,2) % Delta };
        % Labels for the points
 \node[above right, font=\small] at (axis cs:3,4) {$A(3,4)$}; \node[below left, font=\small] at (axis cs:2,1) {$B(2,1)$}; \node[below right, font=\small] at (axis cs:3,-1) {$\Gamma(3,-1)$}; \node[above left, font=\small] at (axis cs:4,2) {$\Delta(4,2)$};
    \end{axis}\end{tikzpicture}
    \end{center}
    \item Να βρείτε τις συντεταγμένες των διανυσμάτων $\vec{AB}$ και $\vec{\Delta\Gamma}$. \monades{9}
    \item Να αποδείξετε ότι το τετράπλευρο $AB\Gamma\Delta$ είναι παραλληλόγραμμο. \monades{8}\end{alist}

\vspace{4mm}

% --- Άσκηση 17075 (17075.tex) ---
\Askhsh{}{2}

Στο παρακάτω σχήμα δίνονται τα διανύσματα $\vec{AB}$ και $\vec{A\Gamma}$ του καρτεσιανού επιπέδου $Oxy$.

\begin{alist}    \item Να αποδείξετε ότι $\vec{AB} = (4,-3)$ και $\vec{A\Gamma} = (0,-2)$. \monades{12}
    \item Να υπολογίσετε το εσωτερικό γινόμενο των διανυσμάτων $\vec{AB}$ και $\vec{A\Gamma}$. \monades{13}\end{alist}

\begin{tikzpicture}\begin{axis}[
    axis lines = middle,
    xlabel = {$x$},
    ylabel = {$y$},
    xmin = -5.5, xmax = 8.5,
    ymin = -2.5, ymax = 4.5,
    xtick = {-5,-4,-3,-2,-1,0,1,2,3,4,5,6,7,8},
    ytick = {-2,-1,0,1,2,3,4},
    grid = both,
    grid style = {gray!30},
    samples = 100,
    width = 10cm, height = 7cm,
    tick label style = {font=\scriptsize},
    every axis x label/.style = {at={(current axis.right of origin)}, anchor=north west},
    every axis y label/.style = {at={(current axis.above origin)}, anchor=south},
    clip = false, 
]
    % Define points
    \coordinate (O) at (0,0);
    \coordinate (A) at (2,3);
    \coordinate (B) at (6,0);
    % The vector AΓ=(0,-2) starting from A=(2,3) ends at (2+0, 3-2) = (2,1)
    \coordinate (Gammaend) at (2,1);
    % Draw vectors
    \draw[->, l3, maincolor] (A) -- (B);
    \draw[->, l3, secondarycolor] (A) -- (Gammaend);
    % Labels for points
    \node[above left] at (A) {A};
    \node[below right] at (B) {B};
    \node[below left] at (O) {O};
    % Label the endpoint of the vector AΓ as Γ
    \node[below left] at (Gammaend) {$\Gamma$}; 
\end{axis}\end{tikzpicture}

\vspace{4mm}

% --- Άσκηση 17076 (17076.tex) ---
\Askhsh{}{4}

Δίνονται τα σημεία $A(-3, -1)$, $B(0, 3)$ και $M(x, y)$ του καρτεσιανού επιπέδου $Oxy$.

\begin{alist}    \item Να βρείτε τις συντεταγμένες των διανυσμάτων $\overrightarrow{AM}$, $\overrightarrow{MB}$ και $\overrightarrow{AB}$. \monades{6}
    \item Να βρείτε τα μέτρα των διανυσμάτων $|\overrightarrow{AM}|$, $|\overrightarrow{MB}|$ και $|\overrightarrow{AB}|$. \monades{6}
    \item Να αποδείξετε ότι $|\overrightarrow{AM}| + |\overrightarrow{MB}| \ge 5$. \monades{6}
    \item Θεωρήστε τον ισχυρισμό: «Υπάρχει ζεύγος πραγματικών αριθμών $(x, y)$ τέτοιο ώστε να ισχύει $\sqrt{(x + 3)^2+(y + 1)^2} + \sqrt{x^2+(y - 3)^2} = 4$.»\\
 Είναι αληθής ή ψευδής ο παραπάνω ισχυρισμός; Να αιτιολογήσετε την απάντησή σας. \monades{7}\end{alist}

\vspace{4mm}

% --- Άσκηση 17077 (17077.tex) ---
\noindent{\titlefont\bfseries\Large ΘΕΜΑ 4}

Στο καρτεσιανό επίπεδο $Oxy$ τα σημεία $A$ και $B$ έχουν διανύσματα θέσεως
\[ \vec{OA} = 2\vec{i} + \lambda\vec{j} \quad \text{και} \quad \vec{OB} = (\lambda + 1)\vec{i} + (\lambda + 3)\vec{j}, \quad \text{με} \quad \lambda \in \mathbb{R}. \]

\begin{alist}    \item Να αποδείξετε ότι $\vec{AB} = (\lambda - 1)\vec{i} + 3\vec{j}$. \monades{6}
    \item Να βρείτε την απόσταση των σημείων $A$ και $B$ ως συνάρτηση του $\lambda$. \monades{7}
    \item Για ποιες τιμές του $\lambda$ η απόσταση των σημείων $A$ και $B$ είναι ίση με $5$; \monades{7}
    \item Θεωρήστε τον ισχυρισμό: «Υπάρχει πραγματικός αριθμός $\lambda$ τέτοιος ώστε η απόσταση των σημείων $A$ και $B$ να παίρνει τη μικρότερη δυνατή τιμή.»\\
 Είναι αληθής ή ψευδής ο παραπάνω ισχυρισμός; Να αιτιολογήσετε την απάντησή σας. \monades{5}\end{alist}

\vspace{4mm}

% --- Άσκηση 17078 (17078.tex) ---
\paragraph{ΘΕΜΑ 4}
Δίνονται τα σημεία $A(3, 2\alpha)$, $B(4,\alpha)$, $\Gamma(\alpha + 1, 1 - \alpha)$ και $\Delta(\alpha, 1)$, με $\alpha \in \mathbb{R}$.

\begin{alist}
    \item Να αποδείξετε ότι:
    \begin{rlist}        \item Η ευθεία που διέρχεται από τα σημεία $A$ και $B$ έχει εξίσωση $y = -ax + 5a$.
        \item Τα σημεία $\Gamma$ και $\Delta$ ανήκουν στην ευθεία $AB$ αν και μόνο αν $\alpha = \frac{5\pm\sqrt{21}}{2}$.  \monades{6}
        \item Το τετράπλευρο $AB\Gamma\Delta$ είναι παραλληλόγραμμο όταν $\alpha \ne \frac{5\pm\sqrt{21}}{2}$.  \monades{7}
 \monades{7}\end{rlist}
    \item Θεωρήστε τον ισχυρισμό: «Υπάρχει πραγματικός αριθμός $\alpha$ ώστε το τετράπλευρο $AB\Gamma\Delta$ να είναι τετράγωνο.» Είναι αληθής ή ψευδής ο παραπάνω ισχυρισμός; Να αιτιολογήσετε την απάντησή σας.
    \monades{5}
\end{alist}

\vspace{4mm}

% --- Άσκηση 17317 (17317.tex) ---
\Askhseis{ΜΑΘ3-25501}\Askhsh{ΜΑΘ325501}{2}
Δίνεται ο κύκλος $C: (x - 1)^2 + (y - 2)^2 = 4$ και η ευθεία $\varepsilon: 3x - 4y = 8$.

\begin{alist}    \item Να βρείτε το κέντρο $K$ του κύκλου $C$ και την ακτίνα του.
    \item Αν $K(1,2)$, να δείξετε ότι η απόσταση του κέντρου του κύκλου $C$ από την ευθεία $\varepsilon$ είναι $d(K, \varepsilon) = \frac{13}{5}$.  \monades{5}
    \item Να αιτιολογήσετε γιατί η ευθεία και ο κύκλος δεν έχουν κανένα κοινό σημείο.  \monades{13}
 \monades{7}\end{alist}

\vspace{4mm}

% --- Άσκηση 17694 (17694.tex) ---
\Askhseis{ΘΕΜΑ 4}

Στο χάρτη μίας πεδινής περιοχής, που είναι εφοδιασμένος με ορθοκανονικό σύστημα συντεταγμένων, δύο κωμοπόλεις Α και Β έχουν συντεταγμένες $A(3,6)$ και $B(7,-2)$.

\begin{alist}    \item Ανάμεσα στις δύο κωμοπόλεις, θα κατασκευαστεί ευθεία σιδηροδρομική γραμμή, κάθε σημείο της οποίας θα ισαπέχει από αυτές. Να βρείτε την εξίσωση της ευθείας, πάνω στην οποία βρίσκεται η σιδηροδρομική γραμμή. \monades{12}
    \item Πάνω στην σιδηροδρομική γραμμή θα κατασκευαστεί σταθμός $\Sigma$, ώστε το εμβαδόν της περιοχής που ορίζεται από τα σημεία $A$, $B$ και $\Sigma$ να ισούται με $20$ τετραγωνικές μονάδες. Να βρείτε τις συντεταγμένες του σταθμού $\Sigma$ στο χάρτη. \monades{13}\end{alist}

\vspace{4mm}

% --- Άσκηση 17695 (17695.tex) ---
\Askhsh{}{4}
Υποθέτουμε, ότι σε ένα επίπεδο που έχουμε εφοδιάσει με ορθοκανονικό σύστημα συντεταγμένων, κινούνται δύο σημεία $A$ και $B$. Κάθε χρονική στιγμή $t$ με $t \ge 0$ η θέση του πρώτου σημείου είναι $A(t-1, 2t-1)$ και του δευτέρου $B(3t-1, -4t-1)$.

\begin{alist}    \item Να βρείτε τις εξισώσεις των γραμμών πάνω στις οποίες κινούνται τα δύο σημεία. \monades{8}
    \item Υπάρχει χρονική στιγμή κατά την οποία τα δύο σημεία ταυτίζονται; \monades{7}
    \item Να υπολογιστεί η απόσταση των δύο σημείων την χρονική στιγμή $t=2$. \monades{5}
    \item Να βρεθεί η χρονική στιγμή $t$ κατά την οποία η απόσταση του σημείου $A$ από την ευθεία $\varepsilon: 4x+3y+7=0$ ισούται με $6$. \monades{5}\end{alist}

\vspace{4mm}

% --- Άσκηση 17805 (17805.tex) ---
\Askhsh{}{2}

Δίνεται το τρίγωνο ΑΟΒ με $A(3,4)$, $B(7,1)$, Ο η αρχή των αξόνων και το σημείο
$\Delta\left(\frac{12}{5}, \frac{16}{5}\right)$ της πλευράς $AO$.

\begin{center}
\begin{tikzpicture}\begin{axis}[
    axis lines = middle,
    xlabel = {},
    ylabel = {},
    xmin = -2.5, xmax = 9.5,
    ymin = -1.5, ymax = 5.5,
    xtick = {-2,-1,0,1,2,3,4,5,6,7,8,9},
    ytick = {-1,0,1,2,3,4,5},
    grid = both,
    grid style = {gray!30, dotted},
    width = 10cm, height = 7cm,
    tick label style = {font=\scriptsize},
    every axis x label/.style = {at={(current axis.right of origin)}, anchor=north west},
    every axis y label/.style = {at={(current axis.above origin)}, anchor=south},
    clip = false,
]
    % Plot points
    \addplot[only marks, mark=*, mark size=1.5pt, maincolor] coordinates {
        (0,0) % O
        (3,4) % A
        (7,1) % B
        (2.4,3.2) % D (12/5, 16/5)
    };
    % Draw segments
    \addplot[very thick, maincolor] coordinates {(0,0) (3,4)}; % OA
    \addplot[very thick, maincolor] coordinates {(3,4) (7,1)}; % AB
    \addplot[very thick, maincolor] coordinates {(7,1) (0,0)}; % OB
    \addplot[very thick, maincolor] coordinates {(2.4,3.2) (7,1)}; % DB
    % Labels for points
    \node[above right, font=\small, maincolor] at (axis cs:3,4) {A};
    \node[below right, font=\small, maincolor] at (axis cs:7,1) {B};
    \node[above left, font=\small, maincolor] at (axis cs:2.4,3.2) {$\Delta$};
\end{axis}\end{tikzpicture}
\end{center}

\begin{alist}    \item Να βρείτε τις συντεταγμένες των διανυσμάτων $\vec{OA}$ και $\vec{\Delta A}$. \monades{7}
    \item Να δείξετε ότι $\vec{\Delta A} = -\frac{1}{5}\vec{OA}$. \monades{9}
    \item Δίνεται ότι $(\text{OAB}) = \frac{25}{2}$ τετραγωνικές μονάδες. Να δείξετε ότι $(\text{A}\Delta\text{B}) = \frac{1}{5}(\text{OAB})$. \monades{9}\end{alist}

\vspace{4mm}

% --- Άσκηση 17942 (17942.tex) ---
\OnPar{ΘΕΜΑ 2}

Δίνεται η κωνική τομή με εξίσωση (C)
\[ \frac{x^2}{4} - \frac{y^2}{9}=1 \]

\begin{alist}    \item Να προσδιορίσετε το είδος της κωνικής τομής και να βρείτε μία εστία της. \monades{12}
    \item Να εξετάσετε αν το σημείο $M(1,2022)$ μπορεί να ανήκει στην (C) . \monades{13}\end{alist}

\vspace{4mm}

% --- Άσκηση 17944 (17944.tex) ---
\Askhsh{HYPB201}{3}
Δίνεται η υπερβολή με εξίσωση της μορφής $(C)$: $\frac{x^2}{\alpha^2} - \frac{y^2}{\beta^2} = 1$, εστιακή απόσταση $EE' = 2\sqrt{7}$ και εκκεντρότητα $\varepsilon = \frac{\sqrt{7}}{2}$.

\begin{alist}    \item Να αποδείξετε ότι $\alpha=2$, $\beta=\sqrt{3}$. \monades{8}
    \item \begin{rlist}
        \item Να βρείτε τις συντεταγμένες των κορυφών $A, A'$ της υπερβολής $(C)$.
        \item Να βρείτε τις εξισώσεις των ασύμπτωτων ευθειών της υπερβολής $(C)$.
    \end{rlist} \monades{8}
    \item Να σχεδιάσετε στο ίδιο σύστημα αξόνων την υπερβολή $(C)$, τις ασύμπτωτές της, τις εστίες της και τις κορυφές της. \monades{9}\end{alist}

\begin{center}
    \begin{tikzpicture}    \begin{axis}[
        axis lines = middle,
        xlabel = {$x$},
        ylabel = {$y$},
        xmin = -4.5, xmax = 4.5,
        ymin = -3.5, ymax = 3.5,
        xtick = {-4,-3,-2,-1,0,1,2,3,4},
        ytick = {-3,-2,-1,0,1,2,3},
        grid = both,
        grid style = {gray!30},
        samples = 200, 
        width = 10cm, height = 7cm,
        tick label style = {font=\scriptsize},
        every axis x label/.style = {at={(current axis.right of origin)}, anchor=north west},
        every axis y label/.style = {at={(current axis.above origin)}, anchor=south east},
    ]
        % Υπερβολή: y = +/- sqrt(3*(x^2/4 - 1))
        \addplot[very thick, maincolor, smooth, domain=-4.5:-2] {sqrt(3*(x^2/4 - 1))};
        \addplot[very thick, maincolor, smooth, domain=2:4.5] {sqrt(3*(x^2/4 - 1))};
        \addplot[very thick, maincolor, smooth, domain=-4.5:-2] {-sqrt(3*(x^2/4 - 1))};
        \addplot[very thick, maincolor, smooth, domain=2:4.5] {-sqrt(3*(x^2/4 - 1))};
        % Ασύμπτωτες: y = +/- (sqrt(3)/2)*x
        \addplot[dashed, secondarycolor, domain=-4.5:4.5] {sqrt(3)/2 * x};
        \addplot[dashed, secondarycolor, domain=-4.5:4.5] {-sqrt(3)/2 * x};
        % Κορυφές A(2,0) και A'(-2,0)
        \addplot[only marks, mark=*, mark size=1.5pt, black] coordinates {(2,0) (-2,0)};
        \node[above right, font=\tiny] at (axis cs:2,0) {$A(2,0)$};
        \node[above left, font=\tiny] at (axis cs:-2,0) {$A'(-2,0)$};
        % Εστίες E(sqrt(7),0) και E'(-sqrt(7),0)
        % (sqrt(7) περίπου 2.646)
        \addplot[only marks, mark=*, mark size=1.5pt, black] coordinates {(sqrt(7),0) (-sqrt(7),0)};
        \node[below right, font=\tiny] at (axis cs:sqrt(7),0) {$E(\sqrt{7},0)$};
        \node[below left, font=\tiny] at (axis cs:-sqrt(7),0) {$E'(-\sqrt{7},0)$};
        % Σημείο αρχής των αξόνων
        \node[below left, font=\tiny] at (axis cs:0,0) {$O$};
    \end{axis}\end{tikzpicture}
\end{center}

\vspace{4mm}

% --- Άσκηση 18236 (18236.tex) ---
\Askhsh{}{2}
Σε τρίγωνο ΑΒΓ είναι $A(-1, 5)$ και $B(2, 1)$. Αν οι πλευρές ΑΓ και ΒΓ βρίσκονται πάνω στις ευθείες $\varepsilon1 : y = -x + 4$ και $\varepsilon2 : y = -\frac{1}{2}x+2$ αντίστοιχα, τότε:

\begin{alist}    \item Να αποδείξετε ότι $\Gamma(4, 0)$.
    \item Να βρείτε:  \monades{12}
    \begin{rlist}
        \item το συντελεστή διεύθυνσης της ευθείας ΑΓ
        \item την εξίσωση του ύψους ΒΔ.  \monades{6}
 \monades{7}
    \end{rlist}\end{alist}

\vspace{4mm}

% --- Άσκηση 18237 (18237.tex) ---
\Askhsh{}{4}
Θεωρούμε τα σημεία $A(-1, 2)$, $B(3, 2)$, $\Gamma(1, 4)$.
\begin{alist}    \item Να αποδείξετε ότι σχηματίζουν τρίγωνο. \monades{6}
    \item Να βρείτε την εξίσωση της μεσοκάθετης της πλευράς $B\Gamma$. \monades{7}\end{alist}
Έστω ότι η μεσοκάθετη της πλευράς $B\Gamma$ είναι η ευθεία $\varepsilon: y=x+1$.
\begin{alist}    \item Να βρείτε σημείο $K$ στην μεσοκάθετη της πλευράς $B\Gamma$ που ισαπέχει από τα $A$, $B$. \monades{7}  [start=3]
    \item Να βρείτε την εξίσωση του περιγεγραμμένου κύκλου του τριγώνου $AB\Gamma$. \monades{5}\end{alist}

\vspace{4mm}

% --- Άσκηση 18238 (18238.tex) ---
\Askhsh{}{2}
Δίνονται τα σημεία $A(1,3)$ και $B(-3,5)$.
\begin{alist}    \item Να βρείτε τις συντεταγμένες του μέσου $K$ του τμήματος $AB$. \monades{7}
    \item Να αποδείξετε ότι $(KA) = \sqrt{5}$. \monades{8}
    \item Να βρείτε την εξίσωση του κύκλου που έχει διάμετρο το ευθύγραμμο τμήμα $AB$. \monades{10}\end{alist}

\vspace{4mm}

% --- Άσκηση 18239 (18239.tex) ---
\Askhsh{}{2}
Δίνεται το σημείο $K(-3,1)$ και η ευθεία $(\varepsilon): 4x-3y+5=0$.
\begin{alist}    \item Να αποδείξετε ότι η απόσταση του σημείου $K$ από την ευθεία $(\varepsilon)$ είναι ίση με $2$. \monades{6}
    \item Να βρείτε την εξίσωση του κύκλου $C$ που έχει κέντρο το σημείο $K$ και εφάπτεται στην ευθεία $(\varepsilon)$. \monades{9}
    \item Να σχεδιάσετε στο ίδιο ορθοκανονικό σύστημα αξόνων τον κύκλο $C$ και την ευθεία $(\varepsilon)$. \monades{10}
    \begin{center}
        \begin{tikzpicture}            \begin{axis}[
 axis lines=middle, xlabel=$x$, ylabel=$y$, xmin=-7, xmax=3, ymin=-4, ymax=4, xtick distance=1, ytick distance=1, grid=both, grid style={gray!30}, width=10cm, height=7cm, tick label style={font=\scriptsize}, every axis x label/.style={at={(current axis.right of origin)}, anchor=north west}, every axis y label/.style={at={(current axis.above origin)}, anchor=south}, set layers, ]
                % Σχεδίαση της ευθείας (ε): 4x-3y+5=0  =>  3y = 4x+5  =>  y = (4x+5)/3
 \addplot[very thick, maincolor, domain=-7:3, samples=100] { (4*x + 5)/3 }; \node[maincolor, above left] at (axis cs:0, 5/3) {$(\varepsilon): 4x-3y+5=0$};
                % Σχεδίαση του κύκλου C: κέντρο K(-3,1) και ακτίνα r=2 (από το α) ερώτημα)
                % Εξίσωση: (x+3)^2 + (y-1)^2 = 2^2 = 4
 \addplot[very thick, secondarycolor, domain=0:360, samples=100, smooth, variable=t] ({ -3 + 2*cos(t) }, { 1 + 2*sin(t) }); \node[secondarycolor, right] at (axis cs:-3+2*cos(45), 1+2*sin(45)) {$C: (x+3)^2+(y-1)^2=4$};
                % Σημείο κέντρου K(-3,1)
 \addplot[only marks, mark=*, mark size=1.5pt, black] coordinates {(-3,1)}; \node[below left] at (axis cs:-3,1) {$K(-3,1)$};
                % Σημείο επαφής T(-1.4,-0.2)
 \addplot[only marks, mark=*, mark size=1.5pt, black] coordinates {(-1.4,-0.2)}; \node[above right] at (axis cs:-1.4,-0.2) {$T$};
            \end{axis}\end{tikzpicture}
    \end{center}\end{alist}

\vspace{4mm}

% --- Άσκηση 18240 (18240.tex) ---
\section*{ΘΕΜΑ 2}

Δίνεται το σημείο $A(1,2)$ και η ευθεία $(\varepsilon): y=x+3$.

\begin{alist}    \item Να βρείτε την απόσταση του σημείου $A$ από την ευθεία $(\varepsilon)$. \monades{7}
    \item Να βρείτε την εξίσωση της ευθείας $(\eta)$ που διέρχεται από το $A$ και είναι παράλληλη στην $(\varepsilon)$. \monades{8}
    \item Να σχεδιάσετε στο ίδιο ορθοκανονικό σύστημα συντεταγμένων τις ευθείες $(\eta), (\varepsilon)$. \monades{10}\end{alist}

\begin{center}
\begin{tikzpicture}\begin{axis}[
    axis lines = middle,
    xlabel = {$x$},
    ylabel = {$y$},
    xmin = -4, xmax = 4,
    ymin = -3, ymax = 7, 
    xtick = {-3,-2,-1,0,1,2,3},
    ytick = {-2,-1,0,1,2,3,4,5,6}, 
    grid = both,
    grid style = {gray!30},
    samples = 100,
    width = 10cm, height = 7cm,
    tick label style = {font=\scriptsize},
    every axis x label/.style = {at={(current axis.right of origin)}, anchor=north west},
    every axis y label/.style = {at={(current axis.above origin)}, anchor=south},
    clip = false, 
]
    % Line (epsilon): y = x+3 (maincolor)
    \addplot[very thick, maincolor, domain=-3.5:3.5] {x + 3} node[pos=0.9, above right, maincolor] {$(\varepsilon): y=x+3$};
    % Line (eta): y = x+1 (secondarycolor) - calculated from problem statement (y-2 = 1(x-1) => y=x+1)
    \addplot[very thick, secondarycolor, domain=-3.5:3.5] {x + 1} node[pos=0.9, above right, secondarycolor] {$(\eta): y=x+1$};
    % Point A(1,2)
    \addplot[only marks, mark=*, mark size=1.5pt, black] coordinates {(1,2)};
    \node[above right] at (axis cs:1,2) {$A(1,2)$};
\end{axis}\end{tikzpicture}
\end{center}

\vspace{4mm}

% --- Άσκηση 18241 (18241.tex) ---
\OnPar{ΘΕΜΑ 2}

Δίνεται ο κύκλος $C$ με εξίσωση $x^2 + y^2 = 25$. Να σχεδιάσετε στο ίδιο ορθοκανονικό σύστημα συντεταγμένων
\begin{alist}    \item τον κύκλο $C$. \monades{9}  [label=\greekalph*)]\end{alist}

\begin{center}
\begin{tikzpicture}\begin{axis}[
    axis lines = middle,
    xlabel = {$x$},
    ylabel = {$y$},
    xmin = -6, xmax = 6,
    ymin = -6, ymax = 6,
    xtick = {-5,-4,...,4,5},
    ytick = {-5,-4,...,4,5},
    grid = both,
    grid style = {gray!30},
    width = 9cm, height = 9cm,
    tick label style = {font=\scriptsize},
    every axis x label/.style = {at={(current axis.right of origin)}, anchor=north west},
    every axis y label/.style = {at={(current axis.above origin)}, anchor=south},
    axis equal image, 
]
    % Plot the circle x^2 + y^2 = 25 (radius 5, center 0,0)
    \addplot[thick, maincolor, domain=0:360, samples=100] ({5*cos(x)}, {5*sin(x)});
    \node[maincolor, above right, font=\small] at (axis cs:5,0) {$C$}; % Label the circle C
    % Mark the center (0,0) and label it O
    \fill (axis cs:0,0) circle (1.5pt); % a small dot for the origin
    \node[below left, font=\tiny] at (axis cs:0,0) {$O$};
    % Mark the intersection points with axes (where tangents will be drawn from)
    \fill (axis cs:5,0) circle (1.5pt); % (5,0)
    \fill (axis cs:-5,0) circle (1.5pt); % (-5,0)
    \fill (axis cs:0,5) circle (1.5pt); % (0,5)
    \fill (axis cs:0,-5) circle (1.5pt); % (0,-5)
\end{axis}\end{tikzpicture}
\end{center}

\begin{alist}    \item τις εφαπτόμενες του $C$ που διέρχονται από τα σημεία τομής του $C$ με τον $yy'$ και να γράψετε τις εξισώσεις τους. \monades{8}  [label=\greekalph*)] \setcounter{enumi}{1} % Continue list from the second item (β)
    \item τις εφαπτόμενες του $C$ που διέρχονται από τα σημεία τομής του $C$ με τον $xx'$ και να γράψετε τις εξισώσεις τους. \monades{8}\end{alist}

\vspace{4mm}

% --- Άσκηση 18242 (18242.tex) ---
\Askhsh{ALG21}{2}
Δίνεται η παραβολή $C$ με εξίσωση $y^2 = 4x$.

\begin{alist}    \item Να βρείτε τις συντεταγμένες της εστίας $E$ και την εξίσωση της διευθετούσας $\delta$ της $C$. \monades{8}
    \item Να γράψετε την εξίσωση της εφαπτομένης $(\varepsilon)$ της $C$ στο σημείο της $M(4,4)$. \monades{8}
    \item Να σχεδιάσετε στο ίδιο ορθοκανονικό σύστημα συντεταγμένων την παραβολή $C$, τη διευθετούσα $\delta$ και την ευθεία $(\varepsilon)$. \monades{9}
    \begin{center}
    \begin{tikzpicture}    \begin{axis}[
 axis lines = middle, xlabel = {$x$}, ylabel = {$y$}, xmin = -2.5, xmax = 9.5, ymin = -6.5, ymax = 6.5, xtick = {-2,-1,0,1,2,3,4,5,6,7,8,9}, ytick = {-6,-4,-2,0,2,4,6}, grid = both, grid style = {gray!30}, samples = 100, width = 10cm, height = 7cm, tick label style = {font=\scriptsize}, every axis x label/.style = {at={(current axis.right of origin)}, anchor=north west}, every axis y label/.style = {at={(current axis.above origin)}, anchor=south east}, ]
        % Parabola C: y^2 = 4x  <=> y = +/- 2*sqrt(x)
 \addplot[very thick, maincolor, domain=0:9] {2*sqrt(x)}; \addplot[very thick, maincolor, domain=0:9] {-2*sqrt(x)}; \node[maincolor, above right] at (axis cs:9,6) {$C: y^2=4x$};
        % Directrix delta: x = -1
 \addplot[very thick, thirdcolor, dashed, domain=-6.5:6.5] coordinates {(-1, -6.5) (-1, 6.5)}; \node[thirdcolor, above left] at (axis cs:-1,6.5) {$\delta: x=-1$};
        % Tangent epsilon: y = 1/2*x + 2
 \addplot[very thick, secondarycolor, domain=-2:9] {0.5*x + 2}; \node[secondarycolor, below right] at (axis cs:9,6.5) {$(\varepsilon): y=\frac{1}{2}x+2$};
        % Point M(4,4)
 \addplot[only marks, mark=*, mark size=1.5pt, black] coordinates {(4,4)}; \node[above right] at (axis cs:4,4) {$M(4,4)$};
        % Focus E(1,0)
 \addplot[only marks, mark=*, mark size=1.5pt, black] coordinates {(1,0)}; \node[below right] at (axis cs:1,0) {$E(1,0)$};
    \end{axis}\end{tikzpicture}
    \end{center}\end{alist}

\vspace{4mm}

% --- Άσκηση 18243 (18243.tex) ---
\section*{\textbf{ΘΕΜΑ 3}}

Θεωρούμε τα διανύσματα $\vec{\alpha}, \vec{\beta}$ με $|\vec{\alpha}|=2$, $|\vec{\beta}|=4$, $(\vec{\alpha}\hat{,}\vec{\beta})=\frac{\pi}{3}$ και τα διανύσματα
$\vec{\gamma} = \vec{\alpha} - \vec{\beta}$ και $\vec{\delta} = 2\vec{\alpha} + \vec{\beta}$.

\begin{alist}    \item Να βρείτε το $\vec{\alpha} \cdot \vec{\beta}$. \monades{5}
    \item Να βρείτε το $\vec{\gamma} \cdot \vec{\delta}$. \monades{7}
    \item Να βρείτε τα $|\vec{\gamma}|, |\vec{\delta}|$. \monades{8}
    \item Να βρείτε τη γωνία $(\vec{\gamma}\hat{,}\vec{\delta})$. \monades{5}\end{alist}

\vspace{4mm}

% --- Άσκηση 18244 (18244.tex) ---
\Askhsh{}{4}

Δίνονται οι ευθείες $\varepsilon1 : y = \sqrt{3}x$ και $\varepsilon2 : y = x$.

\begin{alist}    \item Να σχεδιάσετε τις $\varepsilon1,\varepsilon2$ στο ίδιο ορθοκανονικό σύστημα αξόνων. \monades{6}
    \begin{center}
    \begin{tikzpicture}    \begin{axis}[
 axis lines = middle, xlabel = {$x$}, ylabel = {$y$}, xmin = -3.5, xmax = 3.5, ymin = -3.5, ymax = 3.5, xtick = {-3,-2,-1,1,2,3}, ytick = {-3,-2,-1,1,2,3}, grid = both, grid style = {gray!30}, samples = 100, width = 10cm, height = 10cm, tick label style = {font=\scriptsize}, every axis x label/.style = {at={(current axis.right of origin)}, anchor=north west}, every axis y label/.style = {at={(current axis.above origin)}, anchor=south east}, ]
        % Eutheia e1: y = sqrt(3)*x
 \addplot[very thick, maincolor, domain=-3:3] {sqrt(3)*x}; \node[maincolor, above right, xshift=1mm] at (axis cs:2, sqrt(3)*2) {$\varepsilon1: y = \sqrt{3}x$};
        % Eutheia e2: y = x
 \addplot[very thick, secondarycolor, domain=-3:3] {x}; \node[secondarycolor, above left, xshift=-1mm] at (axis cs:2, 2) {$\varepsilon2: y = x$};
    \end{axis}\end{tikzpicture}
    \end{center}
    \item Να βρείτε τη γωνία που σχηματίζει κάθε μια από τις ευθείες $\varepsilon1$ και $\varepsilon2$ με τον άξονα $xx'$. \monades{6}
    \item Να αιτιολογήσετε γιατί η οξεία γωνία των $\varepsilon1,\varepsilon2$ είναι $15^\circ$. \monades{3}
    \item Να αποδείξετε ότι $\cos 15^\circ = \frac{\sqrt{3}+1}{2\sqrt{2}}$. \monades{10}\end{alist}

\vspace{4mm}

% --- Άσκηση 18245 (18245.tex) ---
\vspace{4mm}
\OnPar{ΘΕΜΑ 4}

Δίνεται η παραβολή $C : y^2 = 4x$ και η εξίσωση $(\lambda^2 -1)x+2\lambda y+\lambda^2 +1=0$ (1), με $\lambda \in \mathbb{R}$.

\begin{alist}    \item Να βρείτε τις συντεταγμένες της εστίας $E$ και την εξίσωση της διευθετούσας $\delta$ της παραβολής $C$. \monades{6}
    \item Να αποδείξετε ότι η (1) για κάθε $\lambda \in \mathbb{R}$ παριστάνει ευθεία $\varepsilon_\lambda$ που δεν διέρχεται από το $O(0,0)$. \monades{6}
    \item Να αποδείξετε ότι η διευθετούσα της παραβολής δεν ανήκει στην οικογένεια ευθειών $\varepsilon_\lambda$. \monades{6}
    \item Έστω $M(\alpha, \beta)$ σημείο του επιπέδου το οποίο δεν ανήκει στην παραπάνω διευθετούσα $\delta$. Αν από το $M$ διέρχεται μόνο μία ευθεία από την οικογένεια ευθειών $\varepsilon_\lambda$, να δείξετε ότι το $M$ ανήκει στον κύκλο που έχει κέντρο την κορυφή της παραβολής $C$ και διέρχεται από την εστία της $E$. \monades{7}\end{alist}

\vspace{4mm}

% --- Άσκηση 18247 (18247.tex) ---
\OnPar{ΘΕΜΑ 4}
Δίνονται τα σημεία $O(0,0)$, $A(\alpha,0)$ και $B(0, \beta)$, όπου $\alpha, \beta >0$.

\begin{alist}    \item Να βρείτε συναρτήσει των $\alpha,\beta$
    \begin{rlist}
        \item τις συντεταγμένες του μέσου $M$ του τμήματος $AB$. \monades{5}
        \item την απόσταση $(OM)$. \monades{5}
    \end{rlist}
    \item Αν $(OM) = \frac{\sqrt{\alpha^2 + \beta^2}}{2}$, τότε:
    \begin{rlist}
        \item να αποδείξετε ότι $(OM) = \frac{(AB)}{2}$. \monades{5}
        \item να γράψετε την πρόταση της Ευκλείδειας Γεωμετρίας που έχει αποδειχθεί. \monades{3}
    \end{rlist}
    \item Να βρείτε την εξίσωση του περιγεγραμμένου κύκλου του τριγώνου $OAB$. \monades{7}\end{alist}

\vspace{4mm}

% --- Άσκηση 18351 (18351.tex) ---
\section*{ΘΕΜΑ 2}
Δίνονται τα σημεία $A(-1,5)$, $B(3,3)$. Να υπολογίσετε:
\begin{alist}    \item Τις συντεταγμένες του μέσου $M$ του τμήματος $AB$. \monades{8}
    \item Τον συντελεστή διεύθυνσης της ευθείας $AB$. \monades{8}
    \item Την εξίσωση της μεσοκαθέτου ($\eta$) του τμήματος $AB$. \monades{9}\end{alist}

\vspace{4mm}

% --- Άσκηση 18372 (18372.tex) ---
\Askhsh{}{4}
Σε καρτεσιανό επίπεδο Oxy θεωρούμε τα σημεία A$(-2,-2)$, B$(0,-4)$ και την παραβολή $y^2 = 4x$.
\begin{alist}    \item Να βρείτε την παράμετρο, την εστία και την διευθετούσα της παραβολής. \monades{09}
    \item Να βρείτε το σημείο M της παραβολής στο οποίο η εφαπτομένη της είναι παράλληλη στην AB. \monades{08}
    \item Αν M$(1,-2)$ και K είναι το σημείο τομής της εφαπτομένης ευθείας του προηγούμενου ερωτήματος με τον άξονα $x'x$, να δείξετε ότι το τετράπλευρο ABMK είναι παραλληλόγραμμο. \monades{08}\end{alist}

\vspace{4mm}

% --- Άσκηση 18415 (18415.tex) ---
\Askhsh{}{4}
Δίνεται η εξίσωση $(x-3\lambda)^2 + (y+2\lambda)^2 = 1$ (1) όπου $\lambda \in \mathbb{R}$ και η ευθεία $\varepsilon: 2x+3y=0$.

\begin{alist}    \item Να αποδείξετε ότι για κάθε $\lambda \in \mathbb{R}$ τα κέντρα των κύκλων που προκύπτουν από την (1) ανήκουν στην ευθεία $\varepsilon$. \monades{6}
    \item Να βρείτε τις εξισώσεις των ευθειών $\varepsilon1,\varepsilon2$ που απέχουν μεταξύ τους 2 μονάδες και έχουν μεσοπαράλληλη την ευθεία $\varepsilon$. \monades{7}
    \item Να αποδείξετε ότι όλοι οι κύκλοι που προκύπτουν από την (1) εφάπτονται σε δύο σταθερές ευθείες. \monades{6}
    \item Να βρείτε το εμβαδόν ενός τετραγώνου του οποίου δύο απέναντι πλευρές ανήκουν στις ευθείες $\varepsilon1,\varepsilon2$ αντίστοιχα. \monades{6}\end{alist}

\vspace{4mm}

% --- Άσκηση 18416 (18416.tex) ---
\Askhsh{}{4}
Δίνεται η εξίσωση:
\[ x(x-4)+y(y-2)=2(x+y-4) \tag{1} \]

\begin{alist}    \item Να δείξετε ότι η εξίσωση (1) παριστάνει κύκλο με κέντρο $K(3,2)$ και ακτίνα $\rho = \sqrt{5}$. \monades{6}
    \item Δίνονται τα σημεία $A(4,4)$ και $B(2,0)$.
    \begin{rlist}
        \item Να δείξετε ότι τα σημεία $A$ και $B$ είναι αντιδιαμετρικά σημεία του κύκλου. \monades{4}
        \item Να βρείτε τις εξισώσεις των εφαπτόμενων του κύκλου οι οποίες είναι παράλληλες στην διάμετρο $AB$. \monades{9}
    \end{rlist}
    \item Να βρείτε την τιμή της παραμέτρου $\lambda$ ώστε η ευθεία ($\eta$) με εξίσωση $y=\lambda x+4$ να τέμνει τον παραπάνω κύκλο σε δύο σημεία $\Gamma$ και $\Delta$ ώστε $(\Gamma\Delta) = \sqrt{20}$. \monades{6}\end{alist}

\vspace{4mm}

% --- Άσκηση 18467 (18467.tex) ---
\Askhsh{004}{4}

Δίνεται η εξίσωση $(x-1)^2+(y+2)^2 = 2(x+3) \quad : (1).$

\begin{alist}    \item Να δείξετε ότι η εξίσωση (1) παριστάνει κύκλο με κέντρο $K(2,-2)$ και ακτίνα $\rho=3$. \monades{06}
    \item Να δείξετε ότι η αρχή $O$ των αξόνων είναι εσωτερικό σημείο του κύκλου. \monades{04}
    \item Να βρείτε την εξίσωση της ευθείας $(\varepsilon)$ η οποία τέμνει τον κύκλο σε δύο σημεία $A$ και $B$ ώστε η αρχή των αξόνων να είναι το μέσο της χορδής $AB$. \monades{08}
    \item Αν η ευθεία $(\varepsilon)$ του προηγούμενου ερωτήματος έχει εξίσωση $y=x$ τότε να βρείτε το εμβαδό του τριγώνου $KAB$. \monades{07}\end{alist}

\vspace{4mm}

% --- Άσκηση 18520 (18520.tex) ---
\Askhsh{}{4}
\begin{alist}    \item Να αποδειχθεί ότι για όλα τα διανύσματα $\vec{a}$, $\vec{\beta}$ ισχύει:
    \item Δίνεται το παραλληλόγραμμο ΟΑΓΒ με $\vec{OA} = \vec{a}$ και $\vec{OB} = \vec{\beta}$.  \[ |\vec{a} + \vec{\beta}|^2 + |\vec{a} - \vec{\beta}|^2 = 2|\vec{a}|^2 + 2|\vec{\beta}|^2 \quad (1) \] \monades{06}
    \begin{rlist}
        \item Να σχεδιάσετε τα διανύσματα $\vec{a} + \vec{\beta}$ και $\vec{a} - \vec{\beta}$.
        \begin{center}
        \begin{tikzpicture} \tkzDefPoint(0,0){O} \tkzDefPoint(4,0){A} % Vector a \tkzDefPoint(2,3){B} % Vector b \tkzDefPointBy[translation=from O to A](B) \tkzGetPoint{G} % Gamma, the fourth vertex OAGB
            % Draw the parallelogram
 \tkzDrawPolygon[l3,maincolor](O,A,G,B)
            % Draw vectors a and b
 \tkzDrawSegment[l4,->,secondarycolor](O,A) \tkzDrawSegment[l4,->,secondarycolor](O,B)
            % Draw a+b
 \tkzDrawSegment[l4,->,thirdcolor](O,G)
            % Draw a-b = BA
 \tkzDrawSegment[l4,->,thirdcolor](B,A) % Vector BA
            % Labels for points
 \tkzLabelPoint[below left](O){O} \tkzLabelPoint[below right](A){A} \tkzLabelPoint[above left](B){B} \tkzLabelPoint[above right](G){$\Gamma$}
            % Labels for vectors
 \tkzLabelSegment[below,pos=0.6](O,A){$\vec{a}$} \tkzLabelSegment[left,pos=0.7](O,B){$\vec{\beta}$} \tkzLabelSegment[above,pos=0.6](O,G){$\vec{a} + \vec{\beta}$} \tkzLabelSegment[below,pos=0.5](B,A){$\vec{a} - \vec{\beta}$} % Label for BA vector
\end{tikzpicture}
        \end{center}

        \item Να δώσετε τη γεωμετρική ερμηνεία της ισότητας (1).  \monades{05}
 \monades{04}
    \end{rlist}
    \item Ένα σώμα σύρεται πάνω σε λείο επίπεδο από δύο ανθρώπους, οι οποίοι εξασκούν πάνω σε αυτό δυνάμεις $\vec{F_1}$ και $\vec{F_2}$ αντίστοιχα, όπως φαίνεται στο παρακάτω σχήμα. Οι δυνάμεις έχουν ίσα μέτρα $10\ N$ (Newton) και η γωνία που σχηματίζουν είναι $60^\circ$. Να σχεδιάσετε την συνισταμένη δύναμη $\vec{F}$ και να βρείτε το μέτρο της.
    \begin{center}
    \begin{tikzpicture}        % Body (small square)
 \tkzDefPoint(0,0){O} \tkzDefPoint(-0.2,-0.2){BL} \tkzDefPoint(0.2,0.2){TR} \tkzDrawRectangle[fill=maincolor!15](BL,TR)
        % Force F1
 \tkzDefPoint(4,0){F1_end} \tkzDrawSegment[l4,->,secondarycolor](O,F1_end) \tkzLabelSegment[above,pos=0.5](O,F1_end){$\vec{F1}$}
        % Force F2
 \tkzDefPoint(4*cos(60),4*sin(60)){F2_end} % 60 degrees \tkzDrawSegment[l4,->,secondarycolor](O,F2_end) \tkzLabelSegment[left,pos=0.5](O,F2_end){$\vec{F2}$}
        % Resultant Force F
 \coordinate (Fresultantend) at ($(F1_end) + (F2_end)$); % Vector addition \tkzDrawSegment[l4,->,thirdcolor](O,Fresultantend) \tkzLabelSegment[above,pos=0.5](O,Fresultantend){$\vec{F}$}
        % Draw parallelogram lines (dashed)
 \tkzDrawSegment[dashed,gray](F1_end,Fresultantend) \tkzDrawSegment[dashed,gray](F2_end,Fresultantend)
        % Angle phi
 \tkzMarkAngle[fill=thirdcolor!20,size=0.8cm](F1_end,O,F2_end) \tkzLabelAngle[pos=1.1](F1_end,O,F2_end){$\varphi = 60^\circ$}
        % Label O (origin)
 \tkzDrawPoint[fill=black](O) \tkzLabelPoint[below left](O){O}\end{tikzpicture}
    \end{center}
 \monades{10}\end{alist}

\vspace{4mm}

% --- Άσκηση 18521 (18521.tex) ---
\Askhsh{}{4}
Δίνονται τα σημεία $A(1,2)$, $B(2,4)$ και $\Gamma(3,1)$.

\begin{alist}    \item Να αποδείξετε ότι η $B\widehat{A}\Gamma = 90^\circ$.
    \item Να βρείτε την εξίσωση του κύκλου $c$, ο οποίος διέρχεται από τα σημεία $A$, $B$ και $\Gamma$.  \monades{06}
    \item Αν ο κύκλος $c$ έχει εξίσωση $\left(x - \frac{5}{2}\right)^2 + \left(y - \frac{5}{2}\right)^2 = \frac{5}{2}$, τότε να βρείτε τις εξισώσεις των εφαπτόμενων του, οι οποίες διέρχονται από την αρχή των αξόνων.  \monades{09}
 \monades{10}\end{alist}

\vspace{4mm}

% --- Άσκηση 18547 (18547.tex) ---
\Askhsh{}{4}
Δίνονται τα σημεία $A(0,-1)$, $B(\lambda, 1)$ και $\Gamma(\lambda-2, \lambda-3)$, όπου $\lambda \in \mathbb{R}$.
\begin{alist}    \item Να βρείτε τις τιμές του $\lambda \in \mathbb{R}$ ώστε:
    \begin{rlist}
        \item Τα σημεία $A$, $B$ και $\Gamma$ να είναι κορυφές τριγώνου. \monades{8}
        \item Το τρίγωνο $AB\Gamma$ να είναι ορθογώνιο με $\hat{A} = 90^\circ$. \monades{7}
    \end{rlist}
    \item Για $\lambda = -2$, να βρείτε:
    \begin{rlist}
        \item Το εσωτερικό γινόμενο $\vec{AB} \cdot \vec{A\Gamma}$. \monades{4}
        \item Το εμβαδό του τριγώνου $AB\Gamma$. \monades{6}
    \end{rlist}\end{alist}

\vspace{4mm}

% --- Άσκηση 18567 (18567.tex) ---
\subsection*{ΘΕΜΑ 4}

Δίνεται ο κύκλος C: $x^2 + y^2 = 4$ και το σημείο A($2\sqrt{2}$, 0).

\begin{alist}    \item
    \begin{rlist}
        \item Να αποδείξετε ότι το σημείο Α είναι εξωτερικό του κύκλου C. \monades{05}
        \item Να βρείτε τις εξισώσεις των εφαπτόμενων του κύκλου C που διέρχονται από το
 σημείο Α και να αποδείξετε ότι είναι μεταξύ τους κάθετες. \monades{12}
    \end{rlist}
    \item Αν Β, Γ τα σημεία επαφής του κύκλου C με τις εφαπτόμενες ευθείες από το σημείο Α, να
 υπολογίσετε το εμβαδό του τετραπλεύρου ΑΒΟ\Gamma. \monades{08}\end{alist}

\vspace{4mm}

% --- Άσκηση 18568 (18568.tex) ---
\Askhsh{}{4}
Δίνονται τα σημεία $A(2,4)$, $B(-1,0)$ και $\Gamma(3,-2)$.

\begin{alist}    \item Να αποδείξετε ότι τα σημεία $A$, $B$, $\Gamma$ αποτελούν κορυφές τριγώνου $AB\Gamma$. \monades{04}
    \item Αν η ευθεία $AB$ τέμνει τον άξονα $y'y$ σε ένα σημείο $\Delta$ και η ευθεία $A\Gamma$ τέμνει τον άξονα $x'x$ σε ένα σημείο $E$, τότε:
    \begin{rlist}
        \item Να βρείτε τις συντεταγμένες των σημείων $\Delta$ και $E$. \monades{10}
        \item Να αποδείξετε ότι $\vec{A\Delta} = 2\vec{\Delta B}$ και $\vec{AE} = 2\vec{E\Gamma}$. \monades{06}
    \end{rlist}
    \item Να αποδείξετε ότι η ευθεία $\Delta E$ είναι παράλληλη της $B\Gamma$. \monades{05}\end{alist}

\vspace{4mm}

% --- Άσκηση 18569 (18569.tex) ---
\Askhsh{}{4}
Δίνεται ο κύκλος $C$: $x^2+y^2=1$.
\begin{alist}    \item Αν $A$ και $A'$ είναι τα σημεία τομής του κύκλου $C$ με τους ημιάξονες $Ox$ και $Ox'$ αντίστοιχα, τότε:
    \begin{rlist}
        \item Να αποδείξετε ότι οι συντεταγμένες των σημείων $A$ και $A'$ είναι $A(1,0)$ και $A'(-1,0)$.
        \item Να βρείτε την εξίσωση της ευθείας $\varepsilon$ που διέρχεται από το $A$ και σχηματίζει με τον άξονα $x'x$ γωνία $150^\circ$.  \monades{05}
 \monades{06}
    \end{rlist}
    \item Αν η ευθεία $\varepsilon$ τέμνει τον κύκλο $C$ και στο σημείο $B$, να αποδείξετε ότι η χορδή $AB$ έχει μήκος $\sqrt{3}$.
    \item Αν η ευθεία $\varepsilon$ έχει εξίσωση $y = -\frac{\sqrt{3}}{3}(x-1)$, να βρείτε την εξίσωση της ευθείας $(\zeta)$ που διέρχεται από τα σημεία $A'$ και $B$.  \monades{08}
 \monades{06}\end{alist}

\vspace{4mm}

% --- Άσκηση 18570 (18570.tex) ---
\Askhsh{}{4}
Δίνεται ο κύκλος με εξίσωση $x^2+y^2-4x-8y-5=0$ και η ευθεία $(\varepsilon): 3x-4y=\mu$, $\mu \in \mathbb{R}$.

\begin{alist}    \item Να βρείτε το κέντρο του κύκλου και την ακτίνα του. \monades{05}
    \item Αν η ευθεία $\varepsilon$ τέμνει τον κύκλο σε δύο διαφορετικά σημεία $A, B$
    \begin{rlist}
        \item Να αποδείξετε ότι $-35<\mu<15$. \monades{07}
        \item Να βρείτε για ποια τιμή του $\mu$ η ευθεία $\varepsilon$ διέρχεται από το κέντρο του. \monades{04}
        \item Να βρεθεί σημείο $\Gamma$ του κύκλου τέτοιο ώστε, το τρίγωνο $\Gamma AB$ να είναι ισοσκελές με βάση τη χορδή $AB$. \monades{09}
    \end{rlist}\end{alist}

\vspace{4mm}

% --- Άσκηση 18700 (18700.tex) ---
\Askhsh{M21}{2} % Example code: M21, Theme number: 2

Δίνεται κύκλος $C$ με κέντρο την αρχή των αξόνων και ακτίνα $R=5$.

\begin{alist}    \item Να γράψετε την εξίσωση του κύκλου $C$ και να τον σχεδιάσετε στο ορθοκανονικό σύστημα αξόνων. \monades{10}
    \item Δίνεται το σημείο $A(3, -4)$.
    \begin{rlist}
        \item Να αποδείξετε ότι το σημείο $A$ ανήκει στον κύκλο $C$. \monades{5}
        \item Να βρείτε την εξίσωση της εφαπτομένης του κύκλου $C$ στο σημείο $A$. \monades{10}
    \end{rlist}\end{alist}

\vspace{4mm}

% --- Άσκηση 18701 (18701.tex) ---
\Askhsh{}{2}
Δίνεται η παραβολή με εξίσωση
\[ y = \frac{1}{2}x^2 \quad \text{(1).} \]

\begin{alist}    \item Να βρείτε την εστία και τη διευθετούσα της παραβολής.
    \item Να βρείτε την εξίσωση της εφαπτομένης της παραβολής στο σημείο $A(2,2)$.  \monades{6}
    \item Να σχεδιάσετε στο ίδιο ορθοκανονικό σύστημα συντεταγμένων την παραβολή (1), την εστία, τη διευθετούσα και την εφαπτομένη της παραβολής.  \monades{10}
    \begin{center}
    \begin{tikzpicture}    \begin{axis}[
 axis lines = middle, xlabel = {$x$}, ylabel = {$y$}, xmin = -4, xmax = 4, ymin = -1.5, ymax = 5, xtick = {-3,-2,-1,0,1,2,3}, ytick = {-1,0,1,2,3,4,5}, grid = both, grid style = {gray!30}, samples = 100, width = 10cm, height = 7cm, tick label style = {font=\scriptsize}, every axis x label/.style = {at={(current axis.right of origin)}, anchor=north west}, every axis y label/.style = {at={(current axis.above origin)}, anchor=south}, clip=false, ]
        % Παραβολή y = 1/2 x^2
 \addplot[very thick, maincolor, domain=-3.5:3.5] {0.5*x^2}; \node[maincolor, above left] at (axis cs:3,4.5) {$y=\frac{1}{2}x^2$};
        % Εστία F(0, 0.5)
 \addplot[only marks, mark=*, mark size=2pt, thirdcolor] coordinates {(0,0.5)}; \node[thirdcolor, right] at (axis cs:0,0.5) {$F(0, \frac{1}{2})$};
        % Διευθετούσα y = -0.5
 \addplot[dashed, thick, secondarycolor, domain=-3.5:3.5] {-0.5}; \node[secondarycolor, above left] at (axis cs:-3.5,-0.5) {$y=-\frac{1}{2}$};
        % Σημείο Α(2,2)
 \addplot[only marks, mark=*, mark size=2pt, black] coordinates {(2,2)}; \node[above right] at (axis cs:2,2) {$A(2,2)$};
        % Εφαπτομένη y = 2x - 2
 \addplot[thick, thirdcolor, domain=-0.5:3] {2*x - 2}; \node[thirdcolor, below right] at (axis cs:2.5,3) {$y=2x-2$};
    \end{axis}\end{tikzpicture}
    \end{center}
 \monades{9}\end{alist}

\vspace{4mm}

% --- Άσκηση 18732 (18732.tex) ---
\Askhsh{}{4}

Σε σύστημα αξόνων δίνονται τα σημεία $A(3,4)$ και $B(7,1)$.

\begin{center}
\begin{tikzpicture}\begin{axis}[
    axis lines=middle,
    xmin=-4, xmax=13,
    ymin=-2, ymax=5,
    xtick distance=1,
    ytick distance=1,
    grid=both,
    grid style={gray!30},
    width=12cm, height=8cm,
    tick label style={font=\scriptsize},
    xlabel style={at={(current axis.right of origin)}, anchor=north west},
    ylabel style={at={(current axis.above origin)}, anchor=south east},
]
    % Points A, B, O
    \addplot[only marks, mark=*, mark size=1.5pt, black] coordinates {(3,4) (7,1) (0,0)};
    \node[above left] at (axis cs:3,4) {A};
    \node[below right] at (axis cs:7,1) {B};
    \node[below left] at (axis cs:0,0) {O}; % Origin label
    % Segments OA, AB (solid)
    \addplot[l3, black] coordinates {(0,0) (3,4)}; % OA
    \addplot[l3, black] coordinates {(3,4) (7,1)}; % AB
    % Segment OB (dashed)
    \addplot[l3, black, dashed] coordinates {(0,0) (7,1)}; % OB
\end{axis}\end{tikzpicture}
\end{center}

\begin{alist}    \item Αν $\Gamma\left(\frac{2}{3}, \frac{8}{3}\right)$ και $\Delta\left(\frac{13}{3}, 3\right)$ να δείξετε ότι:
    \begin{rlist}
        \item $\vec{A\Gamma} = \frac{1}{3}\vec{AO}$ και $\vec{A\Delta} = \frac{1}{3}\vec{AB}$
        \item $\Gamma\Delta \parallel OB$.  \monades{6}
        \item Να δείξετε ότι $(\Alpha\Gamma\Delta) = \left(\frac{1}{3}\right)^2 (\Alpha OB)$.  \monades{5}
 \monades{5}
    \end{rlist}
    \item Γενικεύοντας το παράδειγμα του α) ερωτήματος, αν για τα σημεία $\Gamma$ και $\Delta$ ισχύουν
 \[ \vec{A\Gamma} = \frac{1}{\nu}\vec{AO} \quad \text{και} \quad \vec{A\Delta} = \frac{1}{\nu}\vec{AB}, \] να δείτε ότι $(\Alpha\Gamma\Delta) = \left(\frac{1}{\nu}\right)^2 (\Alpha ABO)$. \monades{9}\end{alist}

\vspace{4mm}

% --- Άσκηση 18733 (18733.tex) ---
\Askhsh{}{2}
Δίνονται τα σημεία $A(4,3)$, $B(1,1)$ και $\Gamma(6,0)$.
\begin{alist}    \item Να βρείτε τις συντεταγμένες των διανυσμάτων $\vec{AB}$ και $\vec{A\Gamma}$. \monades{8}
    \item Να δείξετε ότι τα διανύσματα $\vec{AB}$ και $\vec{A\Gamma}$ είναι κάθετα. \monades{8}
    \item Δίνεται το σημείο $M\left(\frac{7}{2}, \frac{1}{2}\right)$. Να δείξετε ότι $(MA)=(MB)$. \monades{9}\end{alist}

\vspace{4mm}

% --- Άσκηση 18741 (18741.tex) ---
\Askhsh{}{4}
Δίνεται η παραβολή $C: y^2 = \alpha \cdot x$ η οποία διέρχεται από το σημείο $M(16,\alpha+4)$.
\begin{alist}    \item Να αποδείξετε ότι $\alpha = 4$. \monades{5}
    \item Να βρείτε την εστία $E$ και τη διευθετούσα $\delta$ της παραβολής. \monades{7}
    \item Να βρείτε την εξίσωση της εφαπτομένης $\varepsilon1$, της παραβολής $C$ η οποία είναι παράλληλη στην ευθεία $\varepsilon2 : -x+2y+4=0$. \monades{7}
    \item Να βρείτε την εξίσωση κύκλου $C_1$ με κέντρο την κορυφή της παραβολής $C$ ο οποίος εφάπτεται στην ευθεία $\varepsilon1$ του ερωτήματος γ). \monades{6}\end{alist}

\vspace{4mm}

% --- Άσκηση 18745 (18745.tex) ---
\Askhsh{A001}{4}
Επιστήμονες προκειμένου να μελετήσουν υδρόβιο έντομο κατέγραψαν στιγμιότυπα από τους κύκλους με κέντρα τα σημεία A, B, $\Gamma$, $\Delta$ και ακτίνες $3, 2, 1, \frac{1}{2}$ αντίστοιχα, που σχηματίζονται σε κάθε προσγείωση του στο νερό. Η εικόνα από τις εναέριες λήψεις αποτυπώθηκαν σε σύστημα αξόνων όπως φαίνεται στο σχήμα. Το έντομο κινούμενο ευθύγραμμα περνάει από τα σημεία A, B, $\Gamma$, $\Delta$ για να καταγραφεί την στιγμή που καταλήγει στο σημείο M.

\begin{center}
\begin{tikzpicture}\begin{axis}[
    axis lines = middle,
    xlabel = {$x$},
    ylabel = {$y$},
    xmin = -6.5, xmax = 7.5,
    ymin = -2.5, ymax = 4.5,
    xtick = {-6,-5,-4,-3,-2,-1,0,1,2,3,4,5,6,7},
    ytick = {-2,-1,0,1,2,3,4},
    grid = both,
    grid style = {gray!30},
    samples = 100,
    width = 12cm, height = 9cm,
    tick label style = {font=\scriptsize},
    every axis x label/.style = {at={(current axis.right of origin)}, anchor=north west},
    every axis y label/.style = {at={(current axis.above origin)}, anchor=south},
    clip mode=individual, 
]
    % Centers of the circles
    \coordinate (Acenter) at (-2,0);
    \coordinate (Bcenter) at (0,0);
    \coordinate (Gcenter) at (2,0); % Gamma
    \coordinate (Dcenter) at (3.5,0); % Delta
    \coordinate (Mpoint) at (4,1);
    % Draw circles with specified colors and line thickness
    \draw[l3,maincolor] (Acenter) circle (3);
    \draw[l3,secondarycolor] (Bcenter) circle (2);
    \draw[l3,thirdcolor] (Gcenter) circle (1);
    \draw[l3,maincolor!70!white] (Dcenter) circle (0.5); % Lighter shade of maincolor
    % Plot the line (tangent)
    % y = (sqrt(3)/3)x + (3-4sqrt(3))/3
    \addplot[l4, black, domain=-6:7] {(sqrt(3)/3)*x + (3-4*sqrt(3))/3};
    % Draw points and labels for centers and M
    \fill (Acenter) circle (1.5pt); \node[above left] at (Acenter) {A};
    \fill (Bcenter) circle (1.5pt); \node[above left] at (Bcenter) {B};
    \fill (Gcenter) circle (1.5pt); \node[above left] at (Gcenter) {$\Gamma$};
    \fill (Dcenter) circle (1.5pt); \node[above left] at (Dcenter) {$\Delta$};
    \fill (Mpoint) circle (1.5pt); \node[above right] at (Mpoint) {M};
\end{axis}\end{tikzpicture}
\end{center}

\begin{alist}
\item α) Να βρείτε την εξίσωση της πορείας του εντόμου. \monades{4}
\item β)
\begin{rlist}
\item Να δείξετε ότι η ευθεία ($\varepsilon1$): $y = \frac{\sqrt{3}}{3}x + \frac{3-4\sqrt{3}}{3}$ είναι κοινή εφαπτόμενη των τεσσάρων κύκλων. \monades{7}
\item Να βρείτε την εξίσωση της άλλης κοινής εφαπτομένης. \monades{10}
\endrlist
\item γ) Με βάση το μοτίβο που ακολουθούν οι κινήσεις του εντόμου να βρείτε ότι η τελική θέση του εντόμου είναι το σημείο $M(4,1)$. \monades{4}
\end{alist}

Δίνεται ότι $\varepsilon \phi 150^\circ = -\frac{\sqrt{3}}{3}$.

\vspace{4mm}

% --- Άσκηση 18749 (18749.tex) ---
\section*{ΘΕΜΑ 2}

Σε ορθοκανονικό σύστημα συντεταγμένων θεωρούμε τρίγωνο $AB\Gamma$ ώστε $A(5, 6)$, $B(1, 2)$, $\Gamma(12, 2)$ και το ύψος του $A\Delta$, όπου $\Delta$ σημείο της $B\Gamma$, όπως στο παρακάτω σχήμα.

\begin{alist}    \item Να βρείτε τις εξισώσεις των ευθειών $B\Gamma$ και $A\Delta$. \monades{10}
    \item Να βρείτε τις συντεταγμένες του σημείου $\Delta$. \monades{5}
    \item Να βρείτε την εξίσωση του κύκλου με κέντρο το σημείο $A$, ο οποίος εφάπτεται της ευθείας $B\Gamma$ στο σημείο $\Delta$. \monades{10}\end{alist}

\vspace{1cm} % Add some vertical space to separate text from figure

\begin{center}
    \begin{tikzpicture}[scale=0.6]
        % Ορισμός σημείων
        \tkzDefPoint(0,0){O}
        \tkzDefPoint(1,2){B}
        \tkzDefPoint(12,2){C}
        \tkzDefPoint(5,6){A}
        \tkzDefPoint(5,2){D} % Σημείο Δ, πόδι του ύψους από το Α στην ΒΓ
        % Σχεδίαση αξόνων
        \draw[l3, ->] (-1,0) -- (14,0) node[below] {$x$};
        \draw[l3, ->] (0,-1) -- (0,7) node[left] {$y$};
        % Σχεδίαση τριγώνου ΑΒΓ
        \tkzDrawPolygon[l3](A,B,C)
        % Σχεδίαση ύψους ΑΔ
        \tkzDrawSegment[l3,secondarycolor](A,D)
        % Σχεδίαση κύκλου με κέντρο Α και ακτίνα ΑΔ
        \tkzDrawCircle[l3,thirdcolor](A,D)
        % Σημεία
        \tkzDrawPoints(A,B,C,D,O)
        % Ετικέτες σημείων
        \tkzLabelPoint[above right](A){A}
        \tkzLabelPoint[below left](B){B}
        \tkzLabelPoint[below right](C){$\Gamma$}
        \tkzLabelPoint[below](D){$\Delta$}
        \tkzLabelPoint[below left](O){O}
        % Ετικέτες αξόνων (x', y' όπως στο σχήμα)
        \node[below left] at (-0.5,0) {$x'$};
        \node[below left] at (0,-0.5) {$y'$};
        % Σήμανση ορθής γωνίας στο Δ
        \tkzMarkRightAngle[fill=thirdcolor!20,size=0.4](A,D,C)\end{tikzpicture}
\end{center}

\vspace{4mm}

% --- Άσκηση 18781 (18781.tex) ---
\Askhsh{}{4}

Δίνονται δύο κύκλοι με εξισώσεις $C_1: x^2 + y^2 - 6x + 5 = 0$ και $C_2: x^2 + y^2 = 1$.

\begin{alist}    \item Να δείξετε ότι:
    \begin{rlist}
        \item Η εξίσωση του κύκλου $C_1$ γράφεται στη μορφή $(x - 3)^2 + y^2 = 4$. \monades{5}
        \item Οι κύκλοι $C_1, C_2$ εφάπτονται εξωτερικά. \monades{4}
    \end{rlist}
    \item Να βρείτε:
    \begin{rlist}
        \item Το σημείο επαφής των δύο κύκλων $C_1$ και $C_2$. \monades{6}
        \item Την εξίσωση της εσωτερικής κοινής εφαπτομένης των δύο κύκλων $C_1$ και $C_2$. \monades{4}
    \end{rlist}
    \item Αν τα σημεία $M_1, M_2$ διατρέχουν τους κύκλους $C_1, C_2$ αντίστοιχα, να βρείτε τη μέγιστη απόσταση ανάμεσα στα σημεία αυτά. \monades{6}\end{alist}

\vspace{4mm}

% --- Άσκηση 18870 (18870.tex) ---
\Askhsh{}{4}

Στο σχήμα φαίνεται η γραφική παράσταση της παραβολής με εξίσωση $y^2 = 4x$, η εφαπτομένη της $(\varepsilon)$ στο σημείο $A(4,4)$ και η $AK$ κάθετη στην $(\varepsilon)$. Μία φωτεινή ακτίνα $(\zeta)$, ακολουθώντας πορεία παράλληλη προς τον άξονα της παραβολής, προσπίπτουσα στο σημείο $A$ και ανακλώμενη πάνω στην καμπύλη (που αντιστοιχεί σε παραβολικό κάτοπτρο) διέρχεται από το σημείο $M$. Αν γνωρίζετε ότι η γωνία $\theta_1$ που σχηματίζει η προσπίπτουσα φωτεινή ακτίνα $(\zeta)$ με την $(\varepsilon)$ και η γωνία $\theta_2$ που σχηματίζει η ανακλώμενη φωτεινή ακτίνα $AM$ με την $(\varepsilon)$ είναι ίσες, τότε:

\begin{alist}    \item Να βρείτε την εστία και την διευθετούσα της παραβολής. \monades{06}
    \item Να βρείτε την εξίσωση της ευθείας $(\varepsilon)$ και το σημείο $B$ στο οποίο αυτή τέμνει τον άξονα $x'x$. \monades{06}
    \item Να αποδείξετε ότι το τρίγωνο $MAB$ είναι ισοσκελές. \monades{07}
    \item Να αποδείξετε ότι το σημείο $M$ ταυτίζεται με την εστία της παραβολής. \monades{06}\end{alist}

\begin{center}
\begin{tikzpicture}[scale=0.8]
    % Define points based on the problem statement and geometric properties
    \tkzDefPoint(0,0){O} % Origin
    \tkzDefPoint(4,4){A} % Point A on the parabola
    \tkzDefPoint(6,0){K} % K, intersection of normal AK with x-axis
                         % Normal to y^2=4x at (4,4) has slope -2. Eq: y-4 = -2(x-4) => 2x+y-12=0. For y=0, x=6.
    \tkzDefPoint(1,0){M} % M, assumed to be the focus for drawing the reflected ray.
                         % Focus of y^2=4x is (1,0). The reflection property states AM passes through focus.
    \tkzDefPoint(-4,0){B} % B, intersection of tangent (epsilon) with x-axis
                          % Tangent to y^2=4x at (4,4) has slope 1/2. Eq: y-4 = 1/2(x-4) => x-2y+4=0. For y=0, x=-4.
    % Draw axes (x'x and Oy)
    \draw[->, gray!70!black] (-5,0) -- (7.5,0) node[right] {$x$};
    \draw[->, gray!70!black] (0,-2) -- (0,5.5) node[above] {$y$};
    \tkzLabelPoint[below left](O){O}
    % Draw the parabola y^2 = 4x using pgfplots
    \begin{axis}[
        hide axis, 
        xmin=-1, xmax=7, ymin=-4, ymax=5,
        clip=false, 
    ]
        \addplot[thick, thirdcolor, domain=0:7, samples=100] {2*sqrt(x)};
        \addplot[thick, thirdcolor, domain=0:7, samples=100] {-2*sqrt(x)};
    \end{axis}
    % Draw the tangent (ε) through B(-4,0) and A(4,4)
    \tkzDrawLine[add=1.5 and 0.5, l3, maincolor](B,A) % Extends the line segment BA
    \tkzLabelLine[above left](A,B){$(\varepsilon)$}
    % Draw the normal AK
    \tkzDrawSegment[l3, secondarycolor](A,K)
    % Draw incident ray (ζ) parallel to x-axis, goes through A(4,4)
    \coordinate (Zstart) at ($(A) + (-5,0)$); % A point 5 units to the left of A at y=4
    \tkzDrawSegment[add=0 and 0.5, l3, thirdcolor](Zstart,A)
    \tkzLabelLine[above right, pos=0.8](Zstart,A){$(\zeta)$}
    % Draw reflected ray AM
    \tkzDrawSegment[l3, red](A,M)
    % Draw points
    \tkzDrawPoints(A,B,K,M,O)
    % Labels for points
    \tkzLabelPoint[above right](A){A}
    \tkzLabelPoint[below left](B){B}
    \tkzLabelPoint[below right](K){K}
    \tkzLabelPoint[below](M){M}
    % Define auxiliary points for angle drawing
    % A point on tangent (epsilon) to the right of A (slope 1/2)
    \tkzDefPoint(5,4.5){Aextright}
    % A point on incident ray (zeta) to the left of A (y=4)
    \tkzDefPoint(1,4){Zinc}
    % Angle θ1: between incident ray (ζ) and tangent (ε) at A
    \tkzMarkAngle[arc=lmk, size=0.8cm, thirdcolor](Zinc,A,Aextright)
    \tkzLabelAngle[pos=1.1, thirdcolor](Zinc,A,Aextright){$\theta1$}
    % Angle θ2: between reflected ray AM and tangent (ε) at A
    \tkzMarkAngle[arc=lmk, size=0.8cm, red](M,A,Aextright)
    \tkzLabelAngle[pos=1.1, red](M,A,Aextright){$\theta2$}
    % Angle θ3: between tangent (ε) and x-axis at B
    % This is the angle OBA, where O is the origin (0,0) and B is (-4,0)
    \tkzMarkAngle[arc=lmk, size=0.8cm, maincolor](O,B,A)
    \tkzLabelAngle[pos=1.1, maincolor](O,B,A){$\theta3$}\end{tikzpicture}
\end{center}

\vspace{4mm}

% --- Άσκηση 18871 (18871.tex) ---
\OnPar{ΘΕΜΑ 4}

Δίνεται ο κύκλος $(c)$ με κέντρο $O(0,0)$ και ακτίνα $\rho = \sqrt{5}$ και το σημείο $A(3,1)$.

\begin{alist}    \item Να βρείτε την εξίσωση του κύκλου και να αποδείξετε οτι το σημείο $A$ είναι εξωτερικό του κύκλου.
    \item  \monades{07}
    \begin{rlist}
        \item Να αποδείξετε οτι οι ευθείες που εφάπτονται στον κύκλο $(c)$ και διέρχονται από το σημείο $A$ έχουν εξισώσεις $(\varepsilon1): 2x - y = 5$ και $(\varepsilon2): x + 2y = 5$.

        \item Να βρείτε την εξίσωση της διχοτόμου της γωνίας $B\hat{A}\Gamma$, όπου $B$ και $\Gamma$ είναι αντίστοιχα τα σημεία επαφής των ευθειών $(\varepsilon1)$ και $(\varepsilon2)$ με τον κύκλο.  \monades{10}
 \monades{08}
    \end{rlist}\end{alist}

\vspace{4mm}

% --- Άσκηση 18878 (18878.tex) ---
\Askhsh{2}{2}
Ένας εξερευνητής ξεκίνησε από την κατασκήνωσή του (σημείο $O$) τρεις μέρες πριν, για ένα ταξίδι μέσα στη ζούγκλα. Στο τέλος της πρώτης ημέρας έφθασε στο σημείο $A$, στο τέλος της δεύτερης ημέρας έφθασε στο σημείο $B$ και στο τέλος της τρίτης ημέρας έφθασε στο σημείο $T$. Οι τρεις ημέρες του ταξιδιού του μπορούν να περιγραφούν από τα παρακάτω διανύσματα
\[ \vec{OA} = (1, 1), \quad \vec{AB} = (2, 4), \quad \vec{BT} = (2,5\sqrt{3} - 5), \]
όπως φαίνονται στο σχήμα:

\begin{tikzpicture}\begin{axis}[
    axis lines = middle,
    xlabel = {$x$},
    ylabel = {$y$},
    xmin = -1, xmax = 6,
    ymin = -1, ymax = 10,
    xtick = {0,1,2,3,4,5,6},
    ytick = {0,1,2,3,4,5,6,7,8,9,10},
    grid = both,
    grid style = {gray!30},
    width = 10cm, height = 10cm,
    tick label style = {font=\scriptsize},
    every axis x label/.style = {at={(current axis.right of origin)}, anchor=north west},
    every axis y label/.style = {at={(current axis.above origin)}, anchor=south east},
]
    % Calculate coordinates:
    % O = (0,0)
    % A = O + OAvec = (0,0) + (1,1) = (1,1)
    % B = A + ABvec = (1,1) + (2,4) = (3,5)
    % T = B + BTvec = (3,5) + (2, 5*sqrt(3)-5) = (5, 5*sqrt(3))
    % 5*sqrt(3) approx 8.66
    % Plot the path segments
    \addplot[l3, maincolor] coordinates {(0,0) (1,1)}; % OA
    \addplot[l3, maincolor] coordinates {(1,1) (3,5)}; % AB
    \addplot[l3, maincolor] coordinates {(3,5) (5, {5*sqrt(3)})}; % BT
    % Plot the direct vector OT (dashed blue)
    \addplot[l3, dashed, secondarycolor] coordinates {(0,0) (5, {5*sqrt(3)})}; % OT
    % Plot points
    \addplot[only marks, mark=*, mark size=1.5pt, black] coordinates {
        (0,0) (1,1) (3,5) (5, {5*sqrt(3)})
    };
    % Labels for points
    \node[below left] at (axis cs:0,0) {$O$};
    \node[above right] at (axis cs:1,1) {$A$};
    \node[above left] at (axis cs:3,5) {$B$};
    \node[above right] at (axis cs:5, {5*sqrt(3)}) {$T$};
\end{axis}\end{tikzpicture}

Αν οι αποστάσεις εκφράζονται σε χιλιόμετρα, τότε:
\begin{alist}    \item Να αποδείξετε ότι $\vec{OT} = (5,5\sqrt{3})$. \monades{13}
    \item Να υπολογίσετε την απόσταση $(OT)$ του εξερευνητή από την κατασκήνωση στο τέλος της τρίτης ημέρας. \monades{12}\end{alist}

\vspace{4mm}

% --- Άσκηση 18968 (18968.tex) ---
\Askhsh{}{2}
Δίνεται ο κύκλος $C$ με εξίσωση $x^2 + y^2 - 6x - 8y = 0 \quad (1)$.

\begin{alist}    \item Να δείξετε ότι ο κύκλος $C$ έχει κέντρο $K(3,4)$ και ακτίνα $\rho = 5$. \monades{10}
    \item Να δείξετε ότι η απόσταση του κέντρου $K$ του κύκλου από την ευθεία $\varepsilon: 3x + 4y = 0$ ισούται με $5$ μονάδες μήκους. \monades{8}
    \item Να δικαιολογήσετε αν είναι αληθής ή ψευδής ο ισχυρισμός: «Ο κύκλος $C$ και η ευθεία $\varepsilon$ εφάπτονται». \monades{7}\end{alist}

\vspace{4mm}

% --- Άσκηση 18979 (18979.tex) ---
\Askhsh{}{2}

Δίνονται οι ευθείες $\varepsilon1: 2x + 3y = 5$ και $\varepsilon2: 4x + 6y = 8$.

\begin{alist}    \item Να δείξετε ότι οι ευθείες $\varepsilon1, \varepsilon2$ είναι παράλληλες.
    \item Να αποδείξετε ότι το σημείο $A(1,1)$ είναι σημείο της ευθείας $\varepsilon1$.  \monades{10}
    \item Να βρείτε την απόσταση του σημείου $A$ από την ευθεία $\varepsilon2$.  \monades{5}
 \monades{10}\end{alist}

\vspace{4mm}

% --- Άσκηση 19038 (19038.tex) ---
\Askhsh{}{2}
Δίνεται τα διανύσματα $\vec{\alpha} = (2,3)$, $\vec{\beta} = (-1,1)$ και $\vec{\gamma} = (-5,-5)$.
\begin{alist}    \item Να υπολογίσετε τη γωνία που σχηματίζει το διάνυσμα $\vec{\beta}$ με τον άξονα $x'x$. \monades{9}
    \item Να αποδείξετε ότι $|\vec{\gamma}| = 5|\vec{\beta}|$. \monades{8}
    \item Να υπολογίσετε τους πραγματικούς αριθμούς $\lambda$, $\mu$ ώστε το διάνυσμα $\vec{\gamma}$ να γραφεί στη μορφή $\vec{\gamma} = \lambda\vec{\alpha} + \mu\vec{\beta}$. \monades{8}\end{alist}

\vspace{4mm}

% --- Άσκηση 19039 (19039.tex) ---
\Askhsh{2}{2}
Δίνεται η εξίσωση
\[
(x - 1)(x + 3) + (y + 1)(y - 3) = -4 \quad (1)
\]
\begin{alist}    \item Να αποδείξετε ότι η εξίσωση $(1)$ παριστάνει κύκλο με κέντρο $K(-1,1)$ και ακτίνα $R = 2$.
    \item  \monades{9}
    \begin{rlist}
        \item Να βρείτε τα σημεία $A$ και $B$ του κύκλου $(K,R)$ τα οποία έχουν τετμημένη ίση με $-1$.
        \item Να αποδείξετε ότι τα σημεία $A$ και $B$ είναι αντιδιαμετρικά.  \monades{8}
 \monades{8}
    \end{rlist}\end{alist}

\vspace{4mm}

% --- Άσκηση 19047 (19047.tex) ---
\Askhsh{Α.Δ.4}{4}

Στην Golden Gate γέφυρα του San Francisco, το κεντρικό καλώδιο θεωρούμε προσεγγιστικά ότι αποτεύει τμήμα παραβολής. Οι δύο βασικοί πυλώνες απέχουν μεταξύ τους $1280 \text{ m}$, ενώ το ύψος του κάθε πυλώνα σε σχέση με το οδόστρωμα της γέφυρας είναι $160 \text{ m}$. Γνωρίζουμε ότι το κατώτερο σημείο του παραβολικού καλωδίου αγγίζει τη γέφυρα στο μέσο της απόστασης των δύο πυλώνων. Θεωρούμε ορθογώνιο σύστημα αξόνων, όπως στο σχήμα.

\begin{alist}    \item Να αποδείξετε ότι η εξίσωση της παραβολής του κεντρικού καλωδίου σ΄ αυτό το σύστημα των αξόνων είναι $x^2 = 2560y$. \monades{9}
    \item Να βρείτε τις συντεταγμένες της εστίας $E$ και την εξίσωση της διευθετούσας $(\delta)$ της παραβολής. \monades{8}
    \item Η εφαπτομένη της παραβολής στο σημείο $B(640,160)$ τέμνει τον άξονα $y'y$ στο σημείο $\Delta$. Να αποδείξετε ότι $E\Delta = EB$. \monades{8}\end{alist}

\begin{center}
\begin{tikzpicture}    % Scaling for better visualization: 1 unit in TikZ = 1m in real.
    % x ranges from -640 to 640, y from 0 to 160.
    % We use different scales for x and y to make the parabola visible without making the drawing too wide.
    \pgfmathsetmacro{\xunit}{0.01} % 1m in x-direction corresponds to 0.01cm on page
    \pgfmathsetmacro{\yunit}{0.03} % 1m in y-direction corresponds to 0.03cm on page (magnifies y-axis)
    \begin{scope}[xscale=\xunit, yscale=\yunit]
        % Define key coordinates using the actual values in meters
        \coordinate (O) at (0,0);
        \coordinate (PylonLeftBase) at (-640,0);
        \coordinate (PylonRightBase) at (640,0);
        \coordinate (PylonLeftTop) at (-640,160);
        \coordinate (PylonRightTop) at (640,160);
        \coordinate (B) at (640,160); % Point B is the top of the right pylon
        % Axes
        % x-axis (x'Ox)
        \draw[<-, l3, maincolor] (-700,0) node[left, black, font=\small] {$x'$};
        \draw[->, l3, maincolor] (O) -- (700,0) node[right, black, font=\small] {$x$};
        % y-axis (y'Oy)
        \draw[<-, l3, maincolor] (0,-40) node[below, black, font=\small] {$y'$};
        \draw[->, l3, maincolor] (O) -- (0,180) node[above, black, font=\small] {$y$};
        % Origin label
        \node[below left, black] at (O) {$O$};
        % Pylons (drawn as vertical lines in secondarycolor)
        \draw[l3, secondarycolor] (PylonLeftBase) -- (PylonLeftTop);
        \draw[l3, secondarycolor] (PylonRightBase) -- (PylonRightTop);
        % Parabolic cable (drawn using the function y = x^2 / 2560 in thirdcolor)
        \draw[l4, thirdcolor, domain=-640:640, samples=100] plot (\x, {\x*\x / 2560});
        % Mark and label Point B
        \tkzDrawPoint[color=maincolor,fill=maincolor,size=3](B) % Draw a small filled circle for B
        \node[maincolor, above right, font=\small] at (B) {B}; % Label B
        % Label for 1280m distance between pylons
        % Draw a dashed horizontal line slightly below the x-axis, and place the label above it.
        \draw[dashed, l3, black] (-640,-15) -- (640,-15);
        \node[above=1mm, black] at (0,-15) {\small $1280\text{m}$};
    \end{scope}\end{tikzpicture}
\end{center}

\vspace{4mm}

% --- Άσκηση 20090 (20090.tex) ---
\Askhsh{004}{4}
\noindent Δίνεται η παραβολή $y^2 = 4x$ και $M(x_0, y_0)$, με $y_0 > 0$, ένα σημείο της.

\begin{alist}    \item Αν A είναι η προβολή του M στη διευθετούσα της παραβολής,
    \begin{rlist}
        \item Να εκφράσετε τις συντεταγμένες των σημείων M και A συναρτήσει της τεταγμένης $y_0$ του σημείου M. \monades{05}
        \item Αν E είναι η εστία της παραβολής, να βρείτε το σημείο M για το οποίο $(MAE) = \dfrac{5}{8}$ τ.μ. \monades{12}
    \end{rlist}
    \item Αν $M\left(\dfrac{1}{4}, 1\right)$ και $\varepsilon$ η εφαπτομένη της παραβολής στο σημείο M, να αποδείξετε ότι το τετράπλευρο AMEΜ$'$ είναι ρόμβος, όπου E είναι η εστία της παραβολής και M$'$ το σημείο που η ευθεία $\varepsilon$ τέμνει τον άξονα $x'x$. \monades{08}\end{alist}

\vspace{4mm}

% --- Άσκηση 20091 (20091.tex) ---
\Askhsh{}{4}
Τα σημεία $Α(-7, -1)$ και $Β(3, -5)$ είναι σημεία ενός κύκλου $C$ κέντρου $K$. Το σημείο $M$ είναι το μέσο της χορδής $ΑΒ$ και μία ευθεία $\varepsilon$ διέρχεται από τα σημεία $K$ και $M$.

\begin{alist}    \item Να βρείτε:
    \begin{rlist}
        \item Τις συντεταγμένες του σημείου $M$. \monades{04}
        \item Την εξίσωση της ευθείας $KM$. \monades{08}
    \end{rlist}
    \item Αν από το κέντρο $K$ του κύκλου διέρχεται η ευθεία $(\delta): x+y= -12$, τότε:
    \begin{rlist}
        \item Να βρείτε τις συντεταγμένες του σημείου $K$. \monades{07}
        \item Να βρείτε την εξίσωση του κύκλου $C$. \monades{06}
    \end{rlist}\end{alist}

\vspace{4mm}

% --- Άσκηση 20092 (20092.tex) ---
ΘΕΜΑ 4

Δίνεται η παραβολή $y^2 = 4x$, το σημείο της $M\left(\frac{1}{4}, 1\right)$ και η ευθεία $\varepsilon$ του επιπέδου με εξίσωση
\[ \varepsilon: \frac{x}{3} + \frac{y}{4} + 1 = 0. \]

\begin{alist}    \item
    \begin{rlist}
        \item Να δείξετε ότι η ευθεία $\varepsilon$ δεν έχει κοινά σημεία με την παραβολή και να βρείτε την απόστασή του σημείου $M$ από την $\varepsilon$. \monades{07}
        \item Αν η ευθεία $\varepsilon$ τέμνει τους άξονες $x'x$ και $y'y$ στα σημεία $\Gamma$ και $\Delta$ αντίστοιχα, να δείξετε ότι $(\text{M}\Gamma\Delta) = 5$ τ.μ. \monades{05}
    \end{rlist}
    \item
    \begin{rlist}
        \item Να βρείτε την εξίσωση της ευθείας $\zeta$ που εφάπτεται της παραβολής και είναι παράλληλη στην ευθεία $\varepsilon$. \monades{08}
        \item Ποια είναι η απόσταση των ευθειών $\zeta$ και $\varepsilon$; \monades{05}
    \end{rlist}\end{alist}

\vspace{4mm}

% --- Άσκηση 20229 (20229.tex) ---
\Askhsh{}{4} % The problem is ΘΕΜΑ 4. No specific code provided, so first arg is empty.

Δίνεται η εξίσωση: $x^2 + y^2 - (\lambda + 8)x + \lambda y + 7 = 0$ (1), με $\lambda \in \mathbb{R}$.

\begin{alist}    \item Να αποδείξετε ότι για κάθε $\lambda \in \mathbb{R}$ η εξίσωση (1) παριστάνει κύκλο, του οποίου να βρεθεί το κέντρο και η ακτίνα. \monades{6}
    \item Να βρείτε την εξίσωση της γραμμής πάνω στην οποία κινούνται τα κέντρα των κύκλων αυτών. \monades{7}
    \item Να αποδείξετε ότι για κάθε $\lambda \in \mathbb{R}$, όλοι οι παραπάνω κύκλοι, διέρχονται από δύο σταθερά σημεία, τα οποία και να βρεθούν. \monades{7}
    \item Θεωρούμε τον κύκλο που ορίζεται από την (1) για $\lambda=0$. Να βρεθούν τα σημεία του κύκλου αυτού, που απέχουν από την αρχή των αξόνων την ελάχιστη και την μέγιστη απόσταση αντίστοιχα. \monades{5}\end{alist}

\vspace{4mm}

% --- Άσκηση 20235 (20235.tex) ---
\Askhsh{}{2}
Δίνεται η παραβολή C: $y^2 = 8x$.
\begin{alist}    \item Να βρείτε την εστία και την διευθετούσα της παραβολής. \monades{10}
    \item Να αποδείξετε ότι η εφαπτομένη της παραβολής στο σημείο της $\left(\frac{1}{8}, 1\right)$ είναι παράλληλη στην ευθεία $\varepsilon: 8x - 2y + 3 = 0$. \monades{15}\end{alist}

\vspace{4mm}

% --- Άσκηση 20650 (20650.tex) ---
\textbf{ΘΕΜΑ 4}
\begin{alist}    \item Δίνονται οι ευθείες $\varepsilon1:y=x+2$, $\varepsilon2:y=x-2$ και τα σημεία $A(-2,0)$, $B(2,0)$ των ευθειών $\varepsilon1,\varepsilon2$ αντίστοιχα.
    \begin{rlist}
        \item Να αποδειχθεί ότι $\varepsilon1 // \varepsilon2$.
        \item Να βρεθούν οι συντεταγμένες του μέσου $M$, του $AB$.  \monades{04}
        \item Να βρεθεί η εξίσωση της μεσοπαραλλήλου των ευθειών $\varepsilon1,\varepsilon2$.  \monades{02}
 \monades{06}
    \end{rlist}
    \item Ο κύκλος $(K,\rho)$ έχει την ιδιότητα να εφάπτεται των ευθειών $\varepsilon1$ και $\varepsilon2$. Αν το κέντρο $K$ του κύκλου $(K,\rho)$ ανήκει στην ευθεία $(\eta): x=\lambda$, όπου $\lambda \in \mathbb{R}$, τότε:
    \begin{rlist}
        \item Να βρεθούν οι συντεταγμένες του κέντρου $K$, συναρτήσει του $\lambda$.
        \item Να αποδείξετε ότι η ακτίνα $\rho$ είναι ανεξάρτητη του $\lambda$ και να γράψετε την εξίσωση που παριστάνει όλους τους κύκλους $(K,\rho)$, για τις διάφορες τιμές του $\lambda \in \mathbb{R}$.  \monades{05}
 \monades{08}
    \end{rlist}\end{alist}

\vspace{4mm}

% --- Άσκηση 20651 (20651.tex) ---
\Askhsh{}{4}
Θεωρούμε τα σημεία $A(1, 1)$ και $B(2, 4)$.
\begin{alist}    \item Να βρείτε όλα τα σημεία $M$ στον άξονα $y'y$ ώστε το τρίγωνο $MAB$ να είναι ορθογώνιο με υποτείνουσα την $AB$. \monades{8}
    \item Να βρείτε την εξίσωση κύκλου $C$ με διάμετρο $AB$. \monades{7}
    \item Να αποδείξετε ότι ο κύκλος $C$ διέρχεται από τα σημεία $M$ που προσδιορίσατε στο ερώτημα (α). Κατόπιν, να το επιβεβαιώσετε γεωμετρικά. \monades{10}\end{alist}

\vspace{4mm}

% --- Άσκηση 20653 (20653.tex) ---
\Askhsh{}{4}
Κατά τη διάρκεια μιας επιχείρησης εντοπισμού ενός αγνοούμενου σε μια αχανή δασώδη επίπεδη περιοχή, δύο παρατηρητές $M_1$ και $M_2$ βρίσκονται σε διαφορετικά σημεία. Ο αγνοούμενος εκτοξεύει φωτοβολίδες που διαθέτει και οι δύο παρατηρητές σημειώνουν τις χρονικές στιγμές που ακούνε τον ήχο της εκπυρσοκρότησης του όπλου. Είναι γνωστό ότι ο παρατηρητής $M_1$ ακούει σε όλες τις εκρήξεις τον ήχο με διαφορά $4 \text{ sec}$ αργότερα από τον παρατηρητή $M_2$.

\begin{alist}    \item Αν ονομάσουμε $P$ την θέση του αγνοούμενου, να αποδείξετε ότι
    \item Να αποδείξτε ότι η θέση $P$ του αγνοούμενου ανήκει σε έναν κλάδο υπερβολής με εστίες τα σημεία $M_1$ και $M_2$. \monades{8}  \[ (PM_1) - (PM_2) = 1360 \text{ m}. \] Θεωρούμε ότι η ταχύτητα διάδοσης του ήχου είναι $340 \text{ m/sec.}$ \monades{8}
    \item Αν γνωρίζουμε ότι η απόσταση $(M_1M_2)$ είναι $1378 \text{ m}$, να αποδείξετε ότι αυτή η υπερβολή έχει εξίσωση
 \[ \frac{x^2}{680^2} - \frac{y^2}{111^2} = 1, \] θεωρώντας ως άξονα $x'x$ την ευθεία $M_1M_2$ και κέντρο της υπερβολής την αρχή των αξόνων. Δίνεται ότι $37^2 = 1369$. \monades{9}\end{alist}

% TikZ/PGFPlots Figure for the hyperbola
\begin{center}
\begin{tikzpicture}\begin{axis}[
    axis x line=none, 
    axis y line=none, 
    xmin=-4.5, xmax=4.5,
    ymin=-2.5, ymax=2.5,
    width=10cm, height=7cm,
    clip=false 
]
    % Draw custom x-axis and y-axis using the 'l3' style for thickness
    \draw[black, l3] (axis cs:xmin,0) -- (axis cs:xmax,0); % X-axis
    \draw[black, l3] (axis cs:0,ymin) -- (axis cs:0,ymax); % Y-axis
    % Define internal drawing parameters (a, b, c) for visualization purposes
    % These values are chosen to create a visually similar hyperbola to the PDF image,
    % not necessarily matching the problem's specific values (680, 111, 1378).
    \pgfmathsetmacro{\drawA}{3}
    \pgfmathsetmacro{\drawB}{2}
    \pgfmathsetmacro{\drawC}{sqrt(\drawA^2 + \drawB^2)} % Focal distance for drawing
    % Hyperbola branches (using 'thirdcolor' for green as per template)
    % Left branch
    \addplot[thirdcolor, very thick, domain=-\drawA-1.5:-\drawA, samples=100] { \drawB/\drawA * sqrt(x^2 - \drawA^2) };
    \addplot[thirdcolor, very thick, domain=-\drawA-1.5:-\drawA, samples=100] { -\drawB/\drawA * sqrt(x^2 - \drawA^2) };
    % Right branch
    \addplot[thirdcolor, very thick, domain=\drawA:\drawA+1.5, samples=100] { \drawB/\drawA * sqrt(x^2 - \drawA^2) };
    \addplot[thirdcolor, very thick, domain=\drawA:\drawA+1.5, samples=100] { -\drawB/\drawA * sqrt(x^2 - \drawA^2) };
    % Define coordinates for points M1, O, M2, P within the axis system
    \coordinate (M1_pt) at (axis cs:-\drawC,0);
    \coordinate (Opt) at (axis cs:0,0);
    \coordinate (M2_pt) at (axis cs:\drawC,0);
    % Choose P on the right branch; its x-coordinate is arbitrarily chosen for visual representation
    \pgfmathsetmacro{\Px}{\drawA+1.2} % Example x-coordinate for P
    \pgfmathsetmacro{\Py}{\drawB/\drawA * sqrt(\Px^2 - \drawA^2)} % Corresponding y-coordinate for P
    \coordinate (Ppt) at (axis cs:\Px, \Py);
    % Draw segments connecting M1 and M2 to P with dashed lines (using 'l3' style)
    \draw[dashed, black, l3] (M1_pt) -- (Ppt);
    \draw[dashed, black, l3] (M2_pt) -- (Ppt);
    % Draw points M1, O, M2, P and label them
    \fill (M1_pt) circle (1.5pt) node[below left] {$M1$};
    \fill (Opt) circle (1.5pt) node[below] {$O$};
    \fill (M2_pt) circle (1.5pt) node[below right] {$M2$};
    \fill (Ppt) circle (1.5pt) node[above right] {$P$};
\end{axis}\end{tikzpicture}
\end{center}

\vspace{4mm}

% --- Άσκηση 20655 (20655.tex) ---
\textbf{ΘΕΜΑ 4}

Δίνονται τα σημεία $A(2,1)$, $B(3,-1)$ και $\Gamma(-2,0)$.

\begin{alist}    \item
    \begin{rlist}
        \item Να αποδείξετε ότι τα σημεία $A$, $B$ και $\Gamma$ δεν είναι συνευθειακά. \monades{07}
        \item Να αποδείξετε ότι το εμβαδό του τριγώνου $AB\Gamma$ ισούται με $\frac{9}{2}$ τετραγωνικές μονάδες. \monades{03}
    \end{rlist}
    \item Να βρεθεί ο γεωμετρικός τόπος των σημείων $\Delta(x,y)$ για τα οποία ισχύει $(\Delta A\Gamma) = (AB\Gamma)$. \monades{07}
    \item Αν ο γεωμετρικός τόπος των σημείων $\Delta$ του ερωτήματος (β) αποτελείται από τις ευθείες $\varepsilon1:x-4y-7=0$ και $\varepsilon2 : x-4y+11=0$, τότε:
    \begin{rlist}
        \item Να αποδείξετε ότι οι ευθείες $A\Gamma$, $\varepsilon1$ και $\varepsilon2$ είναι παράλληλες. \monades{03}
        \item Να εξετάσετε αν είναι αληθής ή ψευδής ο ισχυρισμός « οι ευθείες $x-4y-7=0$ και $x-4y+11=0$ έχουν ως μεσοπαράλληλο την ευθεία $A\Gamma$ ». \monades{05}
    \end{rlist}\end{alist}

\vspace{4mm}

% --- Άσκηση 20658 (20658.tex) ---
{\OnPar{ΘΕΜΑ 2}}

Δίνεται η έλλειψη \((c)\) με εξίσωση
\[ \frac{x^2}{4} + \frac{y^2}{16} = 1. \]

\begin{alist}    \item Να δικαιολογήσετε ότι $\alpha = 4$, $\beta = 2$ και $\gamma = 2\sqrt{3}$. \monades{08}
    \item Να βρείτε τα μήκη των αξόνων και τις εστίες της έλλειψης \((c)\). \monades{08}
    \item Να σχεδιάσετε την έλλειψη \((c)\) και τον κύκλο $x^2 + y^2 = 16$ στο ίδιο σύστημα αξόνων. \monades{09}\end{alist}

\begin{center}
\begin{tikzpicture}\begin{axis}[
    axis lines=middle,
    xlabel=$x$,
    ylabel=$y$,
    xmin=-5, xmax=5,
    ymin=-5, ymax=5,
    xtick={-4,-2,0,2,4}, 
    ytick={-4,-2,0,2,4}, 
    grid=both,
    grid style={gray!30},
    samples=100,
    width=10cm, height=10cm, 
    tick label style={font=\scriptsize},
    every axis x label/.style={at={(current axis.right of origin)}, anchor=north west},
    every axis y label/.style={at={(current axis.above origin)}, anchor=south east},
]
    % Έλλειψη (c): x^2/4 + y^2/16 = 1 (vertical ellipse, a=4, b=2)
    % Parametric form: x = b*cos(t), y = a*sin(t) => x = 2*cos(t), y = 4*sin(t)
    \addplot[very thick, maincolor, domain=0:360, smooth] ({2*cos(x)}, {4*sin(x)});
    \node[maincolor, above right, font=\scriptsize] at (axis cs:0,4) {$(c)$};
    % Κύκλος: x^2 + y^2 = 16 (radius R=4)
    % Parametric form: x = R*cos(t), y = R*sin(t) => x = 4*cos(t), y = 4*sin(t)
    \addplot[very thick, secondarycolor, domain=0:360, smooth] ({4*cos(x)}, {4*sin(x)});
    \node[secondarycolor, above left, font=\scriptsize] at (axis cs:4,0) {$K$};
\end{axis}\end{tikzpicture}
\end{center}

\vspace{4mm}

% --- Άσκηση 20666 (20666.tex) ---
\textbf{ΘΕΜΑ 4}

Η τροχιά της Γης γύρω από τον Ήλιο είναι μια έλλειψη με μία εστία τον Ήλιο. Η ελάχιστη απόσταση του κέντρου της Γης από το κέντρο του Ήλιου είναι $PH = 147,5$ εκατομμύρια Km και η μέγιστη $AH = 152,5$ εκατομμύρια Km. Στο σχήμα θεωρούμε ότι τα σημεία $H$ και $\Gamma$ είναι τα κέντρα του Ήλιου και της Γης αντίστοιχα. Θεωρούμε ορθογώνιο σύστημα αξόνων με αρχή το μέσο του $HE$ και $x'x$ τον μεγάλο άξονα της έλλειψης, ενώ ο άξονας $y'y$ είναι η μεσοκάθετος του $HE$.

\begin{center}
\begin{tikzpicture}[scale=0.7]
    % Define parameters for drawing, not actual values (schematic representation)
    \pgfmathsetmacro{\aDraw}{5} % Semi-major axis for drawing
    \pgfmathsetmacro{\cDraw}{1} % Focal distance for drawing
    \pgfmathsetmacro{\bDraw}{sqrt(\aDraw^2 - \cDraw^2)} % Semi-minor axis for drawing
    % Origin
    \tkzDefPoint(0,0){O}
    % Foci H and E
    \tkzDefPoint(-\cDraw,0){H}
    \tkzDefPoint(\cDraw,0){E}
    % Vertices P and A
    \tkzDefPoint(-\aDraw,0){P}
    \tkzDefPoint(\aDraw,0){A}
    % Ellipse
    \draw[l3, thirdcolor!80!black] (O) ellipse (\aDraw cm and \bDraw cm);
    % Point Gamma on the ellipse (chosen for visual representation in the first quadrant)
    \pgfmathsetmacro{\xGamma}{\cDraw + 0.5} 
    \pgfmathsetmacro{\yGamma}{\bDraw * sqrt(1 - \xGamma^2 / \aDraw^2)}
    \tkzDefPoint(\xGamma,\yGamma){G} % Gamma
    % Axes (without specific labels, as per image)
    \draw[l3, gray!70] (-\aDraw-1,0) -- (\aDraw+1,0); % x-axis
    \draw[l3, gray!70] (0,-\bDraw-1) -- (0,\bDraw+1); % y-axis
    % Lines from foci to Gamma
    \tkzDrawSegments[l4, maincolor](H,G E,G)
    % Tangent line t't at G
    % Calculate slope of tangent at (xG, yG) for ellipse x^2/a^2 + y^2/b^2 = 1
    \pgfmathsetmacro{\slopeTan}{(-\bDraw^2 * \xGamma) / (\aDraw^2 * \yGamma)}
    % Define two points on the tangent line for drawing
    \pgfmathsetmacro{\xLeft}{\xGamma - 2.5}
    \pgfmathsetmacro{\yLeft}{\yGamma + \slopeTan * (\xLeft - \xGamma)}
    \pgfmathsetmacro{\xRight}{\xGamma + 2.5}
    \pgfmathsetmacro{\yRight}{\yGamma + \slopeTan * (\xRight - \xGamma)}
    \tkzDefPoint(\xLeft,\yLeft){TgLeft}
    \tkzDefPoint(\xRight,\yRight){TgRight}
    \draw[l3, secondarycolor!80!black] (TgLeft) -- (TgRight);
    % Labels for tangent t't (as per image)
    \node[above] at (TgLeft) {$t'$};
    \node[above] at (TgRight) {$t$};
    % Perpendicular dashed line from G to x-axis
    \draw[dashed, thirdcolor!80!black] (G) -- (\xGamma,0);
    % Labels for points
    \tkzLabelPoint[below](O){O}
    \tkzLabelPoint[left](P){P}
    \tkzLabelPoint[below left](H){H}
    \tkzLabelPoint[below right](E){E}
    \tkzLabelPoint[right](A){A}
    \tkzLabelPoint[above right](G){$\Gamma$}
    % Angle marks for t'ΓH and tΓE (as per image)
    \tkzMarkAngle[size=0.8cm,fill=secondarycolor!20,draw=secondarycolor,opacity=0.7](TgLeft,G,H)
    \tkzMarkAngle[size=0.8cm,fill=secondarycolor!20,draw=secondarycolor,opacity=0.7](E,G,TgRight)
    % Draw the points last so they are on top
    \tkzDrawPoints(P,H,O,E,A,G)\end{tikzpicture}
\end{center}

\begin{alist}    \item Να αποδείξετε $(PA) = 300$ εκατομμύρια Km, $(HE) = 5$ εκατομμύρια Km και ότι η εκκεντρότητα της έλλειψης είναι $\varepsilon = \frac{1}{60}$. \monades{10}
    \item Για μια τυχαία θέση της Γης πάνω στην ελλειπτική τροχιά, να υπολογίσετε την περίμετρο του τριγώνου $H\Gamma E$. \monades{8}
    \item Αν ονομάσουμε $t't$ την εφαπτομένη ευθεία της έλλειψης στο $\Gamma$, να αποδείξετε ότι οι γωνίες $\tau' \Gamma H$ και $\tau \Gamma E$ είναι ίσες. \monades{7}\end{alist}

\vspace{4mm}

% --- Άσκηση 20684 (20684.tex) ---
\Askhsh{XXXXX}{4}
\noindent Ένα σημείο $A(x_A, y_A)$ της παραβολής $C: y^2 = 4x$ με $x_A > 0$, $y_A > 0$, έχει την εξής ιδιότητα: η ημιευθεία $AE$ τέμνει την διευθετούσα $(\delta)$ στο σημείο $\Gamma$, έτσι όμως ώστε η εστία $E$ της παραβολής $C$, να είναι το μέσο του τμήματος $A\Gamma$. Επίσης, από το σημείο $A$ φέρνουμε κάθετη στην διευθετούσα $(\delta)$ και έστω $B$ το σημείο τομής, όπως δείχνει το παρακάτω σχήμα.

\begin{alist}    \item Να αποδείξετε ότι το τρίγωνο $ABE$ είναι ισόπλευρο. \monades{7}
    \item Να αποδείξετε ότι $x_A = 3$ και $y_A = 2\sqrt{3}$. \monades{10}
    \item Να βρείτε την εξίσωση του κύκλου που διέρχεται από τις κορυφές του τριγώνου $AB\Gamma$. \monades{8}\end{alist}

\begin{center}
\begin{tikzpicture}\begin{axis}[
    axis lines=middle,
    xmin=-2.5, xmax=4.5,
    ymin=-3.5, ymax=3.5,
    xtick={-1,1,3},
    ytick={-2*sqrt(3), 2*sqrt(3)},
    xticklabels={},
    yticklabels={},
    grid=none,
    width=10cm, height=7cm,
    clip=false,
    axis line style={-Latex, l3, maincolor!80!black},
    tick label style={font=\scriptsize},
]
    % Axes labels (based on the image's placement of x, y, x', y')
    \node[anchor=west, xshift=3mm] at (axis cs:4.5,0) {$x$}; % Positive x-axis label
    \node[anchor=south, yshift=3mm] at (axis cs:0,3.5) {$y$}; % Positive y-axis label
    \node[anchor=east, xshift=-3mm] at (axis cs:-2.5,0) {$x'$}; % Negative x-axis label
    \node[anchor=north, yshift=-3mm] at (axis cs:0,-3.5) {$y'$}; % Negative y-axis label
    % Draw the parabola C: y^2 = 4x (i.e. x = y^2/4)
    \addplot[very thick, blue, domain=-3.5:3.5, samples=100] ({y^2/4},y);
    \node[above right, blue, xshift=5mm, yshift=5mm] at (axis cs:3, 2.5) {$C : y^2 = 4x$};
    % Draw the directrix (δ): x = -1
    \addplot[very thick, red!80!black, domain=-3.5:3.5, dashed] coordinates {(-1,y)};
    \node[above left, red!80!black] at (axis cs:-1, 3.5) {$(\delta)$};
    % Define points using calculated coordinates
    \coordinate (O) at (axis cs:0,0);
    \coordinate (E) at (axis cs:1,0); % Focus of y^2 = 4x is (1,0)
    \coordinate (A) at (axis cs:3, {2*sqrt(3)}); % A(3, 2sqrt(3)) based on problem statement
    \coordinate (B) at (axis cs:-1, {2*sqrt(3)}); % B is projection of A onto directrix x=-1
    \coordinate (G) at (axis cs:-1, {-2*sqrt(3)}); % Gamma is on directrix, E is midpoint of AG
    % Draw construction line A-Gamma (black, l3)
    \draw[l3, black] (A) -- (G);
    % Draw the triangle ABE (thirdcolor, l4)
    \draw[l4, thirdcolor] (A) -- (B) -- (E) -- cycle; % Sides AB, BE, EA
    % Mark the right angle at B, between AB and the directrix.
    % G is on the directrix below B, so (A,B,G) marks the angle at B between BA and BG.
    \tkzMarkRightAngle[fill=gray!30,size=0.2](A,B,G);
    % Draw points
    \tkzDrawPoints[color=black](O,A,B,E,G)
    % Labels for points
    \tkzLabelPoint[below left](O){O}
    \tkzLabelPoint[right](E){E}
    \tkzLabelPoint[above right](A){A}
    \tkzLabelPoint[above left](B){B}
    \tkzLabelPoint[below left](G){$\Gamma$}
\end{axis}\end{tikzpicture}
\end{center}

\vspace{4mm}

% --- Άσκηση 20685 (20685.tex) ---
\OnPar{ΘΕΜΑ 2}
Δίνονται τα διανύσματα $\vec{u} = (1,1)$, $\vec{w} = (-10,2)$ και τα σημεία $A(-1,2)$, $B(\beta, 0)$, $\Gamma(0, \gamma)$. Τα διανύσματα $\vec{u}$, $\vec{AB}$ είναι κάθετα και το διάνυσμα $\vec{w}$ είναι παράλληλο στο διάνυσμα $\vec{A\Gamma}$.

\begin{alist}    \item Να βρείτε τις συντεταγμένες του διανύσματος $\vec{AB}$ και να αποδείξετε ότι $\beta = 1$. \monades{9}
    \item Να βρείτε τις συντεταγμένες του διανύσματος $\vec{A\Gamma}$ και να αποδείξετε ότι $\gamma = \frac{9}{5}$. \monades{9}
    \item Να υπολογίσετε το εσωτερικό γινόμενο $\vec{AB} \cdot \vec{A\Gamma}$. \monades{7}\end{alist}

\vspace{4mm}

% --- Άσκηση 20700 (20700.tex) ---
\Askhsh{XXXX}{4}
Δίνεται το τετράγωνο $MM_1OM_2$ με $M(4,4)$, $M_1(4,0)$, $M_2(0,4)$. Αν $O$ η αρχή των αξόνων του καρτεσιανού συστήματος συντεταγμένων, τότε:
\begin{alist}    \item Να δείξετε ότι ο κύκλος που διέρχεται από τις κορυφές του τετραγώνου $MM_1OM_2$ έχει εξίσωση $C: (x – 2)^2 + (y − 2)^2 = 8$. \monades{8}
    \item Να αποδείξετε ότι η ευθεία $\varepsilon: x + y = 8$ είναι εφαπτομένη του παραπάνω κύκλου $C$. \monades{8}
    \item Να βρείτε το σημείο επαφής της ευθείας $\varepsilon$ με τον κύκλο $C$. \monades{9}\end{alist}

\vspace{4mm}

% --- Άσκηση 20718 (20718.tex) ---
\Askhsh{ELLIPSE01}{2}
Δίνεται η έλλειψη $C$ με εξίσωση
\[ \frac{x^2}{9} + \frac{y^2}{25} = 1. \]
\begin{alist}    \item Να βρείτε τις εστίες της. \monades{8}
    \item Να σχεδιάσετε την έλλειψη $C$ σε ορθοκανονικό σύστημα συντεταγμένων. \monades{8}
    \item Να σχεδιάσετε στο ίδιο ορθοκανονικό σύστημα τις εφαπτόμενες στις κορυφές της $C$ και να γράψετε τις εξισώσεις τους. \monades{9}\end{alist}

\vspace{4mm}

% --- Άσκηση 20721 (20721.tex) ---
\Askhsh{HYPERBOLA001}{2}
Δίνεται η υπερβολή $C$ με εξίσωση
\[ \frac{x^2}{9} - \frac{y^2}{16} = 1. \]
\begin{alist}    \item Να βρείτε τις εστίες της $C$. \monades{8}
    \item Να βρείτε τις εξισώσεις των ασυμπτώτων της $C$. \monades{8}
    \item Να σχεδιάσετε την υπερβολή $C$ και τις ασύμπτωτές της στο ίδιο σύστημα αξόνων. \monades{9}
    \begin{center}
    \begin{tikzpicture}    \begin{axis}[
 axis lines = middle, xlabel = {$x$}, ylabel = {$y$}, xmin = -8, xmax = 8, ymin = -8, ymax = 8, xtick = {-7,-6,-5,-4,-3,-2,-1,0,1,2,3,4,5,6,7}, ytick = {-7,-6,-5,-4,-3,-2,-1,0,1,2,3,4,5,6,7}, grid = both, grid style = {gray!30}, samples = 100, width = 12cm, height = 9cm, tick label style = {font=\scriptsize}, every axis x label/.style = {at={(current axis.right of origin)}, anchor=north west}, every axis y label/.style = {at={(current axis.above origin)}, anchor=south east}, axis on top, ]
        % Hyperbola: y = +/- (4/3)sqrt(x^2 - 9)
 \addplot[very thick, maincolor, domain=3:8, samples=200] {4/3 * sqrt(x^2 - 9)}; \addplot[very thick, maincolor, domain=3:8, samples=200] {-4/3 * sqrt(x^2 - 9)}; \addplot[very thick, maincolor, domain=-8:-3, samples=200] {4/3 * sqrt(x^2 - 9)}; \addplot[very thick, maincolor, domain=-8:-3, samples=200] {-4/3 * sqrt(x^2 - 9)}; \node[maincolor, above] at (axis cs:6, {4/3 * sqrt(6^2-9)}) {$C$}; % Label for hyperbola
        % Asymptotes: y = +/- (4/3)x
 \addplot[dashed, secondarycolor, domain=-8:8] {4/3 * x}; \addplot[dashed, secondarycolor, domain=-8:8] {-4/3 * x}; \node[secondarycolor, above left] at (axis cs:-7.5, 4/3 * 7.5 + 0.5) {$y = -\frac{4}{3}x$}; \node[secondarycolor, above right] at (axis cs:7.5, 4/3 * 7.5 + 0.5) {$y = \frac{4}{3}x$};
        % Foci F'(-5,0) and F(5,0)
 \addplot[only marks, mark=*, mark size=2pt, thirdcolor] coordinates {(-5,0) (5,0)}; \node[thirdcolor, below] at (axis cs:-5,0) {$F'$}; \node[thirdcolor, below] at (axis cs:5,0) {$F$};
        % Vertices A'(-3,0) and A(3,0)
 \addplot[only marks, mark=*, mark size=2pt, black] coordinates {(-3,0) (3,0)}; \node[black, below] at (axis cs:-3,0) {$A'$}; \node[black, below] at (axis cs:3,0) {$A$};
    \end{axis}\end{tikzpicture}
    \end{center}\end{alist}

\vspace{4mm}

% --- Άσκηση 20722 (20722.tex) ---
\textbf{ΘΕΜΑ 4}

Έστω $K(x, y)$ μεταβλητό σημείο του επιπέδου για το οποίο ισχύει $(KE) + (KE') = 10$, όπου $E(3,0)$ και $E'(-3,0)$.

\begin{alist}    \item Να βρείτε το είδος της καμπύλης $C$ πάνω στην οποία κινείται το σημείο $K$ και να γράψετε την εξίσωσή της, αιτιολογώντας την απάντησή σας.
    \item Έστω $C: \frac{x^2}{25} + \frac{y^2}{16} = 1$ και $(\varepsilon): 3x+5y = 25$.  \monades{6}
    \item Να ερμηνεύσετε γραφικά το συμπέρασμα του ερωτήματος β) και να σχεδιάσετε στο ίδιο ορθοκανονικό σύστημα την έλλειψη $C$ και την ευθεία $\varepsilon$.  Να αποδείξετε ότι $C$ και $(\varepsilon)$ έχουν ένα μόνο κοινό σημείο $M$ και να βρείτε τις συντεταγμένες του σημείου $M$. \monades{7}
    \begin{center}
    \begin{tikzpicture}    \begin{axis}[
 axis lines = middle, xlabel = {$x$}, ylabel = {$y$}, xmin = -6, xmax = 6, ymin = -5, ymax = 5, xtick = {-5,-4,-3,-2,-1,0,1,2,3,4,5}, ytick = {-4,-3,-2,-1,0,1,2,3,4}, grid = both, grid style = {gray!30}, samples = 100, width = 10cm, height = 7cm, tick label style = {font=\scriptsize}, every axis x label/.style = {at={(current axis.right of origin)}, anchor=north west}, every axis y label/.style = {at={(current axis.above origin)}, anchor=south}, clip=false, ]
        % Ellipse C: x^2/25 + y^2/16 = 1 (parametric: x=5cos(t), y=4sin(t))
 \addplot[very thick, maincolor, domain=0:360, samples=200, smooth] ({5*cos(x)}, {4*sin(x)}); \node[maincolor, above right] at (axis cs:4,3.5) {$C$};
        % Line epsilon: 3x+5y = 25 => y = (25-3x)/5 = 5 - (3/5)x
 \addplot[very thick, secondarycolor, domain=-6:6] {5 - (3/5)*x}; \node[secondarycolor, above left] at (axis cs:-4,7.4) {$\varepsilon$};
        % Point M(3, 16/5) = M(3, 3.2)
 \addplot[only marks, mark=*, mark size=2.5pt, thirdcolor] coordinates {(3, 3.2)}; \node[thirdcolor, below right] at (axis cs:3,3.2) {$M(3, 3.2)$};
    \end{axis}\end{tikzpicture}
    \end{center}
    \item Να σχεδιάσετε τη διχοτόμο της γωνίας $EME'$ και να βρείτε την εξίσωσή της.  \monades{6}
    \begin{center}
    \begin{tikzpicture}[scale=1]
        % Define points E, E', M
 \tkzDefPoint(3,0){E} \tkzDefPoint(-3,0){E'} \tkzDefPoint(3,3.2){M} % Using 16/5 = 3.2
        % Draw segments ME and ME'
 \tkzDrawSegment[l3,maincolor](M,E) \tkzDrawSegment[l3,maincolor](M,E')
        % Draw points
 \tkzDrawPoints(E,E',M)
        % Labels for points
 \tkzLabelPoint[below](E){$E(3,0)$} \tkzLabelPoint[below left](E'){$E'(-3,0)$} \tkzLabelPoint[above left](M){$M(3, \frac{16}{5})$}
        % Mark the angle EME'
 \tkzMarkAngle[arc=ll,size=0.6cm,fill=maincolor!10,opacity=0.5](E,M,E')
        % Construct and draw the angle bisector (internal)
 \tkzDefLine[bisector](E,M,E') \tkzGetPoint{bisectorptstart} % First point defining the bisector line \tkzGetPoint{bisectorptend} % Second point defining the bisector line \tkzDrawLine[add=0.5 and 0.5, l3, secondarycolor, dashed](bisectorptstart,bisectorptend) % Draw the bisector line \node[secondarycolor, right] at ($(M)!0.8!(bisectorptend)$) {διχοτόμος}; % Label the bisector
\end{tikzpicture}
    \end{center}
 \monades{6}\end{alist}

\vspace{4mm}

% --- Άσκηση 20724 (20724.tex) ---
\Askhsh{MATH}{4} % Αν δεν υπάρχει κωδικός, χρησιμοποιούμε MATH
Η ευθεία $\varepsilon$ με εξίσωση $x+y-1=0$ του παρακάτω σχήματος, αναπαριστά τη γραμμή ενός σιδηροδρομικού δικτύου που εξυπηρετεί τους κατοίκους δύο πόλεων $A(8,1)$, $B(-7,4)$ (για την ακρίβεια $A$, $B$ είναι τα κεντρικά σημεία των πόλεων από τα οποία μετράμε αποστάσεις). Για το λόγο αυτό θα κατασκευαστεί κατά μήκος της γραμμής $(\varepsilon)$, ένας σταθμός σε ένα σημείο $\Sigma$ και μία πεζογέφυρα σε ένα σημείο $\Pi$.
Να βρείτε:
\begin{alist}    \item ποια πόλη από τις $A$, $B$ είναι πλησιέστερα στη γραμμή του τραίνου.
    \item τις συντεταγμένες του $\Pi$, αν είναι γνωστό ότι θα κατασκευαστεί στο πλησιέστερο σημείο της γραμμής στην πόλη $B$.  \monades{6}
    \item τις συντεταγμένες του $\Sigma$ στις παρακάτω περιπτώσεις  \monades{7}
    \begin{rlist}
        \item ο σταθμός $\Sigma$ να ισαπέχει από τις πόλεις $A,B$.
        \item το οδικό δίκτυο που θα συνδέει το σταθμό $\Sigma$ με τις πόλεις $A,B$ να έχει το μικρότερο δυνατό μήκος.  \monades{6}
 \monades{6}
    \end{rlist}\end{alist}

% Γεωμετρικό σχήμα
\begin{tikzpicture}    % Ορισμός σημείων Α και Β
    \tkzDefPoint(8,1){A}
    \tkzDefPoint(-7,4){B}
    % Η εξίσωση της ευθείας είναι $x+y-1=0 \implies y = -x+1$.
    % Ορίζουμε δύο σημεία για να σχεδιάσουμε την ευθεία.
    \tkzDefPoint(-10,11){L1} % y = -(-10)+1 = 11
    \tkzDefPoint(10,-9){L2}  % y = -(10)+1 = -9
    % Σχεδίαση της ευθείας ε
    \tkzDrawLine[add=0.5 and 0.5, color=maincolor, l3](L1,L2)
    % Τοποθέτηση ετικέτας για την ευθεία ε
    \node[below right, maincolor] at (L2) {$\varepsilon$};
    % Σχεδίαση των σημείων Α και Β. Χρησιμοποιούμε maincolor για μπλε κουκκίδες.
    \tkzDrawPoints[color=maincolor, fill=maincolor, size=3pt](A,B)
    % Τοποθέτηση ετικετών για τα σημεία Α και Β
    \tkzLabelPoint[below right](A){A}
    \tkzLabelPoint[above left](B){B}\end{tikzpicture}

\vspace{4mm}

% --- Άσκηση 20726 (20726.tex) ---
\Askhsh{00000}{4} % Κωδικός: 00000, ΘΕΜΑ 4
Ένας κατασκευαστής μπιλιάρδων θέλει να κατασκευάσει ένα ελλειπτικό μπιλιάρδο όπως αυτό του παρακάτω σχήματος (σχήμα 1). Το περίγραμμα του μπιλιάρδου είναι έλλειψη με εστίες τα σημεία $E(3,0)$ και $E'(-3,0)$. Η μοναδική τρύπα του μπιλιάρδου έχει σχήμα κύκλου (ο μαύρος κύκλος στο σχήμα 1) με κέντρο το σημείο $E'$. Για να σχεδιάσει ο κατασκευαστής το περίγραμμα του μπιλιάρδου πάνω σε μία ξύλινη επίπεδη επιφάνεια, τοποθέτησε στα σημεία $E$ και $E'$ δύο καρφιά στα οποία έδεσε τις άκρες ενός σχοινιού μήκους 10 μονάδων μήκους. Στη συνέχεια με ένα μολύβι διατηρούσε το σχοινί τεντωμένο, ώστε αυτό, κατά την κίνησή του, να διαγράψει έλλειψη $C$ όπως φαίνεται στο παρακάτω (σχήμα 2).

\begin{alist}    \item Να βρείτε τα μήκη του μεγάλου και του μικρού άξονα της έλλειψης $C$. \monades{10}
    \item Να γράψετε την εξίσωση της έλλειψης $C$ και να βρείτε την εκκεντρότητά της. \monades{5}
    \item Ένας παίκτης τοποθετεί μια άσπρη μπάλα (ο άσπρος κύκλος στο σχήμα 1) ακριβώς στο σημείο $E$. Σκοπεύει να χτυπήσει την άσπρη μπάλα ώστε αφού αυτή προσκρούσει πρώτα στο ελλειπτικό περίγραμμα του μπιλιάρδου, στη συνέχεια να πέσει στην τρύπα. Αν θεωρήσουμε ότι ο παίκτης θα χτυπήσει με όση δύναμη απαιτείται για να φτάσει η μπάλα στην τρύπα και το χτύπημα θα είναι στο κέντρο της μπάλας ώστε αυτή να κυλά χωρίς να περιστρέφεται, να βρείτε σε ποιο σημείο της έλλειψης $C$ πρέπει να σημαδέψει, ώστε με ένα μόνο χτύπημα η μπάλα να μπει στην τρύπα:
    \begin{rlist}
        \item μόνο στα άκρα του μεγάλου άξονα
        \item μόνο στα άκρα του μικρού άξονα
        \item μόνο στα άκρα του μικρού άξονα και στο ένα άκρο του μεγάλου άξονα
        \item σε οποιοδήποτε σημείο της $C$ εκτός από το ένα άκρο του μεγάλου άξονα
    \end{rlist}
 Επιλέξτε τη μοναδική σωστή απαντήσεις αιτιολογώντας την απάντησή σας. \monades{10}\end{alist}

    \begin{tabular}{cc}
        \begin{tikzpicture}[scale=0.8]
            % Dimensions for the ellipse based on problem statement
            \def\a{5} % semi-major axis (2a = 10 from string length)
            \def\c{3} % focal distance (foci at (c,0) and (-c,0))
            \pgfmathsetmacro{\b}{sqrt(\a*\a - \c*\c)} % semi-minor axis (b^2 = a^2 - c^2)
            % Calculate the bounding box for the border around the ellipse
            \pgfmathsetmacro{\bboxWidth}{2*\a + 0.8} % Example margin
            \pgfmathsetmacro{\bboxHeight}{2*\b + 0.8} % Example margin
            % Outer wooden border (rounded rectangle shape)
            \fill[brown!80!black, rounded corners=3mm] (-\bboxWidth/2, -\bboxHeight/2) rectangle (\bboxWidth/2, \bboxHeight/2);
            % Inner green felt (ellipse)
            \fill[thirdcolor] (0,0) ellipse (\a cm and \b cm);
            % Hole (black circle at E'(-3,0))
            \fill[black] (-\c,0) circle (0.25cm);
            % Ball (white circle at E(3,0))
            \fill[white] (\c,0) circle (0.25cm);
            \draw[black, thin] (\c,0) circle (0.25cm); % outline for white ball
\end{tikzpicture}
        &
        \begin{tikzpicture}[scale=0.8]
            % Ellipse parameters
            \def\a{5}
            \def\c{3}
            \pgfmathsetmacro{\b}{sqrt(\a*\a - \c*\c)}
            % Draw the ellipse
            \draw[l3, secondarycolor] (0,0) ellipse (\a cm and \b cm);
            % Foci E' and E
            \tkzDefPoint(-\c,0){E'}
            \tkzDefPoint(\c,0){E}
            \tkzDrawPoints[fill=black](E,E')
            \tkzLabelPoint[below left](E'){$E'$}
            \tkzLabelPoint[below right](E){$E$}
            % Point M on the ellipse (example point from image, upper left quadrant)
            \pgfmathsetmacro{\Mx}{-2} % example x-coordinate for M
            \pgfmathsetmacro{\My}{\b*sqrt(1-(\Mx/\a)^2)} % Calculate y-coordinate on ellipse
            \tkzDefPoint(\Mx,\My){M}
            \tkzDrawPoint[fill=black](M)
            \tkzLabelPoint[above](M){$M$}
            % String segments (E'M and EM)
            \draw[dashed, maincolor] (E') -- (M);
            \draw[dashed, maincolor] (E) -- (M);
            % Pencil representation (approximating visual from image)
            \coordinate (PencilTop) at ($(M) + (140:1.5)$); % End of pencil body
            \coordinate (PencilLeftCorner) at ($(M) + (50:0.08)$); % Bottom-left corner of body
            \coordinate (PencilRightCorner) at ($(M) + (-130:0.08)$); % Bottom-right corner of body
            % Draw pencil body
            \fill[gray!70!black] (PencilLeftCorner) -- ($(PencilTop) + (50:0.08)$) -- ($(PencilTop) + (-130:0.08)$) -- (PencilRightCorner) -- cycle;
            % Draw pencil tip
            \fill[black] (M) -- (PencilLeftCorner) -- (PencilRightCorner) -- cycle;
\end{tikzpicture}


\vspace{4mm}

% --- Άσκηση 20728 (20728.tex) ---
\Askhsh{}{4}

Δίνονται οι ευθείες $\varepsilon1 : y = \frac{\sqrt{3}}{3}x$ και $\varepsilon2 : y = x$.

\begin{alist}    \item Να σχεδιάσετε τις $\varepsilon1,\varepsilon2$ στο ίδιο ορθοκανονικό σύστημα αξόνων. \monades{6}
    \begin{center}
    \begin{tikzpicture}    \begin{axis}[
 axis lines = middle, xlabel = {$x$}, ylabel = {$y$}, xmin = -1, xmax = 5, ymin = -1, ymax = 5, xtick = {1,2,3,4,5}, ytick = {1,2,3,4,5}, grid = both, grid style = {gray!30}, samples = 100, width = 8cm, height = 8cm, tick label style = {font=\scriptsize}, every axis x label/.style = {at={(current axis.right of origin)}, anchor=north west}, every axis y label/.style = {at={(current axis.above origin)}, anchor=south}, axis line style = {-Latex}, ]
        % Line epsilon1: y = sqrt(3)/3 * x (approx 0.577x)
 \addplot[very thick, maincolor, domain=-1:5] {sqrt(3)/3 * x}; \node[maincolor, anchor=south west] at (axis cs:4, {sqrt(3)/3 * 4}) {$\varepsilon1: y=\frac{\sqrt{3}}{3}x$};
        % Line epsilon2: y = x
 \addplot[very thick, secondarycolor, domain=-1:5] {x}; \node[secondarycolor, anchor=south west] at (axis cs:4, 4) {$\varepsilon2: y=x$};
        % Origin point label
 \node[below left] at (axis cs:0,0) {O};
    \end{axis}\end{tikzpicture}
    \end{center}
    \item Να βρείτε τη γωνία που σχηματίζει κάθε μια με τον άξονα $x'x$. \monades{6}
    \item Να βρείτε το εμβαδόν του τριγώνου ΟΑΒ, όπου $Ο(0,0)$, $A(3,\sqrt{3})$, $B(3,3)$. \monades{6}
    \item Να αποδείξετε ότι $\eta\mu15^\circ = \frac{\sqrt{3}-1}{2\sqrt{2}}$. \monades{7}\end{alist}

(Θυμίζουμε ότι το εμβαδόν ενός τριγώνου δίνεται από το ημιγινόμενο δύο πλευρών του επί το ημίτονο της περιεχόμενης γωνίας τους).

\vspace{4mm}

% --- Άσκηση 20732 (20732.tex) ---
\Askhsh{}{2}
Δίνονται τα διανύσματα $\vec{a} = (2, 1)$ και $\vec{b} = (-8, -4)$.
\begin{alist}    \item Να δείξετε ότι τα διανύσματα $\vec{a}$, $\vec{b}$ είναι αντίρροπα και ότι $|\vec{b}| = 4|\vec{a}|.$
    \item Να βρείτε τη γωνία που σχηματίζουν τα διανύσματα $\vec{a}$, $\vec{b}$ και να δικαιολογήσετε την απάντησή σας.  \monades{12}
    \item Να δείξετε ότι $\vec{a} \cdot \vec{b} < 0.$  \monades{6}
 \monades{7}\end{alist}

\vspace{4mm}

% --- Άσκηση 20733 (20733.tex) ---
\Askhseis{ΘΕΜΑ 2}
Δίνονται τα μη μηδενικά διανύσματα $\vec{a}$, $\vec{b}$, με $|\vec{a}| = |\vec{b}|$ και $\vec{AB} = \vec{a} - \vec{b}$ και $\vec{A\Gamma} = \vec{a} + \vec{b}$.

\begin{alist}    \item Να εκφράσετε το διάνυσμα $\vec{B\Gamma}$ συναρτήσει του διανύσματος $\vec{b}$. \monades{10}
    \item Να υπολογίσετε το εσωτερικό γινόμενο $\vec{AB} \cdot \vec{A\Gamma}$. \monades{10}
    \item Να αιτιολογήσετε γιατί τα διανύσματα $\vec{AB}$ και $\vec{A\Gamma}$ είναι κάθετα. \monades{05}\end{alist}

\vspace{4mm}

% --- Άσκηση 20773 (20773.tex) ---
\Askhsh{}{2}
Δίνονται τα διανύσματα $\vec{\alpha} = (1, -2)$ και $\vec{\beta} = (2, 3)$.

\begin{alist}    \item Να βρείτε το διάνυσμα $\vec{u}=2\vec{\alpha}+\vec{\beta}$. \monades{08}
    \item Αν $\vec{u} = (4,-1)$ να βρείτε την τιμή του $\kappa \in \mathbb{R}$ ώστε το διάνυσμα $\vec{u}$ να είναι κάθετο στο διάνυσμα $\vec{v} = (1, \kappa)$. \monades{09}
    \item Για $\kappa = 4$ να υπολογίσετε το μέτρο του διανύσματος $\vec{v}$ του προηγούμενου ερωτήματος. \monades{08}\end{alist}

\vspace{4mm}

% --- Άσκηση 20861 (20861.tex) ---
\Askhsh{}{4}

Δίνεται το σημείο $M(-2, 2)$.

\begin{alist}    \item Να βρείτε τις εξισώσεις των ευθειών $(\varepsilon)$ που διέρχονται από το σημείο $M$.
    \item  \monades{6}
    \begin{rlist}
        \item Να βρείτε ποιες από τις παραπάνω εξισώσεις ευθειών σχηματίζουν τρίγωνο με τον αρνητικό ημιάξονα $Ox'$ και τον θετικό ημιάξονα $Oy$.

        \item Να βρείτε την εξίσωση της ευθείας $(\varepsilon1)$, η οποία διέρχεται από το σημείο $M$ και σχηματίζει με τον αρνητικό ημιάξονα $Ox'$ και τον θετικό ημιάξονα $Oy$ τρίγωνο, με εμβαδόν $E = 8$.  \monades{4}
 \monades{10}
    \end{rlist}
    \item Αν $(\varepsilon1): y = x + 4$, να βρείτε το μήκος του ύψους του ορθογωνίου τριγώνου, που σχηματίζει η $(\varepsilon1)$ με τους άξονες, το οποίο φέρεται από την κορυφή $O$.
 \monades{5}\end{alist}

\vspace{4mm}

% --- Άσκηση 20862 (20862.tex) ---
\Askhsh{A}{4}
Δίνονται τα σημεία $M(-2, 2)$, $E(0, -\frac{1}{2})$ και η ευθεία $(\zeta)$ με εξίσωση $y = \frac{1}{2}$.

\begin{alist}    \item Να βρείτε την εξίσωση της ευθείας $(\varepsilon1)$ που διέρχεται από το σημείο $M$ και σχηματίζει γωνία $45^\circ$ με τον άξονα $x'x$. \monades{05}
    \item Να βρείτε την εξίσωση, που εκφράζει το σύνολο των σημείων του επιπέδου, που απέχουν ίση απόσταση από το σημείο $E$ και την ευθεία $(\zeta)$. \monades{06}
    \item 
    \begin{rlist}
        \item Να βρείτε την εξίσωση της εφαπτομένης $(\eta)$ της καμπύλης $C: x^2 + 2y = 0$, που είναι παράλληλη στην ευθεία $(\varepsilon1)$, με εξίσωση $y = x + 4$. \monades{07}
        
        \item Να σχεδιάσετε την γραφική παράσταση της καμπύλης $C$ και των ευθειών $(\varepsilon1)$ και $(\eta)$. Με τη βοήθεια του σχήματος (ή με οποιονδήποτε άλλον τρόπο) να αποδείξετε ότι η ελάχιστη απόσταση των σημείων της $C$ από την ευθεία $(\varepsilon1)$ είναι $\frac{7\sqrt{2}}{4}$. \monades{07}
        \begin{center}
        \begin{tikzpicture}        \begin{axis}[
 axis lines = middle, xlabel = {$x$}, ylabel = {$y$}, xmin = -6, xmax = 6, ymin = -6, ymax = 6, xtick = {-5,-4,-3,-2,-1,0,1,2,3,4,5}, ytick = {-5,-4,-3,-2,-1,0,1,2,3,4,5}, grid = both, grid style = {gray!30}, samples = 100, width = 10cm, height = 7cm, tick label style = {font=\scriptsize}, every axis x label/.style = {at={(current axis.right of origin)}, anchor=north west}, every axis y label/.style = {at={(current axis.above origin)}, anchor=south}, ]
            % Καμπύλη C: x^2 + 2y = 0  <=> y = -1/2 x^2
 \addplot[very thick, maincolor, domain=-5:5] {-0.5*x^2}; \node[maincolor, above right] at (axis cs:2,-2) {$C$};
            % Ευθεία ε1: y = x + 4 (derived in part a)
 \addplot[very thick, secondarycolor, domain=-6:2] {x+4}; \node[secondarycolor, above left] at (axis cs:-2,2) {$\varepsilon1$};
            % Ευθεία η: y = x + 1/2 (derived in thought process for tangency with C)
 \addplot[very thick, thirdcolor, domain=-6:5.5] {x+0.5}; \node[thirdcolor, below right] at (axis cs:2,2.5) {$\eta$};
            % Σημεία M(-2,2)
 \addplot[only marks, mark=*, mark size=1.5pt, black] coordinates {(-2,2)}; \node[above left] at (axis cs:-2,2) {$M(-2, 2)$};
            % Σημεία E(0,-1/2)
 \addplot[only marks, mark=*, mark size=1.5pt, black] coordinates {(0,-0.5)}; \node[below right] at (axis cs:0,-0.5) {$E(0, -\frac{1}{2})$};
            % Ευθεία ζ: y = 1/2
 \addplot[dashed, gray, domain=-6:6] {0.5}; \node[gray, right] at (axis cs:6,0.5) {$\zeta$};
            % Tangency point for eta and C (-1, -1/2)
 \addplot[only marks, mark=*, mark size=1.5pt, thirdcolor] coordinates {(-1,-0.5)}; \node[thirdcolor, above left] at (axis cs:-1,-0.5) {$(-1, -\frac{1}{2})$};
        \end{axis}\end{tikzpicture}
        \end{center}
    \end{rlist}\end{alist}

\vspace{4mm}

% --- Άσκηση 20863 (20863.tex) ---
\Askhsh{}{4}
Δίνονται τα σημεία $A(1,0)$ και $B(3,0)$.

\begin{alist}    \item Να βρείτε την εξίσωση της μεσοκάθετης ευθείας ($\zeta$) του ευθύγραμμου τμήματος $AB$.
    \item Αν $K$ είναι ένα τυχαίο σημείο της ευθείας ($\zeta$), να βρείτε την εξίσωση ($c$) όλων των κύκλων, οι οποίοι έχουν κέντρο $K$ και διέρχονται από τα σημεία $A$ και $B$, συναρτήσει μιας παραμέτρου $\lambda \in \mathbb{R}$.  \monades{07}
    \item Αν η εξίσωση $(x - 2)^2 + (y - \lambda)^2 = \lambda^2 + 1, \lambda \in \mathbb{R}$, παριστάνει όλους τους κύκλους ($c$) του ερωτήματος β), τότε:  \monades{08}
    \begin{rlist}
        \item Να σχεδιάσετε τον κύκλο, ο οποίος έχει διάμετρο το ευθύγραμμο τμήμα $AB$.

        \item Να αποδείξετε ότι η ευθεία ($\varepsilon$): $x + \lambda y - 1 = 0$ εφάπτεται σε όλους τους κύκλους ($c$) στο σημείο $A(1,0)$.  \monades{05}
 \monades{05}
    \end{rlist}\end{alist}

\vspace{4mm}

% --- Άσκηση 20864 (20864.tex) ---
\Askhsh{}{2}

Δίνονται οι ευθείες: $\varepsilon1 : 2x + y - 6 = 0$ και $\varepsilon2 : 2x + y + 2 = 0$.

\begin{alist}    \item Να δείξετε ότι οι ευθείες $\varepsilon1$ και $\varepsilon2$ είναι παράλληλες. \monades{12}
    \item
    \begin{rlist}
        \item Να δείξετε ότι το σημείο $A(0,6)$ ανήκει στην ευθεία $\varepsilon1$. \monades{5}
        \item Να υπολογίσετε την απόσταση των ευθειών $\varepsilon1$ και $\varepsilon2$. \monades{8}
    \end{rlist}\end{alist}

\vspace{4mm}

% --- Άσκηση 20865 (20865.tex) ---
\Askhsh{}{2}
Δίνονται οι ελλείψεις $C_1: x^2 + 4y^2 = 4$, $C_2: \frac{x^2}{16} + \frac{y^2}{4} = 1$ και οι γραφικές τους παραστάσεις στο παρακάτω σχήμα.

\begin{center}
\begin{tikzpicture}    % Define the axes
    \draw[->] (-5,0) -- (5,0) node[right] {$x$};
    \draw[->] (0,-2.5) -- (0,2.5) node[above] {$y$};
    % Draw the origin
    \fill (0,0) circle (1.5pt);
    \node[below left] at (0,0) {O};
    % Draw C1 (inner, green)
    % C1: x^2/4 + y^2/1 = 1 => x radius = 2, y radius = 1
    \draw[l3, thirdcolor] (0,0) ellipse (2 and 1);
    \node[thirdcolor, above right] at (0.5, 1.2) {$C1$}; % Adjust position for label
    % Draw C2 (outer, blue)
    % C2: x^2/16 + y^2/4 = 1 => x radius = 4, y radius = 2
    \draw[l3, maincolor] (0,0) ellipse (4 and 2);
    \node[maincolor, above right] at (1.5, 2.2) {$C2$}; % Adjust position for label
\end{tikzpicture}
\end{center}

\begin{alist}    \item Να βρείτε τα μήκη των αξόνων και τις εστίες των δύο ελλείψεων. \monades{14}
    \item Από το σχήμα φαίνεται ότι οι δύο ελλείψεις έχουν την ίδια εκκεντρότητα. Να αποδείξετε ότι αυτό είναι αληθές. \monades{11}\end{alist}

\vspace{4mm}

% --- Άσκηση 20866 (20866.tex) ---
\Askhsh{}{3}
Στο παρακάτω σχήμα φαίνεται η γραφική παράσταση μιας παραβολής $(c)$, που έχει κορυφή την αρχή των αξόνων, άξονα συμμετρίας τον $y'y$ και διέρχεται από το σημείο $A(1,2)$.

\begin{tikzpicture}    \begin{axis}[
        axis lines = middle,
        xlabel = {$x$},
        ylabel = {$y$},
        xmin = -3.5, xmax = 3.5,
        ymin = -1.5, ymax = 6.5,
        xtick = {-3,-2,-1,1,2,3},
        ytick = {-1,1,2,3,4,5,6},
        grid = both,
        grid style = {gray!30},
        samples = 100,
        width = 10cm, height = 7cm,
        tick label style = {font=\scriptsize},
        every axis x label/.style = {at={(current axis.right of origin)}, anchor=north west},
        every axis y label/.style = {at={(current axis.above origin)}, anchor=south},
    ]
        % Parabola (c): y = 2x^2
        \addplot[very thick, maincolor, domain=-1.8:1.8] {2*x^2};
        \node[maincolor, above right] at (axis cs:1.7, 5.5) {$(c)$};
        % Point A(1,2)
        \addplot[only marks, mark=*, mark size=1.5pt, black] coordinates {(1,2)};
        \node[above right] at (axis cs:1,2) {$A$};
    \end{axis}\end{tikzpicture}

\begin{alist}    \item Να βρείτε την εξίσωση, την εστία και την διευθετούσα της παραβολής. \monades{06}
    \item Να βρείτε το συμμετρικό του σημείου $A$ ως προς τον άξονα της παραβολής. \monades{04}
    \item
        \begin{rlist}
            \item Να βρείτε την εξίσωση της εφαπτομένης $(\varepsilon)$ της παραβολής στο σημείο $A'(-1,2)$. \monades{08}
            \item Να βρείτε το σημείο τομής της $(\varepsilon)$ με τον άξονα $y'y$ και στη συνέχεια να την σχεδιάσετε. \monades{07}
        \end{rlist}\end{alist}

\vspace{4mm}

% --- Άσκηση 20868 (20868.tex) ---
\OnPar{ΘΕΜΑ 2}

Στο παρακάτω σχήμα φαίνονται τα σημεία $A(1,1)$, $B(3,2)$ και $\Gamma$ μιας ευθείας $(\varepsilon)$.

\begin{center}
\begin{tikzpicture}[scale=1]
    % Set up the coordinate system dimensions
    \tkzInit[xmin=-2.5,xmax=4,ymin=-1.5,ymax=3]
    % Draw x-axis with label
    \draw[l3,->] (-5,0) -- (5,0) node[below right] {$x$};
    % Draw y-axis with label
    \draw[l3,->] (0,-5) -- (0,5) node[above left] {$y$};
    % Define and label the origin
    \tkzDefPoint(0,0){O}
    \tkzLabelPoint[below left](O){O}
    % Define points A and B
    \tkzDefPoint(1,1){A}
    \tkzDefPoint(3,2){B}
    % Draw points A and B (size and fill are configured in the preamble)
    \tkzDrawPoints(A,B)
    % Label points A and B
    \tkzLabelPoint[above left](A){A}
    \tkzLabelPoint[above right](B){B}
    % Define the line passing through A and B
    \tkzDefLine[through=A,B]\tkzGetPoint{lineA}\tkzGetPoint{lineB}
    % Draw the line, extending it slightly on both ends.
    % 'l3' is a predefined line style from the preamble (line width=0.3mm).
    \tkzDrawLine[l3, add=0.5 and 0.5](lineA,lineB)
    % Label the line (varepsilon)
    \tkzLabelLine[pos=1.05, right](lineA,lineB){$(\varepsilon)$}
    % Find the intersection of the line with the x-axis (point Gamma)
    % tkzPointResultX is a predefined point (1,0) for the x-axis direction
    \tkzInterLL(lineA,lineB)(O,tkzPointResultX)
    \tkzGetPoint{Gamma}
    % Draw point Gamma and label it
    \tkzDrawPoints(Gamma)
    \tkzLabelPoint[below left](Gamma){$\Gamma$}\end{tikzpicture}
\end{center}

\begin{alist}    \item Να βρείτε την κλίση $\lambda$ της ευθείας $(\varepsilon)$. \monades{10}
    \item Να αποδείξετε ότι η εξίσωση της ευθείας $(\varepsilon)$ είναι η $y = \frac{1}{2}x + \frac{1}{2}$. \monades{05}
    \item Να βρείτε τις συντεταγμένες του σημείου $\Gamma$, στο οποίο η ευθεία $(\varepsilon)$ τέμνει τον άξονα $x'x$. \monades{10}\end{alist}

\vspace{4mm}

% --- Άσκηση 20869 (20869.tex) ---
\Askhsh{}{2}
Στο παρακάτω σχήμα φαίνονται η υπερβολή $C: x^2 - y^2 = 1$, η εστία της $E$, η εφαπτομένη της $\zeta$ στο σημείο $A(1,0)$ και το σημείο $\Gamma$ στο οποίο αυτή τέμνει την ασύμπτωτη ευθεία $\varepsilon1$ της υπερβολής.

\begin{tikzpicture}\begin{axis}[
    axis lines = middle,
    xlabel = {$x$},
    ylabel = {$y$},
    xmin = -3.5, xmax = 3.5,
    ymin = -3.5, ymax = 3.5,
    xtick = {-3,-2,-1,1,2,3},
    ytick = {-3,-2,-1,1,2,3},
    xticklabel=\empty, 
    yticklabel=\empty, 
    grid = none,
    samples = 100,
    width = 10cm, height = 8cm,
    tick label style = {font=\scriptsize},
    every axis x label/.style = {at={(current axis.right of origin)}, anchor=north west},
    every axis y label/.style = {at={(current axis.above origin)}, anchor=south east},
    axis line style={-Latex}, 
]
    % Hyperbola C: x^2 - y^2 = 1
    \addplot[very thick, gray, domain=1:3.2] {sqrt(x^2-1)};
    \addplot[very thick, gray, domain=1:3.2] {-sqrt(x^2-1)};
    \addplot[very thick, gray, domain=-3.2:-1] {sqrt(x^2-1)};
    \addplot[very thick, gray, domain=-3.2:-1] {-sqrt(x^2-1)};
    \node[gray, font=\scriptsize] at (axis cs:-2.2,0.8) {C};
    % Asymptote epsilon1: y = x (blue, secondarycolor)
    \addplot[very thick, secondarycolor, domain=-3.2:3.2] {x};
    \node[secondarycolor, above right, font=\scriptsize] at (axis cs:2.8,2.8) {$\varepsilon1$};
    % Asymptote epsilon2: y = -x (gray)
    \addplot[very thick, gray, domain=-3.2:3.2] {-x};
    % Tangent zeta: x = 1 (red, maincolor)
    \draw[very thick, maincolor] (axis cs:1, -3.2) -- (axis cs:1, 3.2);
    \node[maincolor, above left, font=\scriptsize] at (axis cs:1,3.2) {$\zeta$};
    % Points
    \fill (axis cs:0,0) circle (1.5pt) node[below left, font=\scriptsize] {O};
    \fill (axis cs:1,0) circle (1.5pt) node[below right, font=\scriptsize] {A};
    \fill (axis cs:sqrt(2),0) circle (1.5pt) node[below right, font=\scriptsize] {E}; % Focus E
    \fill (axis cs:1,1) circle (1.5pt) node[above left, font=\scriptsize] {$\Gamma$}; % Intersection point
\end{axis}\end{tikzpicture}

\begin{alist}    \item Να βρείτε τις εστίες $E', E$ και τις ασύμπτωτες $\varepsilon1, \varepsilon2$ της υπερβολής. \monades{10}
    \item
    \begin{rlist}
        \item Να βρείτε την εξίσωση της εφαπτομένης $\zeta$. \monades{07}
        \item Να αποδείξετε ότι το σημείο $\Gamma$ έχει συντεταγμένες $(1,1)$. \monades{08}
    \end{rlist}\end{alist}

\vspace{4mm}

% --- Άσκηση 20883 (20883.tex) ---
\Askhsh{}{2}

Δίνεται η εξίσωση της έλλειψης $C: 16x^2 + 25y^2 = 400$.

\begin{alist}    \item Να βρείτε τα μήκη $BB'$, $AA'$ του μικρού και του μεγάλου άξονα της έλλειψης, καθώς και τις εστίες της $E$ και $E'$. \monades{12}
    \item Αν $E'(-3,0)$ και $E(3,0)$, να γράψετε την εξίσωση της παραβολής που έχει εστία το σημείο $E'$ και διευθετούσα την ευθεία που διέρχεται από το $E$ και είναι παράλληλη στον άξονα $y'y$. \monades{13}\end{alist}

\vspace{4mm}

% --- Άσκηση 20885 (20885.tex) ---
\Askhsh{2023BLYKALGEUCLID002}{2}
Η ευθεία $\varepsilon$ διέρχεται από το σημείο $A(-3, -1)$ και σχηματίζει με τον άξονα $x'x$ γωνία $\frac{3\pi}{4}$.

\begin{alist}    \item Να βρείτε την εξίσωση της ευθείας $\varepsilon$. \monades{12}
    \item Να αποδείξετε ότι το εμβαδό του τριγώνου, που σχηματίζει η ευθεία $\varepsilon$ με τους άξονες $x'x$ και $y'y$, είναι: $E = 8$. \monades{13}\end{alist}

\vspace{4mm}

% --- Άσκηση 20888 (20888.tex) ---
\Askhsh{}{2}
Δίνονται τα διανύσματα $\vec{a}$, $\vec{b}$ και $\vec{y}$, για τα οποία ισχύουν: $|\vec{a}|=4$, $|\vec{b}|=5$, $\left(\vec{a},\vec{b}\right) = \frac{2\pi}{3}$ και $\vec{y} = 2\vec{a} + 3\vec{b}$. Να υπολογίσετε:
\begin{alist}    \item το εσωτερικό γινόμενο $\vec{a} \cdot \vec{b}$. \monades{10}
    \item το μέτρο του διανύσματος $\vec{y}$. \monades{15}\end{alist}

\vspace{4mm}

% --- Άσκηση 20890 (20890.tex) ---
\Askhsh{}{2}
Δίνεται το τρίγωνο $AB\Gamma$ με κορυφές τα σημεία $A(3,-3)$, $B(2,-8)$ και $\Gamma(7,-3)$. Να βρείτε:
\begin{alist}    \item την εξίσωση της πλευράς $B\Gamma$. \monades{10}
    \item την εξίσωση του κύκλου που έχει κέντρο το $A$ και εφάπτεται στην πλευρά $B\Gamma$. \monades{15}\end{alist}

\vspace{4mm}

% --- Άσκηση 20914 (20914.tex) ---
\textbf{ΘΕΜΑ 2}

Στο παρακάτω σχήμα, δίνονται τα διανύσματα $\vec{a}, \vec{b}$.

\begin{alist}    \item Να σημειώσετε πάνω στο σχήμα τα διανύσματα $\vec{OA} = \vec{a}$, $\vec{OB} = \vec{b}$ όπου $O$ η αρχή των αξόνων. \monades{08}
    \item Να βρείτε τις συντεταγμένες των διανυσμάτων $\vec{OA} = \vec{a}$, $\vec{OB} = \vec{b}$ και $\vec{AB}$. \monades{09}
    \item Να αποδείξετε ότι το τρίγωνο $OAB$ είναι ορθογώνιο. \monades{08}\end{alist}

\begin{center}
    \begin{tikzpicture}        \begin{axis}[
            axis lines = middle,
            xlabel = {$x$},
            ylabel = {$y$},
            xmin = -5.5, xmax = 3.5,
            ymin = -1.5, ymax = 5.5,
            xtick = {-5,-4,-3,-2,-1,0,1,2,3},
            ytick = {-1,0,1,2,3,4,5},
            grid = both,
            grid style = {gray!30},
            width = 10cm, height = 8cm,
            tick label style = {font=\scriptsize},
            every axis x label/.style = {at={(current axis.right of origin)}, anchor=north west},
            every axis y label/.style = {at={(current axis.above origin)}, anchor=south},
            xticklabel style = {
                /pgf/number format/print zero/.code=\ifnum\pgfmathparse{\tick}==0\relax\else\pgfmathprintnumber{\tick}\fi
            },
            yticklabel style = {
                /pgf/number format/print zero/.code=\ifnum\pgfmathparse{\tick}==0\relax\else\pgfmathprintnumber{\tick}\fi
            },
        ]
            % Vector a
            \addplot[->, l3, maincolor] coordinates {(-5,1) (-4,3)};
            \node[maincolor, above right, font=\small] at (axis cs: -4.5, 2.5) {$\vec{a}$};
            % Vector b
            \addplot[->, l3, secondarycolor] coordinates {(-3,5) (-2,2)};
            \node[secondarycolor, right, font=\small] at (axis cs: -2.5, 3.5) {$\vec{b}$};
            % Origin O
            \addplot[only marks, mark=o, mark size=2.5pt, draw=maincolor, fill=white] coordinates {(0,0)};
            \node[below right, font=\small] at (axis cs:0,0) {$O$};
        \end{axis}\end{tikzpicture}
\end{center}

\vspace{4mm}

% --- Άσκηση 20926 (20926.tex) ---
\Askhsh{}{}
{2}
Δίνεται η ευθεία $ε: x - 2y = 1$ και τα σημεία $A(0,2)$, $B(1,0)$.
\begin{alist}    \item Να αποδείξετε ότι το σημείο Β ανήκει στην ευθεία $ε$ ενώ το σημείο Α δεν είναι σημείο της $ε$. \monades{08}
    \item Να βρείτε την απόσταση του σημείου Α από την ευθεία $ε$. \monades{08}
    \item Να υπολογίσετε την απόσταση του Α από το Β και να αποδείξετε ότι η προβολή του Α στην ευθεία $ε$ είναι το Β. \monades{09}\end{alist}

\vspace{4mm}

% --- Άσκηση 20938 (20938.tex) ---
\begin{center}
    \Askhseis{ΘΕΜΑ 4}
\end{center}

Σε ορθογώνιο σύστημα συντεταγμένων $Oxy$, θεωρούμε τρίγωνο $OAB$, με $\overrightarrow{OA} = \vec{a}$, $\overrightarrow{OB} = (6,8)$. Για το διάνυσμα $\vec{a}$ γνωρίζουμε
ότι $\vec{a} = (\left|\vec{a}\right| - 4, \left|\vec{a}\right| - 2)$.

\begin{alist}    \item Να δείξετε ότι $\left|\vec{a}\right| = 2$. \monades{10}
    \item Να βρείτε σημείο $\Gamma$ έτσι, ώστε το τετράπλευρο $OA\Gamma B$ να αποτελεί παραλληλόγραμμο. \monades{08}
    \item Να βρείτε την γωνία που σχηματίζει η πλευρά $OA$ με τη διαγώνιο $AB$ του παραλληλογράμμου $OA\Gamma B$. \monades{07}\end{alist}

\vspace{4mm}

% --- Άσκηση 20939 (20939.tex) ---
\OnPar{ΘΕΜΑ 4}
Σε ένα ορθοκανονικό σύστημα συντεταγμένων $Oxy$ είναι τοποθετημένα $7$ χωριά ως σημεία του επιπέδου και μια πηγή νερού σε ένα σημείο $\Pi$. Γνωρίζουμε ότι υπάρχουν $6$ αγωγοί νερού που συνδέουν την πηγή με έξι από τα παραπάνω χωριά. Οι αγωγοί αυτοί ανήκουν στις γραμμές με εξισώσεις της μορφής:
\[ (\lambda + 1)x + (\lambda - 1)y + 2 = 0 \]
με $\lambda \in \{0,1,2,3,4,5\}$.

\begin{alist}    \item Να αποδείξετε ότι και οι $6$ γραμμές είναι ευθείες.
    \item Να βρείτε τις συντεταγμένες του σημείου $\Pi$.  \monades{04}
    \item Το έβδομο χωριό βρίσκεται στο σημείο $O(0,0)$. Να αποδείξετε ότι κανένας από τους παραπάνω αγωγούς νερού δεν διέρχεται από το χωριό αυτό.  \monades{06}
    \item Προκειμένου να έχει πρόσβαση στο νερό το χωριό $O$, υπάρχουν δύο επιλογές:  \monades{04}
    \begin{rlist}
        \item $1^{\eta}$ επιλογή: Να συνδέσουμε απευθείας το χωριό $O$ με την πηγή
        \item $2^{\eta}$ επιλογή: Να συνδέσουμε το χωριό $O$ με έναν από τους παραπάνω αγωγούς μέσω της συντομότερης διαδρομής.
    \end{rlist}
 Με δεδομένο ότι το κόστος κατασκευής ανά μονάδα μήκους για κάθε μία από τις παραπάνω επιλογές είναι το ίδιο,
    \begin{rlist}
        \item να βρείτε την τιμή του $\lambda$ για την οποία οι δύο επιλογές οδηγούν στο ίδιο κόστος κατασκευής.
        \item Πως εξηγείται γεωμετρικά το συμπέρασμα;  \monades{08}
 \monades{03}
    \end{rlist}\end{alist}

\vspace{4mm}

% --- Άσκηση 21152 (21152.tex) ---
\Askhsh{}{1}

\begin{alist}    \item  Να χαρακτηρίσετε καθεμιά από τις προτάσεις που ακολουθούν ως \textbf{Σωστή (Σ)} ή \textbf{Λανθασμένη (Λ)}, γράφοντας στην κόλλα σας, δίπλα στον αριθμό που αντιστοιχεί σε καθεμιά από αυτές το γράμμα Σ αν η πρόταση είναι Σωστή, ή το γράμμα Λ αν αυτή είναι Λάθος.
    \begin{rlist}
        \item Κάθε διάνυσμα στον χώρο είναι ίσο με τη διανυσματική ακτίνα του πέρατος μείον τη διανυσματική ακτίνα της αρχής.
        \item Η ευθεία που διέρχεται από το σημείο $A(x_0,y_0)$ και είναι παράλληλη στον άξονα $x'x$ έχει εξίσωση $x = x_0$.
        \item Η ευθεία με εξίσωση $Ax + By + \Gamma = 0$ είναι παράλληλη στο διάνυσμα $\vec{n} = (A,B)$.
        \item Η παραβολή με εξίσωση $y^2 = 4x$ έχει εστία το σημείο $E(1,0)$.
        \item Η εφαπτομένη του κύκλου $x^2 + y^2 = \rho^2$ στο σημείο του $A(x_1,y_1)$ έχει εξίσωση $x x_1 + y y_1 = \rho^2$.
    \end{rlist}
    \item Να αποδείξετε ότι η εξίσωση της ευθείας που διέρχεται από το σημείο $A(x_0,y_0)$ και έχει συντελεστή διεύθυνσης $\lambda$ είναι  \monades{10}
 \[ y - y_0 = \lambda (x - x_0) \] \monades{15}\end{alist}

\vspace{4mm}

% --- Άσκηση 21154 (21154.tex) ---
\Askhsh{}{4}
Δίνεται η εξίσωση
\[ x^2 + y^2 - 4ax - 4ay = 0 \quad (1) \]
όπου $a$ είναι πραγματικός αριθμός.

\begin{alist}    \item Να βρείτε τις τιμές του $a$ για τις οποίες η εξίσωση (1) παριστάνει κύκλο.
    \item Να προσδιορίσετε το κέντρο $K$ και την ακτίνα $R$ των κύκλων ως συνάρτηση του $a$.  \monades{8}
    \item Να βρείτε τον γεωμετρικό τόπο των κέντρων των κύκλων για τις διάφορες τιμές του $a$ του ερωτήματος (α).  \monades{6}
    \item Να εξετάσετε αν υπάρχει τιμή του $a$ ώστε ο αντίστοιχος κύκλος που ορίζεται από την εξίσωση (1) να εφάπτεται στον άξονα $x'x$.  \monades{5}
 \monades{6}\end{alist}

\vspace{4mm}

% --- Άσκηση 21159 (21159.tex) ---
\Askhsh{}{4}
Δίνονται τα σημεία $A(\alpha, 0)$ και $B(0, \beta)$ με $\alpha, \beta > 0$ και $\alpha + \beta = 10$.

\begin{alist}    \item Να αποδείξετε ότι η εξίσωση των κύκλων με διάμετρο την $AB$, για κάθε τιμή των $\alpha$ και $\beta$ είναι $x^2 + y^2 - \alpha x - (10 - \alpha)y = 0$. \monades{8}
    \item Να αποδείξετε ότι όλοι οι κύκλοι με διάμετρο την $AB$, για τις διάφορες τιμές των $\alpha$ και $\beta$ διέρχονται από δύο σταθερά σημεία, την αρχή $O$ των αξόνων και ένα σημείο $P$ του οποίου να βρείτε τις συντεταγμένες. \monades{9}
    \item Να βρείτε τον γεωμετρικό τόπο των κέντρων όλων των κύκλων με διάμετρο την $AB$ για τις διάφορες τιμές των $\alpha$ και $\beta$. \monades{8}\end{alist}

\vspace{4mm}

% --- Άσκηση 21160 (21160.tex) ---
\Askhsh{25026}{4} % The template's exercise command, with a placeholder code (25026) and Theme number (4)
Σε καρτεσιανό σύστημα συντεταγμένων $Oxy$ θεωρούμε το τρίγωνο που ορίζεται από τα σημεία $O(0, 0)$, $B(\kappa, 0)$ και $\Gamma(0,2\kappa)$ όπου $\kappa$ θετικός πραγματικός αριθμός. Εξωτερικά του τριγώνου $OB\Gamma$ κατασκευάζουμε τετράγωνα $OB\Delta E$ και $O\Gamma ZH$, τότε:

% --- START GEOMETRIC FIGURE ---
\begin{center}
\begin{tikzpicture}[scale=0.8]
    % Ορισμός της τιμής του κ για σκοπούς σχεδίασης.
    % Το πρόβλημα ορίζει το κ ως θετικό πραγματικό αριθμό,
    % οπότε για οπτικοποίηση επιλέγουμε μια συγκεκριμένη θετική τιμή (π.χ. κ=1).
    \pgfmathsetmacro{\k}{1} 
    % Αρχικοποίηση αξόνων, πλέγματος και συστήματος συντεταγμένων
    \tkzInit[xmin=-2.5*\k,xmax=1.5*\k,ymin=-1.5*\k,ymax=2.5*\k,xstep=0.5*\k,ystep=0.5*\k]
    \tkzGrid[sub,color=gray!10] % Ελαφρύ πλέγμα για ευκολία ανάγνωσης
    \draw[l3,->] (-5,0) -- (5,0) node[below right] {$x$}; % Άξονας x
    \draw[l3,->] (0,-5) -- (0,5) node[above left] {$y$}; % Άξονας y
    \tkzLabelX[font=\scriptsize] % Ετικέτες στους άξονες
    \tkzLabelY[font=\scriptsize] 
    % Ορισμός των σημείων με βάση το κ
    \tkzDefPoint(0,0){O}
    \tkzDefPoint(\k,0){B}
    \tkzDefPoint(0,2*\k){G} % Σημείο Γ
    % Σημεία για το τετράγωνο ΟΒΔΕ. Κατασκευάζεται εξωτερικά του ΟΒΓ,
    % οπότε οι y-συντεταγμένες είναι αρνητικές σε σχέση με την ΟΒ.
    \tkzDefPoint(\k,-\k){D} % Σημείο Δ
    \tkzDefPoint(0,-\k){E}
    % Σημεία για το τετράγωνο ΟΓΖΗ. Κατασκευάζεται εξωτερικά του ΟΒΓ,
    % οπότε οι x-συντεταγμένες είναι αρνητικές σε σχέση με την ΟΓ.
    \tkzDefPoint(-2*\k,2*\k){Z}
    \tkzDefPoint(-2*\k,0){H}
    % Σχεδίαση του τριγώνου ΟΒΓ
    \tkzDrawPolygon[l3,thirdcolor!70!black](O,B,G) % Τρίγωνο με 'thirdcolor' (πράσινο)
    % Σχεδίαση των τετραγώνων με τα καθορισμένα χρώματα του template
    \tkzDrawPolygon[l3,secondarycolor!70!black](O,B,D,E) % Τετράγωνο ΟΒΔΕ με 'secondarycolor' (ματζέντα)
    \tkzDrawPolygon[l3,maincolor!70!black](O,G,Z,H)     % Τετράγωνο ΟΓΖΗ με 'maincolor' (μπλε)
    % Σχεδίαση των συγκεκριμένων τμημάτων που αναφέρονται στις ερωτήσεις (ΓΔ και ΒΖ)
    % Χρησιμοποιούνται πιο παχιές και ευδιάκριτες γραμμές για να τονιστούν
    \tkzDrawSegment[l4,red!70!black](G,D) % Ευθύγραμμο τμήμα ΓΔ
    \tkzDrawSegment[l4,blue!70!black](B,Z) % Ευθύγραμμο τμήμα ΒΖ
    % Σχεδίαση όλων των καθορισμένων σημείων
    \tkzDrawPoints(O,B,G,D,E,Z,H)
    % Ετικέτες σημείων
    \tkzLabelPoint[below left](O){$O$}
    \tkzLabelPoint[below right](B){$B$}
    \tkzLabelPoint[above left](G){$\Gamma$}
    \tkzLabelPoint[below right](D){$\Delta$}
    \tkzLabelPoint[left](E){$E$}
    \tkzLabelPoint[above left](Z){$Z$}
    \tkzLabelPoint[below left](H){$H$}\end{tikzpicture}
\end{center}
% --- END GEOMETRIC FIGURE ---

\begin{alist}    \item Να βρείτε τις εξισώσεις των ευθειών που ανήκουν τα ευθύγραμμα τμήματα $\Gamma\Delta$ και $BZ$.
    \item Να βρεθεί η εξίσωση του ύψους του τριγώνου $OB\Gamma$ που διέρχεται από το $O$.  \monades{10}
    \item Να αποδείξετε ότι οι ευθείες $\Gamma\Delta$, $BZ$ και το ύψος του β) ερωτήματος διέρχονται από το ίδιο σημείο.  \monades{7}
 \monades{8}\end{alist}

\vspace{4mm}

% --- Άσκηση 21162 (21162.tex) ---
\Askhsh{}{2}
Δίνονται τα σημεία $A(3,2)$ και $B(-1,-6)$. Να βρεθούν:
\begin{alist}    \item Οι συντεταγμένες του μέσου $M$ του ευθύγραμμου τμήματος $AB$. \monades{8}
    \item Ο συντελεστής διεύθυνσης της ευθείας που διέρχεται από τα σημεία $A$ και $B$. \monades{8}
    \item Η εξίσωση της μεσοκαθέτου ευθείας $(\varepsilon)$ του ευθύγραμμου τμήματος $AB$. \monades{9}\end{alist}

\vspace{4mm}

% --- Άσκηση 21165 (21165.tex) ---
\Askhsh{}{2}
Θεωρούμε το παραλληλόγραμμο ΑΒΓΔ και έστω $\vec{AB} = \vec{\alpha}$ και $\vec{A\Delta} = \vec{\beta}$. Τα σημεία Ε και Ζ είναι τέτοια ώστε $\vec{AE} = -\frac{1}{2} \vec{AB}$ και $\vec{AZ} = \frac{1}{3} \vec{A\Delta}$.

\begin{alist}    \item Να αποδείξετε ότι: $\vec{EZ} = \frac{1}{2} \vec{\alpha} + \frac{1}{3} \vec{\beta}$ και $\vec{Z\Gamma} = \vec{\alpha} + \frac{2}{3} \vec{\beta}$. \monades{10}
    \item Να αποδείξετε ότι: $\vec{Z\Gamma} = 2\vec{EZ}$. \monades{9}
    \item Να δείξετε ότι τα σημεία Ζ, Ε και Γ είναι συνευθειακά. \monades{6}\end{alist}

\begin{center}
\begin{tikzpicture}    % Define the vertices of the parallelogram A B Γ Δ
    \tkzDefPoint(0,0){A}
    \tkzDefPoint(4,0){B}
    \tkzDefPoint(1.5,2.5){D} % Coordinates adjusted to resemble the PDF image
    \tkzDefPointBy[translation=from A to D](B) \tkzGetPoint{C} % C = B + vector(AD)
    % Define point E such that AE = -1/2 AB
    % A is (0,0), B is (4,0). Vector AB is (4,0).
    % Vector AE is -1/2 * (4,0) = (-2,0). So E is at (-2,0).
    \tkzDefPoint(-2,0){E}
    % Define point Z such that AZ = 1/3 AD
    % A is (0,0), D is (1.5, 2.5). Vector AD is (1.5, 2.5).
    % Vector AZ is 1/3 * (1.5, 2.5) = (0.5, 2.5/3). So Z is at (0.5, 0.833).
    \tkzDefPointWith[linear, K=1/3](A,D) \tkzGetPoint{Z}
    % Draw the parallelogram A B Γ Δ
    \tkzDrawPolygon[l3](A,B,C,D)
    % Draw the connecting segments EZ and ZΓ
    \tkzDrawSegments[l3](E,Z Z,C)
    % Draw the points
    \tkzDrawPoints(A,B,C,D,E,Z)
    % Labels for the points
    \tkzLabelPoint[below left](E){E}
    \tkzLabelPoint[below](A){A}
    \tkzLabelPoint[below](B){B}
    \tkzLabelPoint[above left](D){$\Delta$}
    \tkzLabelPoint[above right](C){$\Gamma$}
    \tkzLabelPoint[above left](Z){Z}\end{tikzpicture}
\end{center}

\vspace{4mm}

% --- Άσκηση 21218 (21218.tex) ---
\Askhsh{}{2}
Δίνονται οι υπερβολές $(C_1): x^2 - y^2 = 1$, $(C_2): y^2 - x^2 = 1$.

\begin{alist}    \item Να αποδείξετε ότι οι εστίες της $C_1$ είναι οι $E_1(\sqrt{2},0)$, $E'_1(-\sqrt{2},0)$.
    \item Αν $E_2, E'_2$ οι εστίες της $C_2$, τότε να αποδείξετε ότι το $E_1 E_2 E'_1 E'_2$ είναι τετράγωνο.  \monades{12}
 \monades{13}\end{alist}

\vspace{4mm}

% --- Άσκηση 21248 (21248.tex) ---
\Askhsh{}{2}
Δίνεται το σημείο $E(2,0)$, η ευθεία $\delta: x = -2$ και τυχαίο σημείο $M(x,y)$ του επιπέδου.

\begin{alist}    \item
        \begin{rlist}
            \item Να βρείτε την απόσταση $(ME)$ του σημείου $M(x,y)$ από το $E(2,0)$ ως συνάρτηση των $x, y$. \monades{8}
            \item Να βρείτε την απόσταση $d(M,\delta)$ του σημείου $M$ από την ευθεία $\delta$ ως συνάρτηση των $x, y$. \monades{8}
        \end{rlist}
    \item Αν ισχύει $(ME) = d(M,\delta)$ να δείξετε ότι το σημείο $M$ ανήκει στην παραβολή $y^2 = 8x$. \monades{9}\end{alist}

\vspace{4mm}

% --- Άσκηση 21260 (21260.tex) ---
\Askhsh{}{2}
Δίνεται η ευθεία ($\varepsilon$): $y - 2x = 0$ και τα σημεία $B(1,1)$ και $\Gamma(-1,3)$.
\begin{alist}    \item Να δείξετε ότι το σημείο $A(5,10)$ ανήκει στην ευθεία ($\varepsilon$). \monades{5}
    \item Να βρείτε τις συντεταγμένες των διανυσμάτων $\vec{AB}$ και $\vec{A\Gamma}$. \monades{10}
    \item Να βρείτε το εμβαδόν του τριγώνου $\triangle AB\Gamma$. \monades{10}\end{alist}

\vspace{4mm}

% --- Άσκηση 21276 (21276.tex) ---
\Askhsh{XXXX}{4}
Σε μια σύγχρονη πόλη, κατασκευάζεται σιδηροδρομικό δίκτυο που περιλαμβάνει:
\begin{itemize}
    \item τη γραμμή $γ_1$, κάθε σημείο της οποίας στο ορθοκανονικό σύστημα συντεταγμένων είναι της μορφής: $Α (\lambda-1, 2\lambda+1)$, $\lambda \in \mathbb{R}$.
    \item τη γραμμή $γ_2$, που περνάει από το σταθμό $Σ (-4, 2)$ και είναι παράλληλη στο διάνυσμα $\vec{u} = (-1, 3)$.
\end{itemize}

\begin{alist}    \item Να βρείτε τις εξισώσεις των ευθειών πάνω στις οποίες βρίσκονται οι γραμμές $γ_1$ και $γ_2$.
    \item Η είσοδος του αθλητικού σταδίου μιας συνοικίας θα βρίσκεται στο σημείο $K(1, 1)$ του ορθοκανονικού συστήματος συντεταγμένων. Οι κατασκευαστές θέλουν να συνδέσουν την είσοδο του σταδίου απ' ευθείας με κάθετο δρόμο, με μια από τις γραμμές $γ_1$ και $γ_2$. Να βρείτε με ποια από τις δύο γραμμές είναι πιο συμφέρουσα η σύνδεση. Δίνεται ότι το κόστος σύνδεσης ανά μονάδα μήκους, είναι το ίδιο και για τις δύο γραμμές.  \monades{10}
    \item Γύρω από το στάδιο θα δημιουργηθεί κυκλικό πάρκο. Να βρείτε την εξίσωση του κύκλου, που θα ορίζει το πάρκο, αν το κέντρο του είναι το σημείο $K$ και επιπλέον ο κύκλος αυτός εφάπτεται της γραμμής $γ_1$.  \monades{9}
 \monades{6}\end{alist}

\vspace{4mm}

% --- Άσκηση 21306 (21306.tex) ---
\Askhseis{ΘΕΜΑ 2}
Σε καρτεσιανό επίπεδο $Oxy$ δίνεται η παραβολή με άξονα συμμετρίας τον $\chi'\chi$, κορυφή $O(0,0)$ και εστία $E(2,0)$, όπως στο παρακάτω σχήμα. Το σημείο $A$ της παραβολής έχει τετμημένη $3$ και βρίσκεται στο πρώτο τεταρτημόριο του $Oxy$.

\begin{alist}    \item Να αποδείξετε ότι η εξίσωση της παραβολής είναι $y^2 = 8x$ και ότι $A(3, 2\sqrt{6})$.
    \item Να σχεδιάσετε τη διευθετούσα ($\delta$) της παραβολής και να γράψετε την εξίσωσή της.  \monades{10}
    \item Να βρείτε την εξίσωση της εφαπτομένης ($\varepsilon$) της παραβολής στο σημείο $A$.  \monades{06}
 \monades{09}\end{alist}

\begin{center}
\begin{tikzpicture}\begin{axis}[
    axis lines = middle,
    xmin = -3.5, xmax = 4.5,
    ymin = -5.5, ymax = 6.5,
    xtick = {-3,-2,-1,0,1,2,3,4},
    ytick = {-5,-4,-3,-2,-1,1,2,3,4,5,6},
    grid = none,
    width = 8cm, height = 10cm,
    tick label style = {font=\scriptsize},
]
    % Parabola y^2 = 8x -> y = +/- sqrt(8x)
    \addplot[very thick, maincolor, samples=100, domain=0:4.1] {sqrt(8*x)};
    \addplot[very thick, maincolor, samples=100, domain=0:4.1] {-sqrt(8*x)};
    % Point E(2,0)
    \addplot[only marks, mark=*, mark size=1.5pt, thirdcolor] coordinates {(2,0)};
    \node[thirdcolor, below right] at (axis cs:2,0) {$E$};
    % Point A(3, 2*sqrt(6))
    \addplot[only marks, mark=*, mark size=1.5pt, maincolor] coordinates {(3,{2*sqrt(6)})};
    \node[maincolor, above right] at (axis cs:3,{2*sqrt(6)}) {$A$};
    % Dashed line from A to x-axis
    \draw[dashed, gray] (axis cs:3,0) -- (axis cs:3,{2*sqrt(6)});
\end{axis}\end{tikzpicture}
\end{center}

\vspace{4mm}

% --- Άσκηση 21307 (21307.tex) ---
\Askhsh{}{2}
Σε καρτεσιανό επίπεδο $Oxy$ δίνεται η παραβολή με εξίσωση $x^2 = 12y$.

\begin{alist}    \item Να αποδείξετε ότι η εστία της παραβολής είναι το σημείο $E(0,3)$ και να βρείτε τα σημεία της παραβολής που έχουν τεταγμένη 3. \monades{12}
    \item Να αποδείξετε ότι εφαπτομένες $(\varepsilon1)$ και $(\varepsilon2)$ της παραβολής στα σημεία $A(6,3)$ και $B(-6,3)$, αντίστοιχα, έχουν εξισώσεις $y = x - 3$ και $y = -x - 3$. \monades{08}
    \item Να βρείτε το σημείο τομής των $(\varepsilon1)$ και $(\varepsilon2)$. \monades{05}\end{alist}

\vspace{4mm}

% --- Άσκηση 21308 (21308.tex) ---
\Askhsh{ELL101}{2}
Σε καρτεσιανό επίπεδο $Oxy$ δίνεται η έλλειψη
\[
\frac{x^2}{25} + \frac{y^2}{16} = 1
\]
Να βρείτε:
\begin{alist}    \item Τις συντεταγμένες των εστιών $E$ και $E'$ της έλλειψης και την απόστασή τους.
    \item Το μήκος του μικρού άξονα και το μήκος του μεγάλου άξονα της έλλειψης.  \monades{09}
    \item Την εξίσωση της εφαπτομένης $(\epsilon)$ της έλλειψης στο σημείο της $B(0,4)$.  \monades{08}
 \monades{08}\end{alist}

\vspace{4mm}

% --- Άσκηση 21349 (21349.tex) ---
\textbf{ΘΕΜΑ 4}

Σε ορθοκανονικό σύστημα αξόνων με αρχή το σημείο $O$ θεωρούμε κύκλο $(C)$ και ευθεία $(\varepsilon)$ με εξισώσεις
\[ x^2 + y^2 - 9x - 3y + 10 = 0 \quad (1) \]
και
\[ 4x + 3y - 10 = 0 \quad (2) \]
αντίστοιχα.

\begin{alist}    \item
    \begin{rlist}
        \item Να βρείτε το κέντρο $K$ και την ακτίνα $R$ του κύκλου $(C)$. \monades{5}
        \item Να υπολογίσετε την απόσταση του κέντρου $K$ από την ευθεία $(\varepsilon)$ και να αποδείξετε ότι η ευθεία $(\varepsilon)$ τέμνει τον κύκλο $(C)$ σε δύο σημεία. \monades{4}
        \item Να προσδιορίσετε τα σημεία $A$ και $B$ στα οποία η ευθεία $(\varepsilon)$ τέμνει τον κύκλο $(C)$. \monades{5}
    \end{rlist}
    \item Αν είναι $A(1,2)$ και $B(4, -2)$, τότε:
    \begin{rlist}
        \item Να υπολογίσετε το εσωτερικό γινόμενο $\vec{OA} \cdot \vec{OB}$. \monades{5}
        \item Να αποδείξετε ότι ο κύκλος με διάμετρο $AB$ διέρχεται από το σημείο $O$. \monades{6}
    \end{rlist}\end{alist}

\vspace{4mm}

% --- Άσκηση 21647 (21647.tex) ---
\Askhsh{}{2}
Η έλλειψη $C$ έχει εστίες τα σημεία $E(4,0)$, $E'(-4,0)$ και μεγάλο άξονα 10. Να βρείτε:
\begin{alist}    \item την εξίσωση της $C$. \monades{10}
    \item την εκκεντρότητά της $C$. \monades{7}
    \item την εξίσωση της εφαπτομένης της $C$ στο σημείο της $M\left(4, \frac{9}{5}\right)$. \monades{8}\end{alist}

\vspace{4mm}

% --- Άσκηση 21648 (21648.tex) ---
\subsection*{\titlefont{\color{maincolor}ΘΕΜΑ 2}}

Η έλλειψη $C$ έχει εστίες τα σημεία $E(3,0)$, $E'(-3,0)$ και διέρχεται από το σημείο $M\left(4, \frac{12}{5}\right)$.

\begin{alist}    \item Να αποδείξετε ότι το μήκος του μεγάλου άξονα είναι 10. \monades{10}
    \item Να βρείτε την εξίσωση της $C$. \monades{8}
    \item Να βρείτε την εξίσωση της εφαπτομένης της $C$ στο σημείο της $M\left(4, \frac{12}{5}\right)$. \monades{7}\end{alist}

Δίνεται ότι $\sqrt{1369} = 37$.

\vspace{4mm}

% --- Άσκηση 21649 (21649.tex) ---
\textbf{ΘΕΜΑ 2}
\vspace{4mm}
Η υπερβολή $C$ έχει εστίες τα σημεία $E(5,0)$, $E'(-5,0)$ και εκκεντρότητα $\varepsilon = \frac{5}{4}$. Να βρείτε:
\begin{alist}    \item την εξίσωση της $C$. \monades{10}
    \item τις εξισώσεις των ασύμπτωτων της $C$. \monades{8}
    \item την εξίσωση της εφαπτομένης της $C$ στο σημείο της $M\left(5, -\frac{9}{4}\right)$. \monades{7}\end{alist}

\vspace{4mm}

% --- Άσκηση 21651 (21651.tex) ---
\OnPar{ΘΕΜΑ 2}

Η υπερβολή $C$ έχει εστίες τα σημεία $E(5,0)$, $E'(-5,0)$ και διέρχεται από το σημείο $A(4,0)$.

\begin{alist}    \item Να αποδείξετε ότι έχει εκκεντρότητα $\frac{5}{4}$. \monades{10}
    \item Να βρείτε την εξίσωση της $C$. \monades{8}
    \item Να βρείτε την εξίσωση της εφαπτομένης της $C$ στο σημείο $M\left(5, \frac{7}{4}\right)$. \monades{7}\end{alist}

\vspace{4mm}

% --- Άσκηση 21652 (21652.tex) ---
\Askhsh{}{4}

Δίνονται οι εξισώσεις $\lambda x + y = 2\lambda$ (1) και $x + \lambda y = \lambda +1$ (2), όπου $\lambda \in \mathbb{R}$.

\begin{alist}    \item Να δείξετε ότι για κάθε $\lambda \in \mathbb{R}$ οι εξισώσεις (1) και (2) παριστάνουν ευθείες $\varepsilon_\lambda$ και $\eta_\lambda$ αντίστοιχα.
    \item Να βρείτε για ποιες τιμές του $\lambda$ οι ευθείες $\varepsilon_\lambda$ και $\eta_\lambda$ τέμνονται.  \monades{6}
    \item Για $\lambda \in \mathbb{R} \setminus \{-1,1\}$  \monades{6}
    \begin{rlist}
        \item να βρείτε συναρτήσει του $\lambda$ τις συντεταγμένες του σημείου τομής $M$ των $\varepsilon_\lambda$ και $\eta_\lambda$.

        \item αν $M \left( \frac{2\lambda+1}{\lambda+1}, \frac{\lambda}{\lambda+1} \right)$ να αποδείξετε ότι το $M$ κινείται στην ευθεία $\zeta: x-y=1$.  \monades{7}
 \monades{6}
    \end{rlist}\end{alist}

\vspace{4mm}

% --- Άσκηση 21653 (21653.tex) ---
\Askhsh{XXXXX}{4}
Στο παρακάτω ορθοκανονικό σύστημα συντεταγμένων, το 1ο τεταρτημόριο αντιστοιχεί σε μια θαλάσσια περιοχή και τα υπόλοιπα τεταρτημόρια σε στεριά. Οι ημιάξονες $Ox, Oy$ οριοθετούν ένα λιμάνι. Ένα πλοίο ρυμουλκείται στο λιμάνι, δεμένο με δύο συρματόσχοινα στο ίδιο σημείο $\Pi(\kappa,\lambda)$ του πλοίου. Το ένα από τα δύο ρυμουλκά είναι σταθερό στο σημείο $E(2,0)$ και το άλλο κινείται ώστε η θέση να περιγράφεται από το σημείο $P(-2, \lambda)$. Η ρυμούλκηση γίνεται με τέτοιο τρόπο ώστε κάθε χρονική στιγμή της ρυμούλκησης να ισχύει $(PE) = (P\Pi)$.

\begin{alist}    \item Να αποδείξετε ότι το σημείο $P(-2,\lambda)$ κινείται σε σταθερή ευθεία $(\delta)$ της οποίας να βρείτε την εξίσωση.
    \item Να αιτιολογήσετε γιατί κάθε χρονική στιγμή της ρυμούλκησης είναι $P\Pi \perp (\delta)$.  \monades{5}
    \item Να αποδείξετε ότι η πορεία του $\Pi(\kappa,\lambda)$ είναι παραβολή $C$ της οποίας να βρείτε την εξίσωση.  \monades{5}
    \item Αν $y^2 = 8x$ η εξίσωση της παραβολής $C$ να αποδείξετε ότι κάθε χρονική στιγμή η μεσοκάθετος του $EP$ εφάπτεται της παραβολής $C$ στο σημείο $\Pi$.  \monades{7}
 \monades{8}\end{alist}

\begin{center}
\begin{tikzpicture}\begin{axis}[
    axis lines = middle,
    xmin = -3.5, xmax = 7.5,
    ymin = -1.5, ymax = 7.5,
    xtick = {-3,-2,-1,0,1,2,3,4,5,6,7},
    ytick = {-1,1,2,3,4,5,6,7},
    grid = both,
    grid style = {gray!30},
    width = 10cm, height = 8cm,
    tick label style = {font=\scriptsize},
    xlabel={}, ylabel={}
]
    % Define points for drawing based on the diagram
    \coordinate (E) at (axis cs:2,0);
    \coordinate (P) at (axis cs:-2,6);
    \coordinate (Pi) at (axis cs:4,6);
    % Draw segments (lines) as shown in the diagram
    \draw[l3, black] (P) -- (E);
    \draw[l3, black] (P) -- (Pi);
    \draw[l3, black] (Pi) -- (E);
    % Draw points as solid black circles
    \fill[black] (E) circle (1.5pt);
    \fill[black] (P) circle (1.5pt);
    \fill[black] (Pi) circle (1.5pt);
    % Labels for points
    \node[below right] at (E) {E};
    \node[above left] at (P) {P};
    \node[above right] at (Pi) {$\Pi$};
\end{axis}\end{tikzpicture}
\end{center}

\vspace{4mm}

% --- Άσκηση 21655 (21655.tex) ---
\paragraph{ΘΕΜΑ 4}
Στο παρακάτω σχήμα δίνεται η έλλειψη $C$ με κέντρο το $O(0,0)$, εστίες τα σημεία $E, E'$ και κορυφές τα σημεία $A(5,0), A'(-5,0), B, B'$. Αν είναι γνωστό ότι το τετράπλευρο $BEB'E'$ είναι τετράγωνο, να βρείτε:

\begin{alist}    \item τις συντεταγμένες των σημείων $B, B', E, E'$. \monades{8}
    \item την εξίσωση της έλλειψης $C$. \monades{7}
    \item Έστω $M$ τυχαίο σημείο της $C$, που δεν ταυτίζεται με κάποιο από τα $A, A'$.
    \begin{rlist}
        \item να αποδείξετε ότι όλα τα τρίγωνα $EME'$ έχουν την ίδια περίμετρο την οποία να προσδιορίσετε. \monades{5}

        \item να βρείτε τις συντεταγμένες του $M$ για τις οποίες το εμβαδόν του τριγώνου $EME'$ παίρνει τη μέγιστη τιμή του, την οποία και να προσδιορίσετε. \monades{5}
    \end{rlist}\end{alist}

\begin{center}
    \begin{tikzpicture}        % Define ellipse parameters based on problem statement
        % Given A(5,0), A'(-5,0) => a = 5
        \pgfmathsetmacro{\aval}{5}
        % Given BEB'E' is a square, and E, E' are foci, B, B' are vertices on y-axis.
        % This implies the distance from origin to E is equal to distance from origin to B.
        % So, c = b.
        % For ellipse: a^2 = b^2 + c^2.
        % Substituting c=b: a^2 = b^2 + b^2 = 2b^2.
        % (5)^2 = 2b^2 => 25 = 2b^2 => b^2 = 25/2 => b = 5/sqrt(2).
        \pgfmathsetmacro{\bval}{5/sqrt(2)} % Calculated b (approx 3.5355)
        \pgfmathsetmacro{\cval}{5/sqrt(2)} % Calculated c (approx 3.5355)
        \begin{axis}[
            axis lines=middle,
            xlabel={}, 
            ylabel={}, 
            xmin=-7.5, xmax=8.5,
            ymin=-5.5, ymax=4.5,
            xtick={-7,-6,-5,-4,-3,-2,-1,0,1,2,3,4,5,6,7,8}, 
            ytick={-5,-4,-3,-2,-1,0,1,2,3,4}, 
            grid=none, 
            domain=0:360, 
            samples=200, 
            width=10cm, height=7cm,
            tick label style={font=\scriptsize},
            axis line style={l3}, 
            every axis x label/.style={at={(current axis.right of origin)}, anchor=north west},
            every axis y label/.style={at={(current axis.above origin)}, anchor=south},
            clip=false 
        ]
            % Plot the ellipse C parametrically: (a*cos(t), b*sin(t))
            \addplot[very thick, maincolor, forget plot] ({(\aval*cos(x))}, {(\bval*sin(x))});
            % Draw the square BEB'E'
            % Vertices are (0,b), (c,0), (0,-b), (-c,0). Since b=c, this is a square with vertices (0,b), (b,0), (0,-b), (-b,0).
            \fill[thirdcolor!30, draw=thirdcolor, l3]
                (axis cs:0, \bval) -- (axis cs:\cval, 0) -- (axis cs:0, -\bval) -- (axis cs:-\cval, 0) -- cycle;
            % Define points using tkzDefPoint. tkz-euclide works within pgfplots axis.
            \tkzDefPoint(0,0){O}
            \tkzDefPoint(\aval,0){A}
            \tkzDefPoint(-\aval,0){A'}
            \tkzDefPoint(0,\bval){B}
            \tkzDefPoint(0,-\bval){B'}
            \tkzDefPoint(\cval,0){E}
            \tkzDefPoint(-\cval,0){E'}
            % M (example point, roughly matching the position in the original PDF, e.g., at x=2)
            \pgfmathsetmacro{\mxcoord}{2}
            % Calculate y-coordinate for M on the ellipse: y = sqrt(b^2 * (1 - x^2/a^2))
            \pgfmathsetmacro{\mycoord}{sqrt((\bval^2)*(1-(\mxcoord^2)/(\aval^2)))}
            \tkzDefPoint(\mxcoord,\mycoord){M}
            % Draw points using tkzDrawPoints, which respects tkzSetUpPoint (size=2.8,fill=white)
            \tkzDrawPoints(A,A',B,B',E,E',O,M)
            % Label points
            \tkzLabelPoint[below right](A){A}
            \tkzLabelPoint[below left](A'){A'}
            \tkzLabelPoint[above left](B){B}
            \tkzLabelPoint[below left](B'){B'}
            \tkzLabelPoint[below right](E){E}
            \tkzLabelPoint[below left](E'){E'}
            \tkzLabelPoint[below right](O){O} % Label O explicitly next to the origin's 0 tick
            \tkzLabelPoint[above right](M){M}
        \end{axis}\end{tikzpicture}
\end{center}

\vspace{4mm}

% --- Άσκηση 21656 (21656.tex) ---
\Askhsh{}{4}
Έστω υπερβολή $C$ με κέντρο το $(0,0)$, εστίες τα σημεία $E(5,0), E'(-5,0)$ και κορυφές τα σημεία $A(4,0), A'(-4,0)$.
\begin{alist}    \item Να βρείτε:
    \begin{rlist}
        \item τις εξισώσεις των ασυμπτώτων της υπερβολής $C$. \monades{3}
        \item την εξίσωση της υπερβολής $C$. \monades{3}
    \end{rlist}
    \item Να σχεδιάσετε στο ίδιο ορθοκανονικό σύστημα, την υπερβολή $C$, τις ασύμπτωτες της $C$ και το ορθογώνιο βάσης της $C$. \monades{9}
    \begin{center}
    \begin{tikzpicture}    \begin{axis}[
 axis lines = middle, xlabel = {$x$}, ylabel = {$y$}, xmin = -7, xmax = 7, ymin = -5, ymax = 5, xtick = {-6,-4,-2,0,2,4,6}, ytick = {-4,-2,0,2,4}, grid = both, grid style = {gray!30}, samples = 100, width = 10cm, height = 7cm, tick label style = {font=\scriptsize}, every axis x label/.style = {at={(current axis.right of origin)}, anchor=north west}, every axis y label/.style = {at={(current axis.above origin)}, anchor=south east}, clip=false, ]
        % Hyperbola: x^2/16 - y^2/9 = 1  => y = \pm (3/4) * sqrt(x^2 - 16)
        % Right branch
 \addplot[very thick, maincolor, domain=4:7] {3/4 * sqrt(x^2 - 16)}; \addplot[very thick, maincolor, domain=4:7] {-3/4 * sqrt(x^2 - 16)};
        % Left branch
 \addplot[very thick, maincolor, domain=-7:-4] {3/4 * sqrt(x^2 - 16)}; \addplot[very thick, maincolor, domain=-7:-4] {-3/4 * sqrt(x^2 - 16)};
        % Asymptotes: y = \pm (3/4) * x
 \addplot[very thick, secondarycolor, dashed, domain=-7:7] {3/4 * x}; \addplot[very thick, secondarycolor, dashed, domain=-7:7] {-3/4 * x};
        % Fundamental Rectangle: vertices (\pm 4, \pm 3)
 \draw[thirdcolor, thick] (4,3) rectangle (-4,-3);
        % Foci E(5,0), E'(-5,0)
 \fill[black] (5,0) circle (1.5pt) node[below right, font=\tiny] {$E(5,0)$}; \fill[black] (-5,0) circle (1.5pt) node[below left, font=\tiny] {$E'(-5,0)$};
        % Vertices A(4,0), A'(-4,0)
 \fill[black] (4,0) circle (1.5pt) node[above right, font=\tiny] {$A(4,0)$}; \fill[black] (-4,0) circle (1.5pt) node[above left, font=\tiny] {$A'(-4,0)$};
        % Origin (0,0)
 \node[below left, font=\tiny] at (axis cs:0,0) {$O$};
    \end{axis}\end{tikzpicture}
    \end{center}
    \item Αν $M$ τυχαίο σημείο της $C$, να βρείτε την τιμή της παράστασης $(ME)-(ME')$. \monades{5}
    \item Αν $M(\sqrt{80},6)$ σημείο της $C$, να βρείτε την εξίσωση της διχοτόμου της γωνίας $EME'$. \monades{5}\end{alist}

\vspace{4mm}

% --- Άσκηση 21657 (21657.tex) ---
\Askhsh{}{4}
Έστω υπερβολή $C$ με κέντρο το $(0,0)$, εστίες πάνω στον άξονα $xx'$ της οποίας το ορθογώνιο βάσης είναι τετράγωνο.
\begin{alist}    \item Να βρείτε:
    \begin{rlist}
        \item τις εξισώσεις των ασυμπτώτων της $C$. \monades{6}
        \item την εκκεντρότητα της $C$. \monades{6}
    \end{rlist}
    \item Αν η υπερβολή διέρχεται από το σημείο $(2,0)$ και $(\zeta)$ τυχαία ευθεία παράλληλη σε κάποια εκ των ασύμπτωτων της $C$ (που δεν ταυτίζεται με κάποια από αυτές),
    \begin{rlist}
        \item να δείξετε ότι η $(\zeta)$ έχει ένα μόνο κοινό σημείο με την $C$. \monades{8}
        \item είναι η ευθεία $(\zeta)$ εφαπτόμενη της $C$; Αιτιολογείστε την απάντησή σας. \monades{5}
    \end{rlist}\end{alist}

\vspace{4mm}

% --- Άσκηση 21662 (21662.tex) ---
\Askhsh{}{2}
Δίνεται η ευθεία $\varepsilon$: $-x + y - 2 = 0$ και τα σημεία $A(-5,1)$ και $B(-3,5)$.

\begin{alist}    \item Να βρείτε το συμμετρικό του σημείου $A$ ως προς το σημείο $B$. \monades{10}
    \item Να βρείτε:
    \begin{rlist}
        \item την εξίσωση της ευθείας $\varepsilon'$ που διέρχεται από το $B$ και είναι κάθετη στην $\varepsilon$. \monades{5}
        \item το σημείο τομής των ευθειών $\varepsilon$ και $\varepsilon'$. \monades{5}
        \item το συμμετρικό του σημείου $B$ ως προς την ευθεία $\varepsilon$. \monades{5}
    \end{rlist}\end{alist}

\vspace{4mm}

% --- Άσκηση 21681 (21681.tex) ---
\Askhsh{}{2}
Θεωρούμε τα σημεία $A(1, 2)$, $B(-3, 4)$, $\Gamma(2, 5)$.

\begin{alist}    \item Να βρείτε σημείο $\Delta$ ώστε $\vec{\Delta\Gamma} = \vec{AB}$. \monades{10}
    \item Να αιτιολογήσετε γιατί το τετράπλευρο $AB\Gamma\Delta$ είναι παραλληλόγραμμο. \monades{8}
    \item Να βρείτε το κέντρο $O$ του παραλληλογράμμου. \monades{7}\end{alist}

\vspace{4mm}

% --- Άσκηση 21682 (21682.tex) ---
\Askhsh{}{2}
Έστω $\vec{\alpha}$, $\vec{\beta}$ δυο διανύσματα για τα οποία ισχύει $\vec{\alpha}+\vec{\beta}=(11, 2)$ και $\vec{\alpha}-\vec{\beta}=(-5, -10)$.
\begin{alist}    \item Να αποδείξετε ότι $\vec{\alpha}=(3, -4)$ και να βρείτε το διάνυσμα $\vec{\beta}$. \monades{14}
    \item Να αποδείξετε ότι $\vec{\alpha} \perp \vec{\beta}$ και $|\vec{\beta}|=2|\vec{\alpha}|$. \monades{11}\end{alist}

\vspace{4mm}

% --- Άσκηση 21683 (21683.tex) ---
\Askhsh{MATH4}{4}
\subsection*{\textcolor{maincolor}{ΘΕΜΑ 4}}
Στο διπλανό σχήμα φαίνεται το πρώτο τεταρτημόριο του κύκλου $x^2 + y^2 = 4$ και το τυχαίο σημείο του $M(x_1, y_1)$, $0<x_1 <2$ ανάμεσα στα A, B.

\begin{center}
\begin{tikzpicture}[scale=1.2]
    % Define origin and axes points for the arc
    \tkzDefPoint(0,0){O}
    \tkzDefPoint(2,0){A}
    \tkzDefPoint(0,2){B}
    % Example coordinates for M to draw the tangent
    % Let's use x1 = 1.2 for a clear visual
    \pgfmathsetmacro{\xM}{1.2}
    \pgfmathsetmacro{\yM}{sqrt(4-\xM*\xM)} % y1 = sqrt(4-x1^2) since M is on the circle
    \tkzDefPoint(\xM,\yM){M}
    % Define K and L, intersection points of the tangent with axes
    % Tangent equation at M(x1, y1) is x*x1 + y*y1 = 4
    % For K (y=0): xK = 4/x1
    % For L (x=0): yL = 4/y1
    \pgfmathsetmacro{\xK}{4/\xM}
    \pgfmathsetmacro{\yL}{4/\yM}
    \tkzDefPoint(\xK,0){K}
    \tkzDefPoint(0,\yL){L}
    % Initialize drawing area for axes
    % Ensure xmax and ymax are large enough to contain K and L
    \tkzInit[xmin=-0.5,xmax=4,ymin=-0.5,ymax=4] 
    \draw[l3,->] (-5,0) -- (5,0) node[below right] {$x$};
    \draw[l3,->] (0,-5) -- (0,5) node[above left] {$y$};
    % Draw the quarter circle (arc)
    \tkzDrawArc[l3, color=maincolor](O,A)(B) % Center O, start A, end B
    % Draw the tangent line (L--K)
    \tkzDrawSegment[l3, color=secondarycolor](L,K)
    % Draw points
    \tkzDrawPoints(O,A,B,M,K,L)
    % Label points
    \tkzLabelPoint[below left](O){O}
    \tkzLabelPoint[below](A){A}
    \tkzLabelPoint[left](B){B}
    \tkzLabelPoint[above right](M){M} % Adjust label position for M
    \tkzLabelPoint[below right](K){K}
    \tkzLabelPoint[above left](L){L} % Adjust label position for L
\end{tikzpicture}
\end{center}

\begin{alist}    \item Να βρείτε την εξίσωση της εφαπτομένης του στο M και τις συντεταγμένες των σημείων τομής της Κ, Λ με τους άξονες. \monades{6}
    \item Να αποδείξετε ότι το μήκος $d$ του τμήματος ΚΛ είναι $d= \dfrac{8}{x_1y_1}$. \monades{7}
    \item Να βρείτε το μήκος $d_0$ του τμήματος ΚΛ όταν $x_1 = \sqrt{2}$. \monades{5}
    \item Να αποδείξετε ότι, όταν το $M$ κινείται στο τεταρτοκύκλιο, τότε: $d \ge d_0$. \monades{7}\end{alist}

\vspace{4mm}

% --- Άσκηση 21690 (21690.tex) ---
\textbf{ΘΕΜΑ 4}
Δίνεται η παραβολή $C: y^2 = 3x$ και η ευθεία $\varepsilon: 3x+4y+10=0$.

\begin{alist}    \item Να αποδείξετε ότι η ευθεία και η παραβολή δεν έχουν κοινά σημεία και να τις σχεδιάσετε. \monades{8}
    \begin{center}
    \begin{tikzpicture}    \begin{axis}[
 axis lines = middle, xlabel = {$x$}, ylabel = {$y$}, xmin = -2, xmax = 10, ymin = -7, ymax = 7, xtick = {0,1,2,3,4,5,6,7,8,9,10}, ytick = {-6,-5,-4,-3,-2,-1,0,1,2,3,4,5,6}, grid = both, grid style = {gray!30}, samples = 100, width = 10cm, height = 7cm, tick label style = {font=\scriptsize}, every axis x label/.style = {at={(current axis.right of origin)}, anchor=north west}, every axis y label/.style = {at={(current axis.above origin)}, anchor=south}, ]
        % Parabola C: y^2 = 3x  => y = sqrt(3x) and y = -sqrt(3x)
 \addplot[very thick, maincolor, domain=0:10] {sqrt(3*x)}; \addplot[very thick, maincolor, domain=0:10] {-sqrt(3*x)}; \node[maincolor, above right] at (axis cs:3,3) {$C: y^2=3x$};
        % Line epsilon: 3x + 4y + 10 = 0  => y = (-3/4)x - (10/4) = -0.75x - 2.5
 \addplot[very thick, secondarycolor, domain=-2:10] {-0.75*x - 2.5}; \node[secondarycolor, below right] at (axis cs:2,-4) {$\varepsilon: 3x+4y+10=0$};
    \end{axis}\end{tikzpicture}
    \end{center}
    \item Έστω $M(x_0, y_0)$ ένα σημείο της παραβολής. Να αποδείξετε ότι η απόστασή του $d(M, \varepsilon)$ από την ευθεία είναι
    \item Να βρείτε το σημείο της παραβολής που είναι το πιο κοντινό στην ευθεία. \monades{5}  \[ d(M, \varepsilon) = \frac{(y_0+2)^2+6}{5} \] \monades{8}
    \item Να αποδείξετε ότι η εφαπτομένη της παραβολής στο σημείο που βρήκατε στο ερώτημα γ) είναι παράλληλη στην ευθεία $\varepsilon$. \monades{4}\end{alist}

\vspace{4mm}

% --- Άσκηση 21696 (21696.tex) ---
\Askhsh{Β2ΑΛΓ(2)2}{4ο}
Θεωρούμε την εξίσωση $x^2 + y^2 - 2x + 4y = 0$, \label{eq:circle} (1) και την ευθεία $\varepsilon: x - 2y + 3 = 0$.

\begin{alist}    \item Να αποδείξετε ότι η εξίσωση (1) παριστάνει κύκλο $C$ του οποίου να βρείτε το κέντρο $K$ και την ακτίνα $\rho$.
    \item Να εξετάσετε αν η ευθεία $\varepsilon$ έχει κοινά σημεία με τον κύκλο $C$.  \monades{5}
    \item Να βρείτε τις εφαπτόμενες $\varepsilon1, \varepsilon2$ του κύκλου $C$ που είναι κάθετες στην ευθεία $\varepsilon$.  \monades{5}
    \item Να αποδείξετε ότι $d(\varepsilon1, \varepsilon2) = 2\rho$. Πως αιτιολογείται γεωμετρικά το συμπέρασμα αυτό;  \monades{7}
 \monades{8}\end{alist}

\vspace{4mm}

% --- Άσκηση 21883 (21883.tex) ---
\Askhsh{10141}{4}

ΘΕΜΑ 4

Δίνεται η παραβολή $C:x^2=4y$ και η ευθεία $\varepsilon:y=x-2$.

\begin{alist}    \item Να βρείτε την εστία $E$ και τη διευθετούσα $\delta$ της παραβολής.
    \item Να αποδείξετε ότι η ευθεία και η παραβολή δεν έχουν κοινά σημεία. Στη συνέχεια σε ένα ορθοκανονικό σύστημα αξόνων $Oxy$ να σχεδιαστούν οι γραφικές παραστάσεις της παραβολής $C$ και της ευθείας $\varepsilon$.  \monades{5}
 \monades{8}
    \begin{center}
    \begin{tikzpicture}    \begin{axis}[
 axis lines = middle, xlabel = {$x$}, ylabel = {$y$}, xmin = -4, xmax = 6, ymin = -3, ymax = 6, xtick = {-4,-2,0,2,4,6}, ytick = {-2,0,2,4,6}, grid = both, grid style = {gray!30}, samples = 100, width = 10cm, height = 7cm, tick label style = {font=\scriptsize}, every axis x label/.style = {at={(current axis.right of origin)}, anchor=north west}, every axis y label/.style = {at={(current axis.above origin)}, anchor=south}, ]
        % Parabola C: x^2=4y => y = x^2/4
 \addplot[very thick, maincolor, domain=-3.5:5.5] {x^2/4}; \node[maincolor, above left] at (axis cs:4,4) {$C:y=\frac{x^2}{4}$};
        % Line ε: y=x-2
 \addplot[very thick, secondarycolor, domain=-3.5:5.5] {x-2}; \node[secondarycolor, below right] at (axis cs:4,2) {$\varepsilon:y=x-2$};
    \end{axis}\end{tikzpicture}
    \end{center}
    \item Αν $M(x,y)$ είναι σημείο της παραβολής, τότε:
    \begin{rlist}
        \item Να αποδείξετε ότι η απόσταση του $M$ από την ευθεία $\varepsilon$ είναι

        \item Να βρείτε την ελάχιστη απόσταση του σημείου $M$ από την ευθεία $\varepsilon$ καθώς και τις συντεταγμένες του σημείου $M$ της παραβολής που απέχει την ελάχιστη απόσταση από την ευθεία.  \[d(M,\varepsilon) = \frac{x^2-x+2}{2\sqrt{2}}\] \monades{6}
 \monades{6}
    \end{rlist}\end{alist}

\vspace{4mm}

% --- Άσκηση 21885 (21885.tex) ---
\Askhsh{}{4}
Δίνεται τρίγωνο $AB\Gamma$ και $\Delta,E$ σημεία εσωτερικά των πλευρών $AB$ και $A\Gamma$ αντίστοιχα τέτοια ώστε $\vec{AB}=\kappa \cdot \vec{A\Delta}$ και $\vec{A\Gamma} = \lambda \cdot \vec{AE}$, όπου $\kappa$ και $\lambda$ θετικοί πραγματικοί αριθμοί. Αν $\vec{AB}= \vec{\alpha}$ και $\vec{A\Gamma} = \vec{\beta}$, τότε:

\begin{alist}    \item Να εκφράσετε τα διανύσματα $\vec{\Delta E}$ και $\vec{B\Gamma}$ ως γραμμικό συνδυασμό των $\vec{\alpha}$ και $\vec{\beta}$. \monades{8}
    \item % This item corresponds to 'β)'
    \begin{rlist}
        \item Αν $\kappa=\lambda$, να αποδείξετε ότι $\vec{B\Gamma} // \vec{\Delta E}$ και $|\vec{B\Gamma}| = \kappa|\vec{\Delta E}|$. \monades{10}
        \item Αν $\kappa=\lambda=2$, να γράψετε τη σχέση που συνδέει τα διανύσματα $\vec{\Delta E}$ και $\vec{B\Gamma}$ και να διατυπώσετε λεκτικά ποιο γνωστό θεώρημα της Ευκλείδειας Γεωμετρίας έχει αποδειχθεί. \monades{7}
    \end{rlist}\end{alist}

\vspace{4mm}

% --- Άσκηση 21962 (21962.tex) ---
\Askhsh{2}{2} % ΘΕΜΑ 2
Δίνονται τα σημεία $A(0,3)$, $B(3,4)$ και $\Gamma(1,0)$.
\begin{alist}    \item Να αποδείξετε ότι η γωνία $\widehat{B A \Gamma}$ είναι ορθή. \monades{13}
    \item Να βρείτε το μέσο $K$ της υποτείνουσας $B\Gamma$ του ορθογωνίου τριγώνου $AB\Gamma$. \monades{5}
    \item Να βρείτε την εξίσωση του κύκλου που διέρχεται από τα σημεία $A$, $B$ και $\Gamma$. \monades{7}\end{alist}

\vspace{4mm}

% --- Άσκηση 21964 (21964.tex) ---
\Askhsh{}{2}
Δίνονται το σημείο $A(4,-2)$ και η ευθεία $(\varepsilon1)$ με εξίσωση: $x - y + 2 = 0$. Να βρείτε:
\begin{alist}    \item την ευθεία $(\varepsilon2)$ που διέρχεται από το σημείο $A$ και είναι κάθετη στην ευθεία $(\varepsilon1)$. \monades{8}
    \item το σημείο τομής $B$, των ευθειών $(\varepsilon1)$ και $(\varepsilon2): y = -x + 2$. \monades{8}
    \item το συμμετρικό $\Gamma$ του σημείου $A$, ως προς την ευθεία $(\varepsilon1)$. \monades{9}\end{alist}

\vspace{4mm}

% --- Άσκηση 21965 (21965.tex) ---
\Askhsh{Κωδικός Άσκησης}{2}
Δίνονται τα σημεία $A(2, -4)$ και $B(0, -2)$.
\begin{alist}    \item Να βρείτε το μέσο $M$ του τμήματος $AB$. \monades{7}
    \item Να βρείτε την εξίσωση της μεσοκαθέτου $(\zeta)$ του ευθύγραμμου τμήματος $AB$. \monades{8}
    \item Αν $(\zeta): y = x - 4$ και $(\varepsilon): y = 2x - 6$, τότε να βρείτε το σημείο τομής των ευθειών $(\zeta)$ και $(\varepsilon)$. \monades{10}\end{alist}

\vspace{4mm}

% --- Άσκηση 21973 (21973.tex) ---
\Askhsh{}{1}
\begin{alist}    \item Να χαρακτηρίσετε τις προτάσεις που ακολουθούν γράφοντας στην κόλλα σας δίπλα στον αριθμό που αντιστοιχεί σε κάθε πρόταση τη λέξη \textbf{ΣΩΣΤΟ}, αν η πρόταση είναι σωστή ή \textbf{ΛΑΘΟΣ}, αν η πρόταση είναι λανθασμένη.
    \begin{rlist}
        \item Αν ισχύει $\left|\vec{a}\right| = \lambda\left|\vec{b}\right|$ τότε υποχρεωτικά $\vec{a}//\vec{b}$.
        \item Η εφαπτομένη του κύκλου $C: x^2 + y^2 = \rho^2$ σε ένα σημείο του $A (x_1, y_1)$, έχει εξίσωση $xx_1 + yy_1 = \rho^2$.
        \item Η διευθετούσα της παραβολής $y^2 = 2px$, έχει εξίσωση $x = -\frac{p}{2}$.
        \item Η εκκεντρότητα μιας έλλειψης είναι μικρότερη της μονάδας.
        \item Η εξίσωση: $x^2 + y^2 = \alpha^2$ είναι εξίσωση ισοσκελούς υπερβολής.
    \end{rlist}
    \item Να αποδείξετε ότι η εξίσωση της ευθείας που διέρχεται από το σημείο $A(x_0, y_0)$ και έχει συντελεστή διεύθυνσης $\lambda$ είναι $y - y_0 = \lambda(x - x_0)$.  \monades{10}
 \monades{15}\end{alist}

\vspace{4mm}

% --- Άσκηση 22038 (22038.tex) ---
\OnPar{ΘΕΜΑ 2}
Δίνεται το διάνυσμα $\vec{\alpha} = (5,-12)$.
\begin{alist}    \item Να βρείτε ένα διάνυσμα $\vec{\beta}$ που να είναι ομόρροπο στο $\vec{\alpha}$ και να έχει μέτρο 1. \monades{12}
    \item Να βρείτε ένα διάνυσμα $\vec{\gamma}$ που να είναι αντίρροπο στο $\vec{\alpha}$ και να έχει μέτρο 7. \monades{13}\end{alist}

\vspace{4mm}

% --- Άσκηση 22040 (22040.tex) ---
\Askhsh{}{2}

Δίνεται το διάνυσμα $\vec{a} = (-4, 3)$.

\begin{alist}    \item Να βρείτε ένα διάνυσμα $\vec{\beta}$ που να είναι κάθετο στο $\vec{a}$.
    \item Να βρείτε ένα διάνυσμα $\vec{\gamma}$ που να είναι κάθετο στο $\vec{a}$ και να έχει μέτρο 1.  \monades{12}
 \monades{13}\end{alist}

\vspace{4mm}

% --- Άσκηση 22042 (22042.tex) ---
\Askhseis{ΘΕΜΑ 2}

Στο σχήμα φαίνεται ορθογώνιο παραλληλόγραμμο ΑΒΓΔ με κέντρο Ο. Τα σημεία Κ, Λ, Μ
χωρίζουν την πλευρά ΔΓ σε τέσσερα ίσα τμήματα.

\begin{center}
\begin{tikzpicture}[scale=1.0]
    % Ορισμός κορυφών ορθογωνίου παραλληλογράμμου
    \tkzDefPoint(0,0){D} % \Delta
    \tkzDefPoint(4,0){C} % \Gamma
    \tkzDefPoint(0,2){A}
    \tkzDefPoint(4,2){B}
    % Σχεδίαση του περιγράμματος
    \tkzDrawPolygon[l3](A,B,C,D)
    % Ορισμός του κέντρου O (μέσο της διαγωνίου AC)
    \tkzDefMidPoint(A,C) \tkzGetPoint{O}
    % Ορισμός των σημείων K, Λ, M στην πλευρά ΔΓ
    % ΔK = KΛ = ΛM = MΓ = 1/4 ΔΓ
    \tkzDefPoint(1,0){K}
    \tkzDefPoint(2,0){L} % Χρησιμοποιούμε L για το Λ
    \tkzDefPoint(3,0){M}
    % Σημεία
    \tkzDrawPoints(A,B,C,D,O,K,L,M)
    % Ετικέτες σημείων
    \tkzLabelPoint[above left](A){A}
    \tkzLabelPoint[above right](B){B}
    \tkzLabelPoint[below right](C){$\Gamma$}
    \tkzLabelPoint[below left](D){$\Delta$}
    \tkzLabelPoint[above right](O){O}
    \tkzLabelPoint[below](K){K}
    \tkzLabelPoint[below](L){$\Lambda$}
    \tkzLabelPoint[below](M){M}\end{tikzpicture}
\end{center}

Αν $\vec{\Delta\text{K}} = \vec{\alpha}$ και $\vec{\Delta\text{A}} = \vec{\beta}$ να εκφράσετε καθένα από τα ακόλουθα διανύσματα ως
γραμμικούς συνδυασμούς των $\vec{\alpha}$ και $\vec{\beta}$. Να αιτιολογήσετε τις απαντήσεις σας.

\begin{alist}    \item \quad $\vec{\Delta\text{Γ}}$ \monades{8}
    \item \quad $\vec{\text{M}\text{A}}$ \monades{8}
    \item \quad $\vec{\text{O}\Delta}$ \monades{9}\end{alist}

\vspace{4mm}

% --- Άσκηση 22044 (22044.tex) ---
\Askhsh{}{2}
Δίνονται τα σημεία $A(0,0)$, $B(4,0)$ και $\Gamma(5,1)$.

\begin{alist}    \item Να σχεδιάσετε ένα ορθοκανονικό σύστημα αξόνων και να τοποθετήσετε σε αυτό τα σημεία $A$, $B$, $\Gamma$.
    \item Να βρείτε τις συντεταγμένες ενός τέταρτου σημείου $\Delta$ έτσι ώστε το τετράπλευρο $AB\Gamma\Delta$ να είναι παραλληλόγραμμο. Να αιτιολογήσετε την απάντησή σας.  \monades{10}
 \monades{15}\end{alist}

\vspace{4mm}

% --- Άσκηση 22047 (22047.tex) ---
\Askhsh{202302}{2}
\fulltwoc{
Δίνεται το τετράγωνο του σχήματος με κορυφές $A_1$, $A_2$, $A_3$, $A_4$. Έστω $\lambda_1$, $\lambda_2$, $\lambda_3$, $\lambda_4$ οι κλίσεις των ευθειών $A_1A_2$, $A_2A_3$, $A_3A_4$, $A_4A_1$ αντίστοιχα.

\begin{center}
    \begin{tikzpicture}[scale=0.8]
        % Ορισμός των κορυφών του τετραγώνου ώστε να μοιάζει με το σχήμα
        \tkzDefPoint(0,2){A1}
        \tkzDefPoint(2,0){A2}
        \tkzDefPoint(0,-2){A3}
        \tkzDefPoint(-2,0){A4}
        % Σχεδίαση του τετραγώνου
        \tkzDrawPolygon[l3](A1,A2,A3,A4)
        % Σχεδίαση και επισήμανση των σημείων
        \tkzDrawPoints(A1,A2,A3,A4)
        \tkzLabelPoint[above](A1){A1}
        \tkzLabelPoint[right](A2){A2}
        \tkzLabelPoint[below](A3){A3}
        \tkzLabelPoint[left](A4){A4}\end{tikzpicture}
\end{center}

\begin{alist}    \item Να βρείτε όλα τα ζεύγη των παράλληλων πλευρών του τετραγώνου. Ποιά σχέση συνδέει τις κλίσεις κάθε δύο παράλληλων πλευρών; \monades{15}
    \item Να αποδείξετε ότι: \[ \frac{\lambda_1}{\lambda_3} + \frac{\lambda_2}{\lambda_4} = 2 \] \monades{10}\end{alist}
}

\vspace{4mm}

% --- Άσκηση 22049 (22049.tex) ---
\Askhsh{}{2}

Δίνεται το τετράγωνο του σχήματος με κορυφές $A_1$, $A_2$, $A_3$, $A_4$. Έστω $\lambda_1$, $\lambda_2$, $\lambda_3$, $\lambda_4$ οι κλίσεις των ευθειών $A_1A_2$, $A_2A_3$, $A_3A_4$, $A_4A_1$ αντίστοιχα.

\begin{center}
\begin{tikzpicture}[scale=1.5]
    % Ορισμός κορυφών τετραγώνου σε σχήμα ρόμβου
    \tkzDefPoint(0,1){A1}
    \tkzDefPoint(1,0){A2}
    \tkzDefPoint(0,-1){A3}
    \tkzDefPoint(-1,0){A4}
    % Σχεδίαση του περιγράμματος
    \tkzDrawPolygon[l3](A1,A2,A3,A4)
    % Σημεία
    \tkzDrawPoints(A1,A2,A3,A4)
    % Ετικέτες
    \tkzLabelPoint[above](A1){A1}
    \tkzLabelPoint[right](A2){A2}
    \tkzLabelPoint[below](A3){A3}
    \tkzLabelPoint[left](A4){A4}\end{tikzpicture}
\end{center}

\begin{alist}    \item Να βρείτε όλα τα ζεύγη των πλευρών του τετραγώνου που είναι κάθετες μεταξύ τους. Ποιά σχέση συνδέει τις κλίσεις κάθε δύο κάθετων πλευρών; \monades{15}
    \item Να αποδείξετε ότι:
 \[ \lambda_1 \cdot \lambda_2 + \lambda_2 \cdot \lambda_3 + \lambda_3 \cdot \lambda_4 + \lambda_4 \cdot \lambda_1 = -4 \] \monades{10}\end{alist}

\vspace{4mm}

% --- Άσκηση 22051 (22051.tex) ---
\Askhsh{}{2}

Δίνεται η υπερβολή $x^2 - y^2 = 1$. Να αποδείξετε για τις ασύμπτωτες ευθείες $\varepsilon1$ και $\varepsilon2$ της υπερβολής ότι:

\begin{center}
\begin{tikzpicture}\begin{axis}[
    axis lines = middle,
    xlabel = {$x$},
    ylabel = {$y$},
    xmin = -4.5, xmax = 4.5,
    ymin = -4.5, ymax = 4.5,
    xtick = \empty,
    ytick = \empty,
    samples = 100,
    width = 9cm, height = 7cm,
    tick label style = {font=\scriptsize},
    every axis x label/.style = {at={(current axis.right of origin)}, anchor=north west},
    every axis y label/.style = {at={(current axis.above origin)}, anchor=south},
    clip=false 
]
    % Asymptotes ε1 (y=x) and ε2 (y=-x)
    \addplot[dashed, gray, thick, domain=-4.5:4.5] {x};
    \addplot[dashed, gray, thick, domain=-4.5:4.5] {-x};
    % Hyperbola x^2 - y^2 = 1  => y = +- sqrt(x^2 - 1)
    % The image shows the hyperbola branches with a subtle dash pattern that looks almost solid.
    % Right branch (x >= 1)
    \addplot[maincolor, very thick, dash pattern=on 0.8mm off 0.4mm, domain=1:4] {sqrt(x^2-1)};
    \addplot[maincolor, very thick, dash pattern=on 0.8mm off 0.4mm, domain=1:4] {-sqrt(x^2-1)};
    % Left branch (x <= -1)
    \addplot[maincolor, very thick, dash pattern=on 0.8mm off 0.4mm, domain=-4:-1] {sqrt(x^2-1)};
    \addplot[maincolor, very thick, dash pattern=on 0.8mm off 0.4mm, domain=-4:-1] {-sqrt(x^2-1)};
    % Labels
    \node[below left] at (axis cs:0,0) {$O$};
    \node[above right, xshift=-2mm, yshift=2mm] at (axis cs:3.5,3.5) {$\varepsilon1$};
    \node[above left, xshift=2mm, yshift=2mm] at (axis cs:-3.5,3.5) {$\varepsilon2$};
    \node[right] at (axis cs:3,1.5) {$x^2 - y^2 = 1$};
\end{axis}\end{tikzpicture}
\end{center}

\begin{alist}    \item Συμπίπτουν με τη διχοτόμο του $1\tss{ου}$ και $3\tss{ου}$ τεταρτημορίου και τη διχοτόμο του $2\tss{ου}$ και $4\tss{ου}$ τεταρτημορίου, αντίστοιχα.
    \item Είναι ευθείες κάθετες μεταξύ τους.  \monades{13}
 \monades{12}\end{alist}

\vspace{4mm}

% --- Άσκηση 22052 (22052.tex) ---
\Askhseis{ΘΕΜΑ 2}
\vspace{4mm} % Add some vertical space as seen in the PDF
Δίνονται τα διανύσματα $\vec{a} = (-2,5)$ και $\vec{\beta} = (1, -3)$.

\begin{alist}    \item Να αποδείξετε ότι τα διανύσματα $\vec{a}$, $\vec{\beta}$ δεν είναι παράλληλα. \monades{12}
    \item Να εκφράσετε το διάνυσμα $\vec{v} = (8, -21)$ ως γραμμικό συνδυασμό των διανυσμάτων $\vec{a}$ και $\vec{\beta}$. \monades{13}\end{alist}

\vspace{4mm}

% --- Άσκηση 22055 (22055.tex) ---
\Askhsh{2}{2}

Δίνονται δύο τρίγωνα $AB\Gamma$ και $A'B'\Gamma'$ για τα οποία ισχύει $\hat{A} = \hat{A}'$ και $A\Gamma = A'\Gamma'$.

\begin{center}
\begin{tikzpicture}    % Τρίγωνο ΑΒΓ
    \tkzDefPoint(0,0){B}
    \tkzDefPoint(4,0){C}
    \tkzDefPoint(2,2){A}
    \tkzDrawPolygon[l3](A,B,C)
    \tkzLabelPoint[below](B){B}
    \tkzLabelPoint[below](C){$\Gamma$}
    \tkzLabelPoint[above](A){A}
    % Σήμανση γωνίας Α
    \tkzMarkAngle[arc=l3,size=6mm,fill=maincolor!10](B,A,C)
    % Σήμανση πλευράς ΑΓ
    \tkzMarkSegments[mark=s_](A,C)
    % Τρίγωνο Α'Β'Γ'
    \tkzDefPoint(5,0){Bprime}
    \tkzDefPoint(9,0){Cprime}
    \tkzDefPoint(7,2){Aprime}
    \tkzDrawPolygon[l3](Aprime,Bprime,Cprime)
    \tkzLabelPoint[below](Bprime){B$'$}
    \tkzLabelPoint[below](Cprime){$\Gamma'$}
    \tkzLabelPoint[above](Aprime){A$'$}
    % Σήμανση γωνίας Α'
    \tkzMarkAngle[arc=l3,size=6mm,fill=maincolor!10](Bprime,Aprime,Cprime)
    % Σήμανση πλευράς Α'Γ'
    \tkzMarkSegments[mark=s_](Aprime,Cprime)\end{tikzpicture}
\end{center}

\begin{alist}    \item Να εξηγήσετε γιατί:
    \begin{rlist}
        \item το μήκος της πλευράς $BA$ είναι ίσο με το μήκος της πλευράς $B'A'$ και \monades{3}
        \item το μήκος της πλευράς $A\Gamma$ είναι ίσο με το μήκος της πλευράς $A'\Gamma'$. \monades{3}
    \end{rlist}
    \item
    \begin{rlist}
        \item Να αποδείξετε ότι: $B\Gamma = B'\Gamma'$. \monades{10}
        \item Να εξηγήσετε γιατί το μήκος της πλευράς $B\Gamma$ είναι ίσο με το μήκος της πλευράς $B'\Gamma'$. \monades{3}
    \end{rlist}
    \item Θα μπορούσε η ακόλουθη πρόταση να ήταν κριτήριο ισότητας τριγώνων;
 «Αν για δύο τρίγωνα $AB\Gamma$ και $A'B'\Gamma'$ ισχύει $\hat{A} = \hat{A}'$ και $A\Gamma = A'\Gamma'$, τότε τα τρίγωνα $AB\Gamma$ και $A'B'\Gamma'$ είναι ίσα». Να αιτιολογήσετε την απάντησή σας. \monades{6}\end{alist}

\vspace{4mm}

% --- Άσκηση 22056 (22056.tex) ---
\Askhsh{2}{2}

Έστω $\Omega$ το σύνολο όλων των σημείων $(x, y)$ του επιπέδου για τα οποία ισχύει: $x^2 + y^2 \leq 9$.

\begin{alist}    \item Να σχεδιάσετε το σύνολο $\Omega$ σε ένα ορθοκανονικό σύστημα αξόνων Oxy.
 \monades{13}
    \begin{center}
    \begin{tikzpicture}    \begin{axis}[
 axis lines=middle, xlabel={$x$}, ylabel={$y$}, xmin=-4.5, xmax=4.5, ymin=-4.5, ymax=4.5, xtick=\{-3,-2,-1,0,1,2,3\}, ytick=\{-3,-2,-1,0,1,2,3\}, grid=both, grid style={gray!30}, width=9cm, height=9cm, tick label style={font=\scriptsize}, xlabel style={at={(current axis.right of origin)}, anchor=north west}, ylabel style={at={(current axis.above origin)}, anchor=south}, axis equal image, ]
        % Draw the filled disk x^2 + y^2 <= 9 (radius 3)
 \addplot[ draw=maincolor, ] ({3*cos(x)}, {3*sin(x)});
        % Label the origin O
 \node[black, below left] at (axis cs:0,0) {O};
        % Label key points on the axes
 \node[black, below] at (axis cs:3,0) {3}; \node[black, left] at (axis cs:0,3) {3}; \node[black, above] at (axis cs:-3,0) {--3}; \node[black, right] at (axis cs:0,-3) {--3};
    \end{axis}\end{tikzpicture}
    \end{center}
    \item Υπάρχει σημείο $A$ στο σύνολο $\Omega$ τέτοιο ώστε $\left| \text{OA} \right| = 4$, όπου $O$ η αρχή των αξόνων; Να αιτιολογήσετε την απάντησή σας.
 \monades{12}\end{alist}

\vspace{4mm}

% --- Άσκηση 22059 (22059.tex) ---
\textbf{ΘΕΜΑ 2}

Δίνεται το σημείο $A(1, -3)$ και η ευθεία $\varepsilon: 4x + 6y = 1$.
\begin{alist}    \item Να εξηγήσετε γιατί το $A$ δεν είναι σημείο της ευθείας $\varepsilon$.
    \item Να βρείτε την εξίσωση της ευθείας που διέρχεται από το σημείο $A$ και είναι παράλληλη στην ευθεία $\varepsilon$.  \monades{10}
 \monades{15}\end{alist}

\vspace{4mm}

% --- Άσκηση 22060 (22060.tex) ---
\Askhsh{}{2}
Δίνονται τα σημεία $A(0, 2)$, $B(3, 0)$, $\Gamma(6,2)$ και $\Delta(3,4)$.
\begin{alist}    \item Να βρείτε τις συντεταγμένες των διανυσμάτων $\overrightarrow{A\Delta}$, $\overrightarrow{B\Gamma}$, $\overrightarrow{BA}$ και να επιβεβαιώσετε ότι: $\overrightarrow{A\Delta} = \overrightarrow{B\Gamma}$. \monades{12}
    \item Να δείξετε ότι $|\overrightarrow{BA}| = |\overrightarrow{B\Gamma}|$. Ποιο είναι το σχήμα του τετραπλεύρου $AB\Gamma\Delta$; \monades{13}\end{alist}

\vspace{4mm}

% --- Άσκηση 22061 (22061.tex) ---
\Askhsh{}{4}
Δίνεται τετράγωνο ΑΒΓΔ με μήκος πλευράς $a$ ($a > 0$) και κορυφές $A(0, a)$, $B(0,0)$, $Γ(a, 0)$ και $Δ(a, a)$. Το $M$ είναι το μέσο της πλευράς $ΓΔ$ και το τμήμα $BN$ είναι κάθετο στο τμήμα $AM$, όπως φαίνεται στο σχήμα.

\begin{center}
\begin{tikzpicture}[scale=1]
    % Define 'a' for drawing purposes. Let a=4 units for better visualization.
    \pgfmathsetmacro{\aval}{4} % Use \aval instead of 'a' to avoid conflicts.
    % Define points
    \coordinate (A) at (0,\aval);
    \coordinate (B) at (0,0);
    \coordinate (C) at (\aval,0);
    \coordinate (D) at (\aval,\aval);
    % M is the midpoint of Gamma-Delta (C-D)
    \coordinate (M) at (\aval, \aval/2);
    % Calculate N coordinates (xN = 2a/5, yN = 4a/5) for accurate drawing
    % This is for drawing accuracy, not solving the exercise.
    \coordinate (N) at ({2*\aval/5}, {4*\aval/5});
    % Coordinate system setup: Dashed lines as commonly used for axes
    \draw[->, gray, l3, dashed] (-0.5*\aval,0) -- (1.2*\aval,0) node[right] {$x$};
    \draw[->, gray, l3, dashed] (0,-0.5*\aval) -- (0,1.2*\aval) node[above] {$y$};
    % Fill the square with a light green color
    \tkzFillPolygon[fill=thirdcolor!20, opacity=0.3](A,B,C,D)
    % Draw the square A B C D (solid lines as in the image)
    \tkzDrawPolygon[l3](A,B,C,D)
    % Draw segments AM and BN (solid lines as in the image, maincolor for emphasis)
    \tkzDrawSegment[l4, maincolor](A,M)
    \tkzDrawSegment[l4, maincolor](B,N)
    % Draw points (black dots as in default tkz-euclide)
    \tkzDrawPoints[color=black](A,B,C,D,M,N)
    % Add specific square marks for A, N, D as seen in the PDF image
    \node[draw=thirdcolor, fill=thirdcolor!20, inner sep=1.5pt, minimum size=3mm] at (A) {};
    \node[draw=thirdcolor, fill=thirdcolor!20, inner sep=1.5pt, minimum size=3mm] at (N) {};
    \node[draw=thirdcolor, fill=thirdcolor!20, inner sep=1.5pt, minimum size=3mm] at (D) {};
    % Labels for points, positioned to match the PDF
    \tkzLabelPoint[left](A){$A(0,a)$}
    \tkzLabelPoint[below left](B){$B(0,0)$}
    \tkzLabelPoint[below right](C){$Γ(a,0)$}
    \tkzLabelPoint[above right](D){$Δ(a,a)$}
    \tkzLabelPoint[below](M){$M$}
    \tkzLabelPoint[above right](N){$N$} \end{tikzpicture}
\end{center}

\begin{alist}    \item Να βρείτε τις εξισώσεις των ευθειών:
    \begin{rlist}
        \item $AM$ \monades{5}
        \item $BN$ \monades{5}
    \end{rlist}
    \item Να βρείτε τις συντεταγμένες του σημείου $N$. \monades{5}
    \item Να αποδείξετε ότι το σημείο $N$ ανήκει σε κύκλο με κέντρο $Γ$ και ακτίνα ίση με $a$. Να βρείτε την εξίσωση του κύκλου αυτού. \monades{10}\end{alist}

\vspace{4mm}

% --- Άσκηση 22062 (22062.tex) ---
\Askhsh{}{4}
Δίνεται η οικογένεια κύκλων $C_\lambda: (x - \lambda)^2 + (y - \lambda)^2 = \lambda^2$, με $\lambda \neq 0$.

\begin{center}
\begin{tikzpicture}\begin{axis}[
    axis lines = middle,
    xlabel = {$x$},
    ylabel = {$y$},
    xmin = -2.5, xmax = 5.5,
    ymin = -2.5, ymax = 5.5,
    xtick = {-1, 0, 0.5, 1, 2, 3, 4, 5},
    ytick = {-1, 0, 0.5, 1, 2, 3, 4, 5},
    grid = none,
    width = 10cm, height = 8cm,
    tick label style = {font=\scriptsize},
    every axis x label/.style = {at={(current axis.right of origin)}, anchor=north west},
    every axis y label/.style = {at={(current axis.above origin)}, anchor=south east},
]
    % Line y=x
    \addplot[secondarycolor, dotted, thick, domain=-2:5] {x} node[pos=0.8, anchor=south west] {$y=x$};
    % Circles Clambda: center (lambda, lambda), radius |lambda|
    % C4: lambda=4, center (4,4), radius 4
    \draw[l3, maincolor] (axis cs:4,4) circle (4);
    \node[above right, maincolor] at (axis cs:4,4+4) {$C4$};
    \addplot[only marks, mark=*, mark size=1.5pt, black] coordinates {(4,4)};
    % C2: lambda=2, center (2,2), radius 2
    \draw[l3, maincolor] (axis cs:2,2) circle (2);
    \node[above right, maincolor] at (axis cs:2,2+2) {$C2$};
    \addplot[only marks, mark=*, mark size=1.5pt, black] coordinates {(2,2)};
    % C0.5: lambda=0.5, center (0.5,0.5), radius 0.5
    \draw[l3, maincolor] (axis cs:0.5,0.5) circle (0.5);
    \node[above right, maincolor, font=\tiny] at (axis cs:0.5,0.5+0.5) {$C_{0.5}$};
    \addplot[only marks, mark=*, mark size=1.5pt, black] coordinates {(0.5,0.5)};
    % C-1: lambda=-1, center (-1,-1), radius 1
    \draw[l3, maincolor] (axis cs:-1,-1) circle (1);
    \node[below left, maincolor, font=\tiny] at (axis cs:-1,-1-1) {$C_{-1}$};
    \addplot[only marks, mark=*, mark size=1.5pt, black] coordinates {(-1,-1)};
\end{axis}\end{tikzpicture}
\end{center}
\vspace{1mm}
\noindent Οι κύκλοι $C_\lambda$ για $\lambda=-1$, $\lambda=0.5$, $\lambda=2$, $\lambda=4$
\begin{alist}    \item Να βρείτε το κέντρο και την ακτίνα κάθε κύκλου $C_\lambda$, $\lambda \neq 0$. \monades{6}
    \item Να αποδείξετε ότι το κέντρο κάθε κύκλου $C_\lambda$ βρίσκεται στην ευθεία $y = x$. \monades{6}
    \item Να αποδείξετε ότι η ευθεία $x = 0$ εφάπτεται σε όλους τους κύκλους $C_\lambda$, $\lambda \neq 0$. Να εξηγήσετε με συντομία ότι το ίδιο συμβαίνει και για την ευθεία $y = 0$. \monades{6}
    \item Έστω $a \neq 0$. Να αποδείξετε ότι η ευθεία $x = a$ εφάπτεται σε έναν, και μόνο έναν, από τους κύκλους $C_\lambda$. Να εξηγήσετε με συντομία ότι το ίδιο συμβαίνει και για την ευθεία $y = a$. \monades{7}\end{alist}

\vspace{4mm}

% --- Άσκηση 22063 (22063.tex) ---
\Askhsh{K4.001}{4}
\begin{alist}    \item Έστω $\vec{\alpha}$, $\vec{\beta}$ δύο μη μηδενικά διανύσματα. Να αποδείξετε ότι:
    \begin{rlist}
        \item $|\vec{\alpha} + \vec{\beta}| = |\vec{\alpha}| + |\vec{\beta}| \Leftrightarrow \vec{\alpha} \uparrow\uparrow \vec{\beta}.$ \monades{5}
        \item $|\vec{\alpha} + \vec{\beta}| = ||\vec{\alpha}| - |\vec{\beta}|| \Leftrightarrow \vec{\alpha} \uparrow\downarrow \vec{\beta}.$ \monades{5}
    \end{rlist}
    \item Θεωρούμε τρία διανύσματα $\vec{\alpha}$, $\vec{\beta}$, $\vec{\gamma}$ για τα οποία ισχύουν ότι:
 \[ \vec{\alpha} + \vec{\beta} + \vec{\gamma} = \vec{0}, \quad |\vec{\alpha}| = 1, \quad |\vec{\beta}| = 2, \quad |\vec{\gamma}| = 1. \] Να αποδείξετε ότι:
    \begin{rlist}
        \item $\vec{\alpha} \uparrow\uparrow \vec{\gamma}.$ \monades{5}
        \item $\vec{\alpha} \uparrow\downarrow \vec{\beta}.$ \monades{5}
        \item $\vec{\alpha} = \vec{\gamma}$ και $\vec{\beta} = -2\vec{\alpha}.$ \monades{5}
    \end{rlist}\end{alist}

\vspace{4mm}

% --- Άσκηση 22064 (22064.tex) ---
\Askhsh{}{4}

\paragraph{α)} Αν $\vec{AB}$, $\vec{A\Gamma}$, $\vec{A\Delta}$ είναι τρία διανύσματα, τότε οι ποσότητες $\vec{AB} \cdot \vec{A\Gamma}$ και $\vec{AB} \cdot \vec{A\Delta}$ :
\begin{rlist}    \item είναι αριθμοί ή διανύσματα;
    \item μπορούν να συγκριθούν;\end{rlist}
\monades{5}

\paragraph{β)} Για τα διανύσματα του παρακάτω σχήματος να επιλέξετε την σωστή απάντηση:
\begin{rlist}    \item $\vec{AB} \cdot \vec{A\Gamma} > \vec{AB} \cdot \vec{A\Delta}$
    \item $\vec{AB} \cdot \vec{A\Γ} < \vec{AB} \cdot \vec{A\Delta}$
    \item $\vec{AB} \cdot \vec{A\Γ} = \vec{AB} \cdot \vec{A\Delta}$\end{rlist}

\begin{center}
    \begin{tikzpicture}[scale=0.8]
        \tkzDefPoint(0,0){A}
        \tkzDefPoint(5,0){B}
        \tkzDefPoint(2,2.5){G} % Gamma
        \tkzDefPoint(3.5,3.5){D} % Delta
        % Draw segments (vectors)
        \tkzDrawSegments[l3,->](A,G A,D A,B)
        % Projection of Gamma and right angle
        \tkzDefPointWith[projection](G,A,B) \tkzGetPoint{PG}
        \tkzDrawSegment[dashed](G,PG)
        \tkzMarkRightAngle[size=0.2](G,PG,B)
        % Labels
        \tkzLabelPoint[left](A){A}
        \tkzLabelPoint[right](B){B}
        \tkzLabelPoint[above left](G){$\Gamma$}
        \tkzLabelPoint[above right](D){$\Delta$}\end{tikzpicture}
\end{center}
\monades{10}

\paragraph{γ)} Για τα διανύσματα του παρακάτω σχήματος (η διακεκομμένη γραμμή είναι τμήμα κύκλου με κέντρο Α) να επιλέξετε την σωστή απάντηση:
\begin{rlist}    \item $\vec{AB} \cdot \vec{A\Gamma} > \vec{AB} \cdot \vec{A\Delta}$
    \item $\vec{AB} \cdot \vec{A\Γ} < \vec{AB} \cdot \vec{A\Delta}$
    \item $\vec{AB} \cdot \vec{A\Γ} = \vec{AB} \cdot \vec{A\Delta}$\end{rlist}

\begin{center}
    \begin{tikzpicture}[scale=0.8]
        \tkzDefPoint(0,0){A}
        \tkzDefPoint(5,0){B}
        % Define points G and D on a circle of radius 4 centered at A
        \def\myradius{4}
        \tkzDefPointAt[euclid](A,\myradius){Pbase} % Helper point at radius on x-axis
        \tkzDefPointBy[rotation=center A angle 25](Pbase) \tkzGetPoint{G} % Gamma at 25 deg
        \tkzDefPointBy[rotation=center A angle 40](Pbase) \tkzGetPoint{D} % Delta at 40 deg (25+15)
        % Draw vectors
        \tkzDrawSegments[l3,->](A,G A,D A,B)
        % Draw dashed arc from G to D
        \tkzDrawArc[dashed,l3](A,G)(D) % Arc from G to D with center A
        % Mark angles
        \tkzMarkAngle[fill=maincolor!20,size=0.8cm](B,A,G)
        \tkzLabelAngle[pos=0.5](B,A,G){$25^\circ$}
        \tkzMarkAngle[fill=secondarycolor!20,size=0.8cm](G,A,D) % Angle between G and D
        \tkzLabelAngle[pos=0.5](G,A,D){$15^\circ$}
        % Labels
        \tkzLabelPoint[left](A){A}
        \tkzLabelPoint[right](B){B}
        \tkzLabelPoint[above left](G){$\Gamma$}
        \tkzLabelPoint[above right](D){$\Delta$}\end{tikzpicture}
\end{center}
\monades{10}

Σε όλα τα ερωτήματα, να αιτιολογήσετε τις απαντήσεις σας.

\vspace{4mm}

% --- Άσκηση 22065 (22065.tex) ---
\Askhsh{0000}{4} % Placeholder for a specific exercise code, Theme 4
\noindent Δίνονται τα σημεία $A(0,0)$, $B(8,0)$, $\Gamma(10,4)$, $\Delta(2,4)$. Τα σημεία $M$ και $N$ είναι τα μέσα των πλευρών $AB$ και $B\Gamma$ αντίστοιχα, ενώ $K$ και $\Lambda$ είναι τα σημεία που τέμνουν τα τμήματα $\Delta M$ και $\Delta N$ την διαγώνιο $A\Gamma$ αντίστοιχα, όπως φαίνεται στο σχήμα.

\begin{center}
\begin{tikzpicture}[scale=0.8]
    % Ορισμός κορυφών του τετραπλεύρου
    \tkzDefPoint(0,0){A}
    \tkzDefPoint(8,0){B}
    \tkzDefPoint(10,4){C} % Χρησιμοποιούμε C ως μεταβλητή για το Γ
    \tkzDefPoint(2,4){D} % Χρησιμοποιούμε D ως μεταβλητή για το Δ
    % Ορισμός των μέσων M και N
    \tkzDefMidPoint(A,B) \tkzGetPoint{M}
    \tkzDefMidPoint(B,C) \tkzGetPoint{N}
    % Σχεδίαση του τετραπλεύρου
    \tkzDrawSegments[l3](A,B B,C C,D D,A)
    % Σχεδίαση της διαγωνίου AΓ
    \tkzDrawSegment[l3](A,C)
    % Σχεδίαση των τμημάτων ΔM και ΔN
    \tkzDrawSegments[l3](D,M D,N)
    % Εύρεση των σημείων τομής K και Λ
    \tkzInterLL(D,M)(A,C) \tkzGetPoint{K}
    \tkzInterLL(D,N)(A,C) \tkzGetPoint{L} % Χρησιμοποιούμε L ως μεταβλητή για το Λ
    % Σχεδίαση των σημείων
    \tkzDrawPoints(A,B,C,D,M,N,K,L)
    % Ετικέτες των σημείων με τοποθέτηση που να ταιριάζει με την εικόνα
    \tkzLabelPoint[below left](A){$A$}
    \tkzLabelPoint[below right](B){$B$}
    \tkzLabelPoint[above right](C){$\Gamma$}
    \tkzLabelPoint[above left](D){$\Delta$}
    \tkzLabelPoint[below](M){$M$}
    \tkzLabelPoint[right](N){$N$}
    \tkzLabelPoint[left](K){$K$}
    \tkzLabelPoint[right](L){$\Lambda$}\end{tikzpicture}
\end{center}

\begin{alist}    \item Να δείξετε ότι το τετράπλευρο $AB\Gamma\Delta$ είναι παραλληλόγραμμο. \monades{5}
    \item Να βρείτε τις συντεταγμένες των σημείων $M$ και $N$. \monades{5}
    \item Να βρείτε τις εξισώσεις των ευθειών $A\Gamma$, $\Delta M$, $\Delta N$ και στη συνέχεια τις συντεταγμένες των σημείων $K$ και $\Lambda$. \monades{10}
    \item Να δείξετε ότι τα σημεία $K$ και $\Lambda$ τριχοτομούν την διαγώνιο $A\Gamma$, δηλαδή την χωρίζουν σε τρία ίσα τμήματα. \monades{5}\end{alist}

\vspace{4mm}

% --- Άσκηση 22066 (22066.tex) ---
\section*{\full{\OnPar{ΘΕΜΑ 4}}}

Θεωρούμε το σημείο $A = (\sigma\upsilon\nu\theta, \eta\mu\theta)$ και το σημείο $B = (\sigma\upsilon\nu\theta - \eta\mu\theta, \eta\mu\theta + \sigma\upsilon\nu\theta)$, όπου $\theta \in [0,2\pi)$.

\begin{alist}    \item Να αποδείξετε ότι τα σημεία $A$ και $B$ ανήκουν σε δύο κύκλους με κέντρο την αρχή των αξόνων $O(0,0)$. Να βρείτε τις ακτίνες των δύο κύκλων.
    \item Να αποδείξετε ότι: $\overrightarrow{OA} \perp \overrightarrow{AB}$.  \monades{8}
    \item Να βρείτε το μέτρο της γωνίας $\omega$ μεταξύ των διανυσμάτων $\overrightarrow{OA}$ και $\overrightarrow{OB}$.  \monades{8}
 \monades{9}\end{alist}

\begin{center}
\begin{tikzpicture}[scale=1.2]
    % Axes
    \draw[->, gray] (-2.2,0) -- (2.2,0);
    \draw[->, gray] (0,-2.2) -- (0,2.2);
    % Define points
    % For drawing purposes, let's pick a specific theta, e.g., theta = pi/4 (45 degrees)
    % A = (cos(45), sin(45)) = (sqrt(2)/2, sqrt(2)/2) -> polar(1,45)
    % B = (cos(45)-sin(45), sin(45)+cos(45)) = (0, sqrt(2)) -> polar(sqrt(2),90)
    \tkzDefPoint(0,0){O}
    \tkzDefPointByPolar(1,45){A}
    \tkzDefPointByPolar(sqrt(2),90){B}
    % Draw circles (dashed, gray)
    \tkzDrawCircle[dashed, color=gray](O,A)
    \tkzDrawCircle[dashed, color=gray](O,B)
    % Draw segments
    \tkzDrawSegments[l3,maincolor](O,A O,B A,B)
    % Mark right angle at A
    \tkzMarkRightAngle[fill=thirdcolor!20,size=0.25,opacity=0.8](O,A,B)
    % Mark angle omega
    \tkzMarkAngle[fill=thirdcolor!20,size=0.4,opacity=0.8](A,O,B)
    \tkzLabelAngle[pos=0.6](A,O,B){$\omega$}
    % Label points
    \tkzLabelPoint[below left](O){O}
    \tkzLabelPoint[above right](A){A}
    \tkzLabelPoint[above right](B){B}
    % Draw points
    \tkzDrawPoints(O,A,B)\end{tikzpicture}
\end{center}

\vspace{4mm}

% --- Άσκηση 22067 (22067.tex) ---
\Askhsh{}{4}
Θεωρούμε μία ευθεία $\varepsilon: y = \lambda x$ με θετική κλίση $\lambda$.

\begin{center}
\begin{tikzpicture}[>=latex]
    % Coordinate axes lines
    \draw[->, gray!70!black] (-3.5,0) -- (3.5,0) node[below right] {$x$};
    \draw[->, gray!70!black] (0,-2.5) -- (0,3.5) node[above left, yshift=2mm] {$y$};
    % Extend negative x-axis as dotted and label x'
    \draw[dotted, gray!70!black] (-3.5,0) node[below left] {$x'$} -- (0,0);
    % Origin
    \node[fill=white, inner sep=0.5mm] at (0,0) {O}; 
    % Line epsilon: y = lambda x, with lambda > 0
    % Visual representation chosen for a positive slope
    \draw[l3, maincolor] (-2.5,-1.25) -- (3.5,1.75) node[above right] {$\varepsilon$}; 
    % Bisector delta1 (acute angle, dotted, secondarycolor)
    % Visual representation chosen to be below epsilon
    \draw[l3, dotted, secondarycolor] (0,0) -- (3.5,0.7) node[above right] {$\delta1$};
    % Bisector delta2 (obtuse angle, solid, thirdcolor)
    % Visual representation chosen to be in the second quadrant
    \draw[l3, thirdcolor] (0,0) -- (-2.5,1.8) node[above left] {$\delta2$};\end{tikzpicture}
\end{center}

\begin{alist}    \item Αν $\delta_1$ είναι η διχοτόμος της οξείας γωνίας που σχηματίζει η ευθεία $\varepsilon$ με τον $x'x$ άξονα, τότε να αποδείξετε ότι η εξίσωση της διχοτόμου $\delta_1$ είναι:
    \item Αν $\delta_2$ είναι η διχοτόμος της αμβλείας γωνίας που σχηματίζει η ευθεία $\varepsilon$ με τον $x'x$ άξονα, τότε να αποδείξετε ότι η εξίσωση της διχοτόμου $\delta_2$ είναι:  \[ y = \lambda_1 x \quad \text{με} \quad \lambda_1 = \frac{\lambda}{1+\sqrt{1+\lambda^2}} \] \monades{12}
    \item Αν $\lambda = 1$, να εφαρμόσετε τους τύπους του α) ερωτήματος για να αποδείξετε ότι:  \[ y = \lambda_2 x \quad \text{με} \quad \lambda_2 = \frac{\lambda}{1-\sqrt{1+\lambda^2}} \] \monades{7}
 \[ \text{εφ}22,5^\circ = \sqrt{2} - 1 \] \monades{6}\end{alist}

\vspace{4mm}

% --- Άσκηση 22068 (22068.tex) ---
\Askhsh{}{4}

Σε ένα υλικό σημείο O εφαρμόζονται τρεις δυνάμεις $\vec{F_1}$, $\vec{F_2}$, $\vec{F_3}$ οι οποίες σχηματίζουν ανά δύο γωνία $120^\circ$, έτσι ώστε το υλικό σημείο O να ισορροπεί.

\begin{center}
\begin{tikzpicture}[scale=1]
    % Ορισμός του σημείου Ο
    \tkzDefPoint(0,0){O}
    % Ορισμός των σημείων A, B, Γ για τις δυνάμεις
    % Υποθέτουμε ίσες δυνάμεις για τη σχηματική απεικόνιση και τοποθετούμε F2 στον θετικό άξονα x
    \tkzDefPoint(3,0){B} % Πέρατο για την F2
    \tkzDefPoint(120:3){A} % Πέρατο για την F1 (γωνία 120 μοίρες από τον άξονα x)
    \tkzDefPoint(240:3){G} % Πέρατο για την F3 (γωνία 240 μοίρες από τον άξονα x)
    % Σχεδίαση των διανυσμάτων (τμημάτων με βέλη) με το χρώμα maincolor
    \tkzDrawSegment[l3,->,maincolor](O,A)
    \tkzDrawSegment[l3,->,maincolor](O,B)
    \tkzDrawSegment[l3,->,maincolor](O,G)
    % Σχεδίαση των σημείων
    \tkzDrawPoints(O,A,B,G)
    % Ετικέτες για τα σημεία
    \tkzLabelPoint[above left](A){A}
    \tkzLabelPoint[right](B){B}
    \tkzLabelPoint[below left](G){$\Gamma$} % Χρήση \Gamma για το ελληνικό γράμμα
    \tkzLabelPoint[below right](O){O} 
    % Ετικέτες για τις δυνάμεις F1, F2, F3 όπως φαίνονται στο PDF (απλό κείμενο, ελαφρώς μικρότερο μέγεθος)
    \node[above right] at ($(O)!0.6!(A)$) {\footnotesize F1};
    \node[above] at ($(O)!0.6!(B)$) {\footnotesize F2};
    \node[below right] at ($(O)!0.6!(G)$) {\footnotesize F3};\end{tikzpicture}
\end{center}

\begin{alist}    \item Ποια σχέση ανάμεσα στα διανύσματα $\vec{F_1}$, $\vec{F_2}$, $\vec{F_3}$ εκφράζει την συνθήκη ισορροπίας; \monades{5}
    \item Να αποδείξετε ότι τα διανύσματα $\vec{F_1} + \vec{F_2}$ και $\vec{F_3}$ είναι αντίθετα. \monades{5}
    \item Αν A, B, $\Gamma$, $\Delta$ είναι τα πέρατα των διανυσμάτων $\vec{F_1}$, $\vec{F_2}$, $\vec{F_3}$ και $\vec{F_1} + \vec{F_2}$, αντίστοιχα (θεωρούμενων ως διανυσμάτων με αρχή το σημείο O), τότε να αποδείξετε ότι:
    \begin{rlist}
        \item $\widehat{AO\Delta} = \widehat{BO\Delta} = 60^\circ$. \monades{5}
        \item $\widehat{O\Delta B} = 60^\circ$. \monades{5}
    \end{rlist}
    \item Να αποδείξετε ότι: $|\vec{F_1}| = |\vec{F_2}| = |\vec{F_3}|$. \monades{5}\end{alist}

\vspace{4mm}

% --- Άσκηση 22069 (22069.tex) ---
\Askhsh{}{4}
Τα σημεία $A(3, 2)$ και $B(6, 1)$ βρίσκονται πάνω σε έναν κύκλο $C$ από το κέντρο $K$ του οποίου διέρχεται η ευθεία $\varepsilon: y = 2x - 7$. Να βρείτε:

\begin{center}
\begin{tikzpicture}[scale=1.2]
    % Ορισμός του κέντρου K (αυθαίρετα για το σκίτσο, όχι από υπολογισμούς)
    \tkzDefPoint(0,0){K}
    % Ορισμός ενός σημείου στην περιφέρεια για την ακτίνα (αυθαίρετα)
    \tkzDefPoint(2.0,1.5){Rpoint} 
    % Σχεδίαση του κύκλου C
    \tkzDrawCircle[dashed,l3,maincolor](K,Rpoint)
    % Ετικέτα C (ως κόμβος, όχι στην διαδρομή του κύκλου, για να ταιριάζει με το PDF)
    \node[above=0.3cm] at ($(K) + (0, \tkzCalcLength[cm](K,Rpoint)) $) {$C$};
    % Ορισμός σημείων A και B στον κύκλο (αυθαίρετες θέσεις για το σκίτσο)
    \tkzDefPointBy[rotation=center K angle 130](Rpoint){Asketch}
    \tkzDefPointBy[rotation=center K angle 30](Rpoint){Bsketch}
    % Σχεδίαση των σημείων A, B, K
    \tkzDrawPoints(Asketch, Bsketch, K)
    \tkzLabelPoint[left](Asketch){A}
    \tkzLabelPoint[right](Bsketch){B}
    \tkzLabelPoint[above left](K){K}
    % Σχεδίαση της ευθείας ε που διέρχεται από το K
    % Για να μοιάζει με την εικόνα, η γραμμή πηγαίνει από κάτω-αριστερά προς πάνω-δεξιά
    \tkzDefPoint(-2.5,-1.7){P1} 
    \tkzDefPoint(2.5,1.7){P2}
    \tkzDrawLine[add=0.5 and 0.5, l3, secondarycolor](P1,P2)
    \tkzLabelLine[pos=1.1, above right](P1,P2){$\varepsilon$} \end{tikzpicture}
\end{center}

\begin{alist}
    \item τις συντεταγμένες του κέντρου $K$ του κύκλου $C$. \monades{12}
    \item την ακτίνα $R$ του κύκλου $C$. \monades{8}
    \item την εξίσωση του κύκλου $C$. \monades{5}
\end{alist}

\vspace{4mm}

% --- Άσκηση 22071 (22071.tex) ---
\Askhsh{}{2}
Οι πλευρές $AB$ και $A\Delta$ ενός παραλληλογράμμου $AB\Gamma\Delta$ έχουν εξισώσεις $x+2y+1=0$ και $2x+y+5=0$ αντίστοιχα και το κέντρο του παραλληλογράμμου είναι το σημείο $K(1,2)$.
\begin{alist}    \item Να αποδείξετε ότι:
    \begin{rlist}
        \item Η κορυφή $A$ του παραλληλογράμμου έχει συντεταγμένες $A(-3, 1)$. \monades{8}
        \item Η κορυφή $\Gamma$ του παραλληλογράμμου έχει συντεταγμένες $\Gamma(5, 3)$. \monades{7}
    \end{rlist}
    \item Να βρείτε τις εξισώσεις των άλλων δύο πλευρών του $B\Gamma$ και $\Gamma\Delta$. \monades{10}\end{alist}

\vspace{4mm}

% --- Άσκηση 22072 (22072.tex) ---
\OnPar{ΘΕΜΑ 2}

Δίνονται οι εξισώσεις (1): $\lambda x+(\lambda-1)y-4=0$ και (2): $(3\lambda+1)x-2\lambda y-7=0$ με $\lambda \in \mathbb{R}$.

\begin{alist}    \item Να αποδείξετε ότι οι εξισώσεις (1) και (2) παριστάνουν εξισώσεις ευθειών για κάθε $\lambda \in \mathbb{R}$. \monades{15}
    \item Να βρείτε τις τιμές του $\lambda \in \mathbb{R}$, ώστε οι ευθείες με εξισώσεις τις (1) και (2) να είναι μεταξύ τους κάθετες. \monades{10}\end{alist}

\vspace{4mm}

% --- Άσκηση 22073 (22073.tex) ---
\Askhsh{}{4}
Σε χάρτη με καρτεσιανό σύστημα συντεταγμένων η θέση ενός λιμανιού προσδιορίζεται από το σημείο $\Lambda(2,6)$ και η θέση ενός πλοίου με το σημείο $\Pi(\lambda-1, 2+\lambda)$, $\lambda \in \mathbb{R}$.
\begin{alist}    \item
    \begin{rlist}
        \item Αν το πλοίο κινείται ευθύγραμμα, να βρείτε την εξίσωση της τροχιάς του. \monades{7}
        \item Να εξετάσετε αν το πλοίο θα περάσει από το λιμάνι. \monades{5}
    \end{rlist}
    \item Αν τελικά το πλοίο δεν περάσει από το λιμάνι, να βρείτε:
    \begin{rlist}
        \item Ποια είναι η ελάχιστη απόσταση του πλοίου από το λιμάνι; \monades{6}
        \item Το σημείο του καρτεσιανού επιπέδου που βρίσκεται το πλοίο, όταν απέχει την ελάχιστη απόσταση από το λιμάνι. \monades{7}
    \end{rlist}\end{alist}

\vspace{4mm}

% --- Άσκηση 22092 (22092.tex) ---
\section*{ΘΕΜΑ 2}
Δίνεται τετράπλευρο \textAlpha\textBeta\textGamma\textDelta\ με κορυφή $A(1,4)$. Η πλευρά $A\Delta$ έχει εξίσωση $3x - 2y + 5 = 0$ και η διαγώνιος $B\Delta$ έχει εξίσωση $y = x + 2$.

\begin{alist}    \item Να αποδείξετε ότι η κορυφή $\Delta$ έχει συντεταγμένες $\Delta(-1,1)$. \monades{12}
    \item Αν οι διαγώνιοι $A\Gamma$ και $B\Delta$ του τετραπλεύρου τέμνονται κάθετα, να βρείτε την εξίσωση της διαγωνίου $A\Gamma$. \monades{13}\end{alist}

\vspace{4mm}

% --- Άσκηση 22147 (22147.tex) ---
\OnPar{ΘΕΜΑ 2}

Δίνεται η εξίσωση
\[ x^2 + y^2 - x - y - \frac{7}{2} = 0 \quad (1) \]

\begin{alist}    \item Να αποδείξετε ότι η εξίσωση (1) παριστάνει κύκλο με κέντρο $K \left(\frac{1}{2}, \frac{1}{2}\right)$ και ακτίνα $R = 2$.
    \item Να αποδείξετε ότι το σημείο $A \left(\frac{1}{2}, -\frac{3}{2}\right)$ είναι σημείο του κύκλου $(K,R)$.  \monades{9}
    \item Να βρείτε την εξίσωση της εφαπτομένης του κύκλου $(K,R)$ στο $A$.  \monades{8}
 \monades{8}\end{alist}

\vspace{4mm}

% --- Άσκηση 22168 (22168.tex) ---
\Askhseis{ΘΕΜΑ 2}

Δίνονται η παραβολή $y^2 = 2px$ και η έλλειψη $\frac{x^2}{\alpha^2} + \frac{y^2}{\beta^2}=1$.

\begin{alist}    \item Αν η παραβολή διέρχεται από το σημείο $A(1,2)$, να βρείτε:
    \begin{rlist}
        \item Την εξίσωση της παραβολής. \monades{10}
        \item Την εστία $E$ της παραβολής. \monades{05}
    \end{rlist}
    \item Να βρεθεί η εξίσωση της έλλειψης με κέντρο το $O$, αν η μια εστία της είναι το σημείο $E(1,0)$ και ο μεγάλος άξονας της έχει μήκος ίσο με $4$. \monades{10}\end{alist}

\vspace{4mm}

% --- Άσκηση 22169 (22169.tex) ---
ΘΕΜΑ 2

Δίνεται η υπερβολή $\frac{x^2}{\alpha^2} - \frac{y^2}{\beta^2} = 1$ με ασύμπτωτη την $y = \frac{3}{4} x$. Η απόσταση των κορυφών της $A$ και $A'$ είναι $8$.

\begin{alist}    \item 
    \begin{rlist}
        \item Να βρείτε την εξίσωση της υπερβολής. \monades{10}
        \item Ποιες είναι οι εστίες της υπερβολής; \monades{05}
    \end{rlist}
    \item Να βρείτε την εφαπτομένη της $\frac{x^2}{16} - \frac{y^2}{9} = 1$ στο σημείο της $\left(5, \frac{9}{4}\right)$. \monades{10}\end{alist}

\vspace{4mm}

% --- Άσκηση 22170 (22170.tex) ---
ΘΕΜΑ 2

Δίνονται τα διανύσματα $\vec{\alpha} = (-1,3)$, $\vec{\beta} = \left(-2, -\frac{1}{2}\right)$ και $\vec{\nu} = (x^2, x-1)$.

\begin{alist}    \item Να βρείτε τις συντεταγμένες του διανύσματος $\vec{u} = \vec{\alpha} - 2\vec{\beta}$. \monades{07}
    \item Να βρείτε τους αριθμούς $x$ για τους οποίους τα διανύσματα $\vec{u} = (3,4)$ και $\vec{\nu}$ είναι κάθετα. \monades{09}
    \item Να βρείτε τους αριθμούς $x$ για τους οποίους τα διανύσματα $\vec{\nu}$ και $\vec{\beta}$ είναι συγγραμμικά. \monades{09}\end{alist}

\vspace{4mm}

% --- Άσκηση 22171 (22171.tex) ---
\Askhsh{01}{2}
Δίνονται οι ευθείες $\varepsilon1: 3x - y = 5$ και $\varepsilon2: x - y + 1 = 0$.
\begin{alist}    \item Να βρεθεί το σημείο τομής τους $M$. \monades{10}
    \item Να βρεθεί η εξίσωση της ευθείας που διέρχεται από το σημείο $M(3,4)$ και είναι κάθετη στην $(\varepsilon2)$. \monades{10}
    \item Να βρεθεί ένα διάνυσμα παράλληλο στην $(\varepsilon1)$. \monades{05}\end{alist}

\vspace{4mm}

% --- Άσκηση 22172 (22172.tex) ---
\textbf{ΘΕΜΑ 2}

Θεωρούμε την ευθεία $\varepsilon: 3x - 4y = 0$ και το σημείο $A(-2,1)$.

\begin{alist}    \item Να αποδείξετε ότι η απόσταση του σημείου $A$ από την ευθεία είναι 2. \monades{08}
    \item Να βρείτε την εξίσωση ευθείας $(\eta)$ κάθετης στην $(\varepsilon)$ που διέρχεται από το σημείο $A$. \monades{10}
    \item Να βρείτε την εξίσωση του κύκλου που έχει κέντρο το σημείο $A$ και εφάπτεται στην ευθεία $(\varepsilon)$. \monades{07}\end{alist}

\vspace{4mm}

% --- Άσκηση 22173 (22173.tex) ---
\Askhsh{ALGB21234}{2} % Το "ALG-B-2-1234" είναι ένας ενδεικτικός κωδικός άσκησης.

Δίνεται το τετράγωνο ΑΒΓΟ με κορυφές τα σημεία $A(0,4)$, $B(4,4)$, $\Gamma(4,0)$, $O(0,0)$. Στην διαγώνιο ΑΓ παίρνουμε σημείο Ζ, τέτοιο ώστε $AZ = \frac{3}{4} A\Gamma$. Επίσης, θεωρούμε το μέσο $\Delta$ της ΒΓ.

\begin{alist}    \item Να βρείτε:
    \begin{rlist}
        \item Τις συντεταγμένες του σημείου $\Delta$. \monades{07}
        \item Τις συντεταγμένες του σημείου Ζ. \monades{09}
    \end{rlist}
    \item Αν το σημείο $\Delta$ είναι το $(4,2)$ και το σημείο Ζ το $(3,1)$, να αποδείξετε ότι η ευθεία Ζ$\Delta$ είναι κάθετη στην ευθεία ΑΓ. \monades{09}\end{alist}

\begin{center}
\begin{tikzpicture}    % Αρχικοποίηση του συστήματος συντεταγμένων
    \tkzInit[xmin=-1,xmax=5,ymin=-1,ymax=5]
    % Σχεδίαση των αξόνων χωρίς διαγραμμίσεις (ticks)
    \draw[l3,->] (-5,0) -- (5,0) node[below right] {$x$};
    \draw[l3,->] (0,-5) -- (0,5) node[above left] {$y$};
    % Ορισμός των κορυφών του τετραγώνου
    % Χρησιμοποιούμε G για το Γ (Gamma)
    \tkzDefPoint(0,0){O}
    \tkzDefPoint(0,4){A}
    \tkzDefPoint(4,4){B}
    \tkzDefPoint(4,0){G} 
    % Σχεδίαση του τετραγώνου ΑΒΓΟ με το καθορισμένο πάχος γραμμής (l3)
    \tkzDrawPolygon[l3](A,B,G,O)
    % Σχεδίαση της διαγωνίου ΑΓ
    \tkzDrawSegment[l3](A,G)
    % Ορισμός του σημείου Δ ως το μέσο της ΒΓ
    \tkzDefMidPoint(B,G) \tkzGetPoint{D}
    % Ορισμός του σημείου Ζ στην ΑΓ, τέτοιο ώστε ΑΖ = (3/4)ΑΓ
    % Η εντολή tkzDefPointByFCT(P1, P2, factor) ορίζει ένα σημείο P3 στο τμήμα P1P2
    % έτσι ώστε P1P3 = factor * P1P2.
    \tkzDefPointByFCT(A,G,0.75){Z} 
    % Σχεδίαση όλων των σημείων (χρησιμοποιώντας τις καθολικές ρυθμίσεις του tkzSetUpPoint)
    \tkzDrawPoints(O,A,B,G,D,Z)
    % Τοποθέτηση ετικετών στα σημεία όπως στην εικόνα
    \tkzLabelPoint[left](A){A}
    \tkzLabelPoint[right](B){B}
    \tkzLabelPoint[below](G){$\Gamma$}
    \tkzLabelPoint[above right](D){$\Delta$} 
    \tkzLabelPoint[above left](Z){Z}
    % Το σημείο Ο (αρχή των αξόνων) δεν φέρει ξεχωριστή ετικέτα με γράμμα στην εικόνα.
\end{tikzpicture}
\end{center}

\vspace{4mm}

% --- Άσκηση 22174 (22174.tex) ---
\OnPar{ΘΕΜΑ 4}

Πλανήτης κινείται πάνω σε επίπεδο, ελλειπτικά γύρω από τον ήλιο του. Στο καρτεσιανό επίπεδο ο ήλιος βρίσκεται στην εστία της έλλειψης Ε($\gamma$,0), ενώ η άλλη εστία είναι στο Ε΄(-$\gamma$,0). Η εκκεντρότητα της τροχιάς είναι $0,6$ ενώ ο μεγάλος άξονας $10$.

\begin{alist}
\item Να βρεθεί η εξίσωση της τροχιάς. \monades{09}

\item Θεωρούμε ότι ο πλανήτης κινείται πάνω στην $\frac{x^2}{25} + \frac{y^2}{16} = 1$.
    \begin{rlist}
    \item Τη στιγμή που ο πλανήτης βρίσκεται στο σημείο $\Gamma \left(3, \frac{16}{5}\right)$ εκπέμπεται από αυτόν σήμα που κινείται κατά τη διεύθυνση της εφαπτομένης της τροχιάς του προς τη μεριά του άξονα Oy. Να εξετάσετε αν αυτό το σήμα θα περάσει από το σημείο $\Delta(0,5)$. \monades{09}
    \item Κομήτης κινείται στο ίδιο επίπεδο με τον πλανήτη και πάνω στην καμπύλη $\frac{x^2}{25} - \frac{y^2}{16} = 1$ με $x > 0$. Ποια είναι τα σημεία συνάντησης των δύο τροχιών; \monades{07}
    \end{rlist}
\end{alist}

\begin{center}
\begin{tikzpicture}[scale=0.8]
    % Draw the main ellipse (schematic representation)
    \draw[l3, maincolor] (0,0) ellipse (4.5cm and 3.5cm);
    % Define and draw the foci (schematic positions)
    \tkzDefPoint(-2,0){Ep} % E'
    \tkzDefPoint(2,0){E}   % E
    \tkzDrawPoints(E, Ep)
    % Labels for the foci
    \tkzLabelPoint[left](Ep){E'}
    \tkzLabelPoint[right](E){E}
    \node[below] at (E) {ήλιος};
    % Draw horizontal and vertical reference lines (schematic axes)
    \draw[gray!50, dashed] (-5,0) -- (5,0); % x-axis
    \draw[gray!50, dashed] (0,-4) -- (0,4); % y-axis
    % Define and draw point Gamma (schematic position in the upper-right quadrant)
    \tkzDefPoint(2.5,2.5){G}
    \tkzDrawPoint(G)
    \tkzLabelPoint[above right](G){$\Gamma$}
    % Draw a schematic tangent line at Gamma
    % The tangent is drawn as a line segment passing through G with a visual negative slope,
    % consistent with the sketch in the PDF, without revealing the solution to part b.i.
    \coordinate (T1) at ($(G) + (-1, 0.6)$);
    \coordinate (T2) at ($(G) + (1, -0.6)$);
    \draw[l3, secondarycolor] (T1) -- (T2);\end{tikzpicture}
\end{center}

\vspace{4mm}

% --- Άσκηση 22190 (22190.tex) ---
\Askhsh{}{2}
Δίνεται η παραβολή (C) με εξίσωση
\[ y^2 = x \quad (1) \]
\begin{alist}    \item Να βρείτε τις συντεταγμένες της εστίας $E$ και την εξίσωση της διευθετούσας ($\delta$). \monades{12}
    \item Να αποδείξετε ότι το σημείο $A(1,-1)$ είναι σημείο της παραβολής. \monades{05}
    \item Να βρείτε την εξίσωση της εφαπτομένης παραβολής στο σημείο της $A(1,-1)$. \monades{08}\end{alist}

\vspace{4mm}

% --- Άσκηση 22192 (22192.tex) ---
\Askhsh{}{2}
Δίνεται η έλλειψη $(C)$ με εξίσωση
\[ \frac{x^2}{225} + \frac{y^2}{81} = 1 \quad (1) \]
\begin{alist}    \item Να βρείτε τις συντεταγμένες των εστιών $E$ και $E'$. \monades{10}
    \item Να αποδείξετε ότι το σημείο $B(0,9)$ είναι σημείο της έλλειψης. \monades{05}
    \item Να βρείτε την εξίσωση της εφαπτομένης έλλειψης στο σημείο της $B(0,9)$. \monades{10}\end{alist}

\vspace{4mm}

% --- Άσκηση 22196 (22196.tex) ---
\Askhsh{}{2}

Δίνεται η υπερβολή $(C)$ με εξίσωση
\[ x^2 - y^2 = 25 \quad (1) \]

\begin{alist}    \item Να βρείτε τις συντεταγμένες των εστιών $E$ και $E'$. \monades{10}
    \item Να βρείτε τις ασύμπτωτες $(\varepsilon1)$, $(\varepsilon2)$ της υπερβολής. \monades{10}
    \item Τι γωνία σχηματίζουν οι ασύμπτωτες $(\varepsilon1)$, $(\varepsilon2)$; Να αιτιολογήσετε την απάντησή σας. \monades{05}\end{alist}

\vspace{4mm}

% --- Άσκηση 22214 (22214.tex) ---
\paragraph{ΘΕΜΑ 4}
Σε καρτεσιανό επίπεδο $Oxy$ θεωρούμε τα σημεία $M(x,y)$, $A(-1,0)$, $B(1,0)$ για τα οποία ισχύει $AM^2 + BM^2 = 9|AB|$.
\begin{alist}    \item Να αποδείξετε ότι ο γεωμετρικός τόπος των σημείων $M$ είναι ο κύκλος με εξίσωση $x^2 + y^2 = 8$. \monades{10}
    \item Έστω $\Gamma$ και $\Delta$ δύο σημεία του κύκλου τέτοια ώστε $\Gamma\Delta^2 = 32$.
    \begin{rlist}
        \item Να δείξετε ότι τα σημεία $\Gamma$ και $\Delta$ και η αρχή των αξόνων είναι συνευθειακά σημεία. \monades{08}
        \item Αν το σημείο $M$ κινείται στον κύκλο, να υπολογίσετε το $M\Gamma \cdot M\Delta$. \monades{07}
    \end{rlist}\end{alist}

\vspace{4mm}

% --- Άσκηση 22223 (22223.tex) ---
\Askhsh{}{4}
Σε τρίγωνο $\Alpha\Beta\Gamma$ είναι $\vec{\Alpha\Beta} = (\lambda, \lambda + 1)$, $\vec{\Alpha\Gamma} = (3\lambda, \lambda - 1)$ και το σημείο $\Mu$ είναι το μέσο της $\Beta\Gamma$ με $\lambda \in \mathbb{R}$.
\begin{alist}    \item Να αποδείξετε ότι $\vec{\Alpha\Mu} = (2\lambda, \lambda)$ \monades{08}
    \item Δίνεται επιπλέον ότι η γωνία $\widehat{\Beta\Alpha\Gamma} = 90^\circ$. \monades{08}
    \begin{rlist}
        \item Να υπολογίσετε το $\lambda$.
        \item Αν $\lambda = \frac{1}{2}$ και $\Alpha(2, -\frac{3}{2})$ να βρείτε την εξίσωση του περιγεγραμμένου κύκλου του τριγώνου $\Alpha\Beta\Gamma$. \monades{09}
    \end{rlist}\end{alist}

\vspace{4mm}

% --- Άσκηση 22239 (22239.tex) ---
\Askhsh{}{4}
Σε καρτεσιανό σύστημα αξόνων $Oxy$ η εξίσωση $3x + 4y = 25$ περιγράφει τη θέση ενός αγωγού ύδρευσης. Σε αυτό το σύστημα θέλουμε να σχεδιάσουμε ένα κυκλικό σιντριβάνι με κέντρο το $O(0,0)$ και ακτίνα $2$.

\begin{alist}    \item Α)
    \begin{rlist}
        \item Ποια είναι η εξίσωση του κύκλου που περιγράφει την θέση του σιντριβανιού; \monades{04}
        \item Να εξετάσετε αν ο αγωγός ύδρευσης διέρχεται από το κέντρο του σιντριβανιού, προκειμένου να ενωθεί με αυτό. \monades{05}
        \item Αν ο αγωγός ύδρευσης δεν διέρχεται από το κέντρο του σιντριβανιού, ποιο σημείο του αγωγού ύδρευσης πρέπει να ενωθεί με το κέντρο του σιντριβανιού ώστε να έχουμε την μικρότερη δυνατή απόσταση, άρα και οικονομικότερη κατασκευή; \monades{08}
    \end{rlist}
    \item Β) Ο μηχανικός που θέλει να χαράξει έναν ευθύγραμμο δρόμο, κατέληξε στην εξίσωση $\lambda x + y + x - 2 = 0$, με $\lambda \ne 0$. Μπορείτε να τον βοηθήσετε να βρει για ποια τιμή του $\lambda$ ο δρόμος αυτός εφάπτεται του σιντριβανιού; \monades{08}\end{alist}

\vspace{4mm}

% --- Άσκηση 22262 (22262.tex) ---
\Askhsh{1}{4}
Δίνεται τρίγωνο ΑΒΓ με κορυφές τα σημεία $A(-2, 1)$, $B(1, 5)$ και $\Gamma(5, -1)$.
\begin{alist}    \item Να υπολογίσετε το εμβαδόν του τριγώνου ΑΒΓ. \monades{5}
    \item Να βρείτε την εξίσωση της ευθείας ΒΓ. \monades{5}
    \item Να βρείτε την εξίσωση του ύψους του τριγώνου από την κορυφή Α. Στη συνέχεια να βρείτε το σημείο $\Delta$ της ευθείας ΒΓ, από το οποίο, το Α απέχει την ελάχιστη απόσταση. \monades{8}
    \item Να βρείτε το σύνολο των σημείων Μ του επιπέδου για τα οποία ισχύει:
 \[ (\text{MAB}) = \frac{1}{2} (AB\Gamma) \] \monades{7}\end{alist}

\vspace{4mm}

% --- Άσκηση 22264 (22264.tex) ---
\Askhseis{ΘΕΜΑ 4}
Δίνεται η εξίσωση $x^2 + y^2 + \lambda x + \lambda y + \lambda - 1 = 0$ (1), $\lambda \in \mathbb{R}$.
\begin{alist}    \item Να αποδείξετε ότι η εξίσωση (1), παριστάνει κύκλο για κάθε $\lambda \in \mathbb{R}$. \monades{7}
    \item Να βρείτε το κέντρο και την ακτίνα του κύκλου που ορίζεται από την εξίσωση (1), ο οποίος εφάπτεται της ευθείας $\varepsilon: x + y + 2 = 0$. \monades{9}
    \item Για $\lambda = 1$, στον κύκλο που προκύπτει από την εξίσωση (1), να βρείτε τις εξισώσεις των εφαπτομένων του, που διέρχονται από το σημείο $M\left(-\frac{3}{2}, -\frac{1}{2}\right)$. \monades{9}\end{alist}

\vspace{4mm}

% --- Άσκηση 22265 (22265.tex) ---
\Askhsh{MATHBLYKG04}{4}
Στο καρτεσιανό επίπεδο δίνονται τα σημεία $A(1, -1)$, $B(2, 2)$ και $\Gamma(\mu-1, 3\mu-2)$, $\mu\in\mathbb{R}$.

\begin{alist}    \item Να αποδείξετε ότι καθώς το $\mu$ διατρέχει το $\mathbb{R}$, το σημείο $\Gamma$ κινείται στην ευθεία
    \item Να αποδείξετε ότι καθώς το $\mu$ διατρέχει το $\mathbb{R}$, τα σημεία $A$, $B$, $\Gamma$ είναι κορυφές τριγώνου. \monades{6}  $\epsilon: \gamma = 3x +1$. \monades{6}
    \item Να αποδείξετε ότι το εμβαδόν του τριγώνου $AB\Gamma$ είναι σταθερό. \monades{5}
    \item Να βρείτε τις εξισώσεις των ευθειών που διέρχονται από το σημείο $B$ και από τις οποίες το σημείο $A$, απέχει απόσταση ίση με $1$. \monades{8}\end{alist}

\vspace{4mm}

% --- Άσκηση 22266 (22266.tex) ---
\textbf{ΘΕΜΑ 4}

Δίνεται η εξίσωση $(2\lambda+1)x - (\lambda - 2)y + \lambda - 7 = 0$ (Ε) με $\lambda\in \mathbb{R}$ και η ευθεία ($\zeta$) με εξίσωση:
\[
6x - 8y + 3 = 0
\]

\begin{alist}    \item Να αποδείξετε ότι για κάθε $\lambda\in \mathbb{R}$ η εξίσωση (Ε) παριστάνει ευθεία. \monades{6}
    \item Να αποδείξετε ότι όλες οι ευθείες που ορίζονται από την εξίσωση (Ε), για τα διάφορα $\lambda\in \mathbb{R}$, διέρχονται από το ίδιο σημείο, του οποίου να βρείτε τις συντεταγμένες. \monades{7}
    \item Να βρείτε την τιμή του $\lambda\in \mathbb{R}$, ώστε ευθεία ($\varepsilon$) που ορίζεται από την εξίσωση (Ε) να είναι παράλληλη στην ευθεία ($\zeta$). Ποια είναι η εξίσωση της ($\varepsilon$); \monades{7}
    \item Να βρείτε την απόσταση του σημείου $M(1,3)$ από την ευθεία ($\zeta$). \monades{5}\end{alist}

\vspace{4mm}

% --- Άσκηση 22267 (22267.tex) ---
\Askhsh{Parabolh1}{2}
Δίνεται η εξίσωση $y^2 = 4x$ (1).

\begin{alist}    \item Να γράψετε στο τετράδιό σας συμπληρωμένη την παρακάτω πρόταση:
    \item Να βρείτε την εξίσωση της ευθείας $\varepsilon$ που εφάπτεται στην παραπάνω καμπύλη στο σημείο $\text{A}(1, -2)$. \monades{8}  «Τα σημεία του επιπέδου που επαληθεύουν την εξίσωση (1) βρίσκονται σε μια καμπύλη που ονομάζεται \underline{\hspace{4cm}}. Η εστία της $\text{E}$, έχει συντεταγμένες $\text{E}(\underline{\hspace{1cm}}, \underline{\hspace{1cm}})$ και η διευθετούσα έχει εξίσωση \underline{\hspace{3cm}}». \monades{9}
    \item Να αποδείξετε ότι το σημείο τομής της ευθείας $\varepsilon$ με τον άξονα $x'x$ είναι σημείο της διευθετούσας της παραβολής. \monades{8}\end{alist}

\vspace{4mm}

% --- Άσκηση 22268 (22268.tex) ---
\Askhseis{ΘΕΜΑ 2}

Δίνεται η εξίσωση
\begin{flalign*}
& \frac{x^2}{9} + \frac{y^2}{4} = 1 \quad (1). & &
\end{flalign*}

\begin{alist}    \item Να γράψετε στο τετράδιό σας συμπληρωμένη την παρακάτω πρόταση :
    \item Να βρείτε την εξίσωση της ευθείας $\varepsilon$ η οποία εφάπτεται στην καμπύλη που περιγράφει η εξίσωση (1), στο σημείο της $B(0,-2)$. \monades{10}  «Τα σημεία του επιπέδου που επαληθεύουν την εξίσωση (1) βρίσκονται σε μια καμπύλη που ονομάζεται \ldots. Οι εστίες της $E$ και $E'$, έχουν συντεταγμένες $E(\ldots, \ldots)$ και $E'(\ldots, \ldots)$. Το μήκος του μεγάλου άξονα είναι ίσο με \ldots και η εκκεντρότητα της είναι ίση με \ldots.» \monades{15}\end{alist}

\vspace{4mm}

% --- Άσκηση 22269 (22269.tex) ---
\Askhsh{}{2}
Δίνεται η υπερβολή με εξίσωση:
\[ \frac{x^2}{4} - y^2 = 1 \quad (1) \]
\begin{alist}    \item Να προσδιορίσετε δικαιολογώντας την απάντησή σας :
    \begin{rlist}
        \item Τις συντεταγμένες των εστιών της.
        \item Την εκκεντρότητά της.
        \item Τις εξισώσεις των ασύμπτωτων της υπερβολής.
    \end{rlist}
    \item Να βρείτε την εξίσωση της ευθείας $\varepsilon$ που εφάπτεται στην υπερβολή στο σημείο της, $A\left(\sqrt{5}, \frac{1}{2}\right)$.  \monades{15}
 \monades{10}\end{alist}

\vspace{4mm}

% --- Άσκηση 22273 (22273.tex) ---
\Askhsh{Ex4A}{4}
Δίνεται η έλλειψη με εξίσωση
\[ \frac{x^2}{9} + \frac{y^2}{4} = 1 \quad (1) \]
\begin{alist}
    \item Να προσδιορίσετε δικαιολογώντας την απάντησή σας τις συντεταγμένες :
    \begin{rlist}        \item Των σημείων που η έλλειψη τέμνει τους άξονες $x'x$ και $y'y$.
        \item Των εστιών $E$ και $E'$ της έλλειψης.\end{rlist}
    \monades{12}
    \item Να βρείτε τις εξισώσεις των ευθειών που διέρχονται από το σημείο $A(0, 4)$ και εφάπτονται στη καμπύλη που περιγράφει η εξίσωση (1).
    \monades{13}
\end{alist}

\vspace{4mm}

% --- Άσκηση 22275 (22275.tex) ---
\Askhsh{P1G3}{4} % Placeholder code and theme number
Δίνεται η παραβολή \((C)\) που έχει εξίσωση \(y^2 = 4x\) (1).

\begin{alist}    \item Να σχεδιάσετε πρόχειρα την παραπάνω παραβολή και να γράψετε τις συντεταγμένες της εστίας της \(E\) και την εξίσωση της ευθείας της διευθετούσας \(\delta\). \monades{12}
    \begin{center}
    \begin{tikzpicture}    \begin{axis}[
 axis lines = middle, xlabel = {$x$}, ylabel = {$y$}, xmin = -2.5, xmax = 5.5, ymin = -3.5, ymax = 3.5, xtick = {-2,-1,0,1,2,3,4,5}, ytick = {-3,-2,-1,0,1,2,3}, grid = both, grid style = {gray!30}, samples = 100, width = 8cm, height = 6cm, tick label style = {font=\scriptsize}, every axis x label/.style = {at={(current axis.right of origin)}, anchor=north west}, every axis y label/.style = {at={(current axis.above origin)}, anchor=south}, clip=false ]
        % Parabola y^2 = 4x (i.e., y = +-2*sqrt(x))
 \addplot[very thick, maincolor, domain=0:5, samples=50] {2*sqrt(x)}; \addplot[very thick, maincolor, domain=0:5, samples=50] {-2*sqrt(x)}; \node[maincolor, above right] at (axis cs:3,2.5) {$(C): y^2=4x$};
        % Focus E(1,0)
 \addplot[only marks, mark=*, mark size=1.5pt, black] coordinates {(1,0)}; \node[above right, xshift=2mm, yshift=1mm] at (axis cs:1,0) {$E(1,0)$};
        % Directrix x = -1
 \addplot[dashed, secondarycolor, domain=-3.5:3.5] coordinates {(-1, -3.5) (-1, 3.5)}; \node[secondarycolor, above left, rotate=90, xshift=-1mm] at (axis cs:-1,2) {$\delta: x=-1$};
        % Origin
 \node[below left] at (axis cs:0,0) {$O$};
    \end{axis}\end{tikzpicture}
    \end{center}
    \item Να βρείτε τις εξισώσεις των ευθειών που διέρχονται από το σημείο \(A(0, 2)\) και εφάπτονται στην παραβολή που περιγράφει η εξίσωση (1). \monades{13}  Οι συντεταγμένες της εστίας είναι \(E(1,0)\) και η εξίσωση της διευθετούσας είναι \(\delta: x=-1\).\end{alist}

\vspace{4mm}

% --- Άσκηση 22279 (22279.tex) ---
\OnPar{ΘΕΜΑ 2}
Δίνεται η εξίσωση
\[ (y - 1)^2 = (3 + x)(1 - x) \quad (1) \]
Να αποδείξετε ότι:
\begin{alist}    \item Η εξίσωση (1) παριστάνει κύκλο με κέντρο $K(-1,1)$ και ακτίνα $R = 2$. \monades{9}
    \item Η αρχή $O(0,0)$ των αξόνων είναι εσωτερικό σημείο του κύκλου $(K,R)$. \monades{7}
    \item Η ευθεία $(\varepsilon): x + y = 2$ είναι τέμνουσα του κύκλου $(K,R)$. \monades{9}\end{alist}

\vspace{4mm}

% --- Άσκηση 22280 (22280.tex) ---
\Askhsh{}{4}
Σε ορθοκανονικό σύστημα αξόνων με αρχή το σημείο $O(0,0)$ θεωρούμε τους κύκλους $(K, R)$ και $(\Lambda, \rho)$ με εξισώσεις
\begin{align}
x^2 + y^2 - 6x - 8y + 21 &= 0 \tag{1} \\
x^2 + y^2 + 2x - 2y + 1 &= 0 \tag{2}
\end{align}
αντίστοιχα.

\begin{alist}    \item Να βρείτε τα κέντρα και τις ακτίνες των δύο κύκλων. \monades{12}
    \item Να αποδείξετε ότι οι δύο κύκλοι βρίσκονται ο ένας εξωτερικά του άλλου. \monades{08}
    \item Έστω $M$, $N$ τυχαία σημεία των κύκλων $(K, R)$ και $(\Lambda, \rho)$ αντίστοιχα. Να υπολογίσετε την ελάχιστη και την μέγιστη απόσταση των σημείων $M$ και $N$. \monades{05}\end{alist}

\vspace{4mm}

% --- Άσκηση 22465 (22465.tex) ---
\Askhsh{ALGBP04}{4}
Έστω παραβολή $C$ με κορυφή την αρχή των αξόνων $O$ και άξονα συμμετρίας τον $x'x$. Η απόσταση της εστίας $E$ από την διευθετούσα $\delta$ της παραβολής $C$ είναι $4$ και η γραφική της παράσταση φαίνεται στο παρακάτω ορθοκανονικό σύστημα αξόνων.

\begin{alist}    \item Να δικαιολογήσετε ότι η εστία της είναι η $E(2,0)$, η διευθετούσα της είναι η $\delta : x=-2$ και η εξίσωσή της παραβολής είναι $y^2=8x$. \monades{9}
    \item Να αποδείξετε ότι η εξίσωση της εφαπτομένης της παραβολής στο σημείο της $A(2,4)$ είναι η $\varepsilon: y=x+2$. \monades{9}
    \item Να βρείτε την εξίσωση του κύκλου που διέρχεται από την εστία της παραβολής και εφάπτεται στην ευθεία $\varepsilon$ στο σημείο της $A(2,4)$. \monades{7}\end{alist}

\begin{center}
\begin{tikzpicture}[scale=0.8]
    % Σχεδίαση Αξόνων
    \draw[->, l3] (-3.5,0) -- (7.5,0) node[below] {$x$};
    \draw[->, l3] (0,-5.5) -- (0,5.5) node[left] {$y$};
    % Σημείο Αρχής O
    \node[below left, font=\small] at (0,0) {O};
    \fill (0,0) circle (1.5pt); % Κουκκίδα για το O
    % Εστία E
    \node[below right, font=\small] at (2,0) {E};
    \fill (2,0) circle (1.5pt); % Κουκκίδα για το E
    % Διευθετούσα δ: x = -2
    \draw[secondarycolor, thick, dashed] (-2,-5.5) -- (-2,5.5);
    \node[above, secondarycolor, font=\small] at (-2,5.2) {διευθετούσα};
    % Σύμβολο ορθής γωνίας στο (-2,0)
    \draw[black, line width=0.3pt] (-2,0) rectangle (-1.8,0.2);
    % Παραβολή y^2 = 8x => y = +/- sqrt(8x)
    \draw[maincolor, very thick, samples=200, domain=0:7] plot (\x, {sqrt(8*\x)});
    \draw[maincolor, very thick, samples=200, domain=0:7] plot (\x, {-sqrt(8*\x)});\end{tikzpicture}
\end{center}

\vspace{4mm}

% --- Άσκηση 22508 (22508.tex) ---
\Askhsh{}{4}

Οι κορυφές $A$, $\Gamma$ ενός τετραγώνου $AB\Gamma\Delta$ είναι τα σημεία $(1,4)$ και $(3,0)$ αντιστοίχως.

\begin{alist}    \item Να αποδείξετε ότι η εξίσωση της μεσοκαθέτου του ευθυγράμμου τμήματος $A\Gamma$ γράφεται στη μορφή
    \item Να αποδείξετε ότι η εξίσωση του κύκλου, ο οποίος έχει διάμετρο το ευθύγραμμο τμήμα $A\Gamma$ γράφεται στη μορφή  \[ y - 2 = \frac{1}{2} (x - 2). \] \monades{8}
    \item Να υπολογίσετε τις συντεταγμένες των δύο άλλων κορυφών $B, \Delta$ του τετραγώνου.  \[ (x - 2)^2 + (y - 2)^2 = 5. \] \monades{8}
 \monades{9}\end{alist}

\vspace{4mm}

% --- Άσκηση 22554 (22554.tex) ---
ΘΕΜΑ 2

Στο καρτεσιανό επίπεδο $Oxy$, με μοναδιαία διανύσματα των αξόνων $x'x$, $y'y$ τα $\vec{i}$, $\vec{j}$ αντίστοιχα, τα σημεία Α και Β έχουν διανύσματα θέσεως $\overrightarrow{OA} = 3\vec{i} + 2\vec{j}$ και $\overrightarrow{OB} = 6\vec{i} - \vec{j}$.

Έστω Μ ένα σημείο τέτοιο ώστε $\overrightarrow{OM} = \frac{1}{5} (2\overrightarrow{OA} - \overrightarrow{OB})$.

\begin{alist}    \item Να αποδείξετε ότι:
    \begin{rlist}
        \item $\overrightarrow{AB} = 3\vec{i} - 3\vec{j}$ \monades{8}
        \item $\overrightarrow{OM} = \vec{j}$ \monades{8}
    \end{rlist}
    \item Να βρείτε το εσωτερικό γινόμενο $\overrightarrow{AB} \cdot \overrightarrow{OM}$. \monades{9}\end{alist}

\vspace{4mm}

% --- Άσκηση 22556 (22556.tex) ---
\Askhsh{Μ22222}{2}
\textbf{ΘΕΜΑ 2}
\par
Η έλλειψη του παρακάτω σχήματος έχει κορυφές τα σημεία $\text{A}'(-5,0)$, $\text{A}(5,0)$, $\text{B}'(0, -4)$ και $\text{B}(0,4)$.

\begin{alist}    \item Να αποδείξετε ότι:
    \begin{rlist}
        \item Τα μήκη των αξόνων της έλλειψης είναι $(\text{A}'\text{A}) = 10$ και $(\text{B}'\text{B}) = 8$. \monades{10}
        \item Οι εστίες της έλλειψης είναι τα σημεία $\text{E}'(-3,0)$ και $\text{E}(3,0)$. \monades{10}
    \end{rlist}
    \item Έστω $\text{M}$ ένα σημείο της έλλειψης. Να αποδείξετε ότι $(\text{M}\text{E}') + (\text{M}\text{E}) = 10$. \monades{5}\end{alist}

\begin{center}
\begin{tikzpicture}    % Define the center of the ellipse
    \tkzDefPoint(0,0){O}
    % Define vertices
    \tkzDefPoint(-5,0){Ap} % A'
    \tkzDefPoint(5,0){A}
    \tkzDefPoint(0,-4){Bp} % B'
    \tkzDefPoint(0,4){B}
    % Define foci (for drawing purposes, as they are to be proven in part ii)
    \tkzDefPoint(-3,0){Ep} % E'
    \tkzDefPoint(3,0){E}
    % Define point M on the ellipse (x^2/25 + y^2/16 = 1)
    % For x=3, y = sqrt(16 * (1 - 9/25)) = sqrt(16 * 16/25) = 16/5 = 3.2
    \tkzDefPoint(3,3.2){M}
    % Draw the ellipse
    \draw[l3, maincolor] (O) ellipse (5cm and 4cm);
    % Draw x-axis
    \draw[l3, ->, black] (-6,0) -- (6,0) node[below] {$x$};
    % Draw y-axis
    \draw[l3, ->, black] (0,-5) -- (0,5) node[right] {$y$};
    % Draw segments ME' and ME
    \tkzDrawSegment[l4, secondarycolor](M,Ep)
    \tkzDrawSegment[l4, secondarycolor](M,E)
    % Draw points
    \tkzDrawPoints(O,A,Ap,B,Bp,E,Ep,M)
    % Labels for points
    \tkzLabelPoint[below right](O){O}
    \tkzLabelPoint[left](Ap){$\text{A}'$}
    \tkzLabelPoint[right](A){$\text{A}$}
    \tkzLabelPoint[above left](B){$\text{B}$}
    \tkzLabelPoint[below left](Bp){$\text{B}'$}
    \tkzLabelPoint[below left](Ep){$\text{E}'$}
    \tkzLabelPoint[below right](E){$\text{E}$}
    \tkzLabelPoint[above right](M){$\text{M}$}\end{tikzpicture}
\end{center}

\vspace{4mm}

% --- Άσκηση 22557 (22557.tex) ---
\subsection*{ΘΕΜΑ 2}
Το ισοσκελές τρίγωνο $AB\Gamma$ έχει βάση $B\Gamma$ και ύψος $AO$. Η κορυφή $A$ είναι σημείο του θετικού ημιάξονα $Oy$ και οι κορυφές $B$ και $\Gamma$ είναι σημεία του άξονα $x'x$, όπως φαίνεται στο παρακάτω σχήμα. Έστω $(B\Gamma) = 12$, $(AO) = 8$ και $M$ το μέσο της πλευράς $A\Gamma$.

\begin{alist}    \item Να αποδείξετε ότι:
    \begin{rlist}
        \item $A(0,8)$, $B(-6,0)$ και $\Gamma(6,0)$. \monades{9}
        \item $M(3,4)$ \monades{9}
    \end{rlist}
    \item Να βρείτε το μήκος της διαμέσου $BM$. \monades{7}\end{alist}

\begin{center}
    \begin{tikzpicture}[scale=0.7]
        % Define points
        \tkzDefPoint(0,0){O}
        \tkzDefPoint(0,8){A}
        \tkzDefPoint(-6,0){B}
        \tkzDefPoint(6,0){C} % \Gamma
        \tkzDefMidPoint(A,C) \tkzGetPoint{M} % M is midpoint of AC
        % Draw axes
        \draw[->] (-7,0) -- (7,0) node[below] {$x$};
        \draw[->] (0,-1) -- (0,9) node[left] {$y$};
        % Draw the triangle and segments
        \tkzDrawSegments[l3](A,B B,C C,A A,O B,M)
        % Draw points
        \tkzDrawPoints(A,B,C,O,M)
        % Labels for points
        \tkzLabelPoint[above](A){A}
        \tkzLabelPoint[left](B){B}
        \tkzLabelPoint[right](C){$\Gamma$}
        \tkzLabelPoint[below left](O){O}
        \tkzLabelPoint[above right](M){M}
        % Right angle symbol at O (AO perpendicular to BC)
        \tkzMarkRightAngle[fill=thirdcolor!20,size=0.3](A,O,C) % Or A,O,B
\end{tikzpicture}
\end{center}

\vspace{4mm}

% --- Άσκηση 22558 (22558.tex) ---
\Askhsh{ΜΑΘ120022}{\textbf{2}}

\noindent Η έλλειψη του παρακάτω σχήματος έχει εστίες τα σημεία $Ε'(-2,0)$ και $Ε(2,0)$ και μήκος μεγάλου άξονα $(Α'Α) = 8$.

\begin{alist}    \item Να αποδείξετε ότι η έλλειψη έχει εξίσωση $\frac{x^2}{16} + \frac{y^2}{12} = 1$. \monades{12}
    \item Έστω $\Sigma$ και $P$ τα σημεία της έλλειψης που έχουν την ίδια τετμημένη με την εστία $Ε(2,0)$. Επίσης το $\Sigma$ έχει θετική τεταγμένη και το $P$ αρνητική τεταγμένη.
    \begin{rlist}
        \item Να αποδείξετε ότι $\Sigma(2, 3)$ και $P(2,-3)$. \monades{8}
        \item Να βρείτε το μήκος του τμήματος $\Sigma P$. \monades{5}
    \end{rlist}\end{alist}

\vspace{1cm}

\begin{center}
    \begin{tikzpicture}        % Coordinate system with arrows
        \draw[->, l3, black] (-5,0) -- (5,0) node[right] {$x$};
        \draw[->, l3, black] (0,-4) -- (0,4) node[above] {$y$};
        % Define points
        \tkzDefPoint(0,0){O}
        \tkzDefPoint(4,0){A}
        \tkzDefPoint(-4,0){Ap} % A'
        \tkzDefPoint(2,0){E}
        \tkzDefPoint(-2,0){Ep} % E'
        \tkzDefPoint(2,3){Sigma} % Σ
        \tkzDefPoint(2,-3){P}
        % Draw the ellipse
        % Center (0,0), a=4, b=sqrt(12) approx 3.46
        \draw[l4, maincolor] (0,0) ellipse (4cm and {sqrt(12)}cm);
        % Draw points
        \tkzDrawPoints(O,A,Ap,E,Ep,Sigma,P)
        % Labels for points
        \tkzLabelPoint[below left](O){O}
        \tkzLabelPoint[below right](A){A}
        \tkzLabelPoint[below left](Ap){A'}
        \tkzLabelPoint[below right](E){E}
        \tkzLabelPoint[below left](Ep){E'}
        \tkzLabelPoint[above right](Sigma){$\Sigma$}
        \tkzLabelPoint[below right](P){P}
        % Perpendicular segment and right angle mark
        \tkzDrawSegment[l3, thirdcolor](E,Sigma)
        \tkzMarkRightAngle[size=0.2,fill=thirdcolor!30](O,E,Sigma)\end{tikzpicture}
\end{center}

\vspace{4mm}

% --- Άσκηση 22559 (22559.tex) ---
\Askhsh{}{2}
Η υπερβολή στο παρακάτω σχήμα έχει εστίες τα σημεία $E'(-10,0)$ και $E(10,0)$ και κορυφές τα σημεία $A'(-8,0)$ και $A(8,0)$.

\begin{alist}    \item Να αποδείξετε ότι η υπερβολή έχει εξίσωση $\frac{x^2}{64} - \frac{y^2}{36} = 1$. \monades{12}
    \item Έστω $M$ ένα σημείο της υπερβολής.
    \begin{rlist}
        \item Να αποδείξετε ότι $\left|\left(\text{ME}'\right) - \left(\text{ME}\right)\right| = 16$. \monades{8}
        \item Αν $(\text{ME}) = 9$, να βρείτε την απόσταση του σημείου $M$ από την εστία $\text{E}'$. \monades{5}
    \end{rlist}\end{alist}

\begin{center}
\begin{tikzpicture}[scale=0.7]
    % Define a scaling factor for visual representation to fit in the page
    \def\myScale{0.3}
    % Define points based on the given coordinates and scaling factor
    \coordinate (O) at (0,0);
    \coordinate (Eprime) at (\myScale*-10,0);
    \coordinate (Aprime) at (\myScale*-8,0);
    \coordinate (A) at (\myScale*8,0);
    \coordinate (E) at (\myScale*10,0);
    % Point M on the hyperbola. Using x=10, y=4.5 which satisfies the equation x^2/64 - y^2/36 = 1
    % (10^2/64 - 4.5^2/36 = 100/64 - 20.25/36 = 25/16 - 9/16 = 16/16 = 1)
    \coordinate (M) at (\myScale*10, \myScale*4.5);
    % Draw axes
    \draw[->] (\myScale*-12,0) -- (\myScale*12,0) node[below] {$x$};
    \draw[->] (0,\myScale*-7) -- (0,\myScale*7) node[left] {$y$};
    % Draw hyperbola branches (using path with controls to mimic the sketch style)
    % Right branch (through A)
    \draw[l3, maincolor] (A) .. controls (\myScale*9.5, \myScale*2) and (\myScale*11, \myScale*4) .. (\myScale*12, \myScale*6);
    \draw[l3, maincolor] (A) .. controls (\myScale*9.5, \myScale*-2) and (\myScale*11, \myScale*-4) .. (\myScale*12, \myScale*-6);
    % Left branch (through A')
    \draw[l3, maincolor] (Aprime) .. controls (\myScale*-9.5, \myScale*2) and (\myScale*-11, \myScale*4) .. (\myScale*-12, \myScale*6);
    \draw[l3, maincolor] (Aprime) .. controls (\myScale*-9.5, \myScale*-2) and (\myScale*-11, \myScale*-4) .. (\myScale*-12, \myScale*-6);
    % Draw segments ME' and ME
    \draw[l4, secondarycolor] (M) -- (Eprime);
    \draw[l4, secondarycolor] (M) -- (E);
    % Draw points and labels
    \tkzDrawPoints(O, A, Aprime, E, Eprime, M)
    \tkzLabelPoint[below](O){O}
    \tkzLabelPoint[below](A){A}
    \tkzLabelPoint[below](Aprime){A'}
    \tkzLabelPoint[below](E){E}
    \tkzLabelPoint[below](Eprime){E'}
    \tkzLabelPoint[above right](M){M}\end{tikzpicture}
\end{center}

\vspace{4mm}

% --- Άσκηση 22561 (22561.tex) ---
\OnPar{ΘΕΜΑ 2}

Στο παρακάτω σχήμα η υπερβολή $C$ έχει εξίσωση $x^2 - y^2 = 9$, οι ευθείες $\delta_1$ και $\delta_2$ είναι οι ασύμπτωτες της $C$ και η $\epsilon$ είναι η εφαπτομένη της $C$ στο σημείο της $M(5,4)$.

\begin{alist}
\item Να αποδείξετε ότι:
    \begin{rlist}    \item Οι εξισώσεις των ασυμπτώτων είναι $\delta_1: y = x$ και $\delta_2: y = -x$. \monades{8}
    \item Η εξίσωση της εφαπτομένης στο $M$ είναι $\epsilon: 5x - 4y = 9$. \monades{8}\end{rlist}
\item Να βρείτε τις συντεταγμένες του σημείου τομής των ευθειών $\epsilon$ και $\delta_1$ καθώς και τις συντεταγμένες του σημείου τομής των ευθειών $\epsilon$ και $\delta_2$. \monades{9}
\end{alist}

\vspace{1cm} % Add some vertical space before the figure

\begin{tikzpicture}\begin{axis}[
    axis lines = middle,
    xmin = -7, xmax = 7,
    ymin = -6, ymax = 6,
    xtick = \empty, 
    ytick = \empty, 
    grid = none,
    width = 10cm, height = 7cm,
    tick label style = {font=\scriptsize},
]
    % Hyperbola x^2 - y^2 = 9 => y = +-sqrt(x^2 - 9)
    \addplot[thick, maincolor, domain=3:7, samples=100] {sqrt(x^2 - 9)}; % Upper right branch
    \addplot[thick, maincolor, domain=3:7, samples=100] {-sqrt(x^2 - 9)}; % Lower right branch
    \addplot[thick, maincolor, domain=-7:-3, samples=100] {sqrt(x^2 - 9)}; % Upper left branch
    \addplot[thick, maincolor, domain=-7:-3, samples=100] {-sqrt(x^2 - 9)}; % Lower left branch
    % Asymptotes: delta1: y=x, delta2: y=-x
    % Use l3 for line thickness as defined in the preamble.
    \addplot[l3, secondarycolor, dotted, domain=-7:7] {x}; % delta1 (y=x)
    \addplot[l3, secondarycolor, dotted, domain=-7:7] {-x}; % delta2 (y=-x)
    % Tangent: epsilon: 5x - 4y = 9 => y = (5/4)x - (9/4)
    \addplot[l3, thirdcolor, domain=-1:7] {(5/4)*x - (9/4)}; % epsilon
    % Points O and M
    \coordinate (O) at (axis cs:0,0);
    \coordinate (M) at (axis cs:5,4);
    % Draw points and labels
    \fill (O) circle (1.5pt); % Point O (origin)
    \node[below left] at (O) {O};
    \fill (M) circle (2pt); % Point M(5,4)
    \node[above right] at (M) {M(5,4)};
    % Labels for asymptotes and tangent, adjusted positions to match PDF visual
    \node[below left, secondarycolor] at (axis cs: -3.5, -3.5) {$\delta1$}; % Label for y=x
    \node[above left, secondarycolor] at (axis cs: -3.5, 3.5) {$\delta2$}; % Label for y=-x
    \node[below right, thirdcolor] at (axis cs: 0.5, -2.5) {$\epsilon$}; % Label for tangent epsilon
    % Label C on the hyperbola (beyond M(5,4))
    \node[above right, maincolor] at (axis cs: 6.5, {sqrt(6.5^2 - 9)}) {C};
\end{axis}\end{tikzpicture}

\vspace{4mm}

% --- Άσκηση 22564 (22564.tex) ---
\Askhsh{M2ELP1}{2}
Δίνονται οι ελλείψεις $C_1: \frac{x^2}{12} + \frac{y^2}{7} = 1$ και $C_2: \frac{x^2}{6} + \frac{y^2}{12} = 1$.

\begin{alist}    \item Να αποδείξετε ότι τα σημεία $A(2,2)$ και $B(-2,2)$ ανήκουν και στις δύο ελλείψεις.
    \item Να αποδείξετε ότι η εξίσωση της εφαπτομένης $\varepsilon1$ της έλλειψης $C_1$ στο σημείο $A$ και η εξίσωση της εφαπτομένης $\varepsilon2$ της έλλειψης $C_2$ στο σημείο $B$ είναι αντίστοιχα $x + 2y - 6 = 0$ και $-2x + y - 6 = 0$.  \monades{10}
    \item Να αποδείξετε ότι οι εφαπτομένες $\varepsilon1, \varepsilon2$ είναι κάθετες.  \monades{10}
 \monades{5}\end{alist}

\begin{center}
\begin{tikzpicture}\begin{axis}[
    axis lines = middle, 
    xmin = -4.5, xmax = 4.5,
    ymin = -3.5, ymax = 3.5,
    xtick = {-4,-3,-2,-1,0,1,2,3,4},
    ytick = {-3,-2,-1,0,1,2,3},
    tick label style = {font=\scriptsize},
    axis line style={black}, 
    width = 9cm, height = 7cm,
]
    % Έλλειψη C1: x^2/12 + y^2/7 = 1
    % Ημιάξονες: a = sqrt(12), b = sqrt(7)
    \addplot[domain=0:360, samples=100, black, thick] ({sqrt(12)*cos(x)}, {sqrt(7)*sin(x)});
    % Έλλειψη C2: x^2/6 + y^2/12 = 1
    % Ημιάξονες: a = sqrt(6), b = sqrt(12)
    \addplot[domain=0:360, samples=100, black, thick] ({sqrt(6)*cos(x)}, {sqrt(12)*sin(x)});
    % Σημεία A(2,2) και B(-2,2)
    \addplot[only marks, mark=*, mark size=1.5pt, black] coordinates {(2,2) (-2,2)};
    \node[above right, font=\small] at (axis cs:2,2) {$A$};
    \node[above left, font=\small] at (axis cs:-2,2) {$B$};
    % Ετικέτες για τις ελλείψεις C1 και C2
    % Τοποθέτηση βάσει της εικόνας PDF
    \node[left, font=\small] at (axis cs:-3, 2.2) {$C1$};
    \node[below, font=\small] at (axis cs:0, -2.8) {$C2$};
\end{axis}\end{tikzpicture}
\end{center}

\vspace{4mm}

% --- Άσκηση 22566 (22566.tex) ---
\OnPar{ΘΕΜΑ 2}
Δίνεται η υπερβολή με εξίσωση $4x^2 - y^2 = 4$.

\begin{alist}    \item Να αποδείξετε ότι οι συντεταγμένες της κορυφής της υπερβολής είναι $A(1,0)$ και $A'(-1,0)$. \monades{8}
    \item Να αποδείξετε ότι οι ασύμπτωτες της υπερβολής είναι οι $y = 2x$ και $y = -2x$. \monades{9}
    \item Να αποδείξετε ότι η ευθεία που διέρχεται από την κορυφή $A$ και είναι παράλληλη προς την ασύμπτωτη $y = -2x$ έχει εξίσωση $y = -2x + 2$. \monades{8}\end{alist}

\vspace{4mm}

% --- Άσκηση 22567 (22567.tex) ---
{\OnPar{ΘΕΜΑ 2}}

Στο καρτεσιανό επίπεδο \textgreek{Oxy} η υπερβολή C:
\[ \frac{x^2}{\alpha^2} - \frac{y^2}{\beta^2} = 1 \]
τέμνει τον άξονα $x'x$ στα σημεία $A'(-4,0)$ και $A(4,0)$ και έχει ασύμπτωτες τις ευθείες $y = \frac{3}{4}x$ και $y = -\frac{3}{4}x$.

\begin{alist}    \item Να αποδείξετε ότι:
    \begin{rlist}
        \item $\alpha = 4$ και $\beta = 3$, \monades{10}
        \item οι εστίες της C είναι τα σημεία $E'(-5,0)$ και $E(5,0)$. \monades{10}
    \end{rlist}
    \item Να σχεδιάσετε το παρακάτω σχήμα, συμπληρώνοντάς το με την παραπάνω υπερβολή C. \monades{5}\end{alist}

\begin{tikzpicture}    % Set up the coordinate system
    \draw[->, l3, black] (-7,0) -- (7,0) node[below right, black] {$x$};
    \draw[->, l3, black] (0,-4.5) -- (0,4.5) node[above left, black] {$y$};
    \node[below left, black, inner sep=1pt] at (0,0) {O};
    % Asymptotes
    \draw[l3, secondarycolor, domain=-6:6, samples=2] plot (\x, {3/4 * \x}) node[above left, xshift=-2mm, yshift=2mm, font=\scriptsize, rotate=38, text=secondarycolor] {$y = \frac{3}{4}x$};
    \draw[l3, secondarycolor, domain=-6:6, samples=2] plot (\x, {-3/4 * \x}) node[below left, xshift=-2mm, yshift=-2mm, font=\scriptsize, rotate=-38, text=secondarycolor] {$y = -\frac{3}{4}x$};
    % Points A' and A
    \tkzDefPoint(-4,0){Ap}
    \tkzDefPoint(4,0){A}
    \tkzDrawPoints(Ap,A)
    \tkzLabelPoint[below left](Ap){$A'(-4, 0)$}
    \tkzLabelPoint[below right](A){$A(4, 0)$}
    % Hyperbola C: x^2/16 - y^2/9 = 1 => y = +/- (3/4) * sqrt(x^2 - 16)
    % Domain for plotting: x from 4 to 6.5 (right branch) and -6.5 to -4 (left branch)
    \draw[l4, thirdcolor, domain=4:6.5, samples=100] plot (\x, {3/4 * sqrt(\x*\x - 16)});
    \draw[l4, thirdcolor, domain=4:6.5, samples=100] plot (\x, {-3/4 * sqrt(\x*\x - 16)});
    \draw[l4, thirdcolor, domain=-6.5:-4, samples=100] plot (\x, {3/4 * sqrt(\x*\x - 16)});
    \draw[l4, thirdcolor, domain=-6.5:-4, samples=100] plot (\x, {-3/4 * sqrt(\x*\x - 16)});\end{tikzpicture}

\vspace{4mm}

% --- Άσκηση 32206 (32206.tex) ---
\Askhsh{4}{4}
Η υπερβολή $C$ έχει εστίες τα σημεία $E(5,0), E'(-5,0)$ και διέρχεται από το σημείο $M\left(5, \frac{9}{4}\right)$.

\begin{alist}    \item Να αποδείξετε ότι έχει εκκεντρότητα $\frac{5}{4}$. \monades{7}
    \item Να βρείτε την εξίσωση της $C$. \monades{6}
    \item Να βρείτε την εξίσωση της διχοτόμου της γωνίας $EME'$. \monades{6}
    \item Να βρείτε το συνημίτονο της οξείας γωνίας που σχηματίζουν οι ασύμπτωτές της. \monades{6}\end{alist}
Δίνεται ότι $\sqrt{1681} = 41$.

\vspace{4mm}

% --- Άσκηση 33696 (33696.tex) ---
\Askhsh{}{4}
Το κέντρο $K$ ενός κύκλου $(c)$ βρίσκεται στο πρώτο τεταρτημόριο και είναι σημείο της ευθείας $(\varepsilon)$: $y=2x - 1$. Ο κύκλος $(c)$ έχει ακτίνα $\rho=3\sqrt{2}$ και η ευθεία $(\zeta)$: $x + y -2 = 0$ εφάπτεται στον κύκλο στο σημείο $A$.

\begin{alist}
\item Να αποδείξετε ότι η εξίσωση του κύκλου $(c)$ είναι $(x - 3)^2 + (y - 5)^2 = 18$.
\monades{09}

\item Να αποδείξετε ότι:
    \begin{rlist}    \item Η εξίσωση της ευθείας $KA$ είναι η $x - y + 2 = 0$.
    \item Οι συντεταγμένες του $A$ είναι $(0,2)$.  \monades{05}
 \monades{04}\end{rlist}

\item Να υπολογισθεί το εμβαδόν του τριγώνου $A\Lambda M$, όπου $M$ και $\Lambda$ είναι τα σημεία τομής της ευθείας $(\varepsilon)$ με τον κύκλο $(c)$.
\monades{07}
\end{alist}

\vspace{4mm}

\end{document}
